% Integrated revision: source results through twentieth-cycle V72.
% ==================== 00_front.tex ====================
\documentclass[11pt,a4paper,oneside]{book}
\usepackage[T1]{fontenc}
\usepackage[utf8]{inputenc}
\usepackage{lmodern,microtype}
\usepackage[margin=28mm,headheight=15pt]{geometry}
\usepackage{amsmath,amssymb,amsthm,mathtools,bm,mathrsfs}
\usepackage{booktabs,longtable,tabularx,array,enumitem}
\usepackage{xcolor,graphicx,needspace,fancyhdr}
\usepackage[unicode=true,colorlinks=true,linkcolor=black,citecolor=black,urlcolor=black,bookmarksnumbered=true]{hyperref}
\usepackage{bookmark}
\hypersetup{pdftitle={Renewal Standard Model--Einstein Continuum Emergence},pdfauthor={A. Pelissier},pdfsubject={Conditional reconstruction, continuum limits, and coupled dynamics}}
\setlength{\parindent}{1.2em}
\setlength{\parskip}{2pt}
\linespread{1.035}
\setcounter{tocdepth}{1}
\setcounter{secnumdepth}{2}
\setlist{topsep=5pt,itemsep=3pt,parsep=0pt,leftmargin=2.2em}
\setlength{\emergencystretch}{3em}
\allowdisplaybreaks[2]
\pagestyle{fancy}
\fancyhf{}
\fancyhead[L]{\small\nouppercase{\leftmark}}
\fancyfoot[C]{\thepage}
\renewcommand{\headrulewidth}{0.3pt}
\renewcommand{\chaptermark}[1]{\markboth{\thechapter.\ #1}{}}
\fancypagestyle{plain}{\fancyhf{}\fancyfoot[C]{\thepage}\renewcommand{\headrulewidth}{0pt}}
\theoremstyle{plain}
\newtheorem{theorem}{Theorem}[chapter]
\newtheorem{proposition}[theorem]{Proposition}
\newtheorem{lemma}[theorem]{Lemma}
\newtheorem{corollary}[theorem]{Corollary}
\theoremstyle{definition}
\newtheorem{definition}[theorem]{Definition}
\newtheorem{assumption}[theorem]{Hypothesis}
\newtheorem{criterion}[theorem]{Criterion}
\theoremstyle{remark}
\newtheorem{remark}[theorem]{Remark}
\newtheorem{obstruction}[theorem]{Obstruction}
\newcommand{\R}{\mathbb R}
\newcommand{\C}{\mathbb C}
\newcommand{\N}{\mathbb N}
\newcommand{\Z}{\mathbb Z}
\newcommand{\E}{\mathbb E}
\newcommand{\HH}{\mathcal H}
\renewcommand{\AA}{\mathcal A}
\newcommand{\DD}{\mathcal D}
\newcommand{\XX}{\mathcal X}
\newcommand{\LL}{\mathcal L}
\newcommand{\FF}{\mathcal F}
\newcommand{\SSS}{\mathcal S}
\newcommand{\one}{\mathbf 1}
\newcommand{\id}{\mathrm{id}}
\newcommand{\dd}{\mathrm d}
\newcommand{\ii}{\mathrm i}
\newcommand{\Tr}{\operatorname{Tr}}
\newcommand{\tr}{\operatorname{tr}}
\newcommand{\Cov}{\operatorname{Cov}}
\newcommand{\Var}{\operatorname{Var}}
\newcommand{\Cum}{\operatorname{Cum}}
\newcommand{\Ran}{\operatorname{Ran}}
\newcommand{\supp}{\operatorname{supp}}
\newcommand{\rank}{\operatorname{rank}}
\newcommand{\spec}{\operatorname{spec}}
\newcommand{\dist}{\operatorname{dist}}
\newcommand{\Ric}{\operatorname{Ric}}
\newcommand{\vol}{\operatorname{vol}}
\newcommand{\sgn}{\operatorname{sgn}}
\newcommand{\diag}{\operatorname{diag}}
\newcommand{\Ad}{\operatorname{Ad}}
\newcommand{\Pf}{\operatorname{Pf}}
\newcommand{\Hom}{\operatorname{Hom}}
\newcommand{\End}{\operatorname{End}}
\newcommand{\GSM}{G_{\mathrm{SM}}}
\newcommand{\cO}{\mathcal O}
\newcommand{\norm}[1]{\left\lVert#1\right\rVert}
\newcommand{\abs}[1]{\left\lvert#1\right\rvert}
\newcommand{\pair}[2]{\left\langle#1,#2\right\rangle}
% The second argument preserves the original result identifier for traceability.
\newcommand{\src}[2]{\cite{#1}}
\newcommand{\scope}[1]{\par\smallskip\noindent\textit{Scope.} #1\par\smallskip}
\newcommand{\gate}[1]{\par\smallskip\noindent\textit{Remaining input.} #1\par\smallskip}
\newcommand{\proofsketch}{\noindent\textit{Proof mechanism.}\ }
\begin{document}
\frontmatter
\begin{titlepage}
\centering
\vspace*{18mm}
{\small RENEWAL GEOMETRY\par}
\vspace{13mm}
{\Huge\bfseries Renewal Standard Model--Einstein\par}
\vspace{4mm}
{\Huge\bfseries Continuum Emergence\par}
\vspace{13mm}
{\Large Conditional reconstruction, uniform limits,\par coupled fields, and gravitational dynamics\par}
\vspace{17mm}
{\large A.\ P\'elissier\par}
\vfill
\begin{minipage}{0.81\textwidth}
\small
A standalone mathematical synthesis of the Grand Tensor construction and the Standard Model--Einstein continuum programme. Finite reconstruction, analytic transfer theorems, explicit coupled models, and unresolved continuum hypotheses are kept logically distinct.
\end{minipage}
\vspace{14mm}
\par September 2026
\end{titlepage}
\chapter*{Abstract}
\addcontentsline{toc}{chapter}{Abstract}
The renewal construction starts with operational histories rather than a prescribed spacetime or quantum field theory. A finite predictive quotient and its all-order positive word structure organize persistent degrees of freedom. Under the specified local geometric and internal representation hypotheses, the resulting carrier supports a $K_4/A_3$ spatial branch, calibrated Lorentzian kinematics, the faithful Standard Model gauge group, three structural generations, and a common regulated gauge--Higgs--fermion action. These finite results are the starting point, not the conclusion, of continuum emergence.

This monograph develops the conditional passage from that carrier to a physically reconstructed Standard Model coupled to an Einstein metric. It assembles chiral index and determinant-line results, exact shell integration, source-preserving cutoff comparison, full-face rigid coercivity, positive constrained gravitational sectors, locally normal joint states, reflection positivity, faithful detail-clock reconstruction, and distributional stress descent. Stable quartic reference forms and dynamical holonomy coordinates separate boundary nuclearity from phase, interface, and metric-normalization requirements. A positive Higgs--gauge--conformal lattice branch admits volume-uniform local moments and a translation-invariant thermodynamic Gibbs limit at fixed spacing. Differentiated shell remainders, incidence-sharp curvature bounds, and supported Dirichlet responses identify the additional estimates needed for cutoff-uniform transport. The analysis also records the obstructions that determine the admissible routes: the real Euclidean conformal instability, the inadequacy of a purely local anomaly calculation, the loss of physical clocks under annihilating refinement, the distinction between massive and massless sectors, the possibility that uniform mixing produces a Gaussian limit, and the failure of smooth finite charts to control typical fields of an explicit projective gauge law.

A regular-branch composition theorem identifies ten simultaneous requirements for one cofinal family. Under these requirements the limiting physical Schwinger functional is route-independent, its compiled stress is conserved, the limiting critical metric satisfies the distributional Einstein equation, and a separated connected-response excess certifies non-quasi-free interaction. The full set of requirements has not been established for one physical Standard Model--Einstein regulator family. The final part therefore treats exact and perturbative coupled models on their own terms: compact transverse-traceless gravity, York-completed constraints, corrected quartic interactions, standing-wave cosmology, and an assembled plane-symmetric Einstein--Yang--Mills--Higgs--Dirac system with conserved total stress. Its projected matter sectors have explicit resonance selection rules, finite-volume normalization, backreaction-corrected detuning, integrable isolated triads, shared-mode Hamiltonians, and finite oscillator--doublet quantum blocks. A second-order classical kinetic closure, exact expanding-mode reductions, and the nonlinear two-polarization stress map extend this dynamical content within their stated reductions.

Canonical fermion normalization and fixed-parent phase space then distinguish quantum pair conversion from the commuting triad amplitudes.  A projected one-parent Hamiltonian admits a subtraction-defined self-adjoint ultraviolet limit, an explicit spectral measure, threshold tails, and a weak-coupling exponential survival limit.  Common-profile parent channels yield coherent decay and a precisely invariant dark subspace.  At finite density, exact exclusion-leakage and weighted-correlation identities obstruct naive bosonization and entrywise decorrelation.  An exclusion-preserving collective model instead has a global ultraviolet mild flow and a conserved dressed energy.  Restoring all retained momentum pairings gives a matrix-Pauli covariance flow and a fixed-time many-copy transport theorem for empirical polynomial moments.  These are constructions in explicitly specified models and limit orders, not a finite-species quantum kinetic limit or a realization of the complete Standard Model--Einstein continuum.  Differentiated proper-time trace-class transport and a quantitative three-defect estimate sharpen the determinant-to-stress and constraint-to-Einstein interfaces of the conditional chain.

\chapter*{Preface and logical conventions}
\addcontentsline{toc}{chapter}{Preface and logical conventions}
The exposition separates finite reconstruction, analytic transport, and
physical continuum identification.  Its organizing principle is the common
regulated theory: phase choices, source derivatives, probability laws,
refinement maps, and the physical clock must remain compatible throughout
the construction.  Results are placed at the point where their hypotheses
enter this dependency structure.

Four levels of assertion are distinguished.  A \emph{finite result} is a
statement at a fixed regulator, with its algebraic, spectral, geometric, or
probabilistic assumptions specified.  A \emph{uniform transfer theorem}
states what follows when the required estimates hold independently of the
cutoff.  A \emph{conditional continuum theorem} composes such estimates for
one common family.  An \emph{explicit model result} concerns a declared
truncation, perturbative order, or finite-mode system.  A thermodynamic
limit at fixed spacing is also distinguished from an ultraviolet limit;
existence in the first sense need not establish existence in the second.
Within the explicit models, ultraviolet removal, weak-coupling time
rescaling, collective population growth, parameter quadrature, and a
many-copy limit are kept separate.  A theorem in one of these orders
does not authorize an exchange of limits or a change in the retained
field algebra.

The central identities are accompanied by derivations, and the principal
transfer statements retain their essential hypotheses.  Spatial and scale
trivialization of a determinant line, its reflection-compatible structure,
and positivity of a fermionic boundary form are separate conditions.
Likewise, a supported Dirichlet response is not an unrestricted inverse
estimate.  Finite symbolic or sampled checks are not used as substitutes
for the uniform bounds required by the continuum theorem.

The strongest conclusion available is not a single unconditional existence claim. It is a much more informative separation of what has been constructed, what can be transported, what fails for specific candidate routes, and which quantitative estimates would complete a chosen regular branch. In particular, conservation of a source is not the Einstein equation; a compatible cylinder law is not automatically a local quantum field theory; and a positive transfer operator is not automatically the physical Hamiltonian generated by renewal time.

\begingroup
\setlength{\parskip}{0pt}
\linespread{0.97}\selectfont
\tableofcontents
\endgroup
\mainmatter
\part{Finite reconstruction and the common regulated theory}


% ==================== 01_architecture.tex ====================
\chapter{The endpoint and its conditional architecture}
\label{ch:architecture}
\section{What has to emerge together}
The desired endpoint is a limiting physical theory with the Standard Model gauge, Higgs, and chiral-fermion content on a nondegenerate gravitational branch. Its ingredients must come from one compatible regulator system. Constructing a bosonic probability measure on one family, a chiral determinant on a second family, and a gravitational evolution on a third does not construct their joint theory. The comparison maps, source derivatives, null spaces, reflection, and time calibration must agree on overlaps and under refinement.

Let $n$ index a cofinal cutoff sequence, with microscopic length $a_n\downarrow0$. The regulated data will be denoted
\begin{equation}
 \mathfrak C_n=(X_n,\AA_n,S_n,\mu_n,\Theta_n,Q_n,\mathcal R_{n+1,n},\mathcal J_n).
 \label{eq:common-card}
\end{equation}
Here $X_n$ is the geometric carrier; $\AA_n$ is the graded observable algebra; $S_n$ is the complete action including measure and counterterm contributions; $\mu_n$ denotes its positive physical bosonic body when that body exists; $\Theta_n$ is the declared physical reflection; $Q_n$ is the BRST differential when a BRST presentation is used; $\mathcal R_{n+1,n}$ is an actual refinement comparison; and $\mathcal J_n$ is the joint source map, including metric sources. Fermionic source coefficients are graded multilinear forms, not ordinary probability densities on anticommuting variables.

\begin{definition}[Regular-branch endpoint]
\label{def:endpoint}
A regular-branch Standard Model--Einstein continuum consists of a route-independent physical Schwinger functional on a specified nuclear test algebra, a globally compatible chiral determinant/Pfaffian line, a nondegenerate limiting metric class, a reconstructed physical Hilbert space with declared field domains and time evolution, and a compiled limiting stress $T$ such that
\begin{equation}
 \nabla^\mu T_{\mu\nu}=0,
 \qquad G_{\mu\nu}(g)+\Lambda g_{\mu\nu}=\kappa T_{\mu\nu}
 \label{eq:target}
\end{equation}
in the specified distribution topology. At least one positively separated connected response, after common contact subtraction and removal of the declared dynamically dependent source directions, must differ from its complete Gaussian comparator by a nonzero limit.
\end{definition}
This is the regular-branch target used by the continuum closure theorem. In its effective-action implementation, $T$ is the compiled mean source and $g$ is a limiting critical metric. Equation~\eqref{eq:target} is not silently upgraded to an operator identity for a quantum metric, an almost-sure equation for every metric sample, or a nonperturbative construction of all gravitational sectors. \src{R18}{def:v18v51-hypotheses,thm:v18v51-conditional-closure}

\section{The dependency chain}
The construction is usefully divided into two chains meeting at the common action. The finite chain is
\begin{equation}
\begin{aligned}
 &\text{operational histories and complete prediction}\\
 &\quad\longrightarrow\text{persistent carrier and positive word algebra}\\
 &\quad\longrightarrow\text{geometric and internal loadings}\\
 &\quad\longrightarrow\text{one regulated Einstein--matter source system}.
\end{aligned}
\label{eq:finite-chain}
\end{equation}
The analytic chain is
\begin{equation}
\begin{aligned}
 &\text{physical positivity, phase consistency, and uniform estimates}\\
 &\quad\longrightarrow\text{joint compactness and summable transport}\\
 &\quad\longrightarrow\text{physical reconstruction and stress descent}\\
 &\quad\longrightarrow\text{Einstein variational landing and interaction survival}.
\end{aligned}
\label{eq:analytic-chain}
\end{equation}
An arrow is a theorem only with the inputs stated in the relevant chapter. The two chains do not bypass one another: the finite carrier identifies what must be transported, while the analytic estimates determine whether that transport survives cutoff removal.

\begin{theorem}[Composition principle, preliminary form]
\label{thm:composition-preliminary}
Suppose one literal cofinal family has a common action and chiral line, positive physical reflected forms, joint tightness and moment uniform integrability, vanishing inserted Ward/anomaly defects, a convergent metric variational problem with vanishing full critical residual, a nonzero separated non-Gaussian response excess, summable direct and adjacent transport errors, and the regularity, cluster, domain, and symmetry assumptions of physical reconstruction. Then the family has the regular-branch endpoint of Definition~\ref{def:endpoint}.
\end{theorem}
\begin{proof}
Joint compactness gives subsequential laws and graded Schwinger forms. Uniform integrability passes the required products and connected responses to the limit. Reflection positivity is a closed family of finite Gram inequalities. Vanishing inserted Ward defects descends the physical quotient and conserves the compiled stress. The metric variational convergence identifies a critical limiting metric; convergence of both gravitational and matter variations, with vanishing full residual, gives the second equation in~\eqref{eq:target}. Summable comparison tails make all subsequential and admissible cofinal limits coincide. The separated excess excludes the specified Gaussian comparator. Physical reconstruction is then applied with its own domain and cluster assumptions. Chapter~\ref{ch:closure} states the complete hypothesis list.
\end{proof}
\src{R18}{thm:v18v51-conditional-closure}

\section{Three distinctions that cannot be suppressed}
\subsection{Ward conservation and Einstein dynamics}
The Ward identity is the statement that an assembled source annihilates infinitesimal gauge directions. It constrains the divergence of the metric variation. It does not set the transverse metric variation to zero. A divergence-free tensor can be nonzero. Consequently, even perfect anomaly cancellation and exact stress conservation leave a distinct Einstein residual to control. A metric counterterm changes that residual only within the allowed finite-dimensional space of local variations.

\subsection{Kinematic laws and physical dynamics}
A finite positive state can be made stationary and reflection positive by adjoining a reversible refresh mechanism. An exact projective kernel can preserve that state at every refinement. Neither construction identifies the reconstructed generator with the actual local physical clock. The zero modes, detail modes, and rate scaling must be tested independently. This is especially important when a convenient construction introduces a positive gap by hand.

\subsection{Massive channels and massless sectors}
Exponential decay in a reconstructed massive channel may imply a positive spectral threshold in that channel. A theory containing a photon, or a gravitational sector intended to retain massless excitations, cannot be certified by a positive global mass gap. Channelwise massive estimates and polynomial or infrared-controlled massless estimates must coexist within the same physical state. A gapped scalar comparison theorem is therefore not the final theorem for the complete Standard Model--Einstein system.

\section{What is already strong}
The finite structural and action results identify a concrete theory rather than a placeholder field list. The analytic programme contains exact shell, determinant, response, geometric, and comparison identities. Several results are genuine obstructions to particular routes, not merely missing estimates. The regular-branch theorem gives a sufficient endpoint specification. The explicit coupled models demonstrate that gauge, spin, Yukawa, and gravitational effects can be computed in one consistent source convention. The fixed-spacing Gibbs theorem additionally proves nonemptiness of a positive joint lattice branch. Full-face moments remove a specific coarse rotational defect, and faithful detail dynamics provides a quantified alternative to an annihilating clock. The explicit resonant and two-polarization reductions resolve selected conversion, statistical-response, and geometric-source mechanisms in their declared normalizations.  The canonical continuation constructs an actual one-parent ultraviolet Hamiltonian and a spectral decay limit; an exclusion-preserving nonlinear model supplies global mild evolution and dressed energy; and the fixed-mode many-copy limit supplies a controlled matrix-covariance closure.  The differentiated heat theorem and constraint-angle estimate identify quantitative sufficient conditions at two additional analytic interfaces.  These achievements are logically compatible with the fact that no single physical cofinal family has yet been shown to satisfy the entire continuum hypothesis set.


% ==================== 02_finite.tex ====================
\chapter{A brief reconstruction of the Grand Tensor carrier}
\label{ch:finite}
\section{Prediction, persistence, and positive words}
Start with a fixed typed grammar of potential events and the events that actually occur. Histories are identified only when they induce the same complete future statistics in every retained context and intervention. This history--future equivalence reconstructs the minimum persistent predictive carrier. It does not permit a preferred target geometry to be inserted into the equivalence relation.

\begin{definition}[Predictive equivalence]
Two retained histories $h,h'$ are predictively equivalent when every admissible future test $F$, including its context and intervention labels, satisfies
\begin{equation}
 \mathbb P(F\mid h)=\mathbb P(F\mid h').
 \label{eq:predictive-equivalence}
\end{equation}
The persistent state is the corresponding quotient, with all actually read orientation, scheduler, calibration, and boundary marks retained whenever they distinguish future tests.
\end{definition}

For a word algebra generated by the retained operations, let $\omega$ be the compiled functional. Its positive hierarchy is
\begin{equation}
 G^{(m)}_{ij}=\omega(w_i^*w_j)\succeq0
 \qquad\text{for every finite word panel.}
 \label{eq:word-gram}
\end{equation}
The null space $\mathcal N=\{A:\omega(A^*A)=0\}$ is quotiented before Hilbert completion. Finite-rank history--future/Hankel data give a support-minimal realization when the corresponding finite realization hypotheses hold. In the completely positive presentation, the positive causal-comb structure fixes a canonical support-minimal class rather than a unique arbitrary matrix coordinate system.

\begin{proposition}[All-order persistence is stronger than a covariance shadow]
The degree-one Gram matrix is a finite shadow of the predictive word functional. Agreement of degree-one data, or of any fixed finite correlation order without additional closure hypotheses, does not identify the full interacting state. Complete word positivity and the consistent null quotient are the objects transported toward reconstruction.
\end{proposition}
\begin{proof}
Equation~\eqref{eq:word-gram} constrains every finite linear combination of words, including words of arbitrarily high degree. Restricting the set of words restricts the tested quadratic forms and leaves the untested moments undetermined. A representation of the restricted form therefore need not extend uniquely to a representation of the full functional.
\end{proof}
\cite{GT}

\section{The local spatial branch}
The local $K_4$ selection is relative to a specified branch of operational hypotheses. It is not a theorem that every predictive system has four neighbours. The relevant branch requires a homogeneous source--target reversal pair, a nondegenerate alternating structure on the coefficient space, faithful endpoint and signed-edge Grams, a nonzero cycle, and minimal future-readable innovation.

\begin{theorem}[Conditional $K_4/A_3$ selection]
\label{thm:k4}
On the simplex branch, let the endpoint-difference representation be $W_N$, of dimension $N-1$, and suppose $\R t\oplus W_N$ carries the declared nondegenerate alternating structure. With the nonzero-cycle and minimality assumptions above, the minimal admissible value is $N=4$. Its spatial endpoint representation is $W_4\simeq\R^3$, and
\begin{equation}
 \Hom_{S_N}\!\left(\Lambda^2W_N,W_N\otimes\sgn\right)
 \simeq
 \begin{cases}\R,&N=4,\\0,&N\ne4,\end{cases}
 \label{eq:k4-cross}
\end{equation}
on the stated simplex family.
\end{theorem}
\begin{proof}[Proof mechanism]
Nondegeneracy of an alternating form requires even total dimension $N$. The nonzero-cycle requirement excludes $N=2$. Faithful endpoint and edge ranks are $N-1$ and $N(N-1)/2$, respectively, so the minimal remaining candidate is $N=4$. The equivariant exterior-square calculation identifies the unique signed cross-product channel on that branch.
\end{proof}
\src{GT}{thm:dimension-K4-selector}

\begin{corollary}[Cofinal $3+1$ kinematics]
\label{cor:3plus1}
If mixed-Gram corrections are summable, the threshold ranks $3$ and $6$ remain separated, the pair margin stays positive, and source-loss and extra-writer tails vanish, the local $K_4$ branch survives cofinally. An independently reconstructed, calibrated physical clock then gives a $1+3$ kinematic interpretation.
\end{corollary}
\src{GT}{cor:dimension-cofinal-3plus1}
The four vertices describe a local relational carrier, not four points constituting the whole universe. Likewise, the clock is an additional reconstructed structure; a three-dimensional Gram does not create time by declaration.

\section{Metric and field data are loadings, not invented labels}
The conservative readable completion retains the actual root, orientation, speed, scheduler, source-response, and field records needed by the geometric compiler. Metric, coframe, connection, curvature, and ADM data are deterministic outputs of those records. Their identification with the target geometric objects requires the declared response identities and vanishing residuals. A failed source direction may not be discarded solely because it spoils the desired interpretation.

A useful explicit finite family has
\begin{equation}
 \Sigma_n=\Lambda_{A_3}/d_n\Lambda_{A_3},\qquad a_n=d_n^{-1},
 \qquad k_\alpha=\frac1{8a_n^2}
 \label{eq:a3-rates}
\end{equation}
for the twelve oriented roots. The reference calibration gives $B=I_3$, lapse $N=1$, and shift $\beta=0$. The admitted rate and speed perturbations span the full ten-dimensional ADM tangent at the reference point. This establishes an explicit finite geometric branch and a source response, not convergence of its fluctuating metric law. \cite{GT}

\section{The internal Standard Model branch}
The finite Standard Model result assumes the specified geometric short, faithful private plane, operation-native weak reset, positive colour bridge, and typed determinant--Clifford word structure. These hypotheses do real work: they select the internal algebra and the physical representation within the operational carrier.

\begin{theorem}[Finite structural Standard Model]
\label{thm:finite-sm}
Under those internal loading hypotheses, the active relative commutant is
\begin{equation}
 M_3(\C)\oplus M_2(\C)\oplus\C\oplus\C,
\end{equation}
and the faithful gauge group is
\begin{equation}
 \GSM=S(U(3)\times U(2))
 \simeq\frac{SU(3)\times SU(2)\times U(1)}{\Z_6}.
 \label{eq:gsm}
\end{equation}
The physical multiplets $Q,u,d,L,e,\nu,H$ have integral hypercharge labels
\begin{equation}
 6Y=(1,4,-2,-3,-6,0,3).
 \label{eq:charges}
\end{equation}
The source has three future-visible structural endpoint generations, one weak Higgs doublet, and one grading-odd self-adjoint finite Dirac writer with residue matrices $Y_u,Y_d,Y_e,Y_\nu,M_R$ expressed in the same frame.
\end{theorem}
\src{GT}{thm:SM-active-SM-I,thm:SM-GT}

The group quotient is concrete. The map
\begin{equation}
 (g_3,g_2,z)\longmapsto(z^{-2}g_3,z^3g_2)
 \label{eq:gsm-map}
\end{equation}
has kernel $(z^2I_3,z^{-3}I_2,z)$ with $z^6=1$. The quotient is therefore the faithful group, not an arbitrary product presentation with an ignored central kernel. \src{GT}{thm:SM-group}

Writing $C$ and $W$ for the colour and weak seeds, the generation is compiled from
\begin{equation}
 Q=C\otimes W,\quad u=\Lambda^2C^*,\quad d=C,\quad
 L=W^*,\quad e=\Lambda^2W^*,\quad H=W,
 \label{eq:multiplet-writers}
\end{equation}
with the determinant trivialization fixing the hypercharge assignments. The orthogonal cycle triplet is not counted as an extra set of generations. \src{GT}{thm:SM-generation}

\scope{The result identifies representation content and the common finite writer. It does not calculate the numerical gauge couplings, Yukawa eigenvalues, mixing angles, Majorana scale, Higgs expectation, or Newton constant. Three structural generations are not by themselves a dynamical explanation of their mass hierarchy.}

% ==================== 03_action.tex ====================
\chapter{One action, exact identities, and the source interface}
\label{ch:action}
\section{The finite gauge--Higgs--fermion action}
Let $U$ be link variables in $\GSM$, let $H$ be the weak doublet, and let $\Psi$ carry spatial spin, generation, and the finite internal packet. The dual Grassmann field is retained independently. With the discrete Hodge stars and weights supplied by the geometric loading, the regulated Euclidean matter action is
\begin{equation}
\begin{aligned}
 S_{\mathrm{SM},n}^{E}={}&
 \frac12\pair{F_U}{*_2F_U}_{g_1,g_2,g_3}
 +\frac12\pair{D_UH}{*_1D_UH}\\
 &+\sum_{v\in X_n}m_v\lambda_H\bigl(\norm{H_v}^2-v_H^2\bigr)^2\\
 &+\operatorname{Re}\pair{\bar\Psi}
 {*_0\bigl(D_{\mathrm{sp}}^U+\Gamma_{\mathrm{Cl}}\otimes D_F^Y(H)\bigr)\Psi}.
 \label{eq:finite-sm-action}
\end{aligned}
\end{equation}
The field content and the action are distinct achievements. The former fixes what transforms; the latter fixes a common local source system and its variations.

\begin{theorem}[Finite action, Ward identity, and Higgs vacuum]
\label{thm:action}
The action~\eqref{eq:finite-sm-action} is gauge invariant and local on the declared vertex--edge--plaquette neighbourhood. Its Euler derivatives satisfy the exact finite identity
\begin{equation}
 d_U^*E_U+R_H^*E_H+R_\Psi^*E_\Psi+R_{\bar\Psi}^*E_{\bar\Psi}=0.
 \label{eq:finite-ward}
\end{equation}
At $U=I$, $H=v_Hh_0$, and $\Psi=0$, the stabilizer is
\begin{equation}
 U(3)\simeq\frac{SU(3)\times U(1)_{\rm em}}{\Z_3}.
\end{equation}
The gauge mass matrix has rank three. In the radial convention of~\eqref{eq:finite-sm-action}, the radial potential Hessian is $8\lambda_Hv_H^2$.
\end{theorem}
\begin{proof}
Gauge covariance of $F_U,D_UH,D_F^Y(H)$ and invariance of the pairings imply invariance of each term. Differentiate this invariance along an arbitrary finite gauge parameter and collect Euler derivatives to obtain~\eqref{eq:finite-ward}. The vacuum stabilizer fixes the chosen weak vector; the broken gauge directions give the rank-three quadratic form. Differentiating $\lambda_H(r^2-v_H^2)^2$ twice at $r=v_H$ gives the stated radial Hessian.
\end{proof}
\src{GT}{thm:SM-active-SM-II}

\begin{proposition}[Classical handoff is not measure convergence]
\label{prop:classical-handoff}
On a smooth-field refinement family with the declared consistency and quadrature estimates, the discrete action tends to the continuum classical Standard Model action. Where weak field convergence is supplemented by convergence of the corresponding energies, the stated handoff upgrades to strong $L^2$ convergence of curvature and covariant Higgs derivative and strong $L^4$ convergence of the Higgs field.
\end{proposition}
\src{GT}{thm:SM-active-SM-II-continuum}
This controls deterministic smooth configurations. It neither identifies typical configurations of a fluctuating measure nor proves tightness, anomaly-free phase gluing, reflection positivity, or interacting continuum Schwinger functions.

\section{BRST, BV, and exact elimination}
At finite cutoff, the declared BRST rules include
\begin{equation}
 s c=-c^2,\qquad sU=c_tU-Uc_s,\qquad sH=cH,\qquad
 s\Psi=c\Psi,\qquad s\bar c=B,\qquad sB=0.
 \label{eq:brst}
\end{equation}
The complete field/antifield projection must preserve the odd symplectic pairing. Restricting only the field Hessian and then declaring the result to be a BV action misses the dual projection and its defect. On an admissible finite BV fibration, pushforward along the specified Lagrangian fibre preserves the quantum master equation and composes under successive eliminations. This is an exact algebraic statement; convergence of the resulting effective action remains analytic. \cite{R1,R2}

\begin{proposition}[Connected shell expansion]
\label{prop:cumulant-shell}
For a finite shell reference expectation $\E_0$ and interaction $V$, wherever the exponential moments justify differentiation,
\begin{equation}
 -\log\E_0e^{-V}
 =\E_0V-\frac12\Cum_0(V,V)+\frac16\Cum_0(V,V,V)-\cdots.
 \label{eq:shell-cumulants}
\end{equation}
Truncation after second order has a third-cumulant integral remainder evaluated along the actual exponential tilt. If the reference itself moves, its Radon--Nikodym variation is part of $V$ before this expansion is taken.
\end{proposition}
\begin{proof}
Set $f(t)=-\log\E_0e^{-tV}$. Successive derivatives are alternating tilted connected cumulants. Taylor's formula with integral remainder gives the identity and the remainder statement.
\end{proof}
\cite{R1,R4,R16}

\section{The metric variation is a complete variation}
The stress source is obtained from the same action and measure used to define the state. It includes the variation of the kinetic and potential terms, determinant/Pfaffian phase, normalization, moving domain or support, source transport, gauge-fixing and ghost terms where relevant, and local counterterms. A finite field-dependent vacuum shift is not a harmless constant if it depends on the metric.

\begin{proposition}[Normalized response identity]
\label{prop:normalized-response}
For $\mu_t=Z_t^{-1}e^{-S_t}\lambda$ and a differentiable observable $A_t$,
\begin{equation}
 \frac{\dd}{\dd t}\E_tA_t
 =\E_t\dot A_t-\Cov_t(A_t,\dot S_t).
 \label{eq:normalized-response}
\end{equation}
If $\lambda_t$ varies and is absolutely continuous in a common reference, replace $\dot S_t$ by the derivative of the full action including $-\log(\dd\lambda_t/\dd\lambda)$. For singular deterministic lifts, no fictitious finite Radon--Nikodym derivative is introduced; a coupling comparison is required.
\end{proposition}
\begin{proof}
Differentiate numerator and denominator. The denominator derivative removes exactly the mean of $\dot S_t$, producing the covariance. The same calculation applies after incorporating an actual reference density into the action.
\end{proof}
\cite{R2,R16}

\subsection{Nonminimal Higgs coupling}
For $f=H^\dagger H$, the nonminimal term $\xi fR$ has the complete bulk stress contribution
\begin{equation}
 T_\xi^{\mu\nu}=-2\xi\left[fG^{\mu\nu}
 +(g^{\mu\nu}\Box-\nabla^\mu\nabla^\nu)f\right],
 \qquad \nabla_\mu T_\xi^{\mu}{}_{\nu}=\xi R\nabla_\nu f.
 \label{eq:nonminimal}
\end{equation}
Its boundary variation is paired with the bulk term. In the source convention the geometric--Higgs boundary momentum is
\begin{equation}
 P_{gH}^{ab}=(\chi_R+\xi f)\Pi^{ab}-\xi\gamma^{ab}\nabla_nf.
 \label{eq:nonminimal-boundary}
\end{equation}
The normal derivative cannot be dropped while retaining the bulk improvement term. \src{R1}{thm:V150-Ixi-variation}

If $F=\chi_R+\xi f>0$, the Einstein-frame change $\widetilde g=(F/\chi_*)g$ yields the scalar target metric
\begin{equation}
 \widetilde G_{AB}=\frac{\chi_*}{F}G^0_{AB}
 +\frac{3\chi_*}{F^2}\partial_AF\,\partial_BF.
 \label{eq:einstein-frame}
\end{equation}
This is a conditional positive branch. Covariance does not fix $\xi$, does not ensure $F>0$ on a fluctuating field support, and does not justify crossing $F=0$ with the same frame.

\section{Boundary and domain data}
The boundary source system contains oriented port terms and the corresponding Brown--York/Codazzi identities. Opposite caps may cancel after the same orientation and domain conventions are used. An arbitrary boundary projector, however, need not define a Hadamard state or a causal stress response. The boundary Calder\'on or Poisson realization, spectral sections, and common-domain estimates are separate inputs.

Moving support projections provide a particularly simple warning. If $P=P(s,t)$ is a differentiable projector, its Kato transport has curvature governed by $P[\dd P,\dd P]P$. A scalar likelihood or action can have flat parameter transport while this operator support bundle has nonzero curvature. Scalar normalization therefore does not replace the geometric transport of zero modes, domains, or determinant-line frames. \src{R1}{thm:V119-Kato-transport,prop:V119-GM-versus-Kato}

\gate{A continuum source theorem needs uniform control of this complete variation, not only convergence of the unvaried action. The same condition applies to boundary, ghost, phase, and measure terms.}


% ==================== 04_chiral.tex ====================
\chapter{Chirality, anomalies, and determinant lines}
\label{ch:chiral}
\section{Local anomaly cancellation and its precise reach}
For anomaly arithmetic it is convenient to express all fermions as left-handed Weyl fields. One generation then has hypercharges
\begin{equation}
 Q:\tfrac16,\quad u^c:-\tfrac23,\quad d^c:\tfrac13,
 \quad L:-\tfrac12,\quad e^c:1,\quad\nu^c:0.
\end{equation}
The mixed nonabelian and gravitational traces and the cubic abelian trace cancel:
\begin{align}
 2\left(\tfrac16\right)-\tfrac23+\tfrac13&=0,
 &3\left(\tfrac16\right)-\tfrac12&=0,\\
 6\left(\tfrac16\right)+3\left(-\tfrac23\right)+3\left(\tfrac13\right)
 +2\left(-\tfrac12\right)+1&=0,\\
 6\left(\tfrac16\right)^3+3\left(-\tfrac23\right)^3
 +3\left(\tfrac13\right)^3+2\left(-\tfrac12\right)^3+1&=0.
 \label{eq:anomaly-traces}
\end{align}
There are four weak doublets per generation, including colour multiplicity. Yukawa couplings do not change this representation-theoretic local anomaly calculation. \src{R1}{cor:V138-local-cancellation,cor:V138-Yukawa-independence}

\begin{proposition}[Anomaly separation]
\label{prop:anomaly-separation}
Cancellation of~\eqref{eq:anomaly-traces} removes the declared local perturbative anomaly polynomial. It does not by itself control a regulator remainder, fix a determinant-line trivialization, remove torsion or a global holonomy character on a different geometric category, orient a Pfaffian on all components, or prove reflection positivity of the complete fermionic boundary form.
\end{proposition}
\src{R18}{led:v18v1-anomaly}

\section{The faithful smooth branch}
The smooth faithful branch uses $\GSM=S(U(3)\times U(2))$ on ordinary spin geometry, with the specified elliptic realization, collars, and determinant gluing. The source identifies the relevant closed anomaly character as trivial on this branch. The global statement is stronger than local trace cancellation because it concerns the determinant-line holonomy on the declared bordism domain. It is also narrower than a claim about arbitrary non-spin backgrounds, alternative global quotients, rough regulators, or unconstructed chiral measures. \src{R1}{thm:V141-closed-anomaly-zero}

\begin{theorem}[Smooth faithful phase transport]
\label{thm:smooth-phase}
On the smooth faithful family with compact base, the declared elliptic product-collar realization, paired boundary data, determinant gluing, and vanishing boundary $K^1$ obstruction, spectral sections exist. The combined anomaly-free determinant line admits a unit parallel section, unique up to one constant phase. The resulting geometric transport is path independent and has trivial phase cocycle on the specified family.
\end{theorem}
\src{GT}{thm:SMST-smooth-faithful-phase}

The qualification ``specified family'' is indispensable. To descend a line through a gauge quotient, isotropy must act trivially on its fibre. To glue local trivializations, their transition phases must form the required coboundary. Flatness then removes curvature but not a potentially nontrivial flat holonomy unless the global character has also been controlled. Across cutoffs, a further cocycle condition identifies the phase anchors on successive scales. These are distinct checks.

\subsection{Spatial, scale, and reflection compatibility}
The phase problem has three separate directions: motion through a gauge or
geometric parameter space, passage between cutoffs, and reflection.  A local
choice of determinant coefficient does not identify these directions.  Let
$D_i=g_{ij}D_j$ be local coefficients of the complete determinant line on an
open cover.  A global scalarization of the line exists precisely when
$g_{ij}=h_i h_j^{-1}$ for nowhere-zero functions $h_i$, in which case
$d=h_i^{-1}D_i$ is independent of the chart.  Zeros of the determinant section
are permitted; it is the frame, not the section, that must be nonvanishing.
\src{R10}{thm:v97-scalarization}

\begin{theorem}[Coherent determinant scalarization]
\label{thm:coherent-line}
Let $p_{n+1,n}:B_{n+1}\to B_n$ be refinement maps and let
\[
 \Phi_{n+1,n}:L_{n+1}\longrightarrow p_{n+1,n}^{*}L_n
\]
be isomorphisms of the complete determinant lines, including all eliminated
factors.  Assume their compositions agree with the direct line maps.  Given
a scalarizing frame $\tau_0:L_0\to\C$, there is a unique compatible family
\begin{equation}
 \tau_{n+1}=(p_{n+1,n}^{*}\tau_n)\circ\Phi_{n+1,n}.
 \label{eq:coherent-line-frame}
\end{equation}
A section family satisfying
$\Phi_{n+1,n}D_{n+1}=p_{n+1,n}^{*}D_n$ consequently has scalar coefficients
$d_{n+1}=d_n\circ p_{n+1,n}$.  A nonzero parallel initial frame exists on a
connected flat branch if and only if its loop holonomy is trivial.

Suppose, in addition, that $L_n$ carries an antilinear reflection lift
$\mathcal J_n$ and that the scale maps intertwine these lifts.  If
$\tau_0(\mathcal J_0\ell)=\overline{\tau_0(\ell)}$, then every frame in
\eqref{eq:coherent-line-frame} has the same reflection compatibility.
\end{theorem}
\begin{proof}
Composition with an isomorphism preserves a frame.  The scale cocycle makes
the recursive construction independent of the grouping of successive
refinements, and application to the section identity gives scalar
consistency.  Parallel transport from a fixed nonzero fibre vector is
path independent exactly when every loop has identity holonomy.  Finally,
commuting $\mathcal J_n$ through the recursion proves the reflection
identity by induction.
\end{proof}
\src{R10}{thm:v97-flat-holonomy,thm:v97-projective-scalarization,prop:v97-reflection-frame}

Complex scalarization alone does not supply the initial reflection-compatible
frame: the corresponding Real line cocycle must also be trivialized.
Similarly, a gauge-quotient frame presupposes trivial isotropy action.  None
of these line conditions is a substitute for positivity of the reflected
fermionic coefficient form in Chapter~\ref{ch:reconstruction}.


\section{A chiral regulator with an admissible gap}
For an overlap construction, the Wilson kernel and its sign must be uniformly controlled before one differentiates the chiral projectors. The one-species window fixes the Wilson--Bott class. On the quantitative admissible branch recorded in the source, the Wilson parameter choice gives
\begin{equation}
 H_W^2\succeq\frac{1-30\delta}{a^2}I,
 \qquad \norm{H_W}\le\frac7a,
 \qquad 0\le\delta<\frac1{30}.
 \label{eq:wilson-gap}
\end{equation}
The resulting sign operator and overlap projectors are exponentially local with the declared cutoff and volume constants. The Ginsparg--Wilson relation is
\begin{equation}
 \gamma_5D+D\gamma_5=aD\gamma_5D.
 \label{eq:gw}
\end{equation}
\cite{R3,R6}

\begin{theorem}[Index identification on the smooth curved branch]
\label{thm:curved-index}
Under the smooth bundle, ellipticity, locality, and family-approximation assumptions of the Wilson--Bott construction, the regulator carries the curved twisted index
\begin{equation}
 \operatorname{index}D_R^+
 =\int_X\bigl[\widehat A(TX)\operatorname{ch}(E_R)\bigr]_4.
 \label{eq:curved-index}
\end{equation}
The family version retains the corresponding $K$-theory class. A cofinal approximation must additionally control the required parameter derivatives and its range schedule.
\end{theorem}
\src{R3}{thm:v3v92-curved-index,thm:v3v92-global-Thom}
A fixed admissibility gap is not a uniform smoothness theorem for typical quantum fields. Smooth bundle reconstructions from rough links can have derivatives that deteriorate with cutoff even when the curvature energy is controlled.

\section{Physical zero modes must be retained before inversion}
The Wilson-kernel gap and the physical Dirac spectrum play different roles. The first permits the sign construction. The second may contain physical zero modes and crossings. When a mode crosses a spectral partition, it must be moved into the retained finite-dimensional Grassmann sector \emph{before} an inverse is formed. Dividing by a determinant that is then zero is not a legal hard-mode elimination.

\begin{proposition}[Hard/soft Schur factorization]
\label{prop:hard-soft}
For a finite block operator with invertible hard block $D_{hh}$,
\begin{equation}
 \det D=\det D_{hh}\,
 \det\!\left(D_{ss}-D_{sh}D_{hh}^{-1}D_{hs}\right).
 \label{eq:hard-soft}
\end{equation}
The analogous Berezin pushforward retains the full Grassmann polynomial in the soft variables. The identity remains suitable at a physical crossing only when all crossing modes have already been assigned to the soft block.
\end{proposition}
\begin{proof}
Block Gaussian elimination gives triangular factors with unit determinant and the Schur complement. Berezin integration performs the same algebra while retaining soft sources. Neither derivation requires the soft determinant to be nonzero.
\end{proof}
\src{R4}{thm:v4v73-Berezin-Schur}

For spectral sections, the distinction is between a uniformly invertible complement and a finite crossing packet. Local eta expansions and Berry currents must be joined with the crossing transitions; neither an absolute phase bound on one chart nor a local anomaly trace controls the whole line.

\section{Massive Yukawa blocks and their exclusions}
\begin{proposition}[Massive overlap floor]
\label{prop:massive-overlap}
For
\begin{equation}
 D_m=\left(1-\frac{am}{2}\right)D+mI,
 \qquad0<m\le\frac2a,
 \label{eq:massive-overlap}
\end{equation}
on the specified overlap branch, the smallest singular value is at least $m$ and the vectorlike determinant is positive. With $A_D=I-aD/2$, the kinetic and mass terms assemble as $D+mA_D$.
\end{proposition}
\cite{R6}
The identities $D\widehat P_\pm=P_\mp D$ and $A_D\widehat P_\pm=P_\pm A_D$ are what make the chiral Yukawa assembly consistent. Singular values of the Yukawa matrices give massive paired blocks on the nonzero-Higgs chart. This does not produce a uniform mass for an unmatched neutrino, a zero Yukawa singular value, or a configuration where the Higgs leaves that chart. Antiperiodic finite-volume spin conditions are not a volume-uniform physical mass floor.

\section{Relative determinants and rough-field control}
In four dimensions, an order-$-1$ relative Dirac perturbation is naturally at the weak Schatten-$4$ threshold. A fifth modified determinant isolates the trace-class remainder while keeping the first four Born contributions explicit:
\begin{equation}
 \det_{\rm ren}D
 =\exp\bigl(B_{\le4}^{\rm ren}+T_{\rm loc}\bigr)\det_5(I+K).
 \label{eq:det5}
\end{equation}
Metric principal-symbol changes require an order-zero elliptic microlocal normalizer before the residual becomes order $-1$. The normalizer's own ultraviolet local gravitational variation is retained; it is not discarded as a coordinate change. \src{R4}{thm:v4v91-decomp}

\begin{proposition}[What remains after local subtraction]
The renormalized Born packet contains finite nonlocal polarization, including the logarithmic momentum dependence at the marginal order. Curvature $L^2$ control alone does not bound the associated logarithmic moment. A positive bound on a separated absolute fermionic factor need not exist even when the correctly signed combined Gaussian inverse covariance and polarization is controlled.
\end{proposition}
\cite{R3,R4}
For three-dimensional boundary operators, the corresponding order-$-1$ threshold is $p>3$, with a fourth modified determinant and local multiplicative anomaly. Thus bulk and boundary subtraction orders are not interchangeable.

\section{Proper-time subtraction and differentiated trace-class transport}
\label{sec:proper-time-transfer}
The relative determinant and the stress obtained by differentiating it
require different levels of convergence.  This distinction can be stated
on a common Hilbert carrier, without identifying a scalar heat-trace
cancellation with convergence of an operator-valued logarithm.
Let $B_n(u),C_n(u)\succeq m^2 I$, $m>0$, be the positive squared hard
operators on a compact parameter interval $I$.  The parameter $u$ may
represent a metric or source variation.  Write
\[
 \mathcal H_n(s,u)=e^{-sC_n(u)}-e^{-sB_n(u)}.
\]
Whenever the relative trace is defined, its four-dimensional local head
has the form $c_{0,n}s^{-2}+c_{2,n}s^{-1}+c_{4,n}$.  Its integral over
$[\Lambda^{-2},1]$ against $\dd s/s$ is exactly
\begin{equation}
 \frac{c_{0,n}}2(\Lambda^4-1)
 +c_{2,n}(\Lambda^2-1)+2c_{4,n}\log\Lambda.
 \label{eq:proper-time-head}
\end{equation}
These local coefficients are retained in the renormalized action with
their source dependence.  Boundary heat coefficients, when present,
require their own subtraction and are not covered by the bulk head.

\begin{theorem}[Differentiated proper-time transport]
\label{thm:proper-time-c1}
Let $\mathcal P_n(s,u)$ be declared operator-valued subtractions and
suppose the remainders
$\mathcal A_n(s,u)=\mathcal H_n(s,u)-\mathcal P_n(s,u)$ belong to the
trace class $\mathfrak S_1$.  Suppose that, for almost every $s>0$,
$\mathcal A_n(s,\cdot)$ is continuously differentiable in trace norm,
and that for $j=0,1$
\begin{equation}
 \sup_{u\in I}\frac{\norm{\partial_u^j\mathcal A_n(s,u)}_1}{s}
 \le g_j(s),\qquad g_j\in L^1((0,\infty)),
 \label{eq:heat-derivative-majorant}
\end{equation}
uniformly in $n$.  Assume trace-norm convergence of these remainders
and their derivatives, uniformly in $u$, for almost every $s$.
Then the Bochner integrals
\begin{equation}
 \mathcal K_n(u)=\int_0^\infty\mathcal A_n(s,u)\frac{\dd s}{s}
 \label{eq:proper-time-operator}
\end{equation}
converge in $C^1(I;\mathfrak S_1)$, with
$\partial_u\mathcal K(u)=\int_0^\infty\partial_u\mathcal A(s,u)\dd s/s$.
Consequently the renormalized relative logarithm and its first source
variation converge after adding the separately convergent local
subtraction packet.  With only the $j=0$ hypotheses, the conclusion is
trace-norm convergence without a conclusion about stress derivatives.
\end{theorem}
\begin{proof}
The majorants give existence of the Bochner integrals.  For either
$j$, the supremum over $u$ of the trace norm of the difference of the
integrals is bounded by the integral of the corresponding supremum.
Its integrand tends to zero and is bounded by $2g_j$.  Dominated
convergence therefore gives uniform convergence of the integrals and
of their derivatives.  The fundamental theorem of calculus in
$\mathfrak S_1$ identifies the limiting derivative.  Finally
$|\Tr K|\le\norm K_1$ transfers both conclusions to the scalar
renormalized action.  The sign of the relative logarithm is fixed by
$\log B-\log C=\int_0^\infty(e^{-sC}-e^{-sB})\dd s/s$ whenever this
identity is defined before subtraction.
\end{proof}
\src{R15}{thm:v15v37-operator-heat,thm:v15v38-c1-heat}

A sufficient short-time hypothesis is
$\norm{\partial_u^j\mathcal A_n(s,u)}_1\le Cs^\delta$ with
$\delta>0$; an exponentially decreasing large-time bound then completes
\eqref{eq:heat-derivative-majorant}.  Neither hypothesis follows merely
from a scalar asymptotic expansion.  For example,
$B=\diag(1,2)$ and $C=\diag(2,1)$ have identical heat traces, although
$\norm{\log B-\log C}_1=2\log2$.  Likewise uniform convergence of
$\sin(nu)/n$ does not imply convergence of its derivatives.
\src{R15}{thm:v15v37-subtraction-order,cth:v15v37-trace-cancellation}

\scope{This theorem supplies a precise determinant-to-stress interface
on a uniformly massive hard sector.  It does not construct the
operator-valued subtraction on a rough chiral family, remove physical
zero modes, trivialize a determinant phase, or establish reflection
positivity.  Those requirements remain those of the common regulated
line and source system.}


\gate{The remaining chiral task is a single rough, cutoff-compatible, reflection-compatible physical line and source system, with exact zero-mode saturation and uniform derivative bounds. Smooth anomaly cancellation and finite massive charts supply essential pieces but do not supply this entire object.}


% ==================== 05_renormalization.tex ====================
\part{Uniform analytic bridges and their obstructions}
\chapter{Exact elimination, renormalization, and cutoff ancestry}
\label{ch:rg}
\section{Schur reduction with a complete reference measure}
An elimination is useful only if it preserves the theory that was actually defined. For a positive quadratic block
\begin{equation}
 Q=\begin{pmatrix}A&B\\B^*&C\end{pmatrix},\qquad C>0,
\end{equation}
Gaussian integration gives the coarse quadratic form $A-BC^{-1}B^*$ and the determinant factor from $C$. Omitting the latter changes metric and parameter derivatives even if one subsequently renormalizes the probability density. The Grassmann analogue is Proposition~\ref{prop:hard-soft}; the two determinant powers have their respective bosonic and fermionic meanings. \src{R4}{thm:v4v1-Schur-shell}

For a full-rank linear coarse map $P$ and fine covariance $C_f$, the coarse Gaussian covariance is $C_c=PC_fP^*$. The conditional mean lift and conditional covariance are
\begin{equation}
 L=C_fP^*C_c^{-1},\qquad
 C_{\rm det}=C_f-C_fP^*C_c^{-1}PC_f.
 \label{eq:gaussian-lift}
\end{equation}
Sampling $x_c$ with covariance $C_c$ and an independent detail $\eta$ with covariance $C_{\rm det}$ gives $x_f=Lx_c+\eta$ with exactly covariance $C_f$.

\begin{proposition}[No reference defect for the canonical Gaussian lift]
\label{prop:gaussian-lift}
The lift~\eqref{eq:gaussian-lift} reproduces the fine Gaussian reference exactly. After this lift, an interacting cutoff comparison is governed by the Wilsonian action mismatch and the actual observable mismatch, not by an additional invented Gaussian-density error.
\end{proposition}
\begin{proof}
The lifted covariance is $LC_cL^*+C_{\rm det}=C_f$, and its mean is zero. A centered finite Gaussian law is determined by this covariance. The interacting density is then compared on the resulting common probability space.
\end{proof}
\cite{R16}

\section{Differentiated shell remainders in the comparison norm}
An exact Gaussian lift removes the reference mismatch, but it does not make
a formal cumulant expansion a uniform cutoff estimate.  The necessary
quantity is the differentiated remainder in the same localized norm used
for the adjacent comparison.  Write
\begin{equation}
 W_u(x)=-\log\E_\eta e^{-uZ(x,\eta)},\qquad
 X_u=Z-\E_uZ,\qquad D_{j,u}=\nabla_jZ-\E_u\nabla_jZ,
 \label{eq:localized-shell-tilt}
\end{equation}
where $\E_u$ is the normalized tilted detail expectation.  Assume throughout
that the indicated derivatives and moments exist on the retained coarse
field domain.  The normalized differentiation rule is
\begin{equation}
 \nabla_j\E_uF=\E_u\nabla_jF-u\E_u(FD_{j,u}).
 \label{eq:normalized-shell-derivative}
\end{equation}
It retains the derivative of the partition function, which is essential
for the connected expression below.

\begin{theorem}[Localized third-order shell bound]
\label{thm:localized-shell}
For $0\le u\le u_0$, the third strength derivative obeys
\begin{align}
 \partial_u^3W_u&=\E_uX_u^3,\\
 \nabla_j\partial_u^3W_u
 &=3\E_u(X_u^2D_{j,u})
   +3u\E_uX_u^2\,\E_u(X_uD_{j,u})
   -u\E_u(X_u^3D_{j,u}).
 \label{eq:third-shell-gradient}
\end{align}
Suppose the centered mixed moments have uniform bounds $C_{rs,j}$ on
$\E_u(|X_u|^r|D_{j,u}|^s)$ for the indices appearing here; write $C_{20}$
for the derivative-free second moment.  Set
\[
 B_j=3C_{21,j}+3u_0C_{20}C_{11,j}+u_0C_{31,j}.
\]
For the localized seminorm
\[
 \norm{A}_{\mathcal B_{S,n}}
 =\sup_s\sum_j e^{-m_*d_n^{\rm phys}(S,j)}
                \norm{\nabla_jA}_{L^2(\mu_{n,s})},
\]
and $0\le\epsilon\le u_0$, the remainder
\begin{equation}
 R_3=W_\epsilon-\epsilon\E_0Z
                    +\frac{\epsilon^2}{2}\Var_0Z
 \label{eq:third-shell-remainder}
\end{equation}
satisfies
\begin{equation}
 \norm{R_3}_{\mathcal B_{S,n}}
 \le\frac{\epsilon^3}{6}\,M_{3,n}(S),\qquad
 M_{3,n}(S)=\sum_j e^{-m_*d_n^{\rm phys}(S,j)}B_j.
 \label{eq:third-shell-bound}
\end{equation}
\end{theorem}
\begin{proof}
Differentiating the normalized tilt gives
$\partial_uW_u=\E_uZ$ and
$\partial_u^2W_u=-\E_uX_u^2$; the next derivative is the third cumulant.
In \eqref{eq:normalized-shell-derivative}, use
$\nabla_jX_u=D_{j,u}+u\E_u(X_uD_{j,u})$ to obtain
\eqref{eq:third-shell-gradient}.  The mixed-moment bounds control each
term.  Taylor's integral remainder has kernel
$(\epsilon-u)^2/2$; its integral is $\epsilon^3/6$.  Taking the localized
seminorm proves \eqref{eq:third-shell-bound}.
\end{proof}
\src{R17}{lem:v17v1-differentiation,thm:v17v1-gradient,cor:v17v1-banach}

\subsection{Stable polynomial tilts rather than absolute exponential moments}
The hypotheses of Theorem~\ref{thm:localized-shell} can be satisfied by
one-sided stable tilts without imposing an impossible absolute exponential
moment on a quartic Gaussian polynomial.

\begin{proposition}[Degree obstruction and one-sided replacement]
\label{prop:stable-polynomial-tilt}
Let $\eta$ be a finite nondegenerate Gaussian vector.  For a nonconstant
real polynomial $P$, $\E e^{a|P(\eta)|}$ is infinite for every $a>0$ when
$\deg P>2$; for $\deg P\le2$ it is finite for sufficiently small $a>0$.

Instead, let $Z(x,\eta)\ge-K$ on a controlled coarse domain
$\mathcal G$, and suppose the Gaussian raw moments
\[
 H_{rs,j}=\sup_{x\in\mathcal G}
          \E_0\bigl(|Z|^r|\nabla_jZ|^s\bigr)
\]
are finite for $0\le r\le3$, $0\le s\le1$.  For $0\le u\le u_0$ the
normalized density is bounded above by
\begin{equation}
 \frac{e^{-uZ}}{\E_0e^{-uZ}}
 \le Q:=\exp\{u_0(K+H_{10})\}.
 \label{eq:stable-tilt-domination}
\end{equation}
Consequently the centered mixed moments required in
Theorem~\ref{thm:localized-shell} are bounded by finite combinations of
$Q$ and the $H_{rs,j}$.
\end{proposition}
\begin{proof}
A nonzero leading homogeneous part of degree $p>2$ is bounded away from
zero on an open cone.  Along that cone, the growth $ar^p$ dominates the
Gaussian quadratic decay.  Quadratic or linear growth is absorbed by a
sufficiently small change of the Gaussian quadratic form.  For the
one-sided tilt, Jensen's inequality gives
$\E_0e^{-uZ}\ge e^{-u\E_0Z}\ge e^{-uH_{10}}$, while
$e^{-uZ}\le e^{uK}$.  This proves \eqref{eq:stable-tilt-domination}.
Expansion around the tilted means and the same raw-moment bound control
the centered products.
\end{proof}
\src{R17}{thm:v17v2-degree,lem:v17v2-tilt,thm:v17v2-one-sided,cor:v17v2-polynomial}

Stable polynomial interactions with bounded coefficients on $\mathcal G$
therefore give a finite-regulator Taylor estimate.  A cofinal conclusion
requires, in addition,
\begin{equation}
 \sum_n\epsilon_n^3M_{3,n}(S)<\infty
 \label{eq:summable-third-shell}
\end{equation}
and a separate bound on the complement of $\mathcal G$.  The first two
cumulants, local counterterms, and transported observables must still be
matched.  Equation~\eqref{eq:summable-third-shell} supplies one term in
the complete comparison theorem, not its entire hypothesis set.


\section{Local subtraction before absolute estimates}
The shell effective action has the connected expansion~\eqref{eq:shell-cumulants}. Its local relevant and marginal jets must be assigned to the running local action. The remaining shell is
\begin{equation}
 R_j=(I-P_{{\rm loc},j})\,S_{{\rm shell},j}.
 \label{eq:local-projector}
\end{equation}
If the local projector fails to intertwine the Ward complex, the remainder acquires a new symmetry defect. Uniform bounds on an unprojected shell do not establish a symmetry-compatible renormalization.

\begin{proposition}[Marginal subtraction and summable complements]
On the four-dimensional marginal bubble branch of the source, Taylor-jet subtraction converts the nonlocal complement into a summable shell contribution. At the unsubtracted marginal order the logarithmic accumulation remains. The complete BV-coupled derivative theorem requires the same subtraction in all source and ghost components and an invariant local projector.
\end{proposition}
\cite{R2,R16,R18}

The order of operations matters. Signed contractions, Ward cancellations, Born subtractions, and exact common-action combinations are formed first. Only the resulting residual is bounded in absolute value. Taking separate absolute values before cancellation can replace the true small residual by a non-summable bound. Conversely, declaring a cancellation without the common coefficient writers does not establish it.

\section{A stable irrelevant graph needs a base trajectory}
\begin{theorem}[Nonautonomous stable graph]
\label{thm:stable-graph}
Let a bounded sequence of retained relevant/marginal coordinates $u_j$ be given. Suppose the irrelevant recursion is
\begin{equation}
 v_j=F_j(u_{j+1},v_{j+1}),
 \label{eq:stable-recursion}
\end{equation}
where each $F_j$ maps the common ball of radius $R$ into itself, is Lipschitz in the second argument with constant $c<1$, and satisfies $\sup_j\norm{F_j(u_{j+1},0)}\le r_*$ with $r_*+cR\le R$. Then there is a unique bounded solution on that ball and
\begin{equation}
 \sup_j\norm{v_j}\le\frac{r_*}{1-c}.
 \label{eq:stable-bound}
\end{equation}
If the maps are uniformly Lipschitz in the base trajectory with constant $L_u$, then
\begin{equation}
 \norm{v(u)-v(\widetilde u)}_\infty
 \le\frac{L_u}{1-c}\,d(u,\widetilde u).
\end{equation}
\end{theorem}
\begin{proof}
The right side of~\eqref{eq:stable-recursion} defines a contraction of the product ball in the supremum norm. Its fixed point is unique. Comparing it with zero proves~\eqref{eq:stable-bound}; comparing two fixed points proves the Lipschitz estimate.
\end{proof}
\src{R5}{thm:v5v2-Lipschitz}
\scope{This theorem controls irrelevant coordinates above an already supplied bounded base trajectory. It neither constructs that trajectory nor establishes an ultraviolet completion of hypercharge or gravity.}

On the source's Gaussian-boundary convention, the one-loop coefficients are $b_3=7$, $b_2=19/6$, and $b_Y=-41/6$. Their role is to distinguish the nonabelian asymptotically free directions from the hypercharge Landau tendency on that branch. They are not a theorem excluding every non-Gaussian ultraviolet completion. Likewise, a finite Einstein--Hilbert coupling list is not closed under all higher-curvature counterterms. A claimed continuum trajectory must say which larger local action space is being controlled. \cite{R3,R4,R5}

\section{Each occurrence is charged once}
The repeated shell estimates can be organized by ancestry: an insertion has one shell of origin, one actual comparison path, and the corresponding descendants. At a merger of several branches, a single resampled shell occurrence must not be counted as several independent small factors. The exact Doob or reverse-martingale decomposition identifies which covariance actually owns that occurrence. This yields an additive error ledger rather than an artificial product of gains. \cite{R2,R4,R5}

A scale tradeoff illustrating the principle has two bounds for the same error,
\begin{equation}
 \epsilon(a,\Lambda)\lesssim
 \min\{a^b\Lambda^r,\ a^{-h}\Lambda^{-s}\}.
 \label{eq:two-regime}
\end{equation}
Choose $\Lambda=a^{-\gamma}$. Both displayed powers decay for a common $\gamma$ precisely when
\begin{equation}
 \frac hs<\gamma<\frac br,
 \qquad\text{equivalently}\qquad sb>rh,
 \label{eq:scale-window}
\end{equation}
for positive exponents. The interval must be nonempty before a summable refinement schedule is claimed. Counting the same annulus twice can falsely produce such an interval. \cite{R5,R6}

\section{Massless Coulomb terms retain their polynomial grade}
\begin{theorem}[Two-ended Coulomb cancellation]
\label{thm:coulomb}
Let a three-dimensional Green kernel satisfy the source's finite-difference bounds of Coulomb order. Suppose two compactly rooted sources have cancellation orders $r$ and $s$ with controlled moment norms $M_r,M_s$. Then their sandwich satisfies
\begin{equation}
 \abs{\sum_{z,w}f_x(z)G(z,w)h_y(w)}
 \le C M_r(f)M_s(h)(1+|x-y|)^{-1-r-s}.
 \label{eq:coulomb}
\end{equation}
Its absolute row sum is finite when $r+s>2$. At $r+s=2$, the displayed estimate has a logarithmically divergent row sum.
\end{theorem}
\begin{proof}
Subtract the Taylor terms killed by the moments at both ends. The surviving mixed finite difference has $r+s$ extra powers of decay. Summing shells in three dimensions requires $1+r+s>3$.
\end{proof}
\src{R5}{thm:v5v91-two-ended}
Gauge invariance can give neutrality without removing the first multipole. Exponential localization of a massive factor at one end does not erase a surviving Coulomb tail carried by the other part of the expression. The massless sector therefore needs its own cancellation and infrared ledger.

\section{Finite-rank determinants, phase ratios, and zero divisors}
A finite-rank Sylvester factor can retain zeros exactly while the weak-Schatten complement is treated by a modified determinant. This is more informative than a global lower bound that excludes physical crossings. Relative phase ratios on two nearby charts can remain stable even when a global phase mean is very small; normalized local observables and a partition-function lower bound are nevertheless separate questions.

For an invertible moving finite block in a fixed transported frame,
\begin{equation}
 \frac{\det B_1}{\det B_0}
 =\exp\left(\int_0^1\Tr(B_t^{-1}\dot B_t)\,\dd t\right).
 \label{eq:logdet-path}
\end{equation}
A dimension-free estimate uses a trace norm for the trace-class product. Replacing it by an operator norm without a dimension factor is not valid. The local divergent jet must again be subtracted before applying an absolute trace bound. \cite{R6,R7,R11}

\gate{The exact eliminations are established modules. Their physical use still requires a bounded renormalized base trajectory, a common invariant local projector, uniform complete source derivatives, a phase-compatible hard/soft partition, and summable errors on the same cutoff family.}


% ==================== 06_gravity_positive.tex ====================
\chapter{Positive gravitational sectors and the conformal obstruction}
\label{ch:positive-gravity}
\section{Why matter stability does not solve gravity stability}
The real Euclidean Einstein--Hilbert action, with the convention
\begin{equation}
 S_E(g)=-\frac1{2\kappa}\int_X(R_g-2\Lambda)\,\dd\vol_g,
\end{equation}
is not bounded below on the unrestricted real conformal direction. For $g=e^{2\sigma}g_0$ in dimension four, integration by parts gives
\begin{equation}
 S_E(e^{2\sigma}g_0)
 =-\frac1{2\kappa}\int_X
 \left(e^{2\sigma}R_0+6e^{2\sigma}|\nabla\sigma|^2
       -2\Lambda e^{4\sigma}\right)\dd\vol_0.
 \label{eq:conformal-action}
\end{equation}
Boundary conditions are chosen so that the displayed integration by parts is legitimate.

\begin{theorem}[Fixed-volume conformal unboundedness]
\label{thm:conformal}
On the real Euclidean branch of~\eqref{eq:conformal-action}, bounded-amplitude conformal oscillations can drive $S_E$ to $-\infty$, including after a fixed-volume normalization. A matter factor stable in its own variables does not remove a conformal direction on which that factor contributes no compensating high-frequency coercivity.
\end{theorem}
\begin{proof}
Choose a bounded oscillatory $\sigma_k$ whose gradient energy grows like $k^2$. An additive constant enforces fixed volume and remains bounded. The curvature and cosmological terms stay bounded on this family, whereas the gradient contribution in~\eqref{eq:conformal-action} has a negative coefficient and diverges. Holding spectator matter in a bounded sector leaves the same instability.
\end{proof}
\src{R9}{thm:n9v93-conformal-unboundedness,cor:n9v93-no-direct-EH-coercivity}

This is an obstruction to the direct real positive-density route. It is not a claim that every constrained, Lorentzian, complex-contour, or non-Gaussian gravitational construction fails. It specifies what those alternative routes must avoid or replace.

\section{Two inadequate repairs}
\begin{obstruction}[Imaginary rotation is not a real probability law]
\label{obs:rotation}
Rotating a wrong-sign Gaussian variable onto an imaginary contour can make the contour integral convergent, but the resulting response to a real-field characteristic source has the form $\exp(t^2/(2a))$ for $a>0$. Since its modulus exceeds one for $t\ne0$, it is not the characteristic function of a probability measure on that real observable.
\end{obstruction}
\begin{proof}
Every probability characteristic function is the expectation of a unit-modulus function and therefore has modulus at most one. The rotated Gaussian violates this necessary bound.
\end{proof}
\cite{R9,R10}
A contour construction needs an actual admissible contour, relative homology, singular-divisor control, compatible orientations and phase anchors, normalized new-mode integrals, and a physical reflection criterion. Formal convergence along an imaginary line supplies none of the last requirements automatically.

\begin{obstruction}[Higher-derivative positivity is not reflection positivity]
\label{obs:higher-derivative}
The covariance
\begin{equation}
 \frac1{p^2+\varepsilon p^4}
 =\frac1{p^2}-\frac1{p^2+\varepsilon^{-1}}
 \label{eq:hd-cov}
\end{equation}
has a negative massive spectral contribution. Its positive quadratic action therefore does not imply a positive reflected Gram form. Removing the stabilizer while keeping fixed low modes also need not preserve a uniform comparison density.
\end{obstruction}
\cite{R9,R19}
The negative reflected tests in the source are finite-dimensional witnesses of this distinction. They exclude the stated higher-derivative route under its stated identification of physical observables, not every possible enlarged constrained theory.

\section{Eliminate the unstable constraint direction rather than integrate it}
A different route solves a conformal or longitudinal constraint before constructing a positive reduced law. Write the eliminated variable as $s$ and retained variables as $m$. In matched energy spaces, suppose the constraint can be written
\begin{equation}
 C_0s=Bm+N(s,m).
 \label{eq:constraint-graph}
\end{equation}
If $C_0^{-1}N$ is a contraction in $s$ with constant $q<1$, the equation has a local unique graph $s=\sigma(m)$. The response bounds follow from
\begin{equation}
 D\sigma=(C_0-D_sN)^{-1}(B+D_mN).
 \label{eq:constraint-response}
\end{equation}
A coordinate Euclidean norm is not a substitute for the matched operator domain and dual norm. \cite{R9,R10}

\begin{proposition}[Reduced Hessian and the eliminated-source term]
\label{prop:reduced-hessian}
For $S_{\rm red}(m)=S(\sigma(m),m)$ and a tangent $v$,
\begin{equation}
 D^2S_{\rm red}[v,v]
 =D^2S[(D\sigma v,v),(D\sigma v,v)]
   +D_sS\,[D^2\sigma(v,v)].
 \label{eq:reduced-hessian}
\end{equation}
The last term vanishes when the eliminated variable is stationary for this same action along the graph. Solving a different constraint does not justify dropping it. Under the quantified inverse, response, and Hessian-margin assumptions in the source, the reduced form has a positive floor.
\end{proposition}
\begin{proof}
Differentiate the composite map twice. The second derivative of the graph couples exactly to the remaining eliminated Euler derivative. The positive-floor statement follows by bounding this term and the tangent-form perturbations against the declared reference floor.
\end{proof}
\cite{R10}

\section{Constraint surfaces, gauge quotients, and coarea}
For a regular finite-dimensional constraint $C:\R^N\to\R^r$, put
\begin{equation}
 J_C(x)=\sqrt{\det(DC(x)DC(x)^*)}.
\end{equation}
On a regular chart the hard-constraint limit has density $J_C^{-1}\rho$ relative to the induced surface measure. Coordinate Jacobians, this coarea factor, and the gauge-orbit/Faddeev--Popov factor are three different objects. One cannot replace all of them by a single determinant with an unspecified sign. \src{R10}{prop:v44-three-determinants}

\begin{theorem}[Regular positive constraint reduction]
\label{thm:coarea}
Assume the derivative of $C$ has a uniform nonzero normal singular-value floor, the chosen gauge slice is regular modulo its declared stabilizer, and the positive ambient densities have the domination needed for the shrinking-penalty limit. Then the finite hard-constraint law is positive and is given by coarea on the regular physical locus. Uniform normal floors and summable complete log-density comparisons yield the corresponding regular-chart cutoff comparison.
\end{theorem}
\src{R8}{thm:e8v58-coarea,thm:e8v83-second-coarea,thm:e8v99-coarea-ratio}
Singular strata, changing stabilizer rank, and an escaping exceptional set are not covered by this regular-chart theorem. A global quotient requires their additional measure and phase control.

A useful boundary calculation concerns the fourth-order normal constraint operator in three dimensions. Its leading heat coefficient in the source convention is
\begin{equation}
 C_{\rm con}\vol(\Sigma),\qquad
 C_{\rm con}=\frac{\Gamma(3/4)(17+3\sqrt3)}{64\pi^2}.
 \label{eq:coarea-heat}
\end{equation}
Fixed volume cancels this leading comparison term, and the stated closed parity setting removes the odd coefficients. This closes specified regular determinant/coarea rows; it does not control the probability of leaving those rows. \cite{R8}

\section{Retaining the reflection boundary}
\begin{theorem}[Conditional-half reflection square]
\label{thm:half-square}
Let the full cut information be $b$, with positive law $\nu(\dd b)$. Conditional on $b$, suppose the negative and positive halves are reflected copies with independent conditional laws. Then every positive-half test $F$ satisfies
\begin{equation}
 \E[(\Theta F)F]
 =\int\left|\E[F\mid b]\right|^2\nu(\dd b)\ge0.
 \label{eq:half-square}
\end{equation}
\end{theorem}
\begin{proof}
Condition on $b$, use the reflected conjugacy of the two conditional expectations, and integrate their product against the positive boundary law.
\end{proof}
\cite{R10,R14}
All boundary information used by the action, constraints, phase, and observable must be retained for this factorization. A marginal that forgets necessary cut variables need not preserve it. Moreover, a positive bulk transverse-traceless boundary form has Dirichlet-to-Neumann sign $+|k|$, whereas the unreduced conformal boundary form has the opposite sign. Retention of a boundary variable does not repair a negative law for that variable. \src{R9}{thm:n9v97-Einstein-boundary-symbol}

\section{Contours and determinant noncollapse}
A concrete analytic module for a declared complex Gaussian contour is the following. Let $C$ be complex symmetric and trace class, with $\norm{C}_{\rm op}\le q<1$. Finite normalized Gaussian integrals converge to
\begin{equation}
 \E\exp\!\left(\frac12x^TCx\right)=\det(I-C)^{-1/2},
 \qquad
 \abs{\log\det(I-C)^{-1/2}}
 \le\frac{\norm{C}_1}{2(1-q)}.
 \label{eq:complex-det}
\end{equation}
The branch is anchored at $C=0$, and trace-norm perturbations give the corresponding Lipschitz comparison. Wick subtraction gives the Hilbert--Schmidt variant with a second modified determinant. These are noncollapse and comparison results for the declared contour integral; they do not overturn Obstruction~\ref{obs:rotation}. \cite{R11,R12}

\gate{A viable gravitational route must jointly supply positivity on the physical observables, nondegenerate metric control, a complete reduced source, reflection compatibility, and a uniform treatment of singular strata or contour divisors. Matter stability alone supplies no such joint theorem.}


% ==================== 07_cartan.tex ====================
\chapter{Ordered geometry and the Einstein-divergence gate}
\label{ch:cartan}
\section{Noncentral coframes and inverse control}
The relational carrier can have noncommuting coefficient writers. Order must therefore be fixed before classical contractions are used. Let $E$ be an invertible matrix of coefficient writers and $\eta$ the declared reference signature. The metric writer is
\begin{equation}
 G=E^*\eta E.
 \label{eq:ordered-metric}
\end{equation}
A uniform inverse margin is needed to transport this expression and its derivatives. In a Banach algebra,
\begin{equation}
 \dd(E^{-1})=-E^{-1}(\dd E)E^{-1}.
 \label{eq:inverse-coframe}
\end{equation}
Higher derivatives are ordered sums, not commuting scalar formulas. If $E_n\to E$ with a common inverse bound, then $E_n^{-1}\to E^{-1}$; convergence of the required derivatives follows under the corresponding stronger hypotheses. \cite{R13}

\begin{lemma}[Ordered-word transport]
\label{lem:word-telescope}
For uniformly bounded factors in a Banach algebra,
\begin{equation}
\begin{split}
 \norm{\prod_{j=1}^{\ell}X_{j,n}-\prod_{j=1}^{\ell}X_j}
 \le\sum_{r=1}^{\ell}
 \left(\prod_{j<r}\norm{X_{j,n}}\right)
 \norm{X_{r,n}-X_r}
 \left(\prod_{j>r}\norm{X_j}\right).
\end{split}
 \label{eq:word-telescope}
\end{equation}
Thus componentwise convergence transports every fixed ordered contraction word without any commutation assumption.
\end{lemma}
\begin{proof}
Replace the factors one at a time and telescope. Submultiplicativity bounds each ordered summand.
\end{proof}
\src{R13}{lem:v13v96-word}

\section{Exact Cartan identities}
Let $(\Omega^\bullet,\dd)$ be an associative differential graded algebra with $\dd^2=0$ and the graded Leibniz rule. For a connection one-form $\omega$, define
\begin{equation}
 \mathcal R=\dd\omega+\omega^2,\qquad
 \mathcal T=\dd\theta+\omega\theta.
 \label{eq:ordered-curv}
\end{equation}
Here $\theta$ is a coframe column; this notation separates it from the coefficient matrix $E$. The left connection is $\nabla^L=\dd+\omega$, and the adjoint connection on a homogeneous form $A$ is
\begin{equation}
 \nabla^{\rm ad}A=\dd A+\omega A-(-1)^{|A|}A\omega.
\end{equation}

\begin{theorem}[Ordered Cartan--Bianchi identities]
\label{thm:cartan}
Without permuting noncommuting factors,
\begin{equation}
 (\nabla^L)^2X=\mathcal RX,\qquad
 \nabla^{\rm ad}\mathcal R=0,\qquad
 \nabla^L\mathcal T=\mathcal R\theta.
 \label{eq:cartan-bianchi}
\end{equation}
Consequently, torsion freedom implies $\mathcal R\theta=0$.
\end{theorem}
\begin{proof}
Expand $(\dd+\omega)^2X$. The terms $-\omega\dd X$ and $\omega\dd X$ cancel in their existing order, leaving $(\dd\omega+\omega^2)X$. Next,
\begin{align*}
 \dd\mathcal R&=(\dd\omega)\omega-\omega\dd\omega,\\
 \omega\mathcal R-\mathcal R\omega
 &=\omega\dd\omega-(\dd\omega)\omega,
\end{align*}
where the cubic terms cancel by associativity. This proves the second identity. Apply the first identity to $\theta$ to obtain the third.
\end{proof}
\cite{R13}
Frame change by an invertible $u$ gives
\begin{equation}
 \theta^u=u\theta,\qquad
 \omega^u=u\omega u^{-1}-(\dd u)u^{-1},\qquad
 \mathcal R^u=u\mathcal Ru^{-1},\quad \mathcal T^u=u\mathcal T.
\end{equation}
These identities prove frame covariance of curvature and torsion. They do not yet define a Ricci contraction or an Einstein tensor.

\section{Contracted conservation is an additional theorem}
For a proposed ordered contraction $\mathcal C$ and target covariant derivative $\mathfrak D$, define its intertwining defect by
\begin{equation}
 \mathcal J_{\mathcal C}(A)
 =\mathfrak D\mathcal C(A)-\mathcal C(\nabla^{\rm ad}A).
 \label{eq:contraction-defect}
\end{equation}

\begin{theorem}[Exact Einstein-contraction identity]
\label{thm:einstein-contraction}
For $\mathcal G=\mathcal C(\mathcal R)$ in an exact differential calculus,
\begin{equation}
 \mathfrak D\mathcal G=\mathcal J_{\mathcal C}(\mathcal R).
 \label{eq:contraction-gate}
\end{equation}
At finite cutoff, if $\mathcal B_n=\nabla_n^{\rm ad}\mathcal R_n$ is the actual Bianchi residual, then
\begin{equation}
 \mathfrak D_n\mathcal C_n(\mathcal R_n)
 =\mathcal J_{\mathcal C_n}(\mathcal R_n)+\mathcal C_n(\mathcal B_n).
 \label{eq:finite-contraction-gate}
\end{equation}
\end{theorem}
\begin{proof}
Insert the realized curvature into~\eqref{eq:contraction-defect}. In the exact calculus the adjoint derivative of curvature vanishes; at finite cutoff it equals the retained residual.
\end{proof}
\src{R13}{thm:v13v96-gate}

\begin{corollary}[Cofinal contracted conservation]
If both realized-input terms on the right of~\eqref{eq:finite-contraction-gate} tend to zero in the declared source topology, the contracted curvature is asymptotically conserved. A uniform complete bound on $\mathcal C_n$ and $\mathcal B_n\to0$ is sufficient for the second term, but is stronger than controlling $\mathcal C_n(\mathcal B_n)$ directly.
\end{corollary}
\src{R13}{cor:v13v96-cofinal}
This is the exact place where a formal Bianchi argument can otherwise overreach. An uncontracted identity does not prove that a proposed noncentral Ricci writer has a conserved contraction.

If $P$ and $s$ are candidate Ricci and scalar writers, the symmetric expression
\begin{equation}
 \mathcal G_J=P-\frac14(Gs+sG)
 \label{eq:jordan-einstein}
\end{equation}
is self-adjoint when its inputs are. Self-adjointness still does not establish the classical normalization, covariance, or the intertwining identity. Each is checked separately in the construction. Similarly, an ordinary coefficient-matrix trace need not be cyclic over a noncommutative coefficient algebra; the universal graded cyclic quotient is the appropriate setting for the cyclic characteristic identities. \cite{R13}

\section{Metricity and the torsion-free selector}
The flat compatible reference connection is $\Gamma_0=E^{-1}\dd E$. All metric-compatible corrections can be parameterized, in the source convention, by
\begin{equation}
 \Gamma=\Gamma_0+G^{-1}H,\qquad H^*=-H.
 \label{eq:metricity-param}
\end{equation}
The torsion-free equation is then a linear selector equation for the skew-adjoint coefficient $H$:
\begin{equation}
 G^{-1}H\theta=-\mathcal T_0.
 \label{eq:torsion-selector}
\end{equation}

\begin{theorem}[Local torsion-free metric selector]
\label{thm:torsion-selector}
At the central orthonormal anchor, the declared Cartan coefficient map is an isomorphism. On a perturbation chart where the inverse-times-perturbation norm is strictly below one, the torsion equation~\eqref{eq:torsion-selector} has a unique metric-compatible solution with a quantified inverse bound. The selector and its permitted parameter derivatives converge under the corresponding uniform coefficient and inverse bounds.
\end{theorem}
\begin{proof}[Proof mechanism]
The source gives the central coefficient inverse explicitly. Write the perturbed map as $L_0(I+L_0^{-1}\Delta L)$ and invert by the Neumann series. Differentiate the inverse identity and use ordered-word transport.
\end{proof}
\src{R13}{thm:v13v98-selector,thm:v13v99-cartan,prop:v13v99-isomorphism}
\scope{This is a local noncentral Cartan theorem with a quantitative margin. A global family of compatible charts, overlap control, and a cutoff-uniform inverse margin are still required for a global continuum connection.}

\section{Constraint angles and the three components of the Einstein residual}
\label{sec:constraint-angle}
Conservation constrains only part of the gravitational residual.  A
quantitative decomposition shows exactly which additional components
must be controlled.  On a declared Hilbert source screen $\mathcal Z$,
let $C:\mathcal Z\to\mathcal H_C$ be a bounded closed-range constraint
map and $K:\mathcal X\to\mathcal Z$ a bounded closed-range gauge-tangent
map.  Their Moore--Penrose inverses define orthogonal projections
\begin{equation}
 P=I-C^\dagger C,\qquad E=I-P,\qquad
 T=KK^\dagger,\qquad S=I-T.
 \label{eq:constraint-projections}
\end{equation}
Let $R$ project onto the reduced source space
$\mathcal M=\ker C\cap(\Ran K)^\perp$.  In this section $R$ denotes a
projection, not curvature.  Define
\begin{equation}
 \epsilon=\norm{ET},\qquad
 \eta=\norm{C^\dagger}\norm{CK}\norm{K^\dagger}.
 \label{eq:constraint-angle}
\end{equation}
The exact identity $ET=C^\dagger(CK)K^\dagger$ gives
$\epsilon\le\eta$.  In particular, exact constraint propagation
$CK=0$ makes the constraint-normal and gauge-tangent spaces orthogonal.

\begin{theorem}[Stable reduction and the three-defect estimate]
\label{thm:three-defect}
If $\epsilon<1$, then
\begin{equation}
 \norm{PSP-R}\le\epsilon^2,\qquad
 \norm{(PSP)^j-R}\le\epsilon^{2j}\quad(j\ge1).
 \label{eq:constraint-compression}
\end{equation}
Every residual $r\in\mathcal Z$ satisfies
\begin{equation}
 \norm r^2\le\norm{Rr}^2+
 \frac{\norm{C^\dagger}^2\norm{Cr}^2+
       \norm{K^\dagger}^2\norm{K^*r}^2}{1-\epsilon}.
 \label{eq:three-defect}
\end{equation}
For $CK=0$, one has $R=P-T$ and a single compression $PSP$ is exact.
\end{theorem}
\begin{proof}
Consider $\mathfrak Cx=(Ex,Tx)$ with codomain
$\Ran E\oplus\Ran T$.  Its Gram operator is
\[
 \mathfrak C\mathfrak C^*=
 \begin{pmatrix}I&ET\\TE&I\end{pmatrix}\succeq(1-\epsilon)I.
\]
Its kernel is $\mathcal M$, so
$\norm{(I-R)r}^2\le(\norm{Er}^2+\norm{Tr}^2)/(1-\epsilon)$.
The identities $Er=C^\dagger Cr$ and
$\norm{Tr}\le\norm{K^\dagger}\norm{K^*r}$ prove
\eqref{eq:three-defect} by orthogonality.
For the compression, let $A=T|_{\Ran P}$.  On $\Ran T$,
$AA^*=TPT=I-TET\succeq(1-\epsilon^2)I$, since
$\norm{TET}=\norm{ET}^2$.  Thus on the orthogonal complement of
$\ker A=\mathcal M$ in $\Ran P$, the operator
$PSP=I-A^*A$ has spectrum in $[0,\epsilon^2]$.  It is the identity
on $\mathcal M$ and vanishes on $\Ran E$.  This proves both norm
bounds.  When $ET=0$, $T\preceq P$ and the last assertion follows.
\end{proof}
\src{R14}{thm:v14v18-sharp,cor:v14v18-compression,thm:v14v18-joint,thm:v14v18-three}

\begin{corollary}[Cofinal residual control]
\label{cor:cofinal-three-defect}
On a family of localized source screens, assume uniform bounds on
$\norm{C_n^\dagger}$ and $\norm{K_n^\dagger}$ and
$\eta_n\le\eta_0<1$.  If
$\norm{C_nr_n}\to0$, $\norm{K_n^*r_n}\to0$, and
$\norm{R_nr_n}\to0$, then $\norm{r_n}\to0$.
\end{corollary}
\begin{proof}
Apply \eqref{eq:three-defect} with the common denominator
$1-\eta_0>0$.
\end{proof}
\src{R14}{thm:v14v18-cofinal}

For the realized Einstein residual, these are respectively the
constraint, gauge/Ward, and genuinely reduced dynamical components.
Vanishing divergence does not set the last component to zero.  Passage
from localized Hilbert estimates to a distributional Einstein equation
also requires the common test core and source integrability of
Chapter~\ref{ch:closure}.  The earlier ordered Bianchi and contraction
identities provide inputs to this decomposition, rather than a
substitute for all three estimates.\src{R13}{thm:v13v69-split,cor:v13v71-landing}



% ==================== 08_mixing.tex ====================
\chapter{Coercivity, mixing, and physical scales}
\label{ch:mixing}
\section{A finite floor must survive the actual coarse modes}
A fine quadratic energy can be coercive modulo rigid motions while its
coarse restriction loses rotations.  The relevant interface datum is the
whole face mismatch, not only its value at the face centre.  The following
identity identifies the missing term before any Schur complement or
mixing estimate is applied.

\begin{proposition}[Full-face control of rigid mismatches]
\label{prop:full-face-rigid}
Consider adjacent cubes of side $\ell$ in dimension $d\ge2$, with centres
$x_\pm=x_f\pm\ell e_\mu/2$ and common face $f$.  Let
$r_\pm(x)=a_\pm+\Omega_\pm(x-x_\pm)$, where
$\Omega_\pm^{\mathsf T}=-\Omega_\pm$.  Set
\begin{equation}
 J_\mu=a_+-a_--\frac\ell2(\Omega_++\Omega_-)e_\mu,
 \qquad \Delta_\mu\Omega=\Omega_+-\Omega_-,\qquad
 P_\mu=I-e_\mu\otimes e_\mu.
 \label{eq:face-rigid-definitions}
\end{equation}
Then
\begin{equation}
 \int_f|r_+-r_-|^2\,\dd S
 =\ell^{d-1}|J_\mu|^2+\frac{\ell^{d+1}}{12}
                  \norm{\Delta_\mu\Omega P_\mu}_F^2.
 \label{eq:full-face-moment}
\end{equation}
On the periodic cubic array, the associated full-face quadratic form has
only the global translations in its kernel.  In particular, the
checkerboard rotational mode omitted by centre sampling is not a zero
mode of this form.
\end{proposition}
\begin{proof}
Writing $y=x-x_f$ on the face gives
$r_+-r_-=J_\mu+\Delta_\mu\Omega y$.  Symmetry and elementary integration
on the centred cube face yield
\[
 \int_fy\,\dd S=0,
 \qquad \int_fy\otimes y\,\dd S=\frac{\ell^{d+1}}{12}P_\mu.
\]
Expansion of the square proves \eqref{eq:full-face-moment}.  For a Fourier
mode write $s_\mu=e^{\ii k_\mu}-1$ and $c_\mu=e^{\ii k_\mu}+1$.  After
division by the face area its energy is
\begin{equation}
 E_f(k;a,\Omega)=\sum_\mu\left[
 \left|s_\mu a-\frac\ell2c_\mu\Omega e_\mu\right|^2
 +\frac{\ell^2}{12}|s_\mu|^2\norm{\Omega P_\mu}_F^2\right].
 \label{eq:full-face-symbol}
\end{equation}
If $k\ne0$ modulo the reciprocal lattice, some $s_\mu\ne0$.
Zero energy then implies $\Omega P_\mu=0$.  Skew symmetry forces
$\Omega=0$, and the first terms force $a=0$.  At $k=0$, the energy is
$\ell^2\norm\Omega_F^2$, leaving exactly the translation amplitude $a$.
For $k=(\pi,\ldots,\pi)$ and $a=0$, centre sampling gives zero, whereas
\[
 E_f(k;0,\Omega)=\frac{d-1}{3}\ell^2\norm\Omega_F^2.
\]
The last identity follows by summing
$\norm{\Omega P_\mu}_F^2=\norm\Omega_F^2-|\Omega e_\mu|^2$.
\end{proof}
\src{R16}{lem:v16v79-moment,thm:v16v79-face-energy,thm:v16v79-no-hourglass}

For a face weight satisfying $w_-I\preceq W_f\preceq w_+I$, the weighted
energy lies between $w_-$ and $w_+$ times the unweighted one.  This is a
finite geometric control of the rigid sector.  The coarse Schur floor
still requires the stated fine/coarse coupling bounds, and a fixed-array
kernel calculation is not a volume-independent physical gap.  The
resulting positive-law certificate is an input to mixing, not a continuum
theorem.  \src{R16}{prop:v16v79-weighted,thm:v16v82-flat-schur,thm:v16v83-positive-gap}

\section{The componentwise covariance mechanism}
Let $\mu(\dd x)\propto e^{-U(x)}\dd x$ be a finite block law. Suppose each conditional block has the specified Poincar\'e floor $\rho_i>0$ and
\begin{equation}
 \norm{\nabla_i\nabla_jU}_{\rm op}\le\kappa_{ij}\quad(i\ne j).
\end{equation}
Set $D=\diag(\rho_i)$, $K_{ij}=\kappa_{ij}$ off the diagonal, and $\mathcal M=D-K$. Let $b_i(f)$ be the block-gradient $L^2(\mu)$ norms appearing in the Witten-resolvent estimate.

\begin{theorem}[Influence-resolvent covariance bound]
\label{thm:covariance}
Under the source's conditional Poincar\'e and resolvent-domain hypotheses, if $\mathcal M$ is a nonsingular $M$-matrix, then
\begin{equation}
 |\Cov_\mu(f,g)|\le b(f)^T\mathcal M^{-1}b(g).
 \label{eq:covariance-bound}
\end{equation}
A sufficient quantitative condition is
\begin{equation}
 \theta=\norm{D^{-1}K}_{\ell^\infty\to\ell^\infty}<1.
 \label{eq:influence}
\end{equation}
\end{theorem}
\begin{proof}[Proof mechanism]
The differentiated Poisson equation produces a componentwise inequality $Da\le Ka+b$. Positivity of the inverse of the nonsingular $M$-matrix gives $a\le\mathcal M^{-1}b$. Pair the resulting gradient resolvent with the gradient of the second observable.
\end{proof}
\src{R16}{thm:v16v84-sweeping,thm:v16v84-covariance}

If $K$ has range $R_0$ in block distance, its Neumann expansion has no contribution between $i,j$ before at least $\lceil\dist(i,j)/R_0\rceil$ steps. Thus
\begin{equation}
 0\le(\mathcal M^{-1})_{ij}
 \le\frac{\theta^{\lceil\dist(i,j)/R_0\rceil}}
 {\rho_{\min}(1-\theta)}.
 \label{eq:inverse-decay}
\end{equation}
This is a localized covariance estimate. A spectral gap for a configuration-space Markov generator alone need not imply spatial clustering without the locality information used here. \src{R16}{prop:v16v84-inverse-decay,prop:v16v84-gap-no-cluster}

\section{The continuum rate is measured in physical distance}
Let $\ell_n$ be the physical block size. Translating~\eqref{eq:inverse-decay} into physical distance gives
\begin{equation}
 |\Cov_{\mu_n}(f,g)|\le
 \frac{G_n(f)G_n(g)}{\rho_{\min,n}(1-\theta_n)}
 e^{-m_n d_n^{\rm phys}(\supp f,\supp g)},
 \qquad
 m_n=\frac{-\log\theta_n}{\ell_nR_{0,n}}.
 \label{eq:physical-decay}
\end{equation}
A cutoff-independent estimate requires both a positive lower bound on $m_n$ and a uniform bound on the displayed amplitude. Quoting only $\theta_n<1$ at every cutoff leaves both requirements undecided. \src{R16}{thm:v16v85-physical}

\begin{criterion}[Massive physical covariance regime]
\label{crit:massive}
For the selected massive observable family, require
\begin{equation}
 \inf_n\frac{-\log\theta_n}{\ell_nR_{0,n}}>0,
 \qquad
 \sup_n\frac{G_n(f)G_n(g)}{\rho_{\min,n}(1-\theta_n)}<\infty.
 \label{eq:massive-regime}
\end{equation}
Together with the common reflected reconstruction and limit hypotheses, these estimates can be used for a channelwise spectral threshold. They are not imposed as a positive global gap on the full massless Standard Model--Einstein theory.
\end{criterion}

\section{Local update gaps and global gaps differ}
A conditional heat-bath projection has a unit gap on its own conditional complement. The sum of many such updates has a uniform global gap only under an approximate tensorization, influence, or comparable irreducibility estimate. A two-block model makes the obstruction explicit: the relative angle, equivalently maximal correlation of the retained subspaces in the stated setting, controls the slow mode. Local invertibility does not bound that angle uniformly.

Quartic confinement creates another distinction. Conditional distributions can have total-variation influence arbitrarily close to one on remote configurations even when weighted Wasserstein responses remain finite. A contraction proved on a high-probability bounded region is not a global spectral-gap theorem without a return or Lyapunov estimate. Flat gauge Cartan valleys and changing orbit stabilizers likewise prevent a universal scalar-quartic coercivity argument. \cite{R15,R16}

\section{The Higgs--conformal Hessian is a signed response}
For an oriented edge $e=(p,c)$, conformal difference $d_e$, link $U_e$, and Higgs coordinates $y_x$, let
\begin{equation}
 h_e(d_e,y)=\chi_e\norm{e^{-d_e/2}y_c-e^{d_e/2}U_ey_p}^2.
 \label{eq:hopping}
\end{equation}
With $F_H=-\log Z_H$ and a fixed reference measure,
\begin{align}
 q_e&=\chi_e\bigl(e^{d_e}|y_p|^2-e^{-d_e}|y_c|^2\bigr),\\
 T_e&=\chi_e\bigl(e^{d_e}|y_p|^2+e^{-d_e}|y_c|^2\bigr),\\
 (H_H)_{ef}&=\partial_{d_e}\partial_{d_f}F_H
 =\delta_{ef}\E T_e-\Cov(q_e,q_f).
 \label{eq:higgs-hessian}
\end{align}
The first term and covariance must be assembled with their signs before a curvature bound is taken.

\begin{theorem}[Incidence-sharp negative curvature constant]
\label{thm:incidence}
Let $B$ be the oriented incidence matrix, $W$ a positive diagonal edge weight, and $\Pi$ the orthogonal projection onto $\Ran(W^{1/2}B)$. The smallest nonnegative constant bounding the negative Higgs Hessian on conformal incidence directions is
\begin{equation}
 K_{\rm WI}=\max\!\left\{0,
 \lambda_{\max}\!\left(
 \Pi W^{-1/2}(-H_H)W^{-1/2}\Pi
 \big|_{\Ran(W^{1/2}B)}\right)\right\}.
 \label{eq:kwi}
\end{equation}
Cycle components orthogonal to that range vanish from both the conformal quadratic form and its source pairing. In particular, $K_{\rm WI}$ is no larger than a valid whole-edge absolute covariance bound.
\end{theorem}
\begin{proof}
Write the conformal edge variation as $Bx$ and substitute $z=W^{1/2}Bx$. The ratio of its negative Hessian to its weighted norm is the Rayleigh quotient of the compressed matrix in~\eqref{eq:kwi}. The orthogonal complement never occurs in this quotient.
\end{proof}
\src{R17}{thm:v17v99-optimal}\cite{R18}

\subsection{The supported response and its operator domain}
The curvature estimate concerns a configuration-space domain, not merely a
spatial block.  Let $r$ be the mean-zero conformal field and
$\Omega_R=\{r:\abs{(Br)_e}<R\text{ for every }e\}$.  For the effective
conformal potential $U$, the one-form Witten quadratic form has the
potential contribution
\[
 \mathcal V_U(r)=2\kappa_gI+B^*\bigl(W_J(r)+H_H(r)\bigr)B,
\]
where $W_J$ is the edge Hessian of the retained conformal interaction.
Suppose the signed incidence estimate gives
$\mathcal V_U\succeq\mu_{R,L}^{\rm Ward}I$ on $\Omega_R$, with
$\mu_{R,L}^{\rm Ward}>0$.

\begin{theorem}[Supported Dirichlet response]
\label{thm:supported-response}
Let $\mathcal H^{(1)}_{U,R,D}$ be the Friedrichs realization of the
one-form Witten form on the closure of compactly supported smooth
one-forms in $\Omega_R$.  Then
\begin{equation}
 \mathcal H^{(1)}_{U,R,D}\succeq\mu_{R,L}^{\rm Ward}I,
 \qquad
 \norm{(\mathcal H^{(1)}_{U,R,D})^{-1}}
       \le(\mu_{R,L}^{\rm Ward})^{-1}.
 \label{eq:supported-floor}
\end{equation}
Assume the conditional Higgs law has a global Poincar\'e constant
$C_G$ and satisfies $\E|y_x|^2\le q_HC_G$ at every site.  Put
$A_t=\sum_x t_x|y_x|^2$,
$\zeta_{t,e}=-\Cov(A_t,q_e)$, and
$w_e=2\chi_e^2\cosh(2d_e)$, omitting zero-hopping edges from the
weighted sum.  If the real Higgs dimension is $q_H$, the
maximum degree is $D$, and $w_e\le w_R$ on $\Omega_R$, these estimates give
\begin{equation}
 \sum_e\frac{|\zeta_{t,e}|^2}{w_e}
 \le16q_H^2C_G^4D\norm{t}^2.
 \label{eq:supported-higgs-source}
\end{equation}
For $g_{t,R}=\one_{\Omega_R}B^*\zeta_t$, regarded in the Dirichlet
one-form space,
\begin{equation}
 \pair{g_{t,R}}{(\mathcal H^{(1)}_{U,R,D})^{-1}g_{t,R}}
 \le\frac{32q_H^2C_G^4D^2w_R}{\mu_{R,L}^{\rm Ward}}\norm{t}^2.
 \label{eq:supported-response}
\end{equation}
\end{theorem}
\begin{proof}
The derivative part of the Witten form is nonnegative.  The potential
floor therefore bounds the form on compactly supported one-forms and
passes to its Friedrichs closure.  The spectral theorem yields
\eqref{eq:supported-floor}.  Conditional Poincar\'e applied to
$A_t$ and $q_e$, followed by the one-site moment estimate and the
edge-degree sum, gives \eqref{eq:supported-higgs-source}.
Since $\norm{B}^2\le2D$, it implies
$\norm{g_{t,R}}^2\le32q_H^2C_G^4D^2w_R\norm{t}^2$.
Apply the inverse bound.
\end{proof}
\src{R17}{thm:v17v96-dirichlet,lem:v17v96-source,thm:v17v96-central-response}

A projection onto a spatial interior block does not impose support in
$\Omega_R$.  Thus \eqref{eq:supported-floor} does not bound the
unrestricted global inverse, and \eqref{eq:supported-response} is not an
identity for the full environmental variance.  Both the source and the
inverse domain have been restricted.  An unrestricted response theorem
needs an exterior form estimate and control of the localization error;
these inputs cannot be replaced by the central-chart floor.
\src{R17}{prop:v17v96-erratum}


\section{The canonical scaling bottleneck}
The canonical four-dimensional scaling in this branch is
\begin{equation}
\begin{gathered}
 \kappa_g=\bar\kappa_g a^4,\quad b_g=\bar b_ga^4,
 \quad J=\bar Ja^2,\quad y_x=aH(ax),\\
 \chi_e=\bar\chi_e,\qquad m_H^2=\bar m_H^2a^2,
 \qquad\lambda_H=\bar\lambda_H.
\end{gathered}
\label{eq:canonical-scaling}
\end{equation}
The available one-site and absolute covariance constants are order one, while the relevant conformal force is order $a^2$.

\begin{proposition}[Fixed-regulator massless sign reduction]
\label{prop:massless-sign}
For fixed finite graph, conformal data, and links, the Higgs Hessian is analytic in $\mu=m_H^2$ on the stated stable branch, and
\begin{equation}
 K_{\rm WI}(\mu)=K_{\rm WI}(0)+O(|\mu|).
\end{equation}
Along $\mu=O(a^2)$,
\begin{equation}
 K_{\rm WI}(a)=O(a^2)
 \quad\Longleftrightarrow\quad B^TH_H(0)B\succeq0.
 \label{eq:massless-sign}
\end{equation}
A cofinal version additionally requires the massless sign and the normalized mass-derivative bound uniformly over the regulator family.
\end{proposition}
\begin{proof}
Analyticity gives the local Lipschitz estimate in $\mu$; eigenvalue perturbation gives the same estimate for the compressed negative part. If the massless compression is nonnegative, $K_{\rm WI}(0)=0$. Conversely, an $O(a^2)$ bound forces its continuous limit to vanish. Uniformity cannot be obtained by repeating this fixed-matrix argument with unspecified constants.
\end{proof}
\src{R17}{thm:v17v100-massless-expansion,cor:v17v100-massless-gate}\cite{R18}

The differentiated mass identity displays what a uniform version must
control.  For $A=\sum_x|y_x|^2$ and $\mu$-independent edge insertions,
\begin{align}
 \partial_\mu\E_\mu F&=-\Cov_\mu(F,A),\\
 \partial_\mu(H_H)_{ef}
 &=-\delta_{ef}\Cov_\mu(T_e,A)
                  +\Cum_\mu(q_e,q_f,A).
 \label{eq:mass-hessian-response}
\end{align}
The second formula follows by differentiating the covariance in
\eqref{eq:higgs-hessian}; the derivative of a covariance is the negative
third cumulant.  For a $\mu$-independent positive weight $W$,
\begin{equation}
 |K_{\rm WI}(\mu)-K_{\rm WI}(0)|
 \le |\mu|\sup_{|v|\le|\mu|}
 \norm{\Pi W^{-1/2}(\partial_vH_H)W^{-1/2}\Pi}.
 \label{eq:normalized-mass-derivative}
\end{equation}
At fixed regulator the quartic tail justifies these derivatives.  The
insertion $A$ is extensive, however, and a finite bound does not supply a
volume- or cutoff-uniform bound.  A cofinal sign reduction therefore
requires the supremum in \eqref{eq:normalized-mass-derivative} in its
actual incidence norm, together with the massless sign, uniformly on the
retained field branch.  This identifies the signed response to estimate
without manufacturing an $a^2$ gain from an order-one absolute bound.
\src{R17}{lem:v17v100-mass-derivative,thm:v17v100-massless-expansion}

\section{Trees, finite collisions, and cycle restoration}
Tree elimination gives exact radial messages and a controlled way to separate one-leaf response from the exterior covariance. Mixed responses retain the subtraction in~\eqref{eq:higgs-hessian}; they are not products of independent scalar leaf bounds. The source develops a closed radial message class, local positive sectors, a large-boundary Schur separator, stiffness-dependent contraction, and an additive tree response ledger. These are finite-regulator mechanisms for the massless sign problem, not a proof of it on arbitrary cyclic physical regulators. \cite{R18}

\subsection{Exact radial messages on trees}
For a radial $O(N)$-invariant on-site law on a finite tree, all link transports can be removed by vertex rotations of unit Jacobian. Choose the root rotation arbitrarily and define each child rotation by its unique incoming edge; no cycle holonomy obstructs the recursion. The partition function and its conformal Hessian are therefore independent of the tree link variables. The incidence matrix of a tree is surjective onto its edge space, so the incidence condition becomes ordinary edge-Hessian positivity. \src{R18}{thm:v18v1-gauge-erasure,thm:v18v1-incidence-surjective}

For a convex nondecreasing radial potential $\Psi$, define
\begin{equation}
 M_\Psi(d,z)=\int_{\R^N}
 \exp\!\left[-\Psi(|u|^2)-\chi|e^{-d/2}u-e^{d/2}z|^2\right]\dd u,
 \qquad-\log M_\Psi(d,z)=\ell_\Psi(d,|z|^2).
 \label{eq:radial-message}
\end{equation}

\begin{theorem}[Closed radial message class]
\label{thm:radial-closure}
For the finite integral~\eqref{eq:radial-message}, the conditional curvature $\partial_d^2\ell_\Psi(d,q)$ is nonnegative, and $q\mapsto\ell_\Psi(d,q)$ is nondecreasing and convex. For nonconstant $\Psi$, the connecting-edge curvature is strictly positive and the radial message is strictly increasing. Consequently, an effective on-site potential
\begin{equation}
 \Psi_x(t)=\lambda_xt^2+\sum_{c:\,p(c)=x}\ell_c(t)
\end{equation}
stays in the same radial class through the entire finite tree.
\end{theorem}
\begin{proof}[Proof mechanism]
Write the angular integral as the positive radial power series used in the source. Its radial moments
\begin{equation}
 J_s(a)=\int_0^\infty t^{s-1}e^{-\Psi(t)-at}\dd t
\end{equation}
obey the integration-by-parts recurrence $s=m_s(a+G_{s+1})$, where $m_s=J_{s+1}/J_s$ and $G_s$ is the shape-$s$ expectation of $\Psi'$. Monotonicity of these radial expectations gives the conditional curvature and slope inequalities. Sums with a positive quartic on-site term preserve convexity and monotonicity, so induction proves closure.
\end{proof}
\src{R18}{thm:v18v3-closed-kernel,cor:v18v3-tree-closure}

Conditional edge convexity is not yet convexity after integrating the parent. For $A(d,q)=\Phi(q)+\ell_\Psi(d,q)$ and $\widehat F(d)=-\log\int e^{-A(d,|z|^2)}\dd z$,
\begin{equation}
 \widehat F''(d)=\widehat\E\,\ell_{dd}
                 -\Var_{\widehat\mu}(\ell_d).
 \label{eq:parent-covariance}
\end{equation}
The radial Hessian eigenvalue is $R_A=2A_q+4qA_{qq}$. The radial Brascamp--Lieb estimate in the source gives
\begin{equation}
 \widehat F''(d)\ge\widehat\E
 \left[\ell_{dd}-\frac{4q\ell_{dq}^2}{R_A}\right].
 \label{eq:radial-schur}
\end{equation}
Thus $R_A\ell_{dd}\ge4q\ell_{dq}^2$ is a sufficient pointwise Schur condition; the exterior score variance is not omitted. \src{R18}{thm:v18v4-radial-schur}

\subsection{A coefficient-only depth-independent tree margin}
In the source's compact radial coefficient sector, let $a_*,a^*,b_*,\lambda_*,\lambda^*,\chi^*$ be the positive endpoint bounds, let $B^*$ bound the sum of child slope ceilings, and let $\mathfrak b$ bound branching. For the finite allowed set of half-dimensions $\mathfrak A$, set
\begin{equation}
 \mathfrak m_s(c,\lambda)=
 \frac{\int_0^\infty t^s e^{-\lambda t^2-ct}\dd t}
      {\int_0^\infty t^{s-1}e^{-\lambda t^2-ct}\dd t},
\end{equation}
\begin{equation}
 \begin{aligned}
 M_*&=\min_{\alpha\in\mathfrak A}
       \mathfrak m_{\alpha+1}(a^*+B^*,\lambda^*),
 &G_*&=2\lambda_*M_*,\\
 \gamma_*&=\frac{b_*G_*}{a^*+G_*},
 &p_0&=a_*\min_{\alpha\in\mathfrak A}
       \mathfrak m_\alpha(a^*+B^*,\lambda^*).
 \end{aligned}
\end{equation}
Every outgoing message has slope at least $\gamma_*>0$. The radial response with additional stiffness $\sigma$ satisfies $0\le m'(Q)\le\chi^2/(a+\sigma)^2$. Therefore define
\begin{equation}
 \bar\rho=\left(\frac{\chi^*}{a_*+\gamma_*}\right)^2,
 \quad\rho_\partial=\left(\frac{\chi^*}{a_*}\right)^2,
 \quad C_\partial=\max\{1,\rho_\partial/\bar\rho\}.
\end{equation}
\src{R18}{thm:v18v24-stiffness,thm:v18v25-stiffness}

\begin{theorem}[Uniform tree Hessian from coefficient endpoints]
\label{thm:tree-margin}
For this exact radial message family and its recursive Schur transport, suppose $0<\bar\rho<1/\mathfrak b$ and
\begin{equation}
 \mathfrak Q_0=
 \frac{C_\partial^2(a^*)^2}
 {3\lambda_*p_0(1-\mathfrak b\bar\rho)(1-\bar\rho)}<1.
 \label{eq:tree-certificate}
\end{equation}
Then the complete finite-tree conformal Hessian obeys
\begin{equation}
 H_{\mathcal T}(Q)\succeq(1-\mathfrak Q_0)p_0I
 \label{eq:tree-floor}
\end{equation}
for every admitted exterior boundary value, independently of depth and edge count.
\end{theorem}
\begin{proof}[Proof mechanism]
The quartic radial moment comparison supplies the positive outgoing stiffness floor. It bounds each internal response by $\bar\rho$ and each leaf response by $\rho_\partial$. The exact recursive Schur ledger adds the fresh positive edge energies and subtracts the transported response squares once. Summing depth and branching contributions gives the two geometric denominators in~\eqref{eq:tree-certificate}. The remaining fraction of the fresh floor is~\eqref{eq:tree-floor}.
\end{proof}
\src{R18}{thm:v18v19-ledger,thm:v18v25-tree-gate}

\subsection{Restoring cycles}

A particularly concise cyclic certificate uses one-site half-Hessian floors $A_v>0$, symmetric hopping $\chi_{vw}$, and
\begin{equation}
 C_{vw}=\frac{\chi_{vw}}{A_v},\qquad
 S_{vw}=\frac{\chi_{vw}}{\sqrt{A_vA_w}}.
 \label{eq:normalized-hopping}
\end{equation}
The matrices are similar: $S=A^{1/2}CA^{-1/2}$.

\begin{theorem}[One certificate for response and convexity]
\label{thm:hopping}
For the finite connected graph, the existence of positive weights $w$ and $\kappa<1$ with $Cw\le\kappa w$ is equivalent to $\rho(C)<1$, hence to $\lambda_{\max}(S)<1$. If $\lambda_{\max}(S)\le\kappa<1$, the complete field Hessian on every partial cycle-restoration interpolation satisfies
\begin{equation}
 \nabla^2S_s\succeq2(1-\kappa)\diag(A_vI),\qquad0\le s_e\le1.
 \label{eq:hopping-floor}
\end{equation}
\end{theorem}
\begin{proof}
Positive-matrix spectral theory and diagonal similarity give the equivalence. The mixed hopping quadratic form is bounded by the normalized symmetric matrix $S$ against the one-site energy. Subtracting that bound gives~\eqref{eq:hopping-floor}, uniformly when edges are weakened by $s_e\in[0,1]$.
\end{proof}
\src{R18}{thm:v18v33-equivalence,thm:v18v33-convexity}

\begin{theorem}[Cycle-restoration margin]
\label{thm:cycle-restoration}
Suppose the cut reference conformal Hessian has floor $\mu_0>0$, the normalized hopping certificate holds, and the source's uniform quartic moment, local response, collision-block, and two-/three-arm resolvent estimates hold. Let $\mathcal L_{\rm cyc}$ be their complete integrated cycle-insertion bound. If $\mathcal L_{\rm cyc}<\mu_0$, then
\begin{equation}
 H_1\succeq(\mu_0-\mathcal L_{\rm cyc})I>0
 \label{eq:cycle-margin}
\end{equation}
uniformly in volume.
\end{theorem}
\begin{proof}
Integrate the exact Hessian derivative along cycle restoration. The complete compiled bound gives $\norm{H_1-H_0}\le\mathcal L_{\rm cyc}$; the claimed floor follows. Three-point insertions require a three-arm summation and cannot be inferred from a pair-path estimate alone.
\end{proof}
\src{R18}{thm:v18v34-cycle-margin}
\gate{The strict inequalities in this cyclic theorem have not been verified for a physical cofinal Standard Model--Einstein family. Moreover, strong uniform mixing must be checked against the interaction-survival obstruction developed in Chapter~\ref{ch:stress}.}


% ==================== 09_states.tex ====================
\chapter{Joint states, no escape, and cofinal comparison}
\label{ch:states}
\section{Normalized coercivity and raw-field tightness}
A positive finite law needs normalization estimates as well as a confining exponent. Let the reference $\lambda_n$ be a probability measure and
\begin{equation}
 \mu_n=Z_n^{-1}e^{-V_n}\lambda_n.
\end{equation}

\begin{lemma}[Normalized coercivity]
\label{lem:normalized-coercivity}
Suppose $V_n\ge cE_n-C_n$ with $E_n\ge0$, and $Z_ne^{-C_n}\ge z_*>0$. Then, for $0\le\alpha\le c$,
\begin{equation}
 \E_{\mu_n}e^{\alpha E_n}\le z_*^{-1}.
 \label{eq:normalized-coercivity}
\end{equation}
\end{lemma}
\begin{proof}
The numerator is bounded by $e^{C_n}\int e^{-(c-\alpha)E_n}\dd\lambda_n\le e^{C_n}$. Divide by the partition function and use the stated lower bound.
\end{proof}
\src{R9}{thm:n9v93-normalized-coercivity}
The free-energy lower bound cannot be dropped. A cutoff-dependent additive constant in the exponent may disappear from a formal coercivity inequality while its normalization still degenerates.

\begin{proposition}[Nuclear characteristic-function criterion]
\label{prop:minlos}
Suppose the finite bosonic characteristic functionals converge pointwise on a nuclear test space and are equicontinuous at zero. Then their positive normalized limit defines the corresponding generalized random field. A useful sufficient modulus is obtained when
\begin{equation}
 |\Phi_n(f)|\le X_np(f),\qquad
 \sup_n\mathbb P(X_n>R)\le a(R)\downarrow0:
\end{equation}
\begin{equation}
 1-\operatorname{Re}Z_n(f)
 \le\inf_{R>0}\left(\frac{R^2p(f)^2}{2}+2a(R)\right).
 \label{eq:characteristic-modulus}
\end{equation}
\end{proposition}
\begin{proof}
On $X_n\le R$, use $1-\cos t\le t^2/2$; on its complement use $1-\cos t\le2$. The modulus gives continuity of the positive-definite limit at zero, which is the nuclear-space measure criterion used in the source.
\end{proof}
\src{R9}{thm:n9v91-Minlos}

Negative-Sobolev moments give a complementary route: control in a stronger negative Sobolev space and compact inclusion into a weaker one yield tightness on compact regions. Coordinatewise tightness without a topology-level compact-containment estimate does not exclude escape into finer and finer modes. The enhanced stress and metric coordinates may require stronger estimates than the elementary fields. \src{R16}{thm:v16v88-tightness}

\section{A nonempty fixed-spacing Higgs--gauge--conformal branch}
\label{sec:fixed-spacing-gibbs}
The abstract compactness criteria have a concrete intermediate application.
It concerns a positive conformal lattice model at fixed spacing, rather
than a positive realization of the full Euclidean Einstein action.  This
distinction allows an actual thermodynamic existence theorem without
conflating volume removal with ultraviolet removal.

On a finite graph $G$, let $s_x\in\R$, $y_x\in\R^{q_H}$, and let $U_e$
belong to a compact gauge group acting orthogonally on the real Higgs
coordinates.  With $d=Bs$, use the hopping term \eqref{eq:hopping} and set
\begin{equation}
\begin{split}
 S_G(s,y,U)={}&\sum_x
 \bigl(\kappa_gs_x^2+b_gs_x+\mu|y_x|^2+\lambda_H|y_x|^4\bigr)\\
 &+J\sum_e\cosh(2d_e)+\sum_eh_e(d_e,y)+S_K(U).
\end{split}
\label{eq:fixed-spacing-action}
\end{equation}
Here $\kappa_g,\lambda_H>0$, $J\ge0$, $\mu\in\R$, and
$0\le\chi_e\le\chi_*$.  The compact gauge action $S_K$ is continuous and
finite range.  The reference is Lebesgue measure for $s,y$ and normalized
Haar measure for the links.  At fixed volume the quartic and quadratic
onsite terms, together with nonnegative hoppings, give a finite, strictly
positive partition function.

\subsection{The common mode and effective conformal free energy}
Let $n=|G|$ and write $s=t\one+r$, with $\sum_xr_x=0$.  All edge terms
are independent of $t$.  Consequently, under the joint normalized law,
\begin{equation}
 t\sim N\!\left(-\frac{b_g}{2\kappa_g},\frac1{2n\kappa_g}\right),
 \qquad t\text{ is independent of }(r,y,U).
 \label{eq:common-conformal-mode}
\end{equation}
The conditional quotient potential is
\begin{equation}
 \Phi_y(r)=\kappa_g\norm r^2+J\sum_e\cosh(2d_e)
 +\sum_{e=(p,c)}\chi_e
       \bigl(e^{d_e}|y_p|^2+e^{-d_e}|y_c|^2\bigr).
 \label{eq:conditional-conformal-potential}
\end{equation}
The omitted cross term is independent of $r$.  The Hessian of
$\Phi_y$ on the mean-zero subspace is at least $2\kappa_gI$.
\src{R17}{thm:v17v32-factor,thm:v17v32-hessian}

Let $Z_g(y)$ integrate \eqref{eq:conditional-conformal-potential} over
that subspace and define $F_g(y)=-\log[Z_g(y)/Z_g(0)]$.  If
$h_x^\chi=\sum_{e\ni x}\chi_e$, then
\begin{equation}
 0\le F_g(y)\le e^{1/(2\kappa_g)}
                         \sum_xh_x^\chi|y_x|^2.
 \label{eq:conformal-free-energy-bound}
\end{equation}
Indeed, the added diagonal hopping is nonnegative.  Jensen's inequality
bounds its free energy by its expectation in the even, strongly
log-concave $y=0$ law, for which
$\E_0e^{\pm(Br)_e}\le e^{1/(2\kappa_g)}$.
After restoring the cross term, the effective Higgs action obeys
\begin{equation}
 S_{H,\mathrm{eff}}(y,U)
 \ge\frac{\lambda_H}{2}\sum_x|y_x|^4
       -\frac{n(D\chi_*-\mu)_+^2}{2\lambda_H},
 \label{eq:effective-higgs-stability}
\end{equation}
where $D$ bounds the graph degree and a gauge-only additive term is
suppressed.  This is a stability bound, not yet a volume-uniform
normalization estimate.  \src{R17}{thm:v17v39-free,thm:v17v39-coercive}

\subsection{Volume-uniform local moments}
\begin{theorem}[Uniform Higgs moment and homogeneous conformal tightness]
\label{thm:fixed-spacing-moments}
Fix the couplings in \eqref{eq:fixed-spacing-action} and a degree bound
$D$.  For free or periodic Higgs boundary conditions, and
$0<\theta<\lambda_H/2$, there is a constant $L_H<\infty$, independent
of volume, such that
\begin{equation}
 \sup_{G,x}\log\E_G e^{\theta|y_x|^4}\le L_H.
 \label{eq:uniform-local-higgs}
\end{equation}
The estimate is uniform in the conformal and compact-gauge conditioning.
On homogeneous vertex-transitive periodic quotients, set
$M_{H,2}=(L_H/\theta)^{1/2}$.  Then
\begin{equation}
 \sup_{G,x}\E_G r_x^2
 \le\frac1{2\kappa_g}
 +\frac{D\chi_*e^{1/(2\kappa_g)}}{\kappa_g}M_{H,2}.
 \label{eq:uniform-local-conformal}
\end{equation}
Together with \eqref{eq:common-conformal-mode}, this gives a
volume-uniform second moment for the full coordinate $s_x$.
\end{theorem}
\begin{proof}
The conditional Higgs density at a site is a quartic radial density with
quadratic coefficient at least $\mu$ and a linear source $L_x$.  Increasing
the quadratic coefficient decreases the radial exponential moment.  In
the hopping square, the cross term is independent of $d$, and hence
$|L_x|\le2\chi_*\sum_{z\sim x}|y_z|$ uniformly in $s$ and $U$.
The one-site quartic integral ratio therefore gives constants $C_0,C_L$
depending only on the displayed fixed parameters such that
\[
 \log\E_x e^{\theta|y_x|^4}
 \le C_0+C_L|L_x|^{4/3}
 \le C_0+K_H\sum_{z\sim x}|y_z|^{4/3},
 \quad K_H=C_L(2\chi_*)^{4/3}D^{1/3}.
\]
Use $u^{4/3}\le\zeta u^4+2/(3\sqrt{3\zeta})$.  Writing
$C_\zeta=C_0+2K_HD/(3\sqrt{3\zeta})$, choose $\zeta>0$ so that
$\beta_H=DK_H\zeta/\theta<1$.  Conditional expectation and H\"older's
inequality give, for the maximum $M$ of the finite-volume site log
moments,
$M\le C_\zeta+\beta_HM$.  These moments are finite at fixed volume
by quartic confinement.  Thus one may take
$L_H=C_\zeta/(1-\beta_H)$.

For the quotient conformal law, write $\nu_y\propto e^{-T_y}\nu_0$.
Its relative entropy obeys
$\operatorname{Ent}(\nu_y\mid\nu_0)=F_g(y)-\E_yT_y\le F_g(y)$.
The quadratic transport inequality of the $2\kappa_g$-convex reference
and conditional covariance control give
\[
 \E_y\norm r^2\le\frac{n-1}{2\kappa_g}
                              +\frac{F_g(y)}{\kappa_g}.
\]
Average over $y,U$, apply \eqref{eq:conformal-free-energy-bound} and
\eqref{eq:uniform-local-higgs}, and divide by $n$.  Homogeneity identifies
this average with each site's moment, proving
\eqref{eq:uniform-local-conformal}.
\end{proof}
\src{R17}{thm:v17v41-ratio,thm:v17v41-row,thm:v17v41-uniform,thm:v17v42-conditional,cor:v17v42-local}

On an arbitrary nonhomogeneous graph the last argument gives a
site-averaged bound, not a bound at every exceptional vertex.  Arbitrary
fixed Higgs boundary data also require their own boundary moment control.
The normalizable common mode in \eqref{eq:common-conformal-mode} must be
restored when passing to the full-field specification; the mean-zero
quotient alone is not the unconstrained Gibbs field.

\begin{theorem}[Fixed-spacing thermodynamic existence]
\label{thm:fixed-spacing-dlr}
Let the infinite carrier be a countable locally finite transitive lattice
admitting homogeneous periodic quotients whose injectivity radii tend to
infinity.  Use the fixed finite-range couplings of
\eqref{eq:fixed-spacing-action}.  The normalized periodic laws have a
subsequence converging on every finite coordinate set to a
translation-invariant probability law $\mu_\infty$.  It obeys the local
Gibbs, or Dobrushin--Lanford--Ruelle, specification: for a finite region
$\Lambda$, its conditional kernel $\gamma_\Lambda$, bounded local $f$,
and bounded exterior-measurable $g$,
\begin{equation}
 \int gf\,\dd\mu_\infty
 =\int g\,\gamma_\Lambda f\,\dd\mu_\infty.
 \label{eq:fixed-spacing-dlr}
\end{equation}
In particular, this fixed-spacing joint branch is nonempty.
\end{theorem}
\begin{proof}
The moment bounds make every finite collection of $s,y$ coordinates
tight, and the compact link coordinates need no additional tail bound.
A diagonal subsequence gives consistent finite-dimensional limits and
therefore a law on the countable product space.  Translation symmetry
passes to the limit.  On bounded boundary sets the local kernels are
continuous by domination with the onsite Gaussian and quartic tails;
nonnegative hopping terms preserve that domination.  The local kernels
are consequently Feller.  For sufficiently large quotients, a fixed
region and its interaction collar have the infinite-lattice incidence.
The finite Gibbs identity then passes to the limit for continuous local
tests.  Approximation and a monotone-class argument extend it to the
stated bounded tests, giving \eqref{eq:fixed-spacing-dlr}.
\end{proof}
\src{R17}{lem:v17v43-feller,thm:v17v43-limit,thm:v17v43-dlr,cor:v17v43-nonempty}

\scope{Theorem~\ref{thm:fixed-spacing-dlr} removes volume at fixed spacing.
It does not prove uniqueness, clustering, reflection positivity, a chiral
fermion extension, or an Einstein variational equation.  Its constants
need not survive \eqref{eq:canonical-scaling}; in particular,
$e^{1/(2\kappa_g)}$ deteriorates as $\kappa_g\downarrow0$.
The theorem supplies a positive joint lattice law, not an ultraviolet
construction of the physical Standard Model--Einstein endpoint.}


\section{Fermions and local normality}
Fermionic coefficients are controlled as graded source forms or through positive states on the physical CAR/local operator algebra. They are not treated as a second ordinary random-field measure. Weak-star compactness of all algebraic states may produce a singular limit; a physical locally normal construction needs compactness in the local predual.

\begin{theorem}[Local predual projective limit]
\label{thm:predual}
Let the regional physical algebras have normal unital inclusion maps and the corresponding positive contractive preadjoint restrictions. Suppose the local state families are predual compact and their compatibility defects vanish on each fixed regional inclusion. Then every diagonal accumulation family is locally normal and compatible. If the common determining tests and summable comparisons make the limit unique, the whole family converges to this locally normal projective state.
\end{theorem}
\begin{proof}
Extract successively on a countable cofinal regional family in the compact local preduals. Contractivity and vanishing compatibility defects pass the restriction identities to the limit. Normality is retained because convergence occurs in the predual, not merely among all algebraic functionals. A determining family and a Cauchy estimate remove subsequence dependence.
\end{proof}
\src{R12}{prop:v12v91-restriction,thm:v12v91-projective}

\begin{proposition}[Energy compactness via a positive finite-rank channel]
\label{prop:energy-compact}
Let a local positive reference Hamiltonian have compact resolvent, and suppose $\Tr(\rho H_0)\le C$ for the family of density matrices. Its finite-energy spectral projections $P_M$ give a uniformly vanishing tail $\Tr\rho(I-P_M)\le C/M$. Replacing the discarded tail by a fixed finite-rank state gives completely positive trace-preserving finite-dimensional approximations converging uniformly in trace norm. Hence the local family is predual compact.
\end{proposition}
\begin{proof}
The energy bound proves the tail estimate. Positivity bounds the off-diagonal block by the square root of the discarded mass, so compression plus positive replacement differs from $\rho$ by a quantity tending uniformly to zero. Each approximating state set is finite-dimensional and compact.
\end{proof}
\src{R12}{prop:v12v92-finite-rank,cor:v12v92-energy}
A finite trace in a diffuse algebra is not the same as compact resolvent and does not prove this approximation. In infinite volume the required condition is regional energy control; a uniform bound on the total energy is generally not the appropriate hypothesis.

\section{Joint geometry--matter states}
On a compact metric screen $K$ and a local quantum Hilbert space, a state on $C(K)\otimes\widetilde{\mathcal K(\HH)}$ can contain mass escaping into the unitization boundary. Uniform quantum compact containment rules this out. The resulting joint limit disintegrates as
\begin{equation}
 \omega(F)=\int_K\Tr\bigl(\rho_gF(g)\bigr)\,\nu(\dd g),
 \qquad \rho_g\ge0,\quad\Tr\rho_g=1\quad\nu\text{-a.e.}
 \label{eq:joint-state}
\end{equation}
The geometry law and the conditional matter densities are parts of one state. This is stronger than independently extracting a metric law and a matter state whose correlations have been lost. \src{R12}{thm:v12v81-state-space,thm:v12v82-disintegration}

Countable topological sectors require two controls: a uniformly small tail outside finitely many sectors, and compactness within each retained sector. An action penalty exceeding the sector-count entropy is one sufficient tail mechanism on the declared branch. A gap can keep an index constant along a homotopy within a sector; it does not establish the sector tail. Likewise, a local gauge slice with a Faddeev--Popov floor does not by itself control the full quotient or its singular stabilizers. \cite{R11,R12}

\section{BRST-compatible compactness and zero modes}
A compactness projection must be compatible with the physical differential. For a closed nilpotent $Q$, the Hodge operator $\Delta_Q=Q^*Q+QQ^*$ has a positive-sector spectral problem and a harmonic sector. Projection onto a common reducing spectral core and replacement by a physical harmonic state gives an exact BRST-compatible compactness channel under the source's domain hypotheses. If the core does not reduce $Q$, the channel has an explicit commutator defect, which must be bounded in graph norm rather than declared zero. \src{R12}{prop:v12v96-channel,thm:v12v97-exact,thm:v12v98-integrated}

Harmonic state modes and harmonic anomaly cochains are different objects. In a finite anomaly complex, a Hodge decomposition isolates the cohomological obstruction. A small harmonic representative is forced to vanish only if it lies in the actual integral lattice and its norm is below that lattice's systole; torsion requires its own treatment. Neither a small real coefficient nor an approximate Ward identity alone proves integral anomaly cancellation. \src{R12}{thm:v12v19-Hodge}

\section{Adjacent cutoffs on a common space}
For an actual interpolation $S_t=(1-t)S_n+t\widetilde S_{n+1}$ on a common lifted space, let $\Delta S_n=\widetilde S_{n+1}-S_n$. Equation~\eqref{eq:normalized-response} gives
\begin{equation}
 \E_{n+1}\widetilde A-\E_nA
 =\text{observable-transport term}
 -\int_0^1\Cov_t(A_t,\Delta S_n)\,\dd t.
 \label{eq:cutoff-comparison}
\end{equation}
The block covariance estimate controls the second term by a localized weighted block-gradient norm of the full action mismatch. The full mismatch includes action, determinant/phase, coarea, ambient measure transport, and counterterms. Constants disappear from normalized expectations, but metric-dependent constants do not disappear from metric responses. \cite{R8,R16}

\begin{theorem}[Summable local comparison]
\label{thm:summable}
For each member $A$ of a countable determining physical algebra, suppose the common lifts and the complete interpolation satisfy
\begin{equation}
 |\omega_{n+1}(\mathcal R_{n+1,n}A)-\omega_n(A)|
 \le\varepsilon_n(A),\qquad
 \sum_n\varepsilon_n(A)<\infty.
 \label{eq:summable}
\end{equation}
Then its cylinder expectations are Cauchy. With joint tightness and the required moment uniform integrability, these limits define the same physical state obtained by compactness. If the independently retained direct comparisons satisfy a common vanishing tail, the state is independent of the admissible cofinal route.
\end{theorem}
\begin{proof}
Telescope adjacent differences. The tail sum bounds any two sufficiently fine values. Tightness identifies the resulting functional with an actual limit in the declared state space. Direct-comparison consistency identifies different paths rather than assuming that adjacent telescoping alone does so.
\end{proof}
\cite{R8,R12,R16,R18}

A weighted master ledger can include all local tests and their first or higher source derivatives. Summability in this stronger ledger gives a differentiable state/source map. A finite test bank with small errors does not establish the countable determining or derivative statement by itself.

\section{Commuting ultraviolet and volume limits}
\begin{criterion}[Two-regulator route independence]
\label{crit:double-limit}
Suppose a local expectation $a_{n,L}^{\tau}$ has ultraviolet increments bounded uniformly in volume and boundary condition by a summable $\varepsilon_n$. Suppose boundary-condition differences are bounded uniformly in $n$ by $\delta_L\to0$, and one reference boundary family has its thermodynamic limit. Then the corresponding ultraviolet and volume limits agree whenever taken along the admitted cofinal routes, with errors controlled by the ultraviolet tail and $\delta_L$.
\end{criterion}
\begin{proof}
Compare two route points through the reference boundary family. The ultraviolet portion is bounded by the tail sum; the boundary portion by $\delta_L$. Let the two cutoffs tend to their endpoints in either order.
\end{proof}
\cite{R16}
The uniformity is the content of the criterion. A finite-volume convergence proof with constants growing in $L$ cannot be substituted for it.


% ==================== 10_reconstruction.tex ====================
\chapter{Reflection, boundary nuclearity, and the physical clock}
\label{ch:reconstruction}
\section{Closed reflected Gram forms}
Let $\AA_+$ be the declared positive-time physical algebra, with reflection $\Theta$. The finite reflection form is
\begin{equation}
 \mathcal Q_n(F,G)=\omega_n((\Theta F)G).
\end{equation}

\begin{theorem}[Reflection positivity survives convergent panels]
\label{thm:rp-limit}
If every finite reflected Gram matrix is positive semidefinite at cutoff $n$ and its entries converge on a common determining algebra, the limiting reflected Gram matrices are positive semidefinite. With the stated growth bounds, the limiting form extends to the required physical test space.
\end{theorem}
\begin{proof}
For a finite coefficient vector $c$, the quadratic form $c^*G_nc$ is nonnegative and converges to $c^*Gc$. Thus $G\succeq0$. The growth bound gives the continuous extension.
\end{proof}
\src{R16}{thm:v16v87-rp-closure}
This theorem transfers positivity; it does not prove it for an interacting chiral gravitational cutoff law that has not yet been constructed.

\section{Bosonic, gauge, and fermionic reflection mechanisms}
Positive-type compact gauge weights admit a Peter--Weyl or convolution square. An autocorrelation weight may be supported inside an admissibility neighbourhood while retaining this pure-gauge reflection mechanism. That provides an explicit finite gauge module, not full coupled reflection positivity or the desired universality class. \cite{R6,R19}

A Higgs hinge with the correctly matched side variables has a positive Fock-Gram expansion of its cross term. If the metric weights on the two sides are not reflected partners, the same argument need not apply. In a finite Grassmann hinge, positivity is governed by the cross-boundary coefficient matrix: its positive exterior powers generate the reflected Grassmann form under the declared ordering convention. A real determinant, or even a positive determinant, does not replace this matrix criterion. \cite{R19}

Reflection of the determinant line adds another layer. At a fixed locus, the real line can carry a sign obstruction measured by the relevant mod-two spectral-flow/holonomy data. Multiplet cancellation can remove that sign while leaving a negative cross-boundary Gram matrix. Conversely, positivity of one matrix does not glue the phase on all gauge sectors. The finite hybrid transfer construction therefore joins the gauge, Higgs, fermionic, and line conditions rather than using one as a surrogate for the others. \cite{R10,R18,R19}

\begin{proposition}[Exterior-power and hybrid reflected positivity]
\label{prop:hybrid-rp}
Fix the reflected Grassmann ordering and reflection phases.  Suppose its
crossing exponential has the minor expansion
\begin{equation}
 e^{X_C}=\sum_{k,I,J}\det(C_{I,J})\,
                         \Theta(\psi_I)\psi_J,
 \label{eq:grassmann-minor-gram}
\end{equation}
and that the reference Grassmann functional is positive on the associated
reflected-square cone.  If $C\succeq0$, the form induced by $e^{X_C}$ is
positive on the even half-algebra.  Multiplication by
$\Theta(B)B$, with $B$ an even one-side factor, preserves this property.

Let $\HH_{F,C}$ be the finite coefficient space after quotienting its
null form and let $T_B\succeq0$ be the bounded bosonic crossing operator.
For every square-integrable coefficient amplitude $\Psi$,
\begin{equation}
 \pair{\Psi}{(T_B\otimes I_{\HH_{F,C}})\Psi}\ge0.
 \label{eq:hybrid-rp}
\end{equation}
Reflection-matched one-side factors and simultaneous boundary gauge
projection preserve this inequality by conjugation and compression.
\end{proposition}
\begin{proof}
The coefficient matrix of the minors of order $k$ is $\bigwedge^kC$.
For $C=S^*S$ it equals
$(\bigwedge^kS)^*(\bigwedge^kS)$, so
\eqref{eq:grassmann-minor-gram} is a sum of reflected squares.  An even
$B$ introduces no additional ordering sign.  For the hybrid statement,
expand $\Psi$ in an orthonormal basis of $\HH_{F,C}$; the quadratic form
is the sum of the nonnegative forms $\pair{\psi_j}{T_B\psi_j}$.
Finally, $M^*TM\succeq0$ and $P^*TP\succeq0$ whenever $T\succeq0$.
\end{proof}
\src{R19}{thm:v19v7-gram,thm:v19v16-hybrid,thm:v19v16-coupled}

The proposition applies to a finite coupled slab when the gauge crossing
has its positive heat-kernel transfer, the Higgs crossings are balanced
or occur in the specified reflected two-step pairing, and the complete
fermion cross matrix is positive after transport and phase identification.
Remaining interactions must occur as reflection-matched one-side factors,
with integrable bosonic weights.  A direct integral over reflection-fixed
parameters preserves positivity when these hypotheses hold almost
everywhere.  A parameter-dependent fermion matrix requires this
conditional decomposition or a genuinely positive operator-valued
kernel; pointwise positivity of unrelated fibres alone is not a proof
for an arbitrary mixed kernel.  These are finite factorization
conditions, not a consequence of determinant-line triviality.


\section{Osterwalder--Schrader reconstruction}
The physical reconstruction starts by quotienting the reflected null space and completing the positive form. Euclidean time translations induce a contraction semigroup only when they are compatible with that quotient and satisfy the required continuity. Its generator is then a nonnegative self-adjoint Hamiltonian. The local fields require the additional all-order growth, domain, and distributional continuity conditions; reflection positivity alone is not the entire reconstruction theorem. \src{R14}{thm:v14v88-time}

A rational or dense subgroup of Euclidean symmetries is insufficient without a uniform modulus at the identity. Such a modulus permits passage from the dense subgroup to the full group after the limit. Without it, cutoff-dependent rapidly oscillating labels can agree on a selected dense test pattern while failing the desired limiting continuity. \cite{R13,R16}

For a local net, uniform nondegenerate ADM bounds give a common causal cone. Rational-region locality can then be transported to arbitrary separated regions under continuity. Time-slice density, affiliation of field resolvents, and common graph domains remain distinct requirements. Strong commutativity of spacelike affiliated operators is not a consequence of formal commutation on an unspecified core. \cite{R13,R14}

\section{A positive channel gap and a physical spectral gap}
\begin{theorem}[Euclidean decay implies a channel threshold]
\label{thm:channel-gap}
Suppose a reconstructed centered vector $\psi$ has positive spectral measure $\nu_\psi$ for $H\ge0$ and
\begin{equation}
 \pair{\psi}{e^{-tH}\psi}
 =\int_0^\infty e^{-t\lambda}\,\nu_\psi(\dd\lambda)
 \le C_\psi e^{-mt}\qquad(t\ge0).
\end{equation}
Then $\nu_\psi([0,m))=0$. If the same $m>0$ holds on a dense centered physical family and the vacuum is unique, then $H$ has gap at least $m$ on the orthogonal complement of the vacuum.
\end{theorem}
\begin{proof}
Positive mass on $[0,m-\epsilon]$ would give a lower bound proportional to $e^{-(m-\epsilon)t}$, contradicting the upper bound as $t\to\infty$. Density and spectral-projection continuity give the second conclusion.
\end{proof}
\src{R16}{thm:v16v86-channel,thm:v16v86-dense-gap}
For the complete Standard Model--Einstein endpoint, the last clause is used only for a sector for which a positive gap is actually intended. Photon and gravitational massless channels require their own infrared statement.

\section{Refresh is a construction device, not a mass mechanism}
A finite law can be equipped with a stationary reversible Markov--renewal process by adding refresh to a local generator:
\begin{equation}
 L=\rho(R-I)+L_{\rm loc}.
\end{equation}
The refresh term produces a gap $\rho$ on centered observables and, with the retained cut, the conditional square of Theorem~\ref{thm:half-square}. This constructs a nonempty finite reflected stationary slab, including mixed couplings under the stated integrability hypotheses. \src{GT}{thm:RPESM-reflected-slab}

\begin{proposition}[Refresh cannot certify the refresh-free gap]
\label{prop:refresh}
In the declared reversible setting where refresh projects onto the stationary sector, adding refresh shifts the centered spectrum by $\rho$. A comparison with the actual refresh-free physical generator contains the same shift as a comparison cost. Removing refresh preserves a positive gap only if the refresh-free part already has the corresponding independently controlled coercivity.
\end{proposition}
\src{R14}{thm:v14v94-shift,cor:v14v94-remove,thm:v14v94-correspondence}
This also explains why refresh can mask sector zero modes. It is not a physical argument for confinement, particle masses, or graviton disappearance.

\section{Refinement must not annihilate the new clock modes}
For an isometric conditional lift $J$ of old observables, the exact lifted transfer
\begin{equation}
 P_{n+1}=JP_nJ^*
 \label{eq:annihilating-transfer}
\end{equation}
preserves old dynamics and annihilates the orthogonal detail subspace. Its spectrum is the old spectrum together with zero on those details.

\begin{obstruction}[Annihilated details have no finite physical energy]
\label{obs:annihilation}
A transfer of the form~\eqref{eq:annihilating-transfer} with nontrivial detail kernel cannot equal $e^{-aH}$ for a finite self-adjoint Hamiltonian on those detail modes. A scalar laziness correction does not in general produce one cutoff clock on which both old and detail energies have the required finite nonzero limits.
\end{obstruction}
\begin{proof}
For self-adjoint $H$ bounded below, the spectral multiplier $e^{-a\lambda}$ is strictly positive at every finite spectral value, so $e^{-aH}$ has trivial kernel. The zero detail eigenvalue would represent infinite energy, not an ordinary retained physical mode. The laziness issue follows by comparing the distinct old and detail eigenvalue scalings.
\end{proof}
\cite{R15}

\subsection{A faithful completion of the old clock}
The alternative to detail erasure requires additional geometric and
dynamical data.  Suppose the conditional law over a coarse configuration
$x$ is $K_x$ and there is a measurable family of probability-space
isomorphisms $\phi_x:(U,\nu)\to(Y_x,K_x)$, defined modulo null sets.  The
induced unitary
\[
 (\mathcal WF)(x,u)=F(x,\phi_x(u))
\]
identifies the fine space with $L^2(\mu_n)\widehat\otimes L^2(\nu)$ and
sends the coarse lift $Jf$ to $f\otimes1$.  Disintegration alone does not
select $\phi_x$, its reflection covariance, or its spatial locality.

\begin{theorem}[Faithful old and detail evolution]
\label{thm:faithful-clock}
Let $H_n\ge0$ and $B\ge0$ be self-adjoint Markov generators on the two
factors, with constants as their simple kernels and gap lower bounds
$\gamma_n>0$ and $\beta>0$.  Define
\begin{equation}
 X=\mathcal W^{-1}(H_n\otimes I)\mathcal W,
 \qquad Y=\mathcal W^{-1}(I\otimes B)\mathcal W,
 \qquad H_{n+1}^{\rm prod}=X+Y.
 \label{eq:faithful-product-clock}
\end{equation}
For every clock step $\tau>0$, its transfer is injective and obeys
\begin{equation}
 e^{-\tau H_{n+1}^{\rm prod}}J=Je^{-\tau H_n},
 \qquad
 H_{n+1}^{\rm prod}\ge\min\{\gamma_n,\beta\}(I-E_0),
 \label{eq:faithful-intertwining}
\end{equation}
where $E_0$ projects onto constants.  If $\gamma_n$ and $\beta$ are the
actual spectral gaps, the gap in \eqref{eq:faithful-intertwining} is
exactly their minimum.
\end{theorem}
\begin{proof}
The two summands strongly commute after conjugation by $\mathcal W$.
Their heat operator is the conjugate of
$e^{-\tau H_n}\otimes e^{-\tau B}$, which fixes $f\otimes1$ according to
the old transfer and has no kernel.  Its energies are sums of the two
factor energies.  The only joint zero mode is $1\otimes1$, and every
other spectral sector costs at least the indicated minimum.
\end{proof}
\src{R14}{thm:v14v97-faithful-lift,cor:v14v97-commuting}

This construction retains the old clock and supplies a finite clock to
the new detail modes.  It does not identify the product generator with the
physical local dynamics.  A generator local in the reference variable
$u$ can become nonlocal after $y=\phi_x(u)$; that locality question and
compatibility with physical constraints remain part of the construction.

\begin{proposition}[Noncommuting completion and its clock cost]
\label{prop:split-clock}
Let $X,Y\ge0$ be self-adjoint generators with kernel projections $E_X,E_Y$
and common kernel $\Ran E_0$.  Suppose
\[
 X\ge\gamma(I-E_X),\qquad Y\ge\beta(I-E_Y),\qquad
 c=\norm{(E_X-E_0)(E_Y-E_0)}<1.
\]
On the common form domain,
\begin{equation}
 X+Y\ge\Gamma(I-E_0),\qquad
 \Gamma=\min\{\gamma,\beta\}(1-c).
 \label{eq:clock-angle-floor}
\end{equation}
For bounded finite-regulator generators put
$C=\norm{[X,Y]}$ and $S_\tau=e^{-\tau Y/2}e^{-\tau X}e^{-\tau Y/2}$.
Then
\begin{equation}
 \norm{S_\tau-e^{-\tau(X+Y)}}\le\tau^2C.
 \label{eq:split-clock-error}
\end{equation}
If $q_\tau=e^{-\tau\Gamma}+\tau^2C<1$, the logarithmic transfer
Hamiltonian has the centered gap lower bound
\begin{equation}
 \operatorname{gap}\bigl(-\tau^{-1}\log S_\tau\bigr)
 \ge-\tau^{-1}\log q_\tau.
 \label{eq:split-clock-gap}
\end{equation}
\end{proposition}
\begin{proof}
On the centered space put $P=E_X-E_0$ and $Q=E_Y-E_0$.
The two-projection estimate $\norm{P+Q}\le1+\norm{PQ}=1+c$ implies
$2I-P-Q\ge(1-c)I$, which gives \eqref{eq:clock-angle-floor}.
A double Duhamel integral bounds the interchange of a half $Y$ step and
a full $X$ step by $\tau^2C/2$.  The same integral bounds the error of
the Lie product against $e^{-\tau(X+Y)}$ by $\tau^2C/2$.
Their sum proves \eqref{eq:split-clock-error}.  The centered norm of the
positive operator $S_\tau$ is therefore at most $q_\tau$; spectral
calculus proves \eqref{eq:split-clock-gap}.
\end{proof}
\src{R14}{lem:v14v97-angle,thm:v14v97-split}

A uniform positive $\Gamma$ and $\tau\norm{[X,Y]}\to0$ preserve this
massive-channel lower bound in physical clock units.  These estimates do
not impose a gap on the intended massless sectors.  Unbounded continuum
generators require the corresponding domain and commutator estimates,
not a formal substitution into the bounded-operator inequality.

\subsection{What conditional transport actually supplies}
On a Euclidean shell let
$K_x(\dd y)=Z(x)^{-1}e^{-V(x,y)}\dd y$, with $V\in C^3$ and
$\nabla_y^2V\succeq m_*I$, $m_*>0$.  Along a smooth parameter path $x_t$
set $G_t=\dot x_t\cdot\nabla_xV$ and
$\mathcal L_{x_t}=\Delta_y-\nabla_yV\cdot\nabla_y$.  The mean-zero solution
of
\begin{equation}
 -\mathcal L_{x_t}\psi_t=-G_t+\E_{K_{x_t}}G_t
 \label{eq:conditional-poisson-transport}
\end{equation}
supplies a probability-transport velocity $v_t=\nabla_y\psi_t$ satisfying
\begin{equation}
 \partial_t p_{x_t}+\operatorname{div}(p_{x_t}v_t)=0,
 \qquad
 \norm{v_t}_{L^2(K_{x_t})}\le
        \frac{\norm{\nabla_yG_t}_\infty}{m_*}.
 \label{eq:conditional-transport-energy}
\end{equation}
Indeed, the conditional Poincar\'e inequality first bounds
$\norm{\nabla\psi_t}_2$ by $m_*^{-1/2}\sqrt{\Var(G_t)}$ and then bounds
$\Var(G_t)$ by $m_*^{-1}\norm{\nabla_yG_t}_\infty^2$.
Differentiating the normalized density proves the continuity equation.
When the velocity has the regularity needed for an invertible flow, this
constructs the fibre maps.  A change $\delta$ of one old coordinate with
$\norm{\nabla_y\partial_{x_i}V}_\infty\le L_i$ consequently gives
\[
 W_2(K_x,K_{x+\delta e_i})\le\frac{L_i}{m_*}|\delta|.
\]
Reflection-compatible potentials, base points, and paths give reflected
flows by uniqueness of the mean-zero Poisson solution.
\src{R14}{thm:v14v98-moser,cor:v14v98-influence,cor:v14v98-reflection}

The estimate is finite influence, not spatial quasilocality: the Poisson
inverse can spread a localized parameter derivative.  Nor do base-point
flows automatically compose along different cutoff routes.  Their
quasilocal tails and composition holonomy must be controlled separately.
A quartic double well is coercive without being uniformly convex, and a
compact boundaryless gauge manifold cannot admit a globally strictly
convex potential: integrating its strictly positive Laplacian would give
a contradiction.  Thus this sufficient transport mechanism is confined
to its declared convex branch; it is not a global trivialization of the
chiral gauge--Higgs--gravity shell.
\src{R14}{prop:v14v98-quartic,prop:v14v98-compact}

\section{Critical clock factors must be counted before inversion}
For a short-time channel $T_\tau=e^{\tau L}$,
\begin{equation}
 \partial_aT_\tau
 =\int_0^\tau e^{(\tau-s)L}(\partial_aL)e^{sL}\,\dd s.
 \label{eq:duhamel-channel}
\end{equation}
Thus a contractive channel norm gives $\norm{\partial_aT_\tau}\le\tau\norm{\partial_aL}$. Mixed derivatives carry either one factor $\tau$ multiplying $L_{ab}$ or two factors multiplying the first derivatives. If a fixed-point inverse behaves like $(1-q_\tau)^{-1}\sim(m\tau)^{-1}$, these explicit clock factors can cancel the apparent singularity. They must not be dropped before the inverse estimate. \cite{R7,R8}

Conversely, a critical conserved owner can force the absolute one-site influence to be at least one. Replacing the physical clock by a normalized unit-total update may introduce a slow mode whose gap vanishes with volume or refinement. Exact clock scaling and local-update normalization are therefore part of the continuum data, not a harmless convention.

\section{Boundary nuclearity and interacting heat traces}
For positive one-particle energies $\epsilon_j$, the finite or convergent Fock heat traces are
\begin{equation}
 Z_F(\beta)=\prod_j(1+e^{-\beta\epsilon_j}),\qquad
 Z_B(\beta)=\prod_j(1-e^{-\beta\epsilon_j})^{-1}.
 \label{eq:fock-traces}
\end{equation}
Fermionic finiteness is governed by the one-particle heat sum. Bosonic finiteness also requires the specified infrared control and absence of an uncontrolled zero mode. These are physical reduced spaces; ghost Fock factors are not automatically positive physical Hilbert spaces. \src{R15}{thm:v15v89-fermion,thm:v15v89-boson,cor:v15v89-heat}

\subsection{A stable Higgs reference before the interaction estimate}
A quartic field interaction should not be treated as a small quadratic
perturbation merely because each regulator is finite.  An alternative is
to include a stable quartic in the reference form.  The fixed-cutoff
counterterm algebra can be made explicit without identifying that choice
with the physical renormalization trajectory.

\begin{proposition}[Radial Wick separation and stable reference]
\label{prop:stable-wick-reference}
For $r$ real Gaussian coordinates with contraction
$\E(X_aX_b)=C\delta_{ab}$,
\begin{align}
 :|X|^2:_C&=|X|^2-rC,\\
 :|X|^4:_C&=|X|^4-2(r+2)C|X|^2+r(r+2)C^2.
 \label{eq:radial-wick-quartic}
\end{align}
For $\lambda>0$ and $m_R^2\ge0$, set
\begin{equation}
 \nu_C=2\lambda(r+2)C+m_R^2,\qquad
 E_C=\lambda r(r+2)C^2+rm_R^2C.
 \label{eq:stable-wick-counterterms}
\end{equation}
Then the exact polynomial identity
\begin{equation}
 \lambda:|X|^4:_C+\nu_C:|X|^2:_C+E_C
 =\lambda|X|^4+m_R^2|X|^2\ge0
 \label{eq:stable-wick-identity}
\end{equation}
holds.  Let $H_{G,n}\ge0$ be a Gaussian reference with compact resolvent
and finite heat trace at positive times, and let
\[
 U_n=\sum_\ell w_\ell
       \bigl(\lambda|X_\ell|^4+m_R^2|X_\ell|^2\bigr),\qquad w_\ell>0.
\]
On the dense intersection of the two form domains,
$H_{{\rm ref},n}=H_{G,n}\dotplus U_n$ is a closed nonnegative form sum and
\begin{equation}
 \Tr e^{-tH_{{\rm ref},n}}\le\Tr e^{-tH_{G,n}}
 \qquad(t>0).
 \label{eq:stable-reference-trace}
\end{equation}
\end{proposition}
\begin{proof}
Expanding $|X|^4=\sum_{a,b}X_a^2X_b^2$ separates the equal-index Wick
contractions from the unequal-index contractions and gives
\eqref{eq:radial-wick-quartic}.  Substitution gives
\eqref{eq:stable-wick-identity}.  The sum of the two nonnegative closed
forms is closed on their intersection.  Its form norm dominates the
Gaussian one, preserving compact embedding.  The min--max principle
bounds each reference eigenvalue below by the corresponding Gaussian
eigenvalue; summing their heat factors proves
\eqref{eq:stable-reference-trace}.
\end{proof}
\src{R15}{lem:v15v93-wick,thm:v15v93-stable-wick,prop:v15v93-reference}

For a four-real-component Higgs doublet,
$\nu_C=12\lambda C+m_R^2$ and $E_C=24\lambda C^2+4m_R^2C$.
These are identities in the chosen Gaussian reference, not a derivation
of the physical Higgs vacuum or of cutoff-independent counterterms.
They also show why the scalar and quadratic counterterms cannot be
conflated when $C$ grows with the regulator.

Lower polynomial degrees can be absorbed into this reference by the
explicit inequality
\begin{equation}
 x^p\le\delta x^4+
 \left(1-\frac p4\right)
 \left(\frac{p}{4\delta}\right)^{p/(4-p)},
 \qquad 0<p<4,\quad\delta>0,\quad x\ge0.
 \label{eq:quartic-absorption}
\end{equation}
It follows by maximizing $x^p-\delta x^4$.  Applied to the local
coefficients, it yields a finite-regulator bound
$|v_n[\psi]|\le\alpha\pair\psi{U_n\psi}+\beta_n\norm\psi^2$ with any
prescribed $\alpha>0$.  Uniformity still requires uniform control of the
coefficient combinations and weighted volume entering $\beta_n$.
This does not absorb arbitrary derivative, gauge, or Yukawa interactions
without their own form estimates.
\src{R15}{lem:v15v93-young,thm:v15v93-polynomial}

\subsection{Assembling the interacting heat operator}
\begin{theorem}[Conditional interacting boundary transfer]
\label{thm:nuclearity}
Let a finite family of nonnegative sector Hamiltonians $H_{j,n}$ have
common Hilbert-space realizations and trace-norm heat convergence at
every positive time.  Let $H_{0,n}$ be their tensor-sum reference, with
form $q_{0,n}$.  Suppose the symmetric interaction obeys
\begin{equation}
 |v_n[\psi]|\le\alpha q_{0,n}[\psi]+\beta\norm{\psi}^2,
 \qquad 0\le\alpha<1,\quad \beta<\infty,
 \label{eq:relative-form}
\end{equation}
with cutoff-independent constants on the common form domain.  If the
assembled operators $H_n=H_{0,n}\dotplus V_n$ converge in norm resolvent
to $H$ on the stated nonzero common realization, then, for every $t>0$,
\begin{equation}
 e^{-tH_n}\longrightarrow e^{-tH},\qquad
 \frac{e^{-tH_n}}{\Tr e^{-tH_n}}\longrightarrow
 \frac{e^{-tH}}{\Tr e^{-tH}}
 \quad\hbox{in trace norm}.
 \label{eq:assembled-heat-convergence}
\end{equation}
\end{theorem}
\begin{proof}
Finite tensor products preserve trace-norm convergence, giving uniform
reference heat traces at every positive time.  The form inequality
$H_n\ge(1-\alpha)H_{0,n}-\beta$ and min--max comparison imply
\begin{equation}
 \begin{aligned}
 \Tr\bigl(\one_{(R,\infty)}(H_n)e^{-tH_n}\bigr)
 &\le e^{-tR/2}\Tr e^{-tH_n/2}\\
 &\le e^{-tR/2+t\beta/2}
        \Tr e^{-(1-\alpha)tH_{0,n}/2}.
 \end{aligned}
 \label{eq:assembled-heat-tail}
\end{equation}
The right-hand side tends to zero uniformly in $n$ as $R\to\infty$.
Norm-resolvent convergence controls the finite spectral cores at
continuity thresholds; these uniform tails upgrade core convergence to
trace-norm convergence.  The limiting heat trace is finite and positive,
so division by the convergent traces proves the normalized statement.
\end{proof}
\src{R15}{thm:v15v99-assembled}
Every factor must be controlled.  Adding stable matter does not repair a
divergent gravitational heat factor.  Phase convergence and reflection
compatibility remain additional requirements even when all magnitude
estimates are nuclear.

\subsection{Collective coordinates are dynamical degrees of freedom}
\begin{proposition}[Continuous multiplicity and its elliptic resolution]
\label{prop:collective-heat}
Let $(\mathcal Q,\nu)$ be a non-atomic finite measure space with an
infinite-dimensional $L^2$ space.  For a nonzero compact operator $T$,
$I_{L^2(\mathcal Q)}\otimes T$ is not compact.  In particular, finiteness
of the scalar integral $\int_{\mathcal Q}\Tr T\,\dd\nu$ does not make
this decomposable operator trace class.

On a compact regular holonomy quotient $\mathcal Q$ of dimension $r$,
let $0<w_-\le w\le w_+$ with $w\in W^{1,\infty}$, and consider the
closed, densely defined form
\begin{equation}
 \ell[u]=\int_{\mathcal Q}
 \bigl(\pair{\nabla u}{A\nabla u}+V|u|^2\bigr)w\,\dd\vol,
 \qquad a_-I\preceq A\preceq a_+I,\quad V\ge-b.
 \label{eq:collective-form}
\end{equation}
Under the stated regular quotient and form-domain hypotheses, its
operator $L_{\mathcal Q}$ has compact resolvent and
\begin{equation}
 N_{L_{\mathcal Q}}(R)\le C(1+R+b)^{r/2},
 \qquad \Tr e^{-tL_{\mathcal Q}}<\infty\quad(t>0).
 \label{eq:collective-counting}
\end{equation}
If $e^{-tH_{\rm tr}}$ is trace class, so is the heat operator of
$L_{\mathcal Q}\otimes I+I\otimes H_{\rm tr}$, and its trace is the
product of the two factor traces.
\end{proposition}
\begin{proof}
Choose an orthonormal sequence $f_j$ in $L^2(\mathcal Q)$ and a vector
$v$ with $Tv\ne0$.  The images $f_j\otimes Tv$ are orthogonal with
fixed nonzero norm, proving noncompactness.  For
\eqref{eq:collective-form}, the lower ellipticity and weight bounds
control the $H^1$ norm after adding a constant.  Rellich compactness
and min--max comparison with the quotient Laplacian give
\eqref{eq:collective-counting}; exponential summation of the counting
bound gives the heat trace.  The final assertion follows from the
spectral tensor-product identity.
\end{proof}
\src{R15}{thm:v15v98-holonomy,thm:v15v98-direct-integral,prop:v15v98-product}

The continuous holonomy label must therefore be given a kinetic operator,
not merely integrated as an uncounted degeneracy.  At reducible strata,
the density, dimension, or regular geometry may change.  Unless the
minimal operator is essentially self-adjoint, its differential expression
alone does not select the physical heat semigroup.  Interface conditions
must come from the cutoff measure and reflection sewing; choosing an
extension solely because it is available does not establish that
compatibility.  This is why regular-stratum nuclearity is a conditional
input rather than a global solution of the gauge quotient.
\src{R15}{prop:v15v98-interface,thm:v15v98-handoff}

\subsection{Normalized states do not fix the gravitational amplitude}
For a real scalar shift $\mathcal E(g)$,
\begin{equation}
 \begin{aligned}
 e^{-t(H(g)+\mathcal E(g)I)}&=e^{-t\mathcal E(g)}e^{-tH(g)},\\
 \frac{e^{-t(H(g)+\mathcal E(g)I)}}{Z_{H+\mathcal E}(t,g)}
 &=\frac{e^{-tH(g)}}{Z_H(t,g)},\\
 \delta_g[-\log Z_{H+\mathcal E}]&=
     \delta_g[-\log Z_H]+t\,\delta_g\mathcal E.
 \end{aligned}
 \label{eq:scalar-shift-source}
\end{equation}
Thus normalized Gibbs states can agree at every cutoff while the sewing
amplitudes vanish, diverge, or fail to converge.  The metric derivative
still reads the scalar counterterm.  Neither
\eqref{eq:stable-wick-identity} nor heat normalization permits its deletion
from the Einstein source.
\src{R15}{prop:v15v93-vacuum-shift,cnt:v15v93-shift}


% ==================== 11_stress.tex ====================
\part{Interaction, stress, and Einstein closure}
\chapter{Compiled stress and the survival of interaction}
\label{ch:stress}
\section{A positive continuum may still be Gaussian}
Uniform finite-volume mixing is useful for compactness and comparison, but it can also remove the very connected responses that would certify an interacting limit. The issue is quantitative and depends on the normalization of the physical observable.

Let a four-dimensional regulator in a fixed compact volume have $|\Lambda_n|\le C_\Lambda h_n^{-4}$. For centered elementary coordinates $X_{n,z}$, assume
\begin{equation}
 \sup_n\sup_{z_1}\sum_{z_2,\ldots,z_r}
 \abs{\Cum(X_{n,z_1},\ldots,X_{n,z_r})}\le K_r<\infty.
 \label{eq:cumulant-stock}
\end{equation}
The critical pairing on this branch is
\begin{equation}
 Z_n(f)=h_n^2\sum_z f_zX_{n,z}.
 \label{eq:critical-field}
\end{equation}

\begin{theorem}[Four-dimensional cumulant scaling]
\label{thm:gaussianization}
Under~\eqref{eq:cumulant-stock},
\begin{equation}
 \abs{\Cum(Z_n(f_1),\ldots,Z_n(f_r))}
 \le C_\Lambda K_rh_n^{2r-4}\prod_j\norm{f_j}_\infty.
 \label{eq:cumulant-scaling}
\end{equation}
Every fixed connected cumulant of order $r\ge3$ vanishes. If the bounds hold at every fixed order, the required local moments are uniformly controlled, and the covariance forms converge, finite collections of these pairings converge to the centered Gaussian law with that covariance. The corresponding Sobolev trace bound promotes this to distribution-space convergence.
\end{theorem}
\begin{proof}
Expand the cumulant by multilinearity. Sum all but the first site using~\eqref{eq:cumulant-stock}, then use $|\Lambda_n|\le C_\Lambda h_n^{-4}$. This gives~\eqref{eq:cumulant-scaling}. Moment--cumulant relations identify every limiting joint moment with its Gaussian pairing value. The cited construction's moment control and Gaussian moment determinacy identify the finite-dimensional law; the stated trace bound gives tightness in the chosen distribution topology.
\end{proof}
\src{R18}{thm:v18v35-cumulant-scaling,thm:v18v35-gaussianization}

\begin{corollary}[Necessary loss of summability for an elementary vertex]
\label{cor:survival-threshold}
For a bounded nondegenerate amplitude normalization, survival of a nonzero connected elementary response of order $r\ge3$ requires, along a subsequence, a contracted cumulant stock of order at least $h_n^{-(2r-4)}$. Cubic survival therefore requires order $h_n^{-2}$ and quartic survival order $h_n^{-4}$ on this branch.
\end{corollary}
\begin{proof}
Apply~\eqref{eq:cumulant-scaling} with the scale-dependent stock and solve the inequality for that stock when the left side stays above a positive constant.
\end{proof}
\src{R18}{thm:v18v36-survival-threshold}
\scope{This is a no-go for the stated uniformly summable elementary normalization, not for all composite operators or all interacting continuum routes. A protected nonlocal kernel, a renormalized composite, a critical replacement law, or another branch can evade the hypotheses.}

\section{A quadratic composite can be non-Gaussian in a Gaussian theory}
Let $G_n$ be a centered real Gaussian vector with covariance $C_n$ and let $A_n=A_n^*$. On the covariance support, put
\begin{equation}
 B_n=C_n^{1/2}A_nC_n^{1/2},\qquad
 Q_n=\pair{G_n}{A_nG_n}-\tr(A_nC_n).
 \label{eq:quadratic-score}
\end{equation}

\begin{theorem}[Exact quadratic Gaussian cumulants]
\label{thm:quadratic-cumulants}
For every $r\ge2$,
\begin{equation}
 \Cum_r(Q_n)=2^{r-1}(r-1)!\tr(B_n^r).
 \label{eq:quadratic-cumulants}
\end{equation}
For $W_n=Q_n/\sqrt{2\tr(B_n^2)}$, define
\begin{equation}
 r_{\rm eff}(B_n)=\frac{\tr(B_n^2)^2}{\tr(B_n^4)},\qquad
 \rho_n=\frac{\norm{B_n}_{\rm op}}{\sqrt{\tr(B_n^2)}}.
\end{equation}
Then $W_n\Rightarrow N(0,1)$, $\Cum_4(W_n)\to0$, $r_{\rm eff}(B_n)\to\infty$, and $\rho_n\to0$ are equivalent on this second-chaos family.
\end{theorem}
\begin{proof}
Diagonalize $B_n$. For small $t$,
\begin{equation}
 \log\E e^{tQ_n}
 =-t\tr B_n-\frac12\log\det(I-2tB_n)
 =\sum_{r\ge2}\frac{2^{r-1}}r t^r\tr(B_n^r).
\end{equation}
Coefficient comparison proves~\eqref{eq:quadratic-cumulants}. The normalized fourth cumulant is $12/r_{\rm eff}(B_n)$. The largest normalized eigenvalue vanishes exactly when the fourth-power sum vanishes; the resulting triangular sum of centered Gaussian squares obeys the stated Gaussian limit criterion. The converse uses the fixed-chaos moment control recorded in the source.
\end{proof}
\src{R18}{thm:v18v38-quadratic-cumulants,thm:v18v38-spectral-equivalence}
A coherent quadratic eigenmode can therefore survive with a non-Gaussian law even when the elementary field is Gaussian. Such a mode alone is not evidence of a non-quasi-free underlying theory.

\section{Remove only the declared dynamically dependent directions}
The programme calls a source direction a \emph{follower} when it is already generated by the declared dynamically retained variables. The removal is a specified orthogonal short, not an arbitrary projection selected to enhance an apparent interaction.

For quadratic Gaussian scores, let $P_n$ project onto the retained one-particle follower space. The quadratic follower kernels are $P_nAP_n$.

\begin{proposition}[Exact quadratic follower short]
\label{prop:follower}
In the Hilbert--Schmidt kernel metric, the orthogonal projection onto the quadratic follower space is $B\mapsto P_nBP_n$. Thus
\begin{equation}
 B_n^\perp=B_n-P_nB_nP_n,\qquad
 \Var I_2(B_n^\perp)=2\norm{B_n^\perp}_{\rm HS}^2.
 \label{eq:follower-short}
\end{equation}
\end{proposition}
\begin{proof}
For every $A=P_nAP_n$, cyclicity of the finite Hilbert--Schmidt pairing gives $\tr[(B-P_nBP_n)A]=0$. The variance follows from~\eqref{eq:quadratic-cumulants}.
\end{proof}
\src{R18}{thm:v18v39-quadratic-projection}
The projection must be transported with the dynamics. A cutoff-dependent change of follower definition can itself create an artificial residual. Its leakage and comparison errors therefore belong in the common source ledger.

\section{All stress responses come from one partition formula}
Let $S_n(\bm s)$ be the complete action with metric-source parameters and $W_n(\bm s)=-\log Z_n(\bm s)$. For a nonempty index set $B$, write $S_{n,B}=\partial_BS_n(0)$.

\begin{theorem}[Common-action response compiler]
\label{thm:response-compiler}
Whenever the finite differentiated integral is justified,
\begin{equation}
 \partial_1\cdots\partial_rW_n(0)
 =\sum_{\pi\in\mathcal P([r])}(-1)^{|\pi|+1}
 \Cum_n(S_{n,B}:B\in\pi),
 \label{eq:partition-compiler}
\end{equation}
where a one-entry cumulant means expectation. In particular,
\begin{align}
 W_i&=\E S_i,\\
 W_{ij}&=\E S_{ij}-\Cov(S_i,S_j),\\
 W_{ijk}&=\E S_{ijk}-\Cov(S_{ij},S_k)-\Cov(S_{ik},S_j)
 -\Cov(S_{jk},S_i)+\Cum(S_i,S_j,S_k).
 \label{eq:third-response}
\end{align}
\end{theorem}
\begin{proof}[Proof mechanism]
Differentiate the normalized expectation repeatedly using~\eqref{eq:normalized-response}. Each derivative either joins an existing action derivative or creates a new connected insertion. The resulting terms are indexed by set partitions, with the displayed alternating signs.
\end{proof}
\src{R18}{thm:v18v40-response-compiler}
The contact vertices $S_{ij},S_{ijk},\ldots$ are part of the response. They cannot be removed by computing only a covariance of a preselected tensor.

\begin{corollary}[Separated source response]
\label{cor:separated-response}
Assume metric locality through order $r$ in the sense that the mixed action vertices vanish for the positively separated source supports under consideration. Then, at sufficiently fine cutoff,
\begin{equation}
 \partial_1\cdots\partial_rW_n(0)
 =(-1)^{r+1}\Cum_n(S_{n,1},\ldots,S_{n,r}).
 \label{eq:separated-response}
\end{equation}
\end{corollary}
\src{R18}{thm:v18v40-separated-response}
After a field has been integrated out, its nonlocal effective action must first be differentiated completely. Locality of the microscopic theory does not permit one to discard the resulting determinant loops or mixed effective vertices without their exact assembly.

\section{The complete Gaussian scalar comparator includes fourth chaos}
The scalar stress contains derivative/quadratic terms and the quartic Higgs potential. Matching only the elementary covariance and subtracting a quadratic Gaussian stress leaves a spurious ``interaction'' contribution from the fourth Gaussian chaos.

\begin{theorem}[Degree-complete Gaussian Higgs stress]
\label{thm:gaussian-stress}
For the centered matched-Gaussian scalar stress on the source branch,
\begin{equation}
 T_{H,n}^{G}[q]=I_2(B_n^H[q])+I_4(D_n^H[q]),
 \label{eq:gaussian-stress}
\end{equation}
and
\begin{equation}
\begin{split}
 \Cov(T_{H,n}^{G}[q],T_{H,n}^{G}[r])
 ={}&2\tr(B_n^H[q]B_n^H[r])\\
 &+\E[I_4(D_n^H[q])I_4(D_n^H[r])].
\end{split}
\end{equation}
The Gaussian comparison and follower short must retain both degrees and their complete higher connected-response contractions.
\end{theorem}
\begin{proof}
Wick-expand the quadratic and quartic scalar source with the matched covariance. After centering, only the second and fourth homogeneous Gaussian chaoses remain. Orthogonality of different chaoses gives the covariance formula. Higher cumulants are obtained from the full Wick contraction, not from covariance orthogonality alone.
\end{proof}
\src{R18}{thm:v18v42-complete-gaussian,prop:v18v42-degree-complete-short}

\section{Fermionic metric loops and soft collars}
On an invertible fermionic chart, $S_F=-\log\det D$ and $G=D^{-1}$ satisfy
\begin{align}
 (S_F)_i&=-\Tr(GD_i),\\
 (S_F)_{ij}&=\Tr(GD_iGD_j)-\Tr(GD_{ij}),\\
 (S_F)_{ijk}&=-\Tr(GD_iGD_jGD_k)-\Tr(GD_iGD_kGD_j)\nonumber\\
 &\quad+\Tr(GD_{ij}GD_k)+\Tr(GD_{ik}GD_j)
       +\Tr(GD_{jk}GD_i)-\Tr(GD_{ijk}).
 \label{eq:fermion-responses}
\end{align}
Both cubic loop orientations and every mixed vertex are retained before taking norms. These formulas follow from $\partial_iG=-GD_iG$ and trace differentiation. \src{R18}{thm:v18v44-det-responses}

Small singular values require a separate collar. If the complete Ward defect on the hard region has expectation at most $\epsilon_n$, and its absolute expectation on the soft collar is at most $u_n$, then $\epsilon_n+u_n\to0$ gives the desired Ward descent. A small collar probability alone is insufficient when inverse-Dirac insertions diverge there; it must be paired with an inverse-moment or source-integrability estimate. Exact zero modes are treated by the retained Grassmann packet, not by this invertible-chart formula. \src{R18}{thm:v18v45-ward-collar}

\section{Stress tightness and moment landing}
Let $e_k$ be a reference elliptic eigenbasis on a compact four-dimensional region, with eigenvalues $\lambda_k$. Suppose the centered stress coefficients obey
\begin{equation}
 \Var\bigl(\mathsf T_n[e_k]\bigr)\le C(1+\lambda_k)^\alpha,
\end{equation}
and the mean stress is uniformly bounded in $H^{-s_0}$ for $s_0>\alpha+2$.

\begin{theorem}[Distributional stress compactness]
\label{thm:stress-tightness}
Under those assumptions,
\begin{equation}
 \sup_n\E\norm{\mathsf T_n}_{H^{-s_0}}^2<\infty,
\end{equation}
and the stress laws are tight in $H^{-s}$ for every $s>s_0$.
\end{theorem}
\begin{proof}
Sum the variance estimate with weight $(1+\lambda_k)^{-s_0}$. The four-dimensional elliptic counting exponent is two, so $s_0>\alpha+2$ makes the sum finite. Add the mean bound and use compact Sobolev inclusion to obtain tightness in the weaker topology.
\end{proof}
\src{R18}{thm:v18v46-stress-tightness}

Tightness alone does not pass cubic moments. For a chosen three-source panel, the source supplies a sufficient $L^{3+\epsilon}$ bound, including hard/soft fermion contributions, together with convergence of the direct and complete Gaussian laws and compatible follower projections. These assumptions pass the connected cubic excess to the continuum. A nonzero limit is a valid non-quasi-free witness; a zero limit is an honest Gaussian/Wick outcome on that panel, not a theorem that every possible interaction vanishes. \src{R18}{thm:v18v48-cubic-landing}

\section{The surviving observable}
The interaction test used in the regular-branch theorem is
\begin{equation}
 \mathcal E_{3,n}^{\perp}(q_1,q_2,q_3)
 =\mathcal C_{3,n}^{\perp}(q_1,q_2,q_3)
  -\mathcal C_{3,n}^{G,\perp}(q_1,q_2,q_3),
 \label{eq:interaction-excess}
\end{equation}
where the supports are positively separated, both sides are compiled from the same source convention and contact prescription, and the same dynamically defined short is applied. The target is a finite nonzero limit $e_*$. Elementary Gaussianization, a non-Gaussian quadratic score, and a positive stress covariance do not by themselves establish~\eqref{eq:interaction-excess}. This is why interaction survival is a separate row in the final theorem.


% ==================== 12_closure.tex ====================
\chapter{Conserved sources and the regular-branch closure theorem}
\label{ch:closure}
\section{Ward descent is a topology-level estimate}
Let $T_n$ be the compiled finite stress and $\mathcal A_n(\xi)$ the complete diffeomorphism/Ward defect, including phase, collar, domain, and measure terms. The finite identity is tested against compactly supported vector fields $\xi$ in a declared Sobolev duality.

\begin{theorem}[Distributional stress conservation]
\label{thm:ward-descent}
Suppose $\E T_n\to T$ in $H^{-s}_{\rm loc}$ and, on each compact region $K$,
\begin{equation}
 \sup_{\substack{\xi\in C_c^\infty(K;TM)\\
                 \norm{\xi}_{H^{s+1}}\le1}}
 |\E\mathcal A_n(\xi)|\longrightarrow0.
 \label{eq:ward-dual}
\end{equation}
With convergence of the geometric pairing used by the finite Ward identity, the limiting stress satisfies $\operatorname{div}_gT=0$ distributionally.
\end{theorem}
\begin{proof}
The finite Ward identity writes the stress pairing with the infinitesimal metric change as the complete defect. Pass the stress and geometric pairing to the limit and use~\eqref{eq:ward-dual}. The resulting functional vanishes on every compactly supported vector field, which is distributional conservation.
\end{proof}
\src{R18}{thm:v18v41-ward-descent}
Inserted Ward identities are needed for the full physical observable algebra, not merely for the expectation of one tensor. A nonzero trace anomaly is compatible with this diffeomorphism Ward identity; trace and divergence are different variations. \src{R2}{thm:v2v93-trace-conservation}

\section{Weak convergence can leave a stress defect}
At energy endpoints, weak convergence of scalar derivatives, quartic fields, or metric improvements can leave concentration or oscillation defects in the stress. A global quartic moment tail is not the same as spatial uniform integrability. The source isolates defect tensors built from the weak limits of quadratic and quartic energies and the nonminimal improvement term. Conservation of the total defect does not imply that the defect vanishes. An Orlicz or equivalent spatial uniform-integrability estimate can remove the corresponding concentration mechanism. \cite{R1,R2}

This matters for Einstein closure. Passing $H_n\rightharpoonup H$ and declaring $T_n\to T(H)$ can be false even when the finite equations are exact. One must either prove the required strong/uniformly integrable convergence or retain the enhanced defect variables as part of the limiting source packet.

\section{The Einstein residual lives beyond the Ward quotient}
Let $\mathscr V$ be the physical metric-source test space, and let $\mathscr L$ be the allowed local counterterm-variation subspace. The weak Einstein residual defines a class in the corresponding quotient. A summable cutoff cocycle can give a well-defined limiting residual class. A finite local counterterm removes it exactly when the class vanishes; conservation alone does not force this. \cite{R8,R13}

\begin{proposition}[Counterterms cannot erase an arbitrary conserved residual]
\label{prop:residual-class}
If $R$ is a limiting conserved source residual and the permitted finite local changes span $\mathscr L$, then an allowed counterterm sets the residual to zero if and only if $R\in\mathscr L$ with the required actual coefficient realization. A nonzero quotient class remains an obstruction even when every Ward identity holds.
\end{proposition}
\begin{proof}
The allowed changes translate the residual precisely by elements of $\mathscr L$. The claim is the definition of the quotient, supplemented by the requirement that the permitted variation is produced by an actual common-action counterterm.
\end{proof}
This prevents a free choice of the source from being confused with a derived Einstein equation.

\section{Variational landing on a physical slice}
Let $\Gamma_n$ be the complete gauge-fixed effective action on a Hilbert physical-slice chart $\XX$. The metric chart must be nondegenerate, and its topology must be strong enough, or sufficiently enhanced, to define the limiting gravitational variation. It is not enough that a formal metric distribution exists.

\begin{theorem}[Strongly monotone critical-point stability]
\label{thm:critical-stability}
Suppose the gradients on a common ball satisfy
\begin{equation}
 \pair{D\Gamma_n(x)-D\Gamma_n(y)}{x-y}
 \ge\lambda\norm{x-y}_{\XX}^2,\qquad\lambda>0,
 \label{eq:strong-monotonicity}
\end{equation}
and converge uniformly there to $D\Gamma$ with error $\delta_n\to0$. If $x_n$ and $x_*$ are critical points in that ball for $\Gamma_n$ and $\Gamma$, respectively, then $x_*$ is unique and
\begin{equation}
 \norm{x_n-x_*}_{\XX}\le\lambda^{-1}\delta_n.
 \label{eq:critical-stability}
\end{equation}
\end{theorem}
\begin{proof}
Apply~\eqref{eq:strong-monotonicity} to $x_n,x_*$, use $D\Gamma_n(x_n)=D\Gamma(x_*)=0$, and bound $D\Gamma_n(x_*)-D\Gamma(x_*)$ by $\delta_n$. Divide by $\norm{x_n-x_*}$ when it is nonzero. The same monotonicity proves uniqueness.
\end{proof}
\src{R18}{thm:v18v49-critical-convergence}

\begin{theorem}[Mosco landing with a common convexity margin]
\label{thm:mosco}
Suppose $\Gamma_n$ Mosco-converges to $\Gamma$, the source's common strong-convexity inequality holds, the minimizing sequence $x_n$ is bounded, and $\Gamma$ has a finite minimizer. Then its minimizer $x_*$ is unique,
\begin{equation}
 x_n\to x_*\quad\text{strongly in }\XX,
 \qquad \Gamma_n(x_n)\to\Gamma(x_*).
 \label{eq:mosco}
\end{equation}
\end{theorem}
\begin{proof}[Proof mechanism]
The Mosco lower bound identifies every weak accumulation point as a minimizer; a recovery sequence gives the matching energy upper bound. The common convexity margin converts energy convergence relative to the recovery sequence into strong convergence. Uniqueness removes subsequences.
\end{proof}
\src{R18}{thm:v18v50-minimizers}
Convergence of the \emph{sum} of gravitational and matter effective actions is not sufficient to identify the matter stress separately. The matter subgradient must satisfy the energy and graph convergence conditions of the source theorem. Moving physical slices likewise require convergence of the slice projections. These conditions are retained in the closure hypotheses rather than hidden in the notation $\Gamma_n\to\Gamma$.

\begin{theorem}[Effective-action Einstein landing]
\label{thm:einstein-landing}
Assume a limiting critical metric is obtained by the preceding variational mechanism. Suppose the gravitational first variations converge to $G(g)+\Lambda g$, the compiled matter variations converge to $T$, the couplings converge to $\kappa>0$, and the full finite first-variation residual tends to zero in the same source dual topology. Then
\begin{equation}
 G(g)+\Lambda g=\kappa T
 \label{eq:einstein-landing}
\end{equation}
on every physical slice test. Vanishing Ward and Bianchi/contraction defects extend the conserved identity through the declared gauge-equivalence class.
\end{theorem}
\begin{proof}
Write the finite critical equation as the sum of the complete gravitational and matter first variations plus its residual. Pass each term to the limit in the common dual space. The source identifications give~\eqref{eq:einstein-landing}. The Ward and geometric identities control the gauge directions excluded by the slice chart.
\end{proof}
\src{R18}{thm:v18v49-einstein-closure}

\section{Causal Einstein evolution is another conditional handoff}
A Euclidean effective-action critical point and a Lorentzian initial-value problem are not the same construction. The latter requires a retarded source response, causal domain control, compatible initial and boundary constraints, and a well-posed hyperbolic reduction.

The stress polarization generally has fourth-order/nonlocal behaviour, including a $p^4\log p^2$ contribution on the source branch. It cannot simply be treated as a uniformly second-order Lipschitz perturbation at all frequencies. One can instead include it with the Einstein principal operator in a specified graph space. If $E_R$ is the retarded inverse of the reduced Einstein operator and $\Pi_R$ the complete retarded source derivative, the equation
\begin{equation}
 (I-\kappa E_R\Pi_R)h=E_Rf
 \label{eq:causal-einstein}
\end{equation}
is invertible whenever the common-domain operator norm satisfies $\kappa\norm{E_R\Pi_R}<1$. A Volterra estimate can replace a global smallness bound on an appropriate finite time interval. \src{R2}{thm:v2v93-Ward-Neumann}

\begin{proposition}[Conditional constraint propagation]
\label{prop:constraint-propagation}
On the declared harmonic reduction and common boundary domain, the gauge-constraint field obeys a wave equation with the complete Ward defect as source; in the source convention its principal form is
\begin{equation}
 \Box C_\nu+R_\nu{}^\mu C_\mu
 =-2\kappa\,\mathcal W_\nu.
 \label{eq:subsidiary}
\end{equation}
If the Ward defect vanishes, the initial subsidiary data vanish, and the hyperbolic boundary problem is unique, then $C=0$ and the reduced solution solves the corresponding Einstein system.
\end{proposition}
\begin{proof}
Take the covariant divergence of the reduced metric equation, use the Bianchi identity, and insert the assembled source Ward identity. Uniqueness of the homogeneous subsidiary problem gives the result.
\end{proof}
\src{R2}{thm:v2v95-constraint-equation}

Boundary trace lifts, relative cohomology, glancing characteristic directions, and the common transverse-traceless/matter/reservoir domain enter this statement. Private lift margins can protect a common boundary angle; a shared perturbation into a private cokernel can destroy it. A uniform spectral gap without the coefficient and boundary-symbol control does not supply the required Poisson estimates. These are reasons to retain the boundary-domain hypotheses in a local coupled continuation theorem. \cite{R1,R2}

\scope{The causal handoff is a local, branch-dependent well-posedness theorem once its analytic assumptions hold. It is not an assertion of global nonlinear stability, geodesic completeness, or the absence of gravitational collapse.}

\section{Ten simultaneous closure hypotheses}
Let $\AA_0$ be a countable nuclear test algebra containing elementary bosonic tests, gauge-invariant Wilson tests, graded fermionic source tests, the degree-complete scalar stress packet, and the complete matter/gravity source tests. Every condition below refers to the \emph{same} cofinal family~\eqref{eq:common-card}.

\begin{assumption}[Regular-branch closure conditions]
\label{ass:closure}
\begin{enumerate}[label=\textbf{C\arabic*.},leftmargin=3em]
\item \textbf{Common action and line.} Scalar, gauge, fermion, ghost, measure, and metric terms arise from one regulated action. The determinant/Pfaffian charts glue with controlled Berry current and exactly saturated physical zero modes.
\item \textbf{Physical positivity.} The normalized finite Schwinger forms are reflection positive on the declared physical gauge/BRST quotient, with cutoff-uniform polynomial growth on every fixed test panel.
\item \textbf{Joint compactness.} Bosonic fields, nondegenerate metric data, and enhanced stress packets are tight in the specified compatible spaces. Fermionic source coefficients admit the bounds needed for diagonal convergence of their graded forms.
\item \textbf{Moment landing.} The selected interaction panel has the stated $L^{3+\epsilon}$ stock, and every higher fixed moment needed by $\AA_0$ is uniformly integrable.
\item \textbf{Ward and anomaly descent.} The full inserted defects, including hard/soft Dirac collars and phase/domain contributions, vanish in the common dual topology and generate a closed physical Ward graph.
\item \textbf{Metric variational landing.} The effective actions Mosco-converge on compatible physical slices with one strong-convexity margin, bounded minimizers, and the separate energy/graph convergence that identifies the matter subgradient.
\item \textbf{Einstein residual.} The regulated Bianchi and contraction defects vanish, and the full effective-action critical residual tends to zero in the same source topology.
\item \textbf{Interaction survival.} For three positively separated tests, the common-action partition compiler, the degree-complete Gaussian comparator, and the same dynamically defined follower short give $\mathcal E_{3,n}^{\perp}\to e_*\ne0$.
\item \textbf{Route closure.} Adjacent and independently retained direct transports of fields, chiral lines, covariances, source projections, Ward defects, stress, and metric forms have a common summable error tail.
\item \textbf{Physical reconstruction.} The limiting symmetry, continuity, reflection, cluster, locality, and domain conditions suffice for the declared Osterwalder--Schrader reconstruction. Photon and any retained gravitational massless sectors are separated from claims of a positive massive-channel threshold.
\end{enumerate}
\end{assumption}
\src{R18}{def:v18v51-hypotheses}

\begin{theorem}[Regular-branch Standard Model--Einstein continuum]
\label{thm:main}
If one literal cofinal family satisfies Hypothesis~\ref{ass:closure}, its physical Schwinger forms converge to a unique, route-independent regular-branch continuum. The reconstructed physical Hilbert space and declared fields are positive and gauge compatible. The compiled limiting stress is conserved and satisfies
\begin{equation}
 G(g)+\Lambda g=\kappa T
\end{equation}
in the specified distribution space for the limiting critical metric. The separated excess $e_*\ne0$ excludes a quasi-free description modulo the declared contact and follower directions.
\end{theorem}
\begin{proof}
C3 gives joint subsequential compactness and diagonal graded limits. C4 passes the required products and connected functions. C1 and C5 identify one consistent physical line and quotient rather than unrelated chartwise fermion states. C2 passes each reflected Gram inequality; C10 provides the additional continuity, domain, and local-net inputs for reconstruction. Theorem~\ref{thm:ward-descent} gives source conservation. C6 and Theorem~\ref{thm:mosco} identify the limiting critical metric and its actual matter subgradient; C7 and Theorem~\ref{thm:einstein-landing} give the Einstein equation. C8 and the moment/follower convergence pass the nonzero direct-minus-Gaussian response. C9 telescopes all determining fields and source data, identifies independent transports, and removes subsequence and route dependence. Each conclusion thus uses a distinct part of the hypothesis list.
\end{proof}
\src{R18}{thm:v18v51-conditional-closure}

The explicit results in Chapter~\ref{ch:reconstruction} sharpen several
inputs without discharging them globally.  Theorem~\ref{thm:faithful-clock}
and Proposition~\ref{prop:split-clock} specify how detail modes can retain
a clock once the local, reflected, cofinal transport is supplied.
Proposition~\ref{prop:stable-wick-reference} identifies a stable reference,
and Proposition~\ref{prop:collective-heat} identifies a regular collective
operator.  Their physical realization, uniformity, and compatibility
with the other sectors are still required by C1--C3, C9, and C10.
Equation~\eqref{eq:scalar-shift-source} also shows why normalized-state
convergence alone cannot replace the source requirements in C6--C7.
Theorem~\ref{thm:proper-time-c1} makes a corresponding distinction for
the massive determinant: operator-valued subtraction and integrable
first-variation majorants are needed to pass from determinant values to
stress.  Theorem~\ref{thm:three-defect} provides a quantitative route to
C7 on a Hilbert source screen.  With uniform pseudoinverse bounds and a
constraint-angle margin, constraint, Ward, and reduced critical defects
control the full residual.  It does not remove the need to identify those
screens with the common distributional source topology.

\section{The exact remaining boundary}
No supplied construction has established C1--C10 simultaneously on one physical family. In particular, a finite nonempty reflected slab does not establish C2 uniformly after chiral and gravitational coupling; a local conformal or determinant chart does not establish C3 on the full support; a Ward identity does not establish C7; and a non-Gaussian polynomial of a Gaussian field does not establish C8.

There is also a meaningful intermediate result. One may obtain a positive, compatible, locally normal kinematic state without proving the full interaction witness or the Einstein residual. Such a state should be reported with that scope. Conversely, the exact coupled models in the next part can satisfy classical Einstein equations without supplying the joint quantum measure and cutoff comparison of C1--C10. Their role is to test source consistency and dynamics, not to conceal the remaining continuum estimates.


% ==================== 13_gauge_charts.tex ====================
\chapter{Finite gauge charts and a projective-law obstruction}
\label{ch:gauge-charts}
\section{Local curvature reconstruction is a genuine finite theorem}
A finite curvature chart is useful for phase transport, local estimates, and exact pilots. It must still be distinguished from the support of a continuum quantum law. In the small-link logarithmic chart, the nonabelian curvature map has the structure
\begin{equation}
 F(A)=\dd A+N(A),\qquad N(A)=O(\norm A^2),
 \label{eq:curvature-chart}
\end{equation}
with the Baker--Campbell--Hausdorff order fixed by the actual plaquette word. Let $K$ be the declared coexact inverse/homotopy on the admitted linear face data. The nonlinear problem becomes a projected fixed point with the remaining face components fixed by the nonlinear compatibility graph.

\begin{theorem}[Unique small coexact reconstruction]
\label{thm:coexact}
Suppose the coexact radius-$r$ ball is invariant under the source's curvature reconstruction map $\mathcal T_F$, and $\norm K L(r)=c<1$, where $L(r)$ bounds its nonlinear derivative. Then there is a unique fixed point $A_F$ in that ball,
\begin{equation}
 \norm{A^{(j)}-A_F}\le\frac{c^j}{1-c}\norm{A^{(1)}-A^{(0)}},
 \qquad
 \norm{A_F-A_G}\le\frac{\norm K}{1-c}\norm{F-G}
 \label{eq:coexact-stability}
\end{equation}
for admitted compatible curvature data. Any actual coexact solution in the same ball is this fixed point.
\end{theorem}
\begin{proof}
The invariant-ball and derivative assumptions give a contraction. The first estimate is the geometric tail bound for Picard iteration. Subtract the two fixed-point equations to obtain the second estimate.
\end{proof}
\src{R18}{thm:v18v78-unique-reconstruction}

The word ``compatible'' cannot be replaced by a false abelian closure condition on the full logarithmic curvature. If $P$ selects the linear exact/closed face coordinates, the complementary part is supplied by a nonlinear graph, schematically
\begin{equation}
 F=G+(I-P)N(A_G).
 \label{eq:curvature-graph}
\end{equation}
The source obtains height and derivative bounds for that graph. Setting the second term to zero would impose the wrong nonabelian Bianchi relation. \src{R18}{thm:v18v89-curvature-graph,thm:v18v91-graph-height,thm:v18v92-graph-differential}

\section{Equivariance and naturality are separate}
A rooted finite homotopy can be exactly natural under its rooted restrictions. Averaging it with its reflected transform can restore reflection equivariance at a fixed size while destroying that exact restriction naturality. One must therefore compute the actual overlap or restriction defect. It is not legitimate to combine the equivariance of one choice with the exact projectivity of another without a comparison theorem. \src{R18}{thm:v18v88-reflection-average}

\begin{theorem}[Shrinking charts under raw restriction collapse]
\label{thm:shrinking-chart}
Suppose a deterministic family of face cochains is exactly consistent under ordinary cell restriction and
\begin{equation}
 R^{M,N}G_M=G_N\quad(M\ge N),\qquad
 \norm{G_N}_\infty\le\frac C{N^2}.
 \label{eq:shrinking-chart}
\end{equation}
Then $G_N=0$ for every fixed $N$. The same statement holds almost surely for a projective law supported in those balls at every level.
\end{theorem}
\begin{proof}
For fixed $N$, restriction does not increase the supremum norm, so $\norm{G_N}\le C/M^2$ for all $M\ge N$. Let $M\to\infty$. Apply the same argument on the probability-one set where the countably many consistency and support statements hold.
\end{proof}
\src{R18}{thm:v18v94-trivial-projective}
The obstruction is to raw restriction together with a shrinking global radius. It motivates rescaled geometric blocking, not abandonment of finite charts.

\section{Gauge-covariant blocking uses transported faces}
A coarse link is the ordered product of its two fine links. A coarse plaquette is correspondingly a product of four fine plaquettes transported to a common base point. With the source's square labelling,
\begin{equation}
 P^{\rm c}
 =\Ad_{h_{00}}(p_{10})\,
  \Ad_{h_{00}v_{10}}(p_{11})\,
  p_{00}\,\Ad_{v_{00}}(p_{01}).
 \label{eq:four-face}
\end{equation}
This is an exact nonabelian word identity, not the sum of four untransported logarithms. \src{R18}{thm:v18v96-four-face-word}

On the specified small chart, rescaled curvature blocking has the estimate
\begin{equation}
 \norm{\mathcal G_N^{\rm c}-B_2\mathcal G_{2N}}_\infty
 <\frac1{2\,000\,000\,N},
 \label{eq:nonlinear-blocking}
\end{equation}
with the constants and normalization of that source chart. The nonlinear error is summable on dyadic $N=2^j$. The later projected identity also contains a chart/projection comparison whose available bound does not automatically decay. It is therefore not replaced by the nonlinear error alone. \src{R18}{thm:v18v97-rescaled-error,thm:v18v98-projected-identity}

\section{An exact non-Haar projective gauge law}
Let $G$ be a positive-dimensional compact gauge group with symmetric central heat-convolution laws $\mu_t$. On the mesh $1/N$, assign independent link laws $\mu_{\tau/N}$ and then average over all vertex gauges. Denote the result by $\overline\nu_N$.

\begin{theorem}[Exact dyadic projectivity]
\label{thm:heat-projective}
The ordered link-blocking map satisfies
\begin{equation}
 (\mathfrak B_1)_\#\overline\nu_{2N}=\overline\nu_N.
 \label{eq:heat-projective}
\end{equation}
The family is gauge invariant and invariant under the declared coordinate reflections. At finite $N$ it provides the source's explicit non-Haar projective example.
\end{theorem}
\begin{proof}
Two consecutive fine-link heat increments have law $\mu_{\tau/(2N)}*\mu_{\tau/(2N)}=\mu_{\tau/N}$. Distinct coarse links use disjoint fine increments. Gauge averaging commutes with the gauge-covariant blocking map, and heat symmetry gives the reflected law.
\end{proof}
\src{R18}{thm:v18v99-product-projective,thm:v18v99-gauge-projective}
This is a kinematic gauge-law result. Its typical plaquette size determines whether it lies in the previously constructed smooth charts.

\begin{theorem}[Smooth-chart exclusion for the heat-link law]
\label{thm:heat-chart-exclusion}
Let $d=\dim G>0$ and
\begin{equation}
 \mathcal A_N(C)=\left\{U:\norm{\log P(U)}_{\rm op}\le\frac C{N^2}
 \ \text{for every unit plaquette in the principal chart}\right\}.
\end{equation}
For fixed $C$ there are $C',N_0$ such that
\begin{equation}
 \overline\nu_N(\mathcal A_N(C))
 \le\left(C'N^{-3d/2}\right)^{N^4/4}\longrightarrow0.
 \label{eq:chart-exclusion}
\end{equation}
\end{theorem}
\begin{proof}
A plaquette made of four independent link increments has heat law $\mu_{4\tau/N}$. The small-time heat density is bounded by a constant times $N^{d/2}$, while a principal-chart ball of radius $C/N^2$ has Haar volume at most a constant times $N^{-2d}$. The one-plaquette event thus has probability at most $C'N^{-3d/2}$. Choose at least $N^4/4$ plaquettes with pairwise disjoint edge sets; their events are independent before gauge averaging. The event is gauge invariant, so averaging does not change its probability. Multiplication gives~\eqref{eq:chart-exclusion}.
\end{proof}
\src{R19}{thm:v19v1-chart-exclusion}

\scope{This excludes the fixed smooth $O(N^{-2})$ plaquette chart for this explicit independent heat-link scaling. It does not exclude correlated gauge measures, differently scaled laws, or rough continuum constructions. It proves that a nonempty projective law and a nonempty smooth reconstruction chart need not overlap on typical configurations.}

The exact two-dimensional heat-face gluing identity is a separate result. Four-dimensional faces share a more complicated edge and cell constraint structure, including the six-face Haar projection. The two-dimensional factorization cannot be imported unchanged as a four-dimensional construction. \cite{R19}


% ==================== 14_compact_gravity.tex ====================
\part{Explicit coupled sectors and gravitational dynamics}
\chapter{Compact gravity, physical constraints, and corrected quartic forms}
\label{ch:compact}
\section{The transverse-traceless seed and its zero mode}
The explicit compact model starts from finitely many transverse-traceless (TT) Fourier modes on a three-torus. At nonzero momentum, the physical ADM quotient has two TT polarizations. The corresponding quadratic energy is positive. The homogeneous massless coordinate is different: a free coordinate does not have a normalizable oscillator vacuum. Removing the mean, imposing a Dirichlet condition, or adding a mass are distinct constructions and must be named rather than conflated. \cite{R19}

For nonzero frequencies, the reflected Gaussian seed has a conditional-half square and reconstructs the oscillator Hamiltonian. If $\omega_k$ are the retained frequencies, its transfer is
\begin{equation}
 T_\epsilon=\exp\!\left(-\epsilon\sum_k\omega_k\widehat N_k\right)
 \label{eq:tt-transfer}
\end{equation}
on the Gaussian plane-state realization, with the source's normal-ordering convention. This is an exact agreement with the linearized Einstein Fock Hamiltonian on the finite seed; nonlinear constraints and metric dynamics are further steps. \src{R19}{thm:v19v64-consolidated}

\section{Compact constraints obstruct arbitrary linear perturbations}
For vacuum initial data, use
\begin{equation}
 R(g)+\tau^2-|K|_g^2=2\Lambda,
 \qquad\nabla^j(K_{ij}-\tau g_{ij})=0,
 \qquad\tau=\tr_gK.
 \label{eq:adm-constraints}
\end{equation}
At a flat torus with $K=0$ and $\Lambda=0$, a formal linear TT perturbation need not be tangent to an exact constrained family.

\begin{theorem}[Quadratic compact compatibility]
\label{thm:taub}
Let $(g(\varepsilon),K(\varepsilon))$ be a $C^2$ family of exact vacuum data through $(\delta,0)$ at $\Lambda=0$. Let $h=g'(0)$ be TT and $A=K'(0)$, with $A_0=A-(\tr_\delta A)\delta/3$ and $\tau_1=\tr_\delta A$. Then
\begin{equation}
 \frac14\int_{\mathbb T^3}|\nabla h|^2
 +\int_{\mathbb T^3}|A_0|^2
 =\frac23\int_{\mathbb T^3}\tau_1^2.
 \label{eq:taub}
\end{equation}
\end{theorem}
\begin{proof}[Proof mechanism]
Differentiate the Hamiltonian constraint twice and integrate over the compact torus. Linear scalar-divergence terms integrate to zero. The TT curvature variation contributes the Dirichlet term and the extrinsic-curvature variation contributes $|A_0|^2-2\tau_1^2/3$. Their balance gives~\eqref{eq:taub}.
\end{proof}
\src{R19}{thm:v19v26-taub}
Thus a nontrivial wave cannot be added to the flat, time-symmetric, zero-$\Lambda$ compact branch while keeping every compensating averaged variable fixed.

\section{York--Lichnerowicz completion}
Write $g=\psi^4\bar g$ and use a transverse-traceless seed $\sigma$ with a longitudinal York correction when required. In the vacuum CMC convention,
\begin{equation}
 -8\bar\Delta\psi+R(\bar g)\psi
 -|\sigma|_{\bar g}^2\psi^{-7}
 +\left(\frac23\tau^2-2\Lambda\right)\psi^5=0,
 \label{eq:lichnerowicz}
\end{equation}
with the coupled momentum constraint and the prescribed volume normalization. The mean-zero scalar equation is solved separately from the averaged balance. Near a suitable seed, the source obtains analytic branches by the implicit-function theorem with an explicit inverse margin. The full momentum equation determines the York correction; a convenient diagonal ansatz for the momentum is not automatically a solution. \cite{R19}

The Euclidean-sign version has a different bifurcation structure. A uniqueness statement for one Lorentzian constraint branch is not analytically continued through a sign change as a uniqueness statement for the Euclidean problem. Likewise, retaining a constraint surface under a Markov update does not prove that the update generates Einstein evolution.

\section{Reflection planes, time symmetry, and mean curvature}
An equivariant map of the seed preserves reflection positivity only when it is local to the appropriate time half. Reflection-odd link differences have the correct algebraic parity. A continuous odd field has zero sharp-time restriction on the reflection plane, however, so a nonzero sharp extrinsic-curvature seed cannot be inferred from that parity alone. Time-smearing and the actual continuum regularity must be addressed. \cite{R19}

For time-symmetric torus data, the source's rigidity result makes the nontrivial-wave branch incompatible with nonnegative $\Lambda$ in the stated setting. A negative $\Lambda$ can balance the second-order wave Dirichlet energy. The resulting plane is a Lorentzian volume maximum and a Euclidean volume minimum. A reflected construction centered on that plane need not be time-translation invariant; its positive reflection square does not itself define a stationary physical semigroup. \src{R19}{thm:v19v64-consolidated}

\section{Physical states and the averaged constraint}
In the finite TT quantum seed, the second-order normal-ordered averaged constraint and total-momentum condition select physical states. For an energy-sector label $e$, the regular mean-curvature branches satisfy
\begin{equation}
 \frac23\tau_e^2=e+2\Lambda,
 \qquad\tau_e=\pm\sqrt{\frac32(e+2\Lambda)}.
 \label{eq:tau-branches}
\end{equation}
The threshold $e+2\Lambda=0$ is a singular branch containing the classical time-symmetric slice. Under the source's real plane-state symmetry, the two signs carry equal weight; the construction does not choose a time arrow. Number states arise from Hermite-weighted plane laws and coherent states from shifted Gaussian plane laws. \src{R19}{thm:v19v64-consolidated}

The finite zero-point contribution lies in the averaged constraint and can be removed by the stated normal ordering on states. This does not make the same subtraction a positive pathwise constraint on the Euclidean seed. The source exhibits divergent variances and positivity/conformal failures for that pathwise substitution. Normal ordering a state constraint and solving a samplewise random geometric equation are different operations.

\section{York-completed quartic balance}
All formulas in this section use the corrected York completion. Earlier momentum coefficients obtained from a diagonal ansatz that failed the full momentum constraint are not used.

For one wavenumber and two polarizations, write the configuration amplitude as $x=(x_+,x_\times)$ and the TT momentum amplitude as $p=(p_+,p_\times)$. The corrected quartic balance is
\begin{equation}
\begin{split}
 Q(x,p)={}&-\frac{145}{1024}|x|^4
 -\frac{59}{128}|x|^2|p|^2+\frac7{64}|p|^4
 +\frac5{64}(x\cdot p)^2\\
 ={}&-\frac{145}{1024}|x|^4
 -\frac{49}{128}|x|^2|p|^2+\frac7{64}|p|^4
 -\frac5{64}(x\wedge p)^2.
\end{split}
\label{eq:quartic-york}
\end{equation}
The equality of the two displayed forms is the two-dimensional identity $|x|^2|p|^2=(x\cdot p)^2+(x\wedge p)^2$. The source reports the York-completed coefficient calculation and its rotation checks; this is a finite-mode quartic balance, not an exact all-amplitude Hamiltonian. \src{R19}{thm:v19v86-form}

For the source's mode-pair normalization, the corrected mixed coefficients are:
\begin{center}
\begin{tabular}{c r r r}
\toprule
Pair $(1,l)$ & $p^2q^2$ & $b^2p^2$ & $a^2q^2$\\
\midrule
$(1,2)$ & $35/18$ & $-3/2$ & $-1/24$\\
$(1,3)$ & $35/64$ & $-279/256$ & $1/256$\\
$(1,4)$ & $119/450$ & $-74/75$ & $1/200$\\
$(1,5)$ & $91/576$ & $-725/768$ & $1/256$\\
$(1,6)$ & $37/350$ & $-2259/2450$ & $29/9800$\\
\bottomrule
\end{tabular}
\end{center}
Here $a,b$ are the two configuration amplitudes and $p,q$ their momentum amplitudes. At the resonant pair $(1,3)$, the recorded odd mixed coefficients are $-109/256$ for $a^3b$ and $7/16$ for $p^3q$. These finite coefficient results are retained with their original normalization. \src{R19}{thm:v19v84-pairs}

\section{The corrected bare self-energy diagnostic}
Let $v_1,v_l$ denote the oscillator variances in the finite-mode convention. The coefficient linear in the occupation of mode one from the vacuum of mode $l$ is
\begin{equation}
 \sigma_{1l}=2c^{xx}_{1l}v_1v_l
 +\frac{c^{pq}_{1l}}{8v_1v_l}
 +\frac{c^{aq}_{1l}v_1}{2v_l}
 +\frac{c^{bp}_{1l}v_l}{2v_1}.
 \label{eq:self-energy}
\end{equation}
For $c_g=3/4$, the corrected partial sums for $K=2,\ldots,6$ are
\begin{equation}
 \sigma_1(K)=
 \left(\frac{235}{2592},-\frac{62615}{165888},
 -\frac{35209}{30720},-\frac{225029}{103680},
 -\frac{17452877}{5080320}\right).
 \label{eq:self-energy-values}
\end{equation}
Thus the finite bare one-graviton coefficient $E_1(K)=1/(2c_g)+\sigma_1(K)$ becomes negative at $K=4$ in this corrected calculation. The corrected intrinsic cutoff is $K_*(3/4)=4$, not the earlier value. The leading large-$K$ term of the stated plane-family calculation remains $-K^2/(16c_g^2)$; the corrected momentum contributions change the subleading terms. \src{R19}{thm:v19v84-self-energy}

\scope{This is a diagnostic of the bare truncated quartic model and its normalization. A negative perturbative occupation coefficient is not, without a controlled remainder and a constructed interacting operator, a theorem that the full physical quantum-gravity Hamiltonian is unbounded below. Nor does it rule out a renormalized trajectory with additional terms.}

\section{What the compact model establishes}
The strongest coherent statement is the conjunction of finite seed positivity, exact linearized spectral matching, explicit second-order constraint completion, physical mean-curvature branches, and corrected finite-mode quartic source calculations. Each element tests a different part of the continuum architecture. Their conjunction does not yet supply a uniform interacting gravitational Gibbs state, a cofinal renormalized metric law, or a physical cutoff-independent Einstein generator.


% ==================== 15_waves.tex ====================
\chapter{Standing waves, averaged cosmology, and colour exchange}
\label{ch:waves}
\section{Classical wave balance on the compact branch}
The finite-mode gravity calculations become a classical small-amplitude construction when the initial metric is completed by the scalar and momentum constraints. For a TT profile
\begin{equation}
 h(x)=\sum_k a_k\cos(k\cdot x)e_k,
\end{equation}
with the orthogonality and normalization used in the source, the second-order balancing coefficient is
\begin{equation}
 \Lambda_2=-\frac18\langle|\nabla h|^2\rangle
 =-\frac1{16}\sum_k|k|^2a_k^2|e_k|^2.
 \label{eq:wave-balance}
\end{equation}
Here the brackets are the torus average and $\Lambda=\varepsilon^2\Lambda_2+O(\varepsilon^4)$ when $h$ is the fixed profile multiplied by the common small amplitude $\varepsilon$. The coefficient is additive over the orthogonal mode bank. \src{R19}{thm:v19v93-lambda}

\begin{theorem}[Classical fourth-order coefficient continuity]
\label{thm:h2}
For the balanced TT datum $\psi^4(\delta+\varepsilon h)$ with zero-mean conformal corrections, the source's fourth-order coefficient satisfies
\begin{equation}
 |\Lambda_4|\le C\norm h_{H^2(\mathbb T^3)}^4.
 \label{eq:h2-bound}
\end{equation}
Finite-mode profiles converging in $H^2$ have convergent $\Lambda_2$ and $\Lambda_4$. For a nonzero gradient profile,
\begin{equation}
 \abs{\frac{\Lambda_4}{\Lambda_2}}
 \le\frac{8C\norm h_{H^2}^4}{\norm{\nabla h}_{L^2}^2}
\end{equation}
with the same volume convention.
\end{theorem}
\begin{proof}[Proof mechanism]
The perturbative Lichnerowicz equation solves its mean-zero second-order coefficient by the elliptic inverse. In three dimensions the $H^2$ product bounds control the quadratic and quartic coefficient terms. Their integrated fourth-order combination gives~\eqref{eq:h2-bound}; multilinear continuity gives convergence of the coefficient along $H^2$ approximations.
\end{proof}
\src{R19}{thm:v19v94-sobolev}
\scope{This is a classical mode-limit theorem for fixed-order coefficients. It is not ultraviolet convergence of a quantum vacuum, whose typical mode amplitudes need not lie in this $H^2$ class.}

\section{Focusing is a theorem about the normal congruence}
Use unit lapse and zero shift and the convention in which positive $\tau$ describes contraction of the normal volume element.

\begin{theorem}[Focusing of a balanced time-symmetric datum]
\label{thm:focusing}
For a time-symmetric compact vacuum datum with $\Lambda<0$, along each geodesic normal and for as long as the synchronous development exists,
\begin{equation}
 \tau(t)\ge\sqrt{3|\Lambda|}
 \tan\!\left(\sqrt{|\Lambda|/3}\,t\right).
 \label{eq:focusing}
\end{equation}
The normal congruence focuses no later than
\begin{equation}
 t_f=\frac\pi2\sqrt{\frac3{|\Lambda|}}.
 \label{eq:focusing-time}
\end{equation}
For $\Lambda=-\varepsilon^2/8+O(\varepsilon^4)$, this upper bound is $\pi\sqrt6/\varepsilon+O(\varepsilon)$.
\end{theorem}
\begin{proof}
The vacuum evolution and $|K|^2\ge\tau^2/3$ imply $\dot\tau\ge\tau^2/3+|\Lambda|$, with $\tau(0)=0$. Compare with the exact Riccati solution on the right of~\eqref{eq:focusing}. Its finite-time divergence gives the normal focusing bound.
\end{proof}
\src{R19}{thm:v19v92-focusing}
A caustic or breakdown of synchronous normal coordinates is not automatically a curvature singularity or a theorem of maximal-spacetime inextendibility. The theorem is kept in its congruence-level scope. Also, the expanding/contracting $\Lambda=0$ branches of the source are not all time-symmetric turnaround branches; the preceding negative-$\Lambda$ conclusion is not transferred to them.

\section{Averaged radiation and its limitations}
At second order, directional TT waves carry the stated null-radiation averages. When the directional weights satisfy the isotropy condition
\begin{equation}
 \sum_k\Lambda_{2,k}\,\widehat k\otimes\widehat k
 =\frac{\Lambda_2}{3}I,
 \label{eq:isotropic-wave}
\end{equation}
the averaged pressure is $p=\rho/3$ and the homogeneous response agrees with the corresponding radiation--cosmological-constant balance in the source normalization. For a single direction the stress is anisotropic. At higher order, the standing-wave solution retains spatial harmonics; a homogeneous Bianchi-I ansatz can reproduce some transverse mean terms while failing the longitudinal equation. The local solution is therefore not replaced by its average. \cite{R19}

Within the special plane-symmetric massless scalar system, a scalar and a gravitational polarization can be combined by the source's rotation of $(b,\phi/\sqrt2)$. A scalar mass breaks this equivalence. Massive scalar averages have longitudinal pressure-to-density ratio $k^2/(k^2+m^2)$ in the stated mode calculation. The Maxwell standing-wave calculation has a different local stress and no generated gravitational-radiation component at the stated second order, unlike the corresponding scalar/gravitational examples. These are model-specific comparisons of local source tensors, not claims that fields of different spin are generally equivalent. \src{R19}{thm:v19v97-rotation,thm:v19v98-mode,thm:v19v99-wave}

\section{A nonabelian plane-wave source}
For the diagonal plane metric
\begin{equation}
 \dd s^2=-\dd t^2+e^{2a+2b}\dd x^2
 +e^{2a-2b}\dd y^2+e^{2c}\dd z^2,
 \label{eq:plane-metric}
\end{equation}
consider two $SU(2)$ colour components $A^1=f\,\dd x$, $A^2=g\,\dd y$. Their commutator contributes $efg\,\dd x\wedge\dd y$ in the third colour. Put $W_\pm=e^{-2a\pm2b}$. The gauge equations are
\begin{align}
 f_{tt}+f_t(c_t-2b_t)
 -e^{-2c}\bigl[f_{zz}-f_z(2b_z+c_z)\bigr]
 +e^2g^2fW_+&=0,\label{eq:two-colour-f}\\
 g_{tt}+g_t(c_t+2b_t)
 -e^{-2c}\bigl[g_{zz}+g_z(2b_z-c_z)\bigr]
 +e^2f^2gW_-&=0.\label{eq:two-colour-g}
\end{align}
At $e=0$, these reduce to two abelian components. At nonzero $e$, the first nonabelian resonant and third-harmonic corrections are cubic in wave amplitude, while their gravitational response enters the corresponding fourth-order calculation. Relative phase affects those coefficients and cannot be inferred from the averaged energy alone. \cite{R20}

\section{The colour resonant system}
Use the amplitude convention
\begin{equation}
 f=\varepsilon\operatorname{Re}(u e^{\ii t})\cos z+\cdots,
 \qquad g=\varepsilon\operatorname{Re}(v e^{\ii t})\cos z+\cdots.
\end{equation}

\begin{theorem}[Nonabelian resonant exchange]
\label{thm:colour-amplitude}
The resonant projection of the cubic nonabelian terms, in slow time $\tau=e^2\varepsilon^2t$, is
\begin{equation}
 \dot u=\ii\frac3{32}(2|v|^2u+v^2\bar u),\qquad
 \dot v=\ii\frac3{32}(2|u|^2v+u^2\bar v).
 \label{eq:colour-amplitude}
\end{equation}
It is the Hamiltonian system $\dot u=\ii\partial_{\bar u}H_{\rm NA}$, $\dot v=\ii\partial_{\bar v}H_{\rm NA}$ with
\begin{equation}
 H_{\rm NA}=\frac3{32}\bigl(2|u|^2|v|^2+
          \operatorname{Re}(u^2\bar v^2)\bigr),
 \label{eq:colour-ham}
\end{equation}
and it conserves $N=|u|^2+|v|^2$ and $H_{\rm NA}$.
\end{theorem}
\begin{proof}
Project $g^2f$ and $f^2g$ onto the spatial first harmonic and temporal frequency one. The surviving complex terms are precisely those in~\eqref{eq:colour-amplitude}. Differentiating~\eqref{eq:colour-ham} gives the displayed flow. Its common phase symmetry conserves $N$, and Hamiltonian evolution conserves $H_{\rm NA}$.
\end{proof}
\src{R20}{thm:v20v6-averaged}
This is an exact algebraic resonant projection and an exact invariant statement for the reduced system. It does not, without a separate averaging-error theorem, claim a uniform long-time approximation of the full nonlinear Einstein--Yang--Mills PDE.

\section{Gravitational polarization exchange}
With the source's covariant metric-response prescription, the gravitational part of the third-order resonant coefficients has Hamiltonian
\begin{equation}
 H_{\rm grav}=\frac1{16}\bigl(|u|^2|v|^2-
          \operatorname{Re}(\bar u^2v^2)\bigr)
 =\frac18S^2,\qquad S=\operatorname{Im}(\bar uv).
 \label{eq:grav-amplitude}
\end{equation}
In slow time $\varepsilon^2t$ its equations are
\begin{equation}
 \dot u=\frac S8v,\qquad\dot v=-\frac S8u.
 \label{eq:polarization-rotation}
\end{equation}
Both $N$ and $S$ are conserved. A single linear polarization has no self-shift in this prescription, in-phase linear polarizations have no such exchange, and the normalized circular state has the recorded frequency shift $1/8$. \src{R20}{thm:v20v8-hamiltonian}

\begin{theorem}[Combined polarization sphere]
\label{thm:stokes}
Define
\begin{equation}
 s_1=|u|^2-|v|^2,\qquad
 s_2=2\operatorname{Re}(\bar uv),\qquad
 s_3=2\operatorname{Im}(\bar uv).
\end{equation}
For $H=H_{\rm grav}+e^2H_{\rm NA}$,
\begin{equation}
 H=\beta_2s_2^2+\beta_3s_3^2,
 \qquad\beta_2=\frac9{128}e^2,\quad
 \beta_3=\frac1{32}+\frac3{128}e^2,
\end{equation}
and
\begin{equation}
 \dot s=\bigl(4(\beta_3-\beta_2)s_2s_3,
 -4\beta_3s_1s_3,4\beta_2s_1s_2\bigr).
 \label{eq:stokes}
\end{equation}
The flow conserves $|s|=N$ and $H$. The colour-axis linear states are elliptic for $e\ne0$. The circular and $45$-degree linear states exchange their elliptic/hyperbolic character at $e^2=2/3$; at that value $s_1$ is conserved.
\end{theorem}
\begin{proof}
Differentiate the three Stokes variables using the two complex amplitude equations. This gives~\eqref{eq:stokes}; its scalar products with $s$ and $\nabla H$ vanish. Linearize at the coordinate axes. The tangent-plane determinants have the signs of $e^2(3e^2+4)$, $e^2(3e^2-2)$, and $-(3e^2-2)(3e^2+4)$, respectively.
\end{proof}
\src{R20}{thm:v20v9-stokes}
\scope{The threshold $e^2=2/3$ belongs to this normalization and this reduced model. It is not a prediction for a measured Standard Model coupling.}

\section{Higgs mass is read through the same geometry}
In the unitary-Higgs branch $\Phi=(0,v+h)/\sqrt2$, the potential convention is
\begin{equation}
 V(h)=\frac\lambda4\bigl((v+h)^2-v^2\bigr)^2,
 \qquad m_h^2=2\lambda v^2,
 \qquad m_A^2=\frac{e^2v^2}{4}.
 \label{eq:toy-higgs-masses}
\end{equation}
These masses enter both the matter wave equations and the gravitational stress. They are not independent labels placed on the metric. The finite resonant calculations show how a massive Higgs and a massive gauge wave alter the response denominators and phase exchange. A reduced elimination prescription can fail to preserve the naive two-amplitude Hamiltonian form even though the full common Lagrangian remains variational. Finite rational fits of frequency-dependent coefficients are retained as finite checks, not promoted here to all-frequency theorems. \cite{R20}


% ==================== 16_coupled.tex ====================
\chapter{An assembled Einstein--Yang--Mills--Higgs--Dirac system}
\label{ch:coupled}
\section{Field content and normalization}
The final explicit system uses the plane metric~\eqref{eq:plane-metric}, one retained $SU(2)$ colour component $A^1=f\,\dd x$, the unitary Higgs $\Phi=(0,v+h)/\sqrt2$, and a paired spin--isospin doublet. It is a \emph{toy Standard Model sector}: it does not contain the complete $SU(3)\times SU(2)\times U(1)$ field system with all three physical generations. Its value is that the retained fields have one Lagrangian, one stress convention, and exactly compatible Einstein and matter equations.

In this chapter the gravitational coupling is absorbed into the stress, so
\begin{equation}
 G_{\mu\nu}+\Lambda g_{\mu\nu}=T_{\mu\nu}.
 \label{eq:toy-einstein}
\end{equation}
The symbols $a,b,c$ are metric functions, $W_\pm=e^{-2a\pm2b}$, and the Yukawa mass is $M=y(v+h)/\sqrt2$. These conventions differ from a formula retaining an explicit factor $8\pi G$ and from the radial Higgs normalization in Chapter~\ref{ch:action}; the physical content must be compared after translating conventions.

\section{The locked doublet and exact Dirac equations}
Write $\psi_1=(u,0,w,0)$ and let $P(u,0,w,0)=(0,u,0,-w)$. Impose
\begin{equation}
 \psi_2=\ii P\psi_1,
 \qquad u_2=\ii u,\quad w_2=-\ii w.
 \label{eq:locked-pair}
\end{equation}
The retained Dirac equations before substitution of the partner components are
\begin{align}
 \ii u_t+\ii(a_t+\tfrac12c_t)u
 +\ii e^{-c}(w_z+a_zw)-Mu
 +\tfrac12ef e^{-a-b}w_2&=0,\label{eq:dirac-u}\\
 -\ii w_t-\ii(a_t+\tfrac12c_t)w
 -\ii e^{-c}(u_z+a_zu)-Mw
 -\tfrac12ef e^{-a-b}u_2&=0.\label{eq:dirac-w}
\end{align}
The partner equations are obtained by the source's spin--isospin exchange; on~\eqref{eq:locked-pair} they are the corresponding $\ii$ and $-\ii$ multiples of the displayed equations.

\begin{theorem}[Invariant paired sector]
\label{thm:paired}
For arbitrary metric functions, common mass $M$, and colour field $f$, the paired subspace~\eqref{eq:locked-pair} is invariant. Its full transverse off-diagonal stress vanishes, while its gauge current is
\begin{equation}
 J=-2e e^{-a-b}\operatorname{Im}(\bar wu).
 \label{eq:current}
\end{equation}
The source of gauge and gravitational polarization is the same bilinear, evaluated in the same convention.
\end{theorem}
\begin{proof}[Proof mechanism]
Substitute $u_2=\ii u,w_2=-\ii w$ into the two spin--isospin systems. Their partner rows reduce to multiples of~\eqref{eq:dirac-u}--\eqref{eq:dirac-w}. In the Belinfante stress, the two locked components cancel the forbidden off-diagonal entries. The action derivative with respect to $f$ gives~\eqref{eq:current}.
\end{proof}
\src{R20}{thm:v20v18-paired}
The gauged equation still contains $b$ through $e^{-a-b}f$. The cited construction's sentence claiming that the full gauged equation contains no $b$ is not used: its displayed equation itself shows that dependence. What cancels in the paired construction is the forbidden stress and the relevant uncoupled spin-connection dependence, not every occurrence of the polarization metric in the gauge term.

\section{All matter equations from the total action}
\begin{theorem}[Assembled matter system]
\label{thm:toy-matter}
Together with~\eqref{eq:dirac-u}--\eqref{eq:dirac-w} and~\eqref{eq:locked-pair}, the matter equations are
\begin{align}
 &f_{tt}+f_t(c_t-2b_t)
 -e^{-2c}\bigl[f_{zz}-f_z(2b_z+c_z)\bigr]
 +\frac{e^2}{4}(v+h)^2f=e^{2a+2b}J,
 \label{eq:full-f}\\
 &h_{tt}+h_t(2a_t+c_t)
 -e^{-2c}\bigl[h_{zz}+h_z(2a_z-c_z)\bigr]
 +\lambda(v+h)\bigl((v+h)^2-v^2\bigr)\nonumber\\
 &\hspace{15mm}+\frac{e^2}{4}(v+h)f^2W_-
 +\sqrt2\,y(|u|^2-|w|^2)=0.
 \label{eq:full-h}
\end{align}
They are Euler--Lagrange equations of the retained total gauge--Higgs--Dirac Lagrangian.
\end{theorem}
\begin{proof}[Proof mechanism]
The gauge and Higgs kinetic terms give the metric-weighted wave operators. Varying $e^2(v+h)^2f^2W_-/8$ gives the gauge mass and Higgs gauge source. Varying the Yukawa mass $M=y(v+h)/\sqrt2$ gives the final Higgs bilinear. Varying the gauge-covariant Dirac term gives the current~\eqref{eq:current}. The source records the complete symbolic derivation from the common Lagrangian.
\end{proof}
\src{R20}{thm:v20v24-matter}

\section{The complete stress packet}
Define
\begin{equation}
 \rho=T_{tt},\quad q=T_{tz},\quad p_z=e^{-2c}T_{zz},\quad
 p_\perp=\frac12(W_-T_{xx}+W_+T_{yy}),\quad
 \pi=\frac12(W_+T_{yy}-W_-T_{xx}).
 \label{eq:stress-convention}
\end{equation}
The sign of $\pi$ is fixed here and used throughout the gravitational polarization equation. With $V(h)$ from~\eqref{eq:toy-higgs-masses}, the assembled sources are
\begin{align}
 \rho={}&\frac12W_-(f_t^2+e^{-2c}f_z^2)
 +\frac12h_t^2+\frac12e^{-2c}h_z^2+V(h)
 +\frac{e^2}{8}(v+h)^2f^2W_-\nonumber\\
 &-2\operatorname{Im}(\bar uu_t+\bar ww_t),\label{eq:rho}\\
 p_z={}&\frac12W_-(f_t^2+e^{-2c}f_z^2)
 +\frac12h_t^2+\frac12e^{-2c}h_z^2-V(h)
 -\frac{e^2}{8}(v+h)^2f^2W_-\nonumber\\
 &+2e^{-c}\operatorname{Im}(\bar uw_z+\bar wu_z),\label{eq:pz}\\
 q={}&W_-f_tf_z+h_th_z
 +e^c\operatorname{Im}(\bar uw_t+\bar wu_t)
 -\operatorname{Im}(\bar uu_z+\bar ww_z),\label{eq:q}\\
 p_\perp={}&\frac12h_t^2-\frac12e^{-2c}h_z^2-V(h)
 +ef e^{-a-b}\operatorname{Im}(\bar wu),\label{eq:pperp}\\
 \pi={}&\frac12W_-(f_t^2-e^{-2c}f_z^2)
 -\frac{e^2}{8}(v+h)^2f^2W_-
 -ef e^{-a-b}\operatorname{Im}(\bar wu).\label{eq:pi}
\end{align}
\src{R20}{thm:v20v24-einstein}

\begin{theorem}[Ansatz consistency and on-shell conservation]
\label{thm:toy-conservation}
For the paired sector, $T_{tx}=T_{ty}=T_{xy}=T_{xz}=T_{yz}=0$. The total stress $T=T^{\rm YM}+T^\Phi+T^\Psi$ satisfies
\begin{equation}
 \nabla^\mu T_{\mu\nu}=0
\end{equation}
on the matter equations, for arbitrary metric functions $a,b,c$. Hence the diagonal plane ansatz is consistent and the Einstein equations are compatible with the matter evolution through the Bianchi identity.
\end{theorem}
\begin{proof}[Proof mechanism]
The paired spin--isospin substitution cancels the five forbidden off-diagonal components. In the divergence, replace $f_{tt},h_{tt}$ and the first fermion time derivatives by the matter equations. The gauge-current and Yukawa-transfer terms cancel between the three stress components. The supplied calculation certifies the resulting zero for arbitrary $a,b,c$; no Einstein equation is required for this cancellation.
\end{proof}
\src{R20}{thm:v20v24-stress}

\section{The five Einstein equations}
With~\eqref{eq:stress-convention}, the independent combinations of~\eqref{eq:toy-einstein} are
\begin{align}
 &a_t^2+2a_tc_t-b_t^2
 -e^{-2c}(3a_z^2-2a_zc_z+2a_{zz}+b_z^2)-\rho=\Lambda,
 \label{eq:einstein-constraint-plane}\\
 &a_{tz}+a_ta_z-a_zc_t+b_tb_z+\frac q2=0,
 \label{eq:einstein-momentum-plane}\\
 &b_{tt}+b_t(2a_t+c_t)
 -e^{-2c}\bigl[b_{zz}+b_z(2a_z-c_z)\bigr]+\pi=0,
 \label{eq:einstein-polarization-plane}\\
 &a_{tt}+2a_t^2+a_tc_t
 -e^{-2c}(a_{zz}+2a_z^2-a_zc_z)
 =\Lambda+\frac12(\rho-p_z),
 \label{eq:einstein-a-plane}\\
 &c_{tt}+c_t^2+2a_tc_t
 -e^{-2c}(2a_z^2-2a_zc_z+2a_{zz}+2b_z^2)
 =\Lambda+\frac12(\rho+p_z)-p_\perp.
 \label{eq:einstein-c-plane}
\end{align}
These equations, the source packet~\eqref{eq:rho}--\eqref{eq:pi}, and the matter equations form one closed symbolic system. The first two are the retained Hamiltonian and momentum constraints; the others evolve the metric combinations. \src{R20}{thm:v20v24-einstein}

\section{Exact massless and massive Dirac diagnostics}
\begin{theorem}[Massless standing spinor as homogeneous null dust]
\label{thm:dirac-null}
For $a=a(t),c=c(t),b=0$, the massless spinor
\begin{equation}
 \psi=e^{-a-c/2}e^{-\ii\theta(t)}(\cos z,0,\ii\sin z,0),
 \qquad\dot\theta=e^{-c},
 \label{eq:dirac-null}
\end{equation}
solves the uncoupled massless Dirac equation exactly. Its stress is
\begin{equation}
 T_{tt}=e^{-2c}T_{zz}=\rho=e^{-2a-2c},\qquad
 T_{xx}=T_{yy}=T_{tz}=0,
\end{equation}
and the Einstein equations reduce to
\begin{equation}
\begin{aligned}
 \ddot a+\dot a(2\dot a+\dot c)&=\Lambda,\\
 \ddot c+\dot c(2\dot a+\dot c)&=\Lambda+\rho,\\
 \dot a^2+2\dot a\dot c&=\Lambda+\rho.
\end{aligned}
 \label{eq:dirac-bianchi}
\end{equation}
\end{theorem}
\begin{proof}[Proof mechanism]
The prefactor cancels the spin-connection dilution term and $\dot\theta=e^{-c}$ matches the longitudinal derivative. The spinor bilinears combine $\cos^2z+\sin^2z$ into a homogeneous density; the other displayed stresses cancel. Substitute them into the Einstein combinations above.
\end{proof}
\src{R20}{thm:v20v16-nulldust}

For the flat massive wave, put $\omega=\sqrt{1+m^2}$ and $W=\omega-m$. Then
\begin{equation}
 \psi=e^{-\ii\omega t}(\cos z,0,\ii W\sin z,0)
\end{equation}
has
\begin{equation}
 \rho=\omega(\cos^2z+W^2\sin^2z),\qquad p_z=W,
 \qquad p_x=p_y=T_{tz}=0,
\end{equation}
and $\langle p_z\rangle/\langle\rho\rangle=1/\omega^2$. Thus the massive and massless cases differ in both trace and local density even where a simple averaged equation of state is available. \src{R20}{thm:v20v17-wave}

\section{The corrected gauge--fermion polarization}
For the massless paired seed, $\operatorname{Im}(\bar wu)=-\tfrac12\sin2z$ at leading unit amplitude. The induced colour condensate is
\begin{equation}
 f=\varepsilon^2\frac e4\sin2z+O(\varepsilon^4).
 \label{eq:condensate}
\end{equation}
Its recorded colour self-energy is $-e^2/16$ in the frequency coefficient, and it induces a third-wavenumber fermion response. Those results are unaffected by the later stress-sign correction. \cite{R20}

\begin{proposition}[Corrected fourth-order anisotropic source]
\label{prop:corrected-polarization}
In the convention~\eqref{eq:stress-convention}, the combined colour and fermion fourth-order polarization source is
\begin{equation}
 \pi_4=-\frac{e^2}{8}\cos^22z+\frac{e^2}{8}\sin^22z
 =-\frac{e^2}{8}\cos4z.
 \label{eq:corrected-pi}
\end{equation}
It has zero mean, and the static fourth-order polarization is
\begin{equation}
 b_4=\frac{e^2}{128}\cos4z.
 \label{eq:corrected-b}
\end{equation}
There is no uniform secular term $e^2t^2/16$ from this source.
\end{proposition}
\begin{proof}
The gauge part of~\eqref{eq:pi} contributes $-e^2\cos^22z/8$ and the fermion-current part contributes $+e^2\sin^22z/8$. Their sum is~\eqref{eq:corrected-pi}. At this order the static polarization equation is $-\partial_z^2b_4+\pi_4=0$, which gives~\eqref{eq:corrected-b}.
\end{proof}
\src{R20}{prop:v20v24-corrigendum}
The withdrawn uniform-growth conclusion came from combining the two stress contributions with opposite conventions. Retaining the corrected modulation is necessary for any cosmological interpretation of this example.

\section{Yukawa displacement and the negative frequency shift}
For the massive paired mode of amplitude $\varepsilon$,
\begin{equation}
 \bar\Psi\Psi
 =2\varepsilon^2W(m+\omega\cos2z).
 \label{eq:yukawa-bilinear}
\end{equation}
The second-order static Higgs response is
\begin{equation}
 h_2=-\frac{2mW}{v}
 \left(\frac m{m_h^2}+\frac{\omega\cos2z}{4+m_h^2}\right).
 \label{eq:higgs-displacement}
\end{equation}
Its mean is negative for $m>0$, lowering the local Higgs expectation in this perturbative solution. The mode-parallel mass perturbation gives
\begin{equation}
 \sigma_Y=-\frac{m^2}{v^2}W
 \left(\frac{2m^2}{\omega m_h^2}+\frac\omega{4+m_h^2}\right)<0,
 \label{eq:yukawa-shift}
\end{equation}
so the frequency becomes $\omega+\varepsilon^2\sigma_Y$ with no amplitude growth from that parallel component. \src{R20}{thm:v20v21-displacement,thm:v20v21-shift}

\begin{proof}[Derivation]
The identities $1-W^2=2mW$ and $1+W^2=2\omega W$ give~\eqref{eq:yukawa-bilinear}. Invert the static Higgs operator on its zero and wavenumber-two components to obtain~\eqref{eq:higgs-displacement}. Project the Hermitian Dirac perturbation $(m/v)h_2\diag(1,-1)$ onto the original normalized massive mode. The zero and second harmonics give the two terms in~\eqref{eq:yukawa-shift}.
\end{proof}

\subsection{The partner level and the frozen-background detuning}
The shift of the original level does not by itself determine the resonance:
the partner level must be evaluated in the same induced background.  For
$m_h>0$, retaining the pair-sourced static colour field
$f_2=(eW/4)\sin2z$ and \eqref{eq:higgs-displacement} gives the first-order
frequency corrections
\begin{align}
 \sigma_c^{\pm}&=\mp\frac{e^2W}{16\omega},\\
 \sigma_Y^+&=-\frac{m^2W}{v^2}
     \left(\frac{2m^2}{\omega m_h^2}+\frac{\omega}{4+m_h^2}\right),\\
 \sigma_Y^-&=\frac{m^2W}{v^2}
     \left(\frac{2m^2}{\omega m_h^2}-\frac{\omega}{4+m_h^2}\right).
 \label{eq:paired-level-shifts}
\end{align}
Here the unperturbed levels are $+\omega$ and $-\omega$ and the actual
correction is $\varepsilon^2\sigma$.  The resulting gap correction is
\begin{equation}
 \Delta_{\rm gap}
 =-\frac{e^2W}{8\omega}
   -\frac{4m^4W}{\omega v^2m_h^2},\qquad
 \delta=2\omega+\varepsilon^2\Delta_{\rm gap}-m_h.
 \label{eq:frozen-gap-detuning}
\end{equation}
In particular, the second spatial Higgs harmonic shifts both levels
equally and cancels from the gap; the homogeneous Higgs depression does
not cancel.  These identities follow by projecting the colour and mass
perturbations onto the two orthogonal free modes, using
$1-W^2=2mW$ and $1+W^2=2\omega W$.
\src{R20}{thm:v20v27-shifts}

\scope{Equation~\eqref{eq:frozen-gap-detuning} is a frozen-background,
first-order perturbative detuning.  The fields used in it are sourced by
the original paired mode.  It does not recompute the background during
conversion, include all gravitational frequency shifts, or supply a
self-consistent nonperturbative resonance condition.}


\section{Unequal masses reveal the limitation of the diagonal ansatz}
The paired invariant sector requires a common mass. For unequal masses, the relative phase beats at $\Delta=\omega_2-\omega_1$ and produces an off-diagonal stress. A cross-metric mode must then be retained. In the source's weak-field mode calculation, its forced equation has the form
\begin{equation}
 d_{tt}-d_{zz}=-efS=2T_{xy},
\end{equation}
with spatial wavenumbers one and three and temporal frequencies $\Delta\pm1$. Resonance occurs at the specified gaps $\Delta=2$ and, for the third harmonic, $\Delta=4$. At $\Delta=2$, the mass relation is $m_2^2=m_1^2+4\omega_1+4$; for $m_1=0$ it gives $m_2=2\sqrt2$. These are mode-matching statements in the model, not predictions for Standard Model fermion masses. \cite{R20}

\subsection{An exact geometric completion with two polarizations}
\label{sec:two-polarizations}
The off-diagonal source can be accommodated without treating the cross
polarization only as an additive linear perturbation.  Replace the
transverse metric by $e^{2a}\mathsf M(P,Q)$, where
\begin{equation}
 \mathsf M(P,Q)=
 \begin{pmatrix}
 e^P&e^P Q\\ e^P Q&e^P Q^2+e^{-P}
 \end{pmatrix},\qquad
 g=-\dd t^2+e^{2a}\mathsf M_{ij}\dd x^i\dd x^j+e^{2c}\dd z^2.
 \label{eq:two-polarization-metric}
\end{equation}
Here $i,j\in\{x,y\}$, all fields depend on $(t,z)$, and $P,Q$ are local
polarization coordinates, not the BRST differential or an orthogonal
projection.  One has $\det\mathsf M=1$, $\sqrt{-g}=e^{2a+c}$, and
$g_{xy}=e^{2a+P}Q$.  The original diagonal metric is $Q=0$, $P=2b$.

Define the weighted base wave operator
\begin{equation}
 \mathscr Df=f_{tt}+(2a_t+c_t)f_t
 -e^{-2c}\bigl[f_{zz}+(2a_z-c_z)f_z\bigr]
 \label{eq:polarization-base-operator}
\end{equation}
and the two polarization expressions
\begin{align}
 \mathcal E_P&=\mathscr DP-e^{2P}(Q_t^2-e^{-2c}Q_z^2),\label{eq:polarization-P}\\
 \mathcal E_Q&=\mathscr DQ+2(P_tQ_t-e^{-2c}P_zQ_z).
 \label{eq:polarization-Q}
\end{align}

\begin{theorem}[Transverse Einstein equations as a weighted sigma model]
\label{thm:polarization-sigma}
For \eqref{eq:two-polarization-metric}, let
$X_G=(G^x{}_x-G^y{}_y)/2$ and $Y_G=G^x{}_y$, with the transverse index
raised using $e^{-2a}\mathsf M^{-1}$.  The exact geometric identities are
\begin{equation}
 \begin{aligned}
 X_G&=\frac12\mathcal E_P-\frac12e^{2P}Q\mathcal E_Q,\\
 Y_G&=Q\mathcal E_P+\frac12(1-e^{2P}Q^2)\mathcal E_Q.
 \end{aligned}
 \label{eq:polarization-einstein-map}
\end{equation}
The coefficient matrix has determinant $(1+e^{2P}Q^2)/4>0$.
The vacuum polarization equations are therefore equivalent to
$\mathcal E_P=\mathcal E_Q=0$.  They are the Euler--Lagrange equations,
with their nonzero scalar normalization factors removed, of
\begin{equation}
 L_{\rm pol}=e^{2a+c}\left[
 -(P_t^2+e^{2P}Q_t^2)+e^{-2c}(P_z^2+e^{2P}Q_z^2)
 \right].
 \label{eq:polarization-sigma-lagrangian}
\end{equation}
Its target metric is $\dd P^2+e^{2P}\dd Q^2$, the hyperbolic-plane
metric of curvature $-1$.
\end{theorem}
\begin{proof}[Derivation]
The transverse block has inverse
\[
 \mathsf M^{-1}=
 \begin{pmatrix}e^{-P}+e^P Q^2&-e^P Q\\-e^P Q&e^P\end{pmatrix}.
\]
For a block metric with base indices $A,B\in\{t,z\}$, the mixed
Christoffel symbols are
$\Gamma^i{}_{Aj}=a_A\delta^i_j+\tfrac12(\mathsf M^{-1}\partial_A\mathsf M)^i{}_j$
and
$\Gamma^A{}_{ij}=-\tfrac12g^{AB}\partial_B(e^{2a}\mathsf M_{ij})$.
Substitution in the traceless transverse Ricci tensor, using
$\Tr(\mathsf M^{-1}\partial_A\mathsf M)=0$, gives
\eqref{eq:polarization-einstein-map}; the scalar-curvature term drops
out of these combinations.  The displayed determinant then proves
invertibility.  Direct variation of \eqref{eq:polarization-sigma-lagrangian}
gives $2e^{2a+c}\mathcal E_P$ and $2e^{2a+c+2P}\mathcal E_Q$.
Finally, the curvature of $\dd P^2+f(P)^2\dd Q^2$ is $-f''/f$;
$f=e^P$ gives $-1$.
\end{proof}
\src{R20}{thm:v20v38-metric,thm:v20v38-sigma,thm:v20v38-limits}

At $Q=0$, $P=2b$, the equation $\mathcal E_P/2=-\pi$ is precisely
\eqref{eq:einstein-polarization-plane}.  On a flat background the cross
equation linearizes to $Q_{tt}-Q_{zz}=2T_{xy}$, recovering the
unequal-mass diagnostic.  The nonlinear terms in
\eqref{eq:polarization-P}--\eqref{eq:polarization-Q} specify how the two
polarizations couple after the cross component has been retained.

\subsection{The direct matter sources of both polarizations}
\begin{proposition}[Invertible stress map and bosonic sources]
\label{prop:polarization-sources}
Put $X_T=(T^x{}_x-T^y{}_y)/2$ and $Y_T=T^x{}_y$.  In the stress
normalization \eqref{eq:toy-einstein}, the traceless transverse equations
are $\mathcal E_P=S_P$, $\mathcal E_Q=S_Q$, with
\begin{equation}
 \begin{aligned}
 S_P&=\frac{2[(1-e^{2P}Q^2)X_T+e^{2P}QY_T]}{1+e^{2P}Q^2},\\
 S_Q&=\frac{2[Y_T-2QX_T]}{1+e^{2P}Q^2}.
 \end{aligned}
 \label{eq:polarization-source-map}
\end{equation}
For the gauge kinetic stress of the retained colour field $A^1=f\,\dd x$,
write $\Pi_f=e^{-2a-P}(f_t^2-e^{-2c}f_z^2)$.  Its contributions are
\begin{equation}
 S_P^{\rm kin}=-(1-e^{2P}Q^2)\Pi_f,\qquad
 S_Q^{\rm kin}=2Q\Pi_f.
 \label{eq:colour-polarization-sources}
\end{equation}
A plane scalar's kinetic and scalar-potential stress is transversely
isotropic and contributes zero to both $S_P$ and $S_Q$.  For a general
stress at $Q=0$,
\begin{equation}
 S_P=e^{-2a}(e^{-P}T_{xx}-e^PT_{yy})=-2\pi,
 \qquad S_Q=2e^{-2a-P}T_{xy}.
 \label{eq:diagonal-polarization-sources}
\end{equation}
\end{proposition}
\begin{proof}
Invert the matrix in \eqref{eq:polarization-einstein-map}; the
cosmological term cancels from the mixed traceless combinations.
For the colour kinetic stress, the two nonzero field-strength entries
are $F_{tx}=f_t$ and $F_{zx}=f_z$.  They give
$X_T=-\tfrac12(1+e^{2P}Q^2)\Pi_f$ and $Y_T=0$, proving
\eqref{eq:colour-polarization-sources}.  A scalar with no transverse
derivatives has $T_{ij}$ proportional to $g_{ij}$ in the transverse
block, so $X_T=Y_T=0$.  Setting $Q=0$ in the inverse metric gives
\eqref{eq:diagonal-polarization-sources}.
\end{proof}
\src{R20}{thm:v20v39-map,thm:v20v39-general}

The kinetic colour source is proportional to $Q$ in the cross equation:
it does not seed a cross polarization from identically zero cross data,
but changes the evolution of a cross polarization already present.
Since $\Pi_f$ is signed, this statement is not a universal growth or
instability theorem.  Within the unequal-mass weak-field calculation,
the retained off-diagonal doublet stress supplies the seed.  The scalar
has no direct traceless transverse source, but it still affects the
geometry through the remaining Einstein equations and its couplings.

\scope{The metric identities and the displayed bosonic stress contributions
are exact on the stated plane ansatz.  The full gauge--Higgs mass term and
fermionic stress must enter through the general source map, not through
the gauge kinetic formula alone.  A Dirac tetrad calculation and all
remaining Einstein equations for $Q\ne0$ have not been assembled here.
Consequently this nonlinear geometric extension is not a completed
non-diagonal version of Theorem~\ref{thm:toy-conservation}.}

\subsection{A projected colour-conversion channel}
The failure of the diagonal Einstein ansatz for unequal masses does not
prevent a well-defined flat, projected matter calculation.  Let
$\omega_j=\sqrt{1+m_j^2}$, $W_j=\omega_j-m_j$, and use the unit-normalized
modes on $0\le z\le2\pi$
\begin{equation}
 \phi_1=\frac{(\cos z,\ii W_1\sin z)}{\sqrt{2\pi\omega_1W_1}},
 \qquad
 \phi_2=\frac{(\ii\cos z,W_2\sin z)}{\sqrt{2\pi\omega_2W_2}}.
 \label{eq:colour-normalized-modes}
\end{equation}
Write $\Psi_j=\varepsilon_Nc_j\phi_je^{-\ii\omega_jt}$ and retain
\[
 f=\frac12\bigl(F e^{-\ii\omega_At}+\bar F e^{\ii\omega_At}\bigr)
                    \sin2z,\qquad \omega_A=\sqrt{4+m_A^2}.
\]
The subscript on $\varepsilon_N$ records unit-normalized spinors; it is
not the unnormalized amplitude $\varepsilon$ in
\eqref{eq:yukawa-bilinear}.  At $\omega_2=\omega_1+\omega_A$, projection
of the colour current and the Dirac coupling gives
\begin{equation}
 \ii\dot c_1=-g_c\bar F c_2,\qquad
 \ii\dot c_2=-g_cF c_1,\qquad
 \ii\dot F=-\kappa_c\bar c_1c_2,
 \label{eq:colour-three-wave}
\end{equation}
where
\begin{equation}
 g_c=\frac{e(W_1+W_2)}{16\sqrt{\omega_1W_1\omega_2W_2}},
 \qquad
 \kappa_c=\frac{2\varepsilon_N^2}{\pi\omega_A}g_c.
 \label{eq:colour-couplings}
\end{equation}
The current contains only the retained $\sin2z$ harmonic.  The factors
in \eqref{eq:colour-couplings} are fixed by the spinor norms and the
$\pi$ norm-square of this sine mode.  \src{R20}{thm:v20v28-threewave}

Direct differentiation of \eqref{eq:colour-three-wave} gives
\begin{align}
 N_c&=|c_1|^2+|c_2|^2,\qquad
 M_c=|F|^2+\frac{\kappa_c}{g_c}|c_2|^2,\\
 K_c&=-g_c(\bar F\bar c_1c_2+Fc_1\bar c_2),\\
 E_c&=\omega_1|c_1|^2+\omega_2|c_2|^2
                  +\omega_A\frac{g_c}{\kappa_c}|F|^2
 \label{eq:colour-invariants}
\end{align}
as conserved quantities when the channel is nonzero.  The last identity
uses the resonance relation.  With
$r_c=2\varepsilon_N^2/(\pi\omega_A)$ and $F=\sqrt{r_c}\eta$, the
canonical coupling magnitude is
\begin{equation}
 G_c=\sqrt{\kappa_cg_c}
 =\frac{|e|\varepsilon_N(W_1+W_2)}
 {8\sqrt{2\pi\omega_A\omega_1W_1\omega_2W_2}}.
 \label{eq:canonical-colour-coupling}
\end{equation}
The relabeling $c_0=c_2$, $c_-=c_1$, followed by a constant phase change,
brings this system to the same canonical form as the Yukawa channel.

\scope{Equations~\eqref{eq:colour-three-wave}--\eqref{eq:colour-invariants}
are resonant equations for selected flat matter modes.  The cross-metric
response produced by unequal masses has been omitted.  Their exact
conservation laws are therefore conservation laws of the projected
system, not a proof that it is an invariant subsystem of all five
Einstein equations.}


\section{An exact two-mode fermion reduction and a classical Higgs resonance}
For a spatially homogeneous mass modulation, the two-dimensional spinor space spanned by
\begin{equation}
 \psi_0=(\cos z,\ii W\sin z),\qquad
 \psi_-=(W\cos z,-\ii\sin z)
\end{equation}
is invariant. The free frequencies are $\pm\omega$, the states are orthogonal, and both squared $L^2$ norms are $2\pi\omega W$. Writing the amplitudes as $\alpha=(\alpha_0,\alpha_-)^T$,
\begin{equation}
 \ii\dot\alpha=
 \left[\omega\sigma_3+\frac mvh(t)B_0\right]\alpha,
 \qquad
 B_0=\frac1\omega\begin{pmatrix}m&1\\1&-m\end{pmatrix}.
 \label{eq:two-mode-dirac}
\end{equation}
The invariance is exact for the prescribed homogeneous modulation. Coupling to the retained zero-wavenumber, linearized Higgs mode gives
\begin{equation}
 \ddot h+m_h^2h
 =-\frac{2m\varepsilon^2\omega W}{v}\,\alpha^\dagger B_0\alpha.
 \label{eq:two-mode-higgs}
\end{equation}
This is a declared projected Higgs subsystem; the nonlinear Higgs self-potential and nonzero spatial Higgs modes of the full assembled system are not secretly included in~\eqref{eq:two-mode-higgs}.

\begin{theorem}[Conserved quantities of the projected exchange system]
\label{thm:yukawa-exchange}
Equations~\eqref{eq:two-mode-dirac}--\eqref{eq:two-mode-higgs} conserve
\begin{equation}
 N=|\alpha_0|^2+|\alpha_-|^2
\end{equation}
and
\begin{equation}
 E=\frac12\dot h^2+\frac12m_h^2h^2
 +2\varepsilon^2\omega W\,
 \alpha^\dagger\left[\omega\sigma_3+\frac mvhB_0\right]\alpha.
 \label{eq:yukawa-energy}
\end{equation}
\end{theorem}
\begin{proof}
The two-mode Hamiltonian is Hermitian, so $N$ is constant. Its state-dependent energy derivative is $(m/v)\dot h\,\alpha^\dagger B_0\alpha$; the remaining Schr\"odinger terms cancel. Multiplication by $2\varepsilon^2\omega W$ cancels the derivative of the Higgs oscillator energy by~\eqref{eq:two-mode-higgs}.
\end{proof}
\src{R20}{thm:v20v23-subspace,thm:v20v23-coupled}

At $m_h=2\omega$, the resonant amplitude reduction takes the source's three-wave form
\begin{equation}
 \ii\dot c_0=g_Y Hc_-,\qquad
 \ii\dot c_-=g_Y\bar Hc_0,\qquad
 \ii\dot H=\kappa_Yc_0\bar c_-,
 \label{eq:yukawa-resonance}
\end{equation}
where $g_Y=m/(2\omega v)$ and $\kappa_Y=m\varepsilon^2W/(\omega v)$ in its amplitude convention. It has the stated fermion norm and Manley--Rowe invariant $|H|^2-2\varepsilon^2W|c_-|^2$. With an undepleted constant Higgs amplitude $H_0$, the two-level conversion time is $\pi/(2g_Y|H_0|)$. This is a classical coherent-mode resonance; it is not a derivation of a quantum decay rate or a proof of particle production in the full continuum. \cite{R20}

\section{Canonical conversion, detuning, and finite quantum blocks}
\label{sec:canonical-conversion}
The three-wave system retains pump evolution and therefore distinguishes
exact projected conversion from an undepleted external drive.  Assume a
nonzero Yukawa channel and set
\begin{equation}
 H=\sqrt{\kappa_Y/g_Y}\,\eta,\qquad
 G_Y=\sqrt{\kappa_Yg_Y}
     =\frac{m\varepsilon\sqrt W}{\sqrt2\,\omega v}.
 \label{eq:canonical-yukawa-coupling}
\end{equation}
In a phase convention with $G=G_Y>0$,
\begin{equation}
 \ii\dot c_0=G\eta c_-,\qquad
 \ii\dot c_-=G\bar\eta c_0,\qquad
 \ii\dot\eta=Gc_0\bar c_-.
 \label{eq:canonical-triad}
\end{equation}
These are Hamilton's equations in complex canonical coordinates for
\begin{equation}
 \mathcal K=G(\eta\bar c_0c_-+\bar\eta c_0\bar c_-).
 \label{eq:triad-hamiltonian}
\end{equation}
The colour channel has the same form with coupling magnitude $G_c$ and
the relabeling given above.  \src{R20}{prop:v20v26-canonical}

\begin{theorem}[Population quadrature and complete coherent conversion]
\label{thm:triad-quadrature}
The canonical triad conserves
\begin{equation}
 N=|c_0|^2+|c_-|^2,\qquad
 M=|\eta|^2-|c_-|^2,\qquad \mathcal K.
 \label{eq:canonical-triad-invariants}
\end{equation}
Writing $P=|c_-|^2$, its population obeys
\begin{equation}
 \dot P^2=4G^2P(N-P)(M+P)-\mathcal K^2.
 \label{eq:population-cubic}
\end{equation}
For $c_0(0)=1$, $c_-(0)=0$, and real $\eta(0)=\eta_0>0$, put
$k^2=(1+\eta_0^2)^{-1}$ and $u=G\sqrt{1+\eta_0^2}\,t$.  Then
\begin{equation}
 P(t)=\frac{(1-k^2)\operatorname{sn}^2(u,k)}
                      {\operatorname{dn}^2(u,k)}.
 \label{eq:elliptic-conversion}
\end{equation}
It reaches $P=1$ after half of the population period
\begin{equation}
 T=\frac{2\mathbf K(k)}{G\sqrt{1+\eta_0^2}}
  =\frac1G\int_0^1\frac{\dd P}
                  {\sqrt{P(1-P)(P+\eta_0^2)}},
 \label{eq:elliptic-period}
\end{equation}
where $\mathbf K$ is the complete elliptic integral of the first kind.
In the large-drive limit $TG\eta_0\to\pi$.
\end{theorem}
\begin{proof}
Differentiate the three expressions in
\eqref{eq:canonical-triad-invariants} using
\eqref{eq:canonical-triad}.  Their derivatives cancel.  If
$Z=\eta\bar c_0c_-$, then
$\dot P=-2G\operatorname{Im}Z$, $\mathcal K=2G\operatorname{Re}Z$, and
$|Z|^2=P(N-P)(M+P)$, which proves \eqref{eq:population-cubic}.
For the stated initial condition, $N=1$, $M=\eta_0^2$, and
$\mathcal K=0$.  The Jacobi identities
$\operatorname{dn}^2=1-k^2\operatorname{sn}^2$ and
$(\operatorname{sn})'=\operatorname{cn}\operatorname{dn}$ verify
\eqref{eq:elliptic-conversion} on the increasing branch.  Reflection at
the turning points completes the solution.  The first complete transfer
occurs at $u=\mathbf K(k)$; the return gives
\eqref{eq:elliptic-period}.  As $\eta_0\to\infty$, $k\to0$ and
$\mathbf K(k)\to\pi/2$.
\end{proof}
\src{R20}{thm:v20v29-quadrature,thm:v20v29-elliptic}

The empty-seed configuration $c_-=\eta=0$ is an exact classical fixed
point even when $c_0\ne0$.  The positive-seed elliptic solution does not
assert spontaneous motion from that point.  Conversely, a pure boson
configuration with both fermion amplitudes zero is also fixed.  These
facts distinguish classical seeding from the quantum block below.

\subsection{Detuned coherent transfer}
With the frozen detuning \eqref{eq:frozen-gap-detuning} and the rotating
coordinate $c_-=d\,e^{-\ii\delta t}$, the resonant equations take the form
\begin{equation}
 \ii\dot c_0=g_YHd,\qquad
 \ii\dot d=-\delta d+g_Y\bar Hc_0,\qquad
 \ii\dot H=\kappa_Yc_0\bar d.
 \label{eq:detuned-triad}
\end{equation}
They conserve $|c_0|^2+|d|^2$,
$|H|^2-(\kappa_Y/g_Y)|d|^2$, and
\[
 \mathcal K_\delta
 =g_Y(H\bar c_0d+\bar Hc_0\bar d)-\delta|d|^2.
\]
For an undepleted real amplitude $H_0$, write $\Omega=g_YH_0$.
Starting from $c_0=1$, $d=0$, diagonalization of
$\left(\begin{smallmatrix}0&\Omega\\\Omega&-\delta\end{smallmatrix}\right)$
gives
\begin{equation}
 |d(t)|^2=\frac{\Omega^2}{\Omega^2+\delta^2/4}
             \sin^2\!\left(\sqrt{\Omega^2+\delta^2/4}\,t\right).
 \label{eq:detuned-rabi}
\end{equation}
The detuning bounds the conversion depth as well as changing its frequency.
Equation~\eqref{eq:detuned-rabi} is an undepleted-drive solution, unlike
\eqref{eq:elliptic-conversion}.  \src{R20}{thm:v20v27-classical}

\subsection{The projected quantum oscillator--doublet model}
\begin{proposition}[Finite excitation blocks]
\label{prop:quantum-conversion}
Quantize the retained oscillator and a single two-level sector, with
$[a,a^\dagger]=1$, upper state $|0\rangle$, lower state $|-\rangle$,
and $\sigma_-=|-\rangle\langle0|$.  On the finite-occupation core, the rotating-wave Hamiltonian
\begin{equation}
 \widehat H_\delta
 =G(a^\dagger\sigma_-+a\sigma_+)
                        -\delta|-\rangle\langle-|,
 \label{eq:projected-jc}
\end{equation}
commutes with
$\widehat N=a^\dagger a+|0\rangle\langle0|$.
On $\operatorname{span}\{|0,n\rangle,|-,n+1\rangle\}$ it is
\begin{equation}
 \begin{pmatrix}0&G\sqrt{n+1}\\G\sqrt{n+1}&-\delta\end{pmatrix},
 \qquad
 E_{n,\pm}=-\frac\delta2
       \pm\sqrt{\delta^2/4+G^2(n+1)}.
 \label{eq:quantum-doublet-spectrum}
\end{equation}
Starting in $|0,n\rangle$, the transition probability is
\begin{equation}
 \mathbb P_{n\to n+1}(t)
 =\frac{G^2(n+1)}{G^2(n+1)+\delta^2/4}
       \sin^2\!\left(\sqrt{G^2(n+1)+\delta^2/4}\,t\right).
 \label{eq:quantum-conversion-probability}
\end{equation}
The lower vacuum $|-,0\rangle$ is a decoupled eigenstate.  At resonance,
$|0,0\rangle$ converts coherently to $|-,1\rangle$ with probability
$\sin^2(Gt)$.
\end{proposition}
\begin{proof}
Each interaction term raises the oscillator occupation exactly when it
lowers the two-level occupation, or conversely.  Thus $\widehat N$ is
conserved.  The oscillator matrix element is $\sqrt{n+1}$, yielding the
displayed block.  Subtracting $-\delta I/2$ leaves a traceless matrix
whose square is $(G^2(n+1)+\delta^2/4)I$; its exponential gives the
transition probability.  Both interaction terms annihilate $|-,0\rangle$.
\end{proof}
\src{R20}{thm:v20v26-jc,thm:v20v27-jc}

For large coherent occupation, $G_Y\sqrt n=g_YH_0$ under
$H_0^2=(\kappa_Y/g_Y)n$, recovering the classical drive normalization.
The colour channel uses $G_c$ and the corresponding heavy/light levels.
These are quantizations of the isolated rotating-wave models, not
constructions of the complete fermion Fock field, a chiral measure, or
an irreversible decay rate.  In particular, the amplitude inherited
from the classical spatial reduction has not thereby acquired the full
continuum one-particle normalization.


\section{General-wavenumber selection for the retained vertices}
\label{sec:general-triads}
The elementary resonances are members of explicit mode families.  Their
spatial selection and temporal matching are separate tests: a frequency
identity does not ensure that the projected vertex is nonzero.  All
spinors in this section are unit normalized on $[0,2\pi]$ and have
amplitude $\varepsilon_N$.

\subsection{Yukawa sum and difference harmonics}
For an integer $k\ge1$, write
\begin{equation}
 \omega_k=\sqrt{k^2+m^2},\qquad W_k=\omega_k-m,\qquad
 \mathcal N_k=\frac{2\pi\omega_kW_k}{k^2}.
 \label{eq:general-spinor-normalization}
\end{equation}
The free modes at frequencies $\pm\omega_k$ are
\begin{equation}
 \phi_k^+=\mathcal N_k^{-1/2}
       \left(\cos kz,\ii\frac{W_k}{k}\sin kz\right),\qquad
 \phi_k^-=\mathcal N_k^{-1/2}
       \left(\frac{W_k}{k}\cos kz,-\ii\sin kz\right).
 \label{eq:general-yukawa-modes}
\end{equation}
For a Higgs mode
$h_q=\tfrac12(H e^{-\ii\Omega_qt}+\bar H e^{\ii\Omega_qt})\cos qz$,
let $\Omega_q=\sqrt{q^2+m_h^2}$, with $q\ge0$.

\begin{proposition}[Yukawa triad selection and coupling]
\label{prop:general-yukawa-triad}
For a positive-frequency mode $k_0$ and a negative-frequency mode $k_-$,
the projected Yukawa vertex vanishes unless
\begin{equation}
 q=|k_0-k_-|\quad\hbox{or}\quad q=k_0+k_-.
 \label{eq:yukawa-spatial-selection}
\end{equation}
For $q\ge1$ on these branches, its coupling in the amplitude convention
of \eqref{eq:yukawa-resonance} is
\begin{equation}
 g_Y(k_0,k_-;q)
 =\frac{m}{8v}
  \frac{W_{k_-}k_0+\sigma_qW_{k_0}k_-}
       {\sqrt{\omega_{k_0}W_{k_0}\omega_{k_-}W_{k_-}}},
 \label{eq:general-yukawa-coupling}
\end{equation}
where $\sigma_q=+1$ for the difference harmonic and $-1$ for the sum
harmonic.  For $q=0$, necessarily $k_0=k_-=k$, and
\begin{equation}
 g_Y(k,k;0)=\frac{mk}{2v\omega_k}.
 \label{eq:zero-yukawa-coupling}
\end{equation}
The backreaction coefficient is $\kappa_Y=r_Yg_Y$, where
\begin{equation}
 r_Y=\begin{cases}
 4\varepsilon_N^2/(\pi\Omega_q),&q\ge1,\\
 2\varepsilon_N^2/(\pi\Omega_0),&q=0.
 \end{cases}
 \label{eq:general-yukawa-backreaction}
\end{equation}
Temporal resonance further requires
\begin{equation}
 \omega_{k_0}+\omega_{k_-}=\Omega_q,
 \qquad
 m_h^2=(\omega_{k_0}+\omega_{k_-})^2-q^2.
 \label{eq:yukawa-triad-surface}
\end{equation}
\end{proposition}
\begin{proof}
Before normalization, the spatial mass matrix element is the integral of
\[
 \cos qz\left(
 \frac{W_{k_-}}{k_-}\cos(k_-z)\cos(k_0z)
 +\frac{W_{k_0}}{k_0}\sin(k_-z)\sin(k_0z)\right).
\]
Product-to-sum identities give the two harmonics in
\eqref{eq:yukawa-spatial-selection}.  Their integrals are
$\tfrac\pi2(W_{k_-}/k_-+\sigma_qW_{k_0}/k_0)$ when $q\ge1$,
and $2\pi W_k/k$ when $q=0$.  Division by the two spinor norms and
multiplication by $m/(2v)$ gives
\eqref{eq:general-yukawa-coupling}--\eqref{eq:zero-yukawa-coupling}.
Projection of the reciprocal scalar source uses
$\int_0^{2\pi}\cos^2qz\,\dd z=\pi$ for $q\ge1$, but $2\pi$ for
$q=0$; the oscillator factor $2\Omega_q$ then gives
\eqref{eq:general-yukawa-backreaction}.  The remaining time phase is
stationary precisely under \eqref{eq:yukawa-triad-surface}.
\end{proof}
\src{R20}{thm:v20v31-modes,thm:v20v31-selection,thm:v20v31-coupling,prop:v20v31-manifold}

Writing $\kappa_Y=r_Yg_Y$, rather than dividing by $g_Y$, also includes
zero vertices.  In particular, at equal wavenumbers the sum-harmonic
numerator vanishes even when its frequency equation holds.  For the
homogeneous $k=1$ channel,
$\varepsilon_N^2=2\pi\omega W\varepsilon^2$ and $\Omega_0=2\omega$
reduce \eqref{eq:general-yukawa-backreaction} to
$r_Y=2\varepsilon^2W$, exactly the earlier convention.

\subsection{Finite-volume normalization and kinematic resonances}
\label{sec:finite-volume-triads}
The volume dependence is visible before a many-mode limit is attempted.
On a circle of length $L$, take $k_n=2\pi n/L$ and $q_j=2\pi j/L$.
For $n\ge1$, replace the spinor norm-square in
\eqref{eq:general-spinor-normalization} by
\begin{equation}
 \mathcal N_{k,L}=\frac{L\omega_kW_k}{k^2}.
 \label{eq:finite-circle-spinor-norm}
\end{equation}
The spatial selection rule becomes $j=|n_0-n_-|$ or $j=n_0+n_-$.
At fixed physical wavenumbers, the coefficient $g_Y$ in
\eqref{eq:general-yukawa-coupling} or
\eqref{eq:zero-yukawa-coupling} is unchanged: the length in the mass
matrix element cancels the length in the product of the spinor norms.
The reciprocal scalar projection instead gives
\begin{equation}
 \kappa_{Y,L}=r_{Y,L}g_Y,\qquad
 r_{Y,L}=\begin{cases}
 8\varepsilon_N^2/(L\Omega_q),&j\ge1,\\
 4\varepsilon_N^2/(L\Omega_0),&j=0.
 \end{cases}
 \label{eq:finite-circle-yukawa}
\end{equation}
Indeed, the scalar norm-square is $L/2$ for $j\ge1$ and $L$ for $j=0$;
the remaining oscillator and source normalization is the one used in
\eqref{eq:general-yukawa-backreaction}.  Thus the symmetric coupling is
\begin{equation}
 G_{Y,L}=\sqrt{\kappa_{Y,L}g_Y}=|g_Y|\sqrt{r_{Y,L}}
             \propto L^{-1/2}
 \label{eq:finite-circle-coupling-scale}
\end{equation}
for fixed physical momenta, masses, and $\varepsilon_N$.  Holding the
integer mode labels fixed instead changes the physical momenta and does
not give this pure scaling statement.  At $L=2\pi$,
\eqref{eq:finite-circle-yukawa} recovers the earlier convention.
\src{R20}{thm:v20v41-scalings}

\begin{proposition}[Finite-size resonance does not imply a nonzero vertex]
\label{prop:finite-circle-resonance}
For $n_0=n_-=n$ and $j=2n$, the Yukawa frequency condition holds for
every $L$ precisely when $m_h=2m$.  The sum-harmonic vertex of this
family is nevertheless zero.  For massless fermions on the difference
branch, the frequency condition is
\begin{equation}
 m_h=\frac{4\pi\sqrt{n_0n_-}}{L}.
 \label{eq:massless-finite-circle-resonance}
\end{equation}
These are kinematic identities, not nonzero conversion rates.
\end{proposition}
\begin{proof}
For the equal-wavenumber sum branch, squaring
$\sqrt{4k_n^2+m_h^2}=2\sqrt{k_n^2+m^2}$ gives $m_h^2=4m^2$.
The numerator of \eqref{eq:general-yukawa-coupling} is
$W_{k_n}k_n-W_{k_n}k_n=0$.  At $m=0$, subtracting
$(k_{n_0}-k_{n_-})^2$ from $(k_{n_0}+k_{n_-})^2$ gives
$4k_{n_0}k_{n_-}$, proving \eqref{eq:massless-finite-circle-resonance}.
\end{proof}
\src{R20}{thm:v20v41-resonances,thm:v20v31-coupling}

In the minimal Yukawa normalization used here, $m=yv/\sqrt2$ with
$v>0$ also makes the displayed Yukawa coupling vanish when $m=0$.
The massless resonance relation therefore cannot by itself establish
an active Yukawa channel in this model.

\subsection{Canonical box normalization and the second standing parity channel}
\label{sec:canonical-box}
The preceding finite-circle amplitudes are commuting classical
coordinates.  To determine a one-quantum matrix element, introduce an
independent canonical $1+1$-dimensional benchmark on the same circle.
Let $M>0$ denote the scalar mass, $m>0$ the fermion mass, and
$E_p=\sqrt{m^2+p^2}$, $\Omega_q=\sqrt{M^2+q^2}$.  For
$p,q\in(2\pi/L)\Z$, impose the canonical anticommutation relations
(CAR) on $c_p,d_p$ and the canonical commutation relations on $a_q$.
With $\mathsf B=\diag(1,-1)$ define
\begin{equation}
 A_p=\sqrt{\frac{E_p+m}{2E_p}},\quad
 B_p=\frac{p}{\sqrt{2E_p(E_p+m)}},\quad
 U(p)=\binom{A_p}{B_p},\quad V(p)=\binom{B_p}{A_p}.
 \label{eq:canonical-spinors}
\end{equation}
The equal-time fields are
\begin{equation}
 \begin{split}
 \Psi(z)&=\frac1{\sqrt L}\sum_p
       \bigl[c_pU(p)e^{\ii pz}+d_p^\dagger V(p)e^{-\ii pz}\bigr],\\
 H(z)&=\sum_q\frac{a_qe^{\ii qz}+a_q^\dagger e^{-\ii qz}}
                         {\sqrt{2\Omega_qL}},\qquad
 H_I=y\int_0^L H:\!\Psi^\dagger\mathsf B\Psi\!:\,\dd z.
 \end{split}
 \label{eq:canonical-box-fields}
\end{equation}
In particular, $H_0=(a_0+a_0^\dagger)/\sqrt{2ML}$ contains one
homogeneous oscillator, not two independently normalized standing modes.

\begin{proposition}[Canonical vertex and parity completion]
\label{prop:canonical-box-vertex}
For normalized states $|q\rangle=a_q^\dagger|0\rangle$ and
$|p,r\rangle=c_p^\dagger d_r^\dagger|0\rangle$,
\begin{equation}
 \langle p,r|H_I|q\rangle
 =\delta_{q,p+r}\frac{y(A_pB_r-B_pA_r)}{\sqrt{2\Omega_qL}}.
 \label{eq:canonical-box-vertex}
\end{equation}
For $p>0$ the rest-parent pair-creation term is
\begin{equation}
 -h_{L,p}a_0(c_p^\dagger d_{-p}^\dagger
                    -c_{-p}^\dagger d_p^\dagger),\qquad
 h_{L,p}=\frac{yp}{E_p\sqrt{2ML}}.
 \label{eq:canonical-rest-vertex}
\end{equation}
For either species, set
$b_C=(b_p+b_{-p})/\sqrt2$ and
$b_S=(b_p-b_{-p})/(\ii\sqrt2)$.  Then
\begin{equation}
 c_p^\dagger d_{-p}^\dagger-c_{-p}^\dagger d_p^\dagger
 =\ii(c_C^\dagger d_S^\dagger-c_S^\dagger d_C^\dagger).
 \label{eq:standing-parity-completion}
\end{equation}
The two final channels are orthogonal and have equal squared vertex
norm $h_{L,p}^2$.  Retaining only $c_C^\dagger d_S^\dagger$ omits half
of this rest-parent vertex norm and does not define an invariant
subspace for this interaction.
\end{proposition}
\begin{proof}
The pair-creation coefficient in the normal-ordered product is
$U(p)^*\mathsf B V(r)=A_pB_r-B_pA_r$.  The spatial integral gives the
Kronecker momentum constraint.  At $r=-p$, evenness of $A_p$ and
oddness of $B_p$ give $-2A_pB_p=-p/E_p$, proving
\eqref{eq:canonical-rest-vertex}.  The change from traveling to standing
operators is unitary in the CAR algebra.  Expanding its inverse gives
\eqref{eq:standing-parity-completion}; the two occupation states have
unit norm and are orthogonal.  The omitted channel is therefore an
explicit nonzero component of $(I-P)H_IP$ for the one-channel projection.
\end{proof}
\src{R20}{thm:v20v43-vertex,thm:v20v43-leak,cor:v20v43-nonclosure}

For a canonical fermion mode, $b^\dagger b$ is a projection, with
occupation between zero and one.  Rescaling $b$ by $\varepsilon$ changes
its CAR to $\{\varepsilon b,\overline\varepsilon b^\dagger\}
=|\varepsilon|^2I$.  Thus the parameter $\varepsilon_N$ in the
commuting spinor ansatz cannot be selected by an arbitrary rescaling
of a canonical annihilation operator.  Fermion-parity invariant states
also have vanishing one-point spinor expectation.  Equations
\eqref{eq:canonical-box-fields}--\eqref{eq:canonical-box-vertex} supply
a quantum normalization for this declared benchmark; they do not
identify its Hilbert space with the earlier classical amplitude phase
space.\src{R20}{prop:v20v43-epsilon}


\subsection{The simple-root mode-counting Jacobian}
\label{sec:triad-density}
Use a signed harmonic coordinate $q=k_0-s k_-$, where $s=+1$ denotes
the difference branch and $s=-1$ the sum branch.  Since the scalar
frequency depends on $q^2$, this includes either orientation of the
cosine harmonic.  At fixed $k_0$ define
\begin{equation}
 \delta_Y(k_-)=\sqrt{q^2+m_h^2}-\omega_{k_0}-\omega_{k_-}.
 \label{eq:circle-detuning}
\end{equation}
Direct differentiation yields
\begin{equation}
 \delta_Y'(k_-)=-s v_q-v_-,\qquad
 |\delta_Y'(k_-)|=|v_q-v_{\rm phys}|,
 \quad v_q=\frac q{\Omega_q},\quad
 v_{\rm phys}=-s\frac{k_-}{\omega_{k_-}}.
 \label{eq:relative-velocity-jacobian}
\end{equation}
For a simple resonance root, the local continuum mode-counting Jacobian
associated with the spacing $2\pi/L$ is consequently
\begin{equation}
 \varrho_L^{\rm loc}(0)=
 \frac{L}{2\pi|v_q-v_{\rm phys}|}.
 \label{eq:triad-local-density}
\end{equation}
Distinct simple roots contribute additively when the retained family
contains more than one.  This identifies the inverse-relative-velocity
factor and the explicit cancellation of $L$ between $G_{Y,L}^2$ and
$\varrho_L^{\rm loc}(0)$.
\src{R20}{thm:v20v41-density,thm:v20v41-scalings}

\scope{At finite $L$ the detuning distribution is a discrete counting
measure, not a smooth density.  Equation~\eqref{eq:triad-local-density}
is the Jacobian for its local continuum or smoothed approximation.
It does not control multiple roots, equal-velocity thresholds, the order
of long-time and volume limits, or the one-quantum normalization of
$\varepsilon_N$.  The explicit cancellation of the length factor is
therefore not yet a proved infinite-volume decay width.}

\subsection{Fixed-parent two-body phase space}
\label{sec:fixed-parent-density}
A decay experiment fixes the parent momentum rather than a daughter
momentum.  This changes the derivative in the density-of-states
calculation.  Assume $M>2m>0$, fix $q\in\R$, and put
\begin{equation}
 \Omega=\sqrt{M^2+q^2},\quad k_* =\sqrt{M^2/4-m^2},\qquad
 D_q(p)=\Omega-E_p-E_{q-p}.
 \label{eq:fixed-parent-detuning}
\end{equation}

\begin{theorem}[Fixed-parent roots and sequential density limit]
\label{thm:fixed-parent-density}
The two simple roots and their daughter energies are
\begin{equation}
 p_\pm=\frac q2\pm\frac\Omega M k_*,\qquad
 E_\pm=\frac\Omega2\pm\frac qM k_*.
 \label{eq:fixed-parent-roots}
\end{equation}
Their Jacobian obeys
\begin{equation}
 D_q'(p)=-\frac p{E_p}+\frac{q-p}{E_{q-p}},\qquad
 E_+E_-|D_q'(p_\pm)|=Mk_*.
 \label{eq:fixed-parent-jacobian}
\end{equation}
For a smooth compactly supported test function $\phi$ and a smooth
approximate identity $\eta_h$,
\begin{equation}
 \lim_{h\downarrow0}\lim_{L\to\infty}\frac{2\pi}{L}
 \sum_{p\in(2\pi/L)\Z}\phi(p)\eta_h(D_q(p))
 =\sum_{\sigma=\pm}\frac{\phi(p_\sigma)}{|D_q'(p_\sigma)|}.
 \label{eq:fixed-parent-sequential}
\end{equation}
The same sequential statement allows box momenta $q_L\to q$.
\end{theorem}
\begin{proof}
Boosting the rest-frame daughter energy--momentum pairs
$(M/2,\pm k_*)$ gives \eqref{eq:fixed-parent-roots}.  Conversely
energy and momentum conservation, followed by one squaring, yield
exactly these two roots.  Substitution into the derivative gives
\eqref{eq:fixed-parent-jacobian}.  For fixed $h$, the sum is a Riemann
sum with compact support.  After this limit, a change of variable
through each simple root gives the approximate-identity limit.  No
exchange of these two limits is used.
\end{proof}
\src{R20}{thm:v20v42-roots,thm:v20v42-limit}

With the canonical $1+1$-dimensional normalization
\begin{equation}
 \Phi_2=\int_{\R^2}\frac{\dd p\,\dd r}{4E_pE_r}
       \delta(q-p-r)\delta(\Omega-E_p-E_r)=\frac1{2Mk_*},
 \label{eq:canonical-two-body-phase-space}
\end{equation}
a squared invariant amplitude $\mathcal A^2$ equal on both roots gives
the leading-order benchmark
\begin{equation}
 \Gamma_{\rm can}(q)=\frac{\mathcal A^2}{4\Omega Mk_*}
 =\frac M\Omega\Gamma_{\rm can}(0).
 \label{eq:canonical-width}
\end{equation}
For the single two-component Dirac species and interaction in
\eqref{eq:canonical-box-fields}, the spinor contraction gives
$\mathcal A^2=y^2(M^2-4m^2)=4y^2k_*^2$.  Hence
\begin{equation}
 \Gamma_{\rm can}(q)=\frac{y^2k_*}{M\Omega},\qquad
 \Gamma_{\rm can}(0)=\frac{y^2k_*}{M^2},\qquad
 |G_{L,p}|^2=\frac{\mathcal A^2}{8\Omega E_pE_{q-p}L}
 \quad\hbox{on shell}.
 \label{eq:canonical-yukawa-width}
\end{equation}
These factors also follow directly by substituting
\eqref{eq:canonical-box-vertex} into the two-root box sum.
\src{R20}{thm:v20v42-width,thm:v20v42-yukawa,prop:v20v42-box}

\scope{Equation~\eqref{eq:relative-velocity-jacobian} remains correct for
its fixed-daughter scan, whereas \eqref{eq:fixed-parent-jacobian} is
the fixed-parent decay Jacobian.  The widths here concern a canonical
$1+1$-dimensional free-state benchmark, not a four-dimensional Standard
Model width or an exact finite-coupling exponential decay law.  The
simple-root theorem excludes the threshold $M=2m$.  An actual continuum
survival law for a specified projected quantum Hamiltonian is
constructed in Chapter~\ref{ch:renormalized-conversion}.}


\subsection{Colour sum and difference harmonics}
Let $k_1,k_2\ge1$,
$\omega_j=\sqrt{k_j^2+m_j^2}$, and $W_j=\omega_j-m_j$.
The unit-normalized light spinor has components
$(\cos k_1z,\ii(W_1/k_1)\sin k_1z)$; the heavy spinor has components
$(\ii\cos k_2z,(W_2/k_2)\sin k_2z)$, each divided by the square root of its norm-square
factor in \eqref{eq:general-spinor-normalization}.  Retain a colour
$\sin qz$ mode with $q\ge1$ and
$\Omega_{A,q}=\sqrt{q^2+m_A^2}$.  The zero sine mode is identically zero
and is excluded.

\begin{proposition}[Colour triad selection and coupling]
\label{prop:general-colour-triad}
The retained colour vertex is zero outside the following branches.  On
them, define
\begin{equation}
 A_{12}(q)=\begin{cases}
 W_2k_1+W_1k_2,&q=k_1+k_2,\\
 W_2k_1-W_1k_2,&q=k_2-k_1>0,\\
 W_1k_2-W_2k_1,&q=k_1-k_2>0.
 \end{cases}
 \label{eq:colour-selection-numerator}
\end{equation}
The coefficients in \eqref{eq:colour-three-wave} are
\begin{equation}
 g_c(k_1,k_2;q)=\frac{eA_{12}(q)}
                         {16\sqrt{\omega_1W_1\omega_2W_2}},
 \qquad
 \kappa_c=\frac{2\varepsilon_N^2}{\pi\Omega_{A,q}}g_c.
 \label{eq:general-colour-coupling}
\end{equation}
Their temporal resonance requires
\begin{equation}
 \omega_2=\omega_1+\Omega_{A,q},\qquad
 m_A^2=(\omega_2-\omega_1)^2-q^2\ge0.
 \label{eq:colour-triad-surface}
\end{equation}
Thus a positive energy gap and a nonzero spatial matrix element must both
be imposed.
\end{proposition}
\begin{proof}
The two sine--cosine products in the colour current have unnormalized
integrals $\tfrac\pi2(W_2/k_2+W_1/k_1)$ on the sum branch,
$\tfrac\pi2(W_2/k_2-W_1/k_1)$ on the positive $k_2-k_1$ branch, and
the reversed difference on the remaining branch.  Orthogonality kills
all other $q$.  Spinor normalization, the colour-generator coefficient,
and the real-mode factor $1/2$ give
\eqref{eq:general-colour-coupling}.  The current projection uses the sine
norm-square $\pi$.  Temporal phase cancellation gives
\eqref{eq:colour-triad-surface}.
\end{proof}
\src{R20}{thm:v20v32-modes,thm:v20v32-coupling}

Each isolated, nonzero matched triad has the canonical invariants and
quantum excitation blocks of Section~\ref{sec:canonical-conversion}.
The selection statements classify the two displayed Yukawa and colour
vertices within the chosen mode families.  They do not exclude every
other cubic vertex in the full action, establish simultaneous closure
of a network of resonances, or justify an infinite-mode kinetic limit.


\section{Shared-mode resonances and the limit of isolated-triad integrability}
\label{sec:shared-triads}
The explicit vertex list permits finite networks of resonances to be
assembled without treating a shared mode as an independent variable in
each triad.  For a finite set of retained matched triads, sum their
Hamiltonians in the common canonical normalization,
\begin{equation}
 \mathcal K_{\mathcal T}
 =\sum_{\tau\in\mathcal T}G_\tau
  \bigl(\eta_\tau\bar c_{0,\tau}c_{-,\tau}
       +\bar\eta_\tau c_{0,\tau}\bar c_{-,\tau}\bigr).
 \label{eq:many-triad-hamiltonian}
\end{equation}
An amplitude shared by several terms is differentiated in every such
term.  This is a finite rotating-wave Hamiltonian, not a claim that its
retained Fourier space is invariant under the unrestricted field system.
\src{R20}{thm:v20v33-hamiltonian}

\begin{proposition}[The smallest shared-mode system]
\label{prop:shared-triad-invariants}
For two triads sharing $c_0$,
\begin{equation}
 \mathcal K_{ab}=G_a(\eta_a\bar c_0c_a+\bar\eta_a c_0\bar c_a)
               +G_b(\eta_b\bar c_0c_b+\bar\eta_b c_0\bar c_b),
 \label{eq:shared-triad-hamiltonian}
\end{equation}
Hamilton's equations are
\begin{equation}
 \begin{aligned}
 \ii\dot c_0&=G_a\eta_ac_a+G_b\eta_bc_b,\\
 \ii\dot c_j&=G_j\bar\eta_jc_0,\qquad
 \ii\dot\eta_j=G_jc_0\bar c_j,\qquad j\in\{a,b\}.
 \end{aligned}
 \label{eq:shared-triad-flow}
\end{equation}
They preserve $\mathcal K_{ab}$ and the three independent quadratic
charges
\begin{equation}
 N=|c_0|^2+|c_a|^2+|c_b|^2,\qquad
 M_a=|\eta_a|^2-|c_a|^2,\qquad
 M_b=|\eta_b|^2-|c_b|^2.
 \label{eq:shared-triad-charges}
\end{equation}
At regular charge values where the three phase symmetries act freely,
the reduced phase space has real dimension four.  The invariant subspace
$c_b=\eta_b=0$ recovers the isolated triad.
\end{proposition}
\begin{proof}
Differentiate \eqref{eq:shared-triad-hamiltonian} with respect to the
conjugate amplitudes to obtain \eqref{eq:shared-triad-flow}.
The phase transformation
\[
 (c_0,c_a,c_b,\eta_a,\eta_b)\longmapsto
 (e^{\ii\theta_0}c_0,e^{\ii\theta_a}c_a,e^{\ii\theta_b}c_b,
 e^{\ii(\theta_0-\theta_a)}\eta_a,
 e^{\ii(\theta_0-\theta_b)}\eta_b)
\]
leaves the Hamiltonian invariant.  Its three moment maps are
\eqref{eq:shared-triad-charges}; direct differentiation also verifies
conservation.  Fixing and quotienting the three independent charges
reduces ten real dimensions to four.  The two $b$-equations vanish on the
stated subspace.
\end{proof}
\src{R20}{thm:v20v33-invariants,prop:v20v33-reduced}

At $G_a=1$, $G_b=2$, the phase-invariant polynomial first integrals of
total degree at most four are spanned by $1$, $N$, $M_a$, $M_b$,
$\mathcal K_{ab}$, and the six products of two quadratic charges.
This is a finite coefficient statement: the 27 phase-neutral monomials
of degree at most four give a Poisson-derivative matrix of rank 16,
hence nullity 11, and the displayed 11 integrals are independent.
The unrestricted degree-at-most-two coefficient system has dimension
66 and rank 62, leaving the constant and the three quadratic charges.
Thus the tested two-degree-of-freedom reduction has no additional
phase-invariant polynomial integral of degree at most four.  This does
not exclude a higher-degree or nonpolynomial integral, establish the
same rank for every coupling, or prove non-integrability.  It shows why
the single-triad quadrature alone does not solve the shared-mode problem.
\src{R20}{thm:v20v33-invariants,prop:v20v33-reduced}

\section{Second-order statistical conversion and kinetic closure}
\label{sec:kinetic-conversion}
A finite coherent triad and an irreversible statistical rate are different
objects.  The connection begins with a controlled coefficient in a
fixed-time expansion; extending it to a kinetic time scale requires
additional assumptions.

\begin{proposition}[Random-phase second-order response]
\label{prop:random-phase-triad}
Consider the detuned canonical equations
\begin{equation}
 \begin{aligned}
 \ii\dot c_0&=G\eta c_-e^{-\ii\delta t},\\
 \ii\dot c_-&=G\bar\eta c_0e^{\ii\delta t},\qquad
 \ii\dot\eta=Gc_0\bar c_-e^{\ii\delta t}.
 \end{aligned}
 \label{eq:detuned-phase-triad}
\end{equation}
Take fixed initial squared amplitudes $n_0,n_-,n_\eta$ and independent
uniform initial phases.  Define
\begin{equation}
 \mathcal C=n_\eta n_--n_0(n_-+n_\eta),\qquad
 K_t(\delta)=\frac{1-\cos(\delta t)}{\delta^2},\quad K_t(0)=\frac{t^2}{2}.
 \label{eq:triad-collision-coefficient}
\end{equation}
At fixed time, the phase-averaged occupations satisfy
\begin{equation}
 \begin{pmatrix}
 \E|c_0(t)|^2\\ \E|c_-(t)|^2\\ \E|\eta(t)|^2
 \end{pmatrix}
 =\begin{pmatrix}n_0\\n_-\\n_\eta\end{pmatrix}
 +2G^2\mathcal C K_t(\delta)
       \begin{pmatrix}1\\-1\\-1\end{pmatrix}+O(G^3).
 \label{eq:phase-averaged-triad}
\end{equation}
On the exact resonant flow, without averaging,
\begin{equation}
 \frac{\dd^2}{\dd t^2}|c_0|^2
 =2G^2\bigl[|\eta|^2|c_-|^2-|c_0|^2(|c_-|^2+|\eta|^2)\bigr].
 \label{eq:exact-triad-acceleration}
\end{equation}
\end{proposition}
\begin{proof}
Integrate each amplitude equation once and substitute the first-order
amplitudes back into the right-hand sides.  Phase averaging removes
monomials with unequal powers of any initial amplitude and its
conjugate.  The surviving terms use
$|\int_0^t e^{\ii\delta s}\dd s|^2=2K_t(\delta)$ and give
\eqref{eq:phase-averaged-triad}.  Alternatively, on resonance set
$Z=\eta\bar c_0c_-$.  The equations imply
$\dot{|c_0|^2}=2G\operatorname{Im}Z$ and
$\dot Z=\ii G\mathcal C(t)$, proving
\eqref{eq:exact-triad-acceleration}.  The two other occupation changes
also follow from the exact quadratic invariants.
\end{proof}
\src{R20}{thm:v20v34-second-order}

The elementary resonance integral is
\begin{equation}
 \int_{\R}K_t(\delta)\,\dd\delta=\pi|t|.
 \label{eq:resonance-integral}
\end{equation}
For a bounded detuning density $\varrho$ continuous at zero,
$\int K_t(\delta)\varrho(\delta)\dd\delta\sim\pi t\varrho(0)$ as
$t\to+\infty$, by the change of variables $u=\delta t$ and dominated
convergence.  Applied to the second-order coefficient, this produces
the candidate local rate
\begin{equation}
 R=2\pi G^2\varrho(0).
 \label{eq:kinetic-local-rate}
\end{equation}
For a nonconstant coupling or initial occupation profile, the product
$G^2\mathcal C\varrho$ must be evaluated with its own continuity and
integrability assumptions near resonance.  The finite-volume origin
and limitations of the Yukawa mode-counting factor are given in
Section~\ref{sec:triad-density}.
\src{R20}{thm:v20v34-rate}

\begin{proposition}[Conservation and entropy of the closed kinetic model]
\label{prop:kinetic-entropy}
Declare the random-phase, continuum-detuning closure
\begin{equation}
 \dot n_0=R\mathcal C,\qquad
 \dot n_-=-R\mathcal C,\qquad
 \dot n_\eta=-R\mathcal C,\qquad R\ge0.
 \label{eq:closed-triad-kinetics}
\end{equation}
It preserves the nonnegative occupation region, $N=n_0+n_-$, and
$M=n_\eta-n_-$.  Whenever the level energies obey
$\epsilon_0=\epsilon_-+\Omega$, it also preserves
$E=\epsilon_0n_0+\epsilon_-n_-+\Omega n_\eta$.
For positive occupations the entropy functional
$S=\log n_0+\log n_-+\log n_\eta$ satisfies
\begin{equation}
 \dot S=R\frac{\mathcal C^2}{n_0n_-n_\eta}\ge0.
 \label{eq:kinetic-entropy-production}
\end{equation}
For $R>0$, an interior equilibrium obeys
$1/n_0=1/n_-+1/n_\eta$.  Rayleigh--Jeans occupations
$n_i=T/(\epsilon_i-\mu_i)$ satisfy this relation when the frequencies
and chemical potentials are additive across the triad and all
occupations are positive.
\end{proposition}
\begin{proof}
At $n_0=0$, the first derivative is nonnegative; at either daughter
boundary, its occupation derivative is also nonnegative.  The two
charge derivatives cancel.  The energy derivative is
$R\mathcal C(\epsilon_0-\epsilon_--\Omega)$.  Finally,
\[
 \frac1{n_0}-\frac1{n_-}-\frac1{n_\eta}
 =\frac{\mathcal C}{n_0n_-n_\eta},
\]
which gives \eqref{eq:kinetic-entropy-production}.  Substitution of the
Rayleigh--Jeans formula proves the equilibrium statement.
\end{proof}
\src{R20}{thm:v20v34-H}

\scope{Equation~\eqref{eq:phase-averaged-triad} controls a second-order
coefficient at fixed time.  It does not bound the perturbative remainder
uniformly at times of order $G^{-2}$, propagate independent phases, or
justify a many-triad Markov closure.  The entropy identity is exact for
\eqref{eq:closed-triad-kinetics}, not a proof of that closure or of
convergence to equilibrium in the full field system.  The occupations
are classical amplitudes; no quantum Bose enhancement, Pauli blocking,
or spontaneous collision term has been derived here.}

\begin{remark}[Signed levels and physical particle occupations]
\label{rem:signed-kinetic-levels}
For the heavy/light colour channel, the resonance has the positive-energy
form $\omega_2=\omega_1+\Omega_A$.  For the Yukawa spinors in
\eqref{eq:general-yukawa-modes}, the corresponding classical level
energies are instead $\epsilon_0=\omega_{k_0}$ and
$\epsilon_-=-\omega_{k_-}$; their difference is
$\Omega_q=\omega_{k_0}+\omega_{k_-}$.  The canonical equations conserve
the signed level energy, not a sum obtained by replacing the negative
level by its positive magnitude while leaving the amplitudes unchanged.
A negative-frequency classical amplitude is not already a positive-energy
antifermion population.  That interpretation needs a common quantized
field, vacuum, and particle--hole normalization.  The isolated
oscillator--doublet model does not supply them.  Consequently the
positive-energy stress mixture below is not asserted for the Yukawa
amplitudes merely by relabeling their frequencies.
\end{remark}
\src{R20}{thm:v20v31-modes,thm:v20v34-H,thm:v20v36-mixture}

\section{Canonical occupation transfer and Pauli-blocked closure}
\label{sec:canonical-pauli}
The classical collision polynomial in
\eqref{eq:triad-collision-coefficient} is not a canonical fermion
collision term.  The distinction follows from finite operator algebra
before any limit is taken.  Let $a$ be a boson and $c,d$ two CAR modes,
put $N=a^\dagger a$, $F=c^\dagger c$, $D=d^\dagger d$,
$P^\dagger=c^\dagger d^\dagger$, $P=dc$, and $X=aP^\dagger$.
Consider the projected Hamiltonian
\begin{equation}
 H=\Omega N+E_cF+E_dD+gX+\bar gX^\dagger,\qquad
 \Delta=E_c+E_d-\Omega.
 \label{eq:canonical-single-pair}
\end{equation}

\begin{theorem}[Exact CAR transfer and its correlation remainder]
\label{thm:canonical-blocking}
On each finite-excitation sector,
\begin{equation}
 [X^\dagger,X]=\mathcal B,
 \quad \mathcal B=N(1-F)(1-D)-(N+1)FD
              =N-NF-ND-FD.
 \label{eq:canonical-blocking-operator}
\end{equation}
For $n=\langle N\rangle$, $f=\langle F\rangle$,
$d=\langle D\rangle$, and $z=g\langle X\rangle$,
\begin{equation}
 \dot f=\dot d=-\dot n=2\operatorname{Im}z,\qquad
 \dot z=\ii\Delta z+\ii|g|^2\langle\mathcal B\rangle.
 \label{eq:canonical-coherence-equations}
\end{equation}
The number operators $N+F$, $N+D$, and $F-D$ are conserved.
The factorized blocking expression
\begin{equation}
 \mathcal B_{\rm fac}=n(1-f)(1-d)-(n+1)fd
                    =n(1-f-d)-fd
 \label{eq:pauli-collision-factor}
\end{equation}
has the exact error
\begin{equation}
 \langle\mathcal B\rangle-\mathcal B_{\rm fac}
 =-\Cov(N,F)-\Cov(N,D)-\Cov(F,D).
 \label{eq:canonical-correlation-defect}
\end{equation}
\end{theorem}
\begin{proof}
The CAR give $PP^\dagger=(1-F)(1-D)$ and $P^\dagger P=FD$;
$aa^\dagger=N+1$ then gives \eqref{eq:canonical-blocking-operator}.
The Heisenberg commutators yield \eqref{eq:canonical-coherence-equations}
and the three conserved operators.  Expanding the four expectations
in $N-NF-ND-FD$ proves \eqref{eq:canonical-correlation-defect}.
\end{proof}
\src{R20}{thm:v20v44-algebra,thm:v20v44-correlations}

Even number-diagonal initial data do not determine the transfer from
occupations alone.  With $N=1$, the equal mixture of pair states
$00,11$ and the equal mixture of $10,01$ both have $f=d=1/2$ and
$z=0$, but their $\langle\mathcal B\rangle$ are respectively $-1/2$
and $0$.  Their identical factorized value is $-1/4$.  The exact
accelerations $\ddot f=2\Delta\operatorname{Re}z+
2|g|^2\langle\mathcal B\rangle$ therefore differ.

For disjoint fermion pairs sharing the same parent, let
$X_i=aP_i^\dagger$.  In addition to each channel's blocking term,
\begin{equation}
 [X_j^\dagger,X_i]=-P_i^\dagger P_j\quad(i\ne j)
 \label{eq:shared-parent-cross-coherence}
\end{equation}
enters its conversion equation.  A closure must control these
coupling-weighted cross coherences as well as the number covariances.
\src{R20}{thm:v20v44-cross}

\begin{proposition}[Entropy of the declared canonical rate model]
\label{prop:pauli-entropy}
Declare, rather than infer from the finite Hamiltonian,
\begin{equation}
 \dot f=\dot d=-\dot n=R(A-C),\qquad
 A=n(1-f)(1-d),\quad C=(n+1)fd,\quad R\ge0.
 \label{eq:canonical-rate-model}
\end{equation}
This model preserves $n\ge0$, $0\le f,d\le1$ and the above number
charges.  In the interior the Bose--Fermi entropy
\[
 \mathcal S=(n+1)\log(n+1)-n\log n
 -f\log f-(1-f)\log(1-f)-d\log d-(1-d)\log(1-d)
\]
satisfies
\begin{equation}
 \dot{\mathcal S}=R(A-C)\log(A/C)\ge0.
 \label{eq:pauli-entropy}
\end{equation}
\end{proposition}
\begin{proof}
The vector field points inward at each occupation boundary.  Charge
conservation is immediate.  Differentiation of the entropy gives the
displayed logarithm, and monotonicity of $\log$ gives its sign.
\end{proof}
\src{R20}{prop:v20v44-entropy}

\scope{The factors $1-f,1-d$ enforce exclusion and $(n+1)$ is bosonic
stimulation.  The entropy identity establishes consistency of the
specified rate model, not a kinetic limit of the quantum dynamics.
Sections~\ref{sec:weighted-coherence} and
\ref{sec:many-copy-transport} give, respectively, an obstruction to
naive decorrelation and a distinct fixed-time limit in which a specified
empirical correlation remainder does vanish.}


\section{Occupation stresses and anisotropic gravitational response}
\label{sec:kinetic-gravity}
The geometry distinguishes conversion channels through their stresses,
not only their conserved energy.  A unit-normalized positive-frequency
plane Dirac mode has energy weight $\omega_k$ and longitudinal pressure
weight $k^2/\omega_k$, with zero averaged transverse pressure.  For the
quadratic Higgs mode $h=A\cos(qz)\cos(\Omega_qt)$, spatial and temporal
averaging gives
\begin{equation}
 \langle\rho_h\rangle=\frac{c_qA^2\Omega_q^2}{2},\qquad
 \langle p_{z,h}\rangle=\frac{c_qA^2q^2}{2},\qquad
 \langle p_{\perp,h}\rangle=0,
 \label{eq:mode-averaged-stresses}
\end{equation}
where $c_0=1$ and $c_q=1/2$ for a nonzero cosine harmonic.  The scalar
identities follow from the averaged squared time derivative, spatial
derivative, and mass term, together with $\Omega_q^2=q^2+m_h^2$.
The spinor identity follows from the free Dirac bilinears and
$\omega_k^2=k^2+m^2$.  Thus the longitudinal equation-of-state ratio is
$k^2/\omega_k^2$ for either positive-energy mode family.
\src{R20}{thm:v20v36-eos}

For positive-energy occupations in a fixed normalization volume
$\mathcal V$, the corresponding averaged mixture has
\begin{equation}
 \rho=\frac1{\mathcal V}\sum_i n_i\epsilon_i,
 \qquad p_z=\frac1{\mathcal V}\sum_i n_i\frac{k_i^2}{\epsilon_i},
 \qquad p_\perp=0.
 \label{eq:positive-occupation-stress}
\end{equation}
If the retained channel has the positive-energy orientation
$\epsilon_0=\epsilon_-+\Omega$ and the common normalization of
\eqref{eq:closed-triad-kinetics}, then
\begin{equation}
 \dot\rho=0,\qquad
 \dot p_z=\frac{R\mathcal C}{\mathcal V}
 \left(\frac{k_0^2}{\epsilon_0}-\frac{k_-^2}{\epsilon_-}
                         -\frac{q^2}{\Omega}\right).
 \label{eq:kinetic-pressure-drift}
\end{equation}
This follows by differentiation.  Energy-preserving conversion can
therefore change the anisotropic pressure.  The signed-level distinction
in Remark~\ref{rem:signed-kinetic-levels} is essential before using this
identity for a fermionic particle interpretation.
\src{R20}{thm:v20v36-mixture}

\begin{proposition}[Constant-equation-of-state Bianchi--I branch]
\label{prop:anisotropic-power-law}
For $\Lambda=0$, $p_\perp=0$, and $p_z=w\rho$ with fixed $0\le w<1$,
the nonzero-density power-law solution of the homogeneous plane Einstein
equations has, up to constant spatial rescalings and a time origin,
\begin{equation}
 \begin{aligned}
 a(t)&=\alpha\log(t/t_0),&
 \alpha&=\frac{2(1-w)}{w^2+3},\\
 c(t)&=\gamma\log(t/t_0),&
 \gamma&=\frac{2(1+w)}{w^2+3},\\
 \rho(t)&=\frac{\rho_0}{t^2},&
 \rho_0&=\frac{4(1-w)(w+3)}{(w^2+3)^2}.
 \end{aligned}
 \label{eq:anisotropic-power-law}
\end{equation}
At $w=0$ this is the isotropic dust solution.  As $w\uparrow1$, the
family tends to the vacuum Kasner exponents $(\alpha,\gamma)=(0,1)$
with $\rho_0\to0$, not to a nonzero null-dust solution of this power-law
class.
\end{proposition}
\begin{proof}
The constraint and the two homogeneous evolution equations become
\[
 \alpha^2+2\alpha\gamma=\rho_0,\qquad
 \alpha(2\alpha+\gamma-1)=\frac{1-w}{2}\rho_0,\qquad
 \gamma(2\alpha+\gamma-1)=\frac{1+w}{2}\rho_0.
\]
For $\rho_0>0$ and $w<1$, these fix
$\gamma/\alpha=(1+w)/(1-w)$ and give
\eqref{eq:anisotropic-power-law}.  They also imply
$2\alpha+(1+w)\gamma=2$, which is the conservation law
$\dot\rho+(2\dot a+\dot c)\rho+\dot c\,p_z=0$.
The two limits follow directly from the coefficients.
\end{proof}
\src{R20}{thm:v20v36-kasner}

This exact constant-$w$ solution supplies a gravitational comparison for
a statistical matter mixture.  Substituting an instantaneous drifting
$w(t)$ into its exponents would be an adiabatic approximation, not an
exact solution of the kinetic--Einstein system.  Likewise,
\eqref{eq:kinetic-pressure-drift} is a fixed-background identity and
contains no cosmological dilution or redshift term.

\section{Expanding modes and the drift of resonance}
\label{sec:expanding-resonances}
The next step can be made exact at the one-mode level on a prescribed
homogeneous metric
$-\dd t^2+e^{2a(t)}(\dd x^2+\dd y^2)+e^{2c(t)}\dd z^2$.
It extends the massless diagnostic without claiming a self-consistent
statistical solution of the Einstein equations.

\begin{proposition}[Exact massive two-level evolution in an expanding metric]
\label{prop:expanding-dirac-mode}
For constant mass $m$ and comoving wavenumber $k>0$, write
\[
 u=e^{-a-c/2}U(t)\cos(kz),\qquad
 w=e^{-a-c/2}\ii V(t)\sin(kz).
\]
The uncoupled Dirac equation is equivalent to
\begin{equation}
 \ii\frac{\dd}{\dd t}\begin{pmatrix}U\\V\end{pmatrix}
 =\begin{pmatrix}m&\kappa_k(t)\\\kappa_k(t)&-m\end{pmatrix}
       \begin{pmatrix}U\\V\end{pmatrix},
 \qquad \kappa_k(t)=ke^{-c(t)}.
 \label{eq:expanding-dirac-two-level}
\end{equation}
The comoving norm $|U|^2+|V|^2$ is conserved.  Away from a degenerate
zero-energy point, the instantaneous eigenvalues are
$\pm\omega_k(t)=\pm\sqrt{\kappa_k(t)^2+m^2}$.  Real normalized
eigenvectors can be chosen so that
\begin{equation}
 \langle-|\partial_t+\rangle=\frac{m\dot\kappa_k}{2\omega_k^2},
 \qquad \langle+|\partial_t+\rangle=0.
 \label{eq:expanding-branch-coupling}
\end{equation}
The branch-mixing coefficient vanishes for $m=0$ and also in the
spatially homogeneous zero-momentum sector treated separately.
\end{proposition}
\begin{proof}
The prefactor cancels $a_t+c_t/2$ in the spin connection.
The cosine and sine derivatives give the off-diagonal entries in
\eqref{eq:expanding-dirac-two-level}.  Hermiticity gives norm
conservation.  Write the matrix as
$\omega_k(\cos\vartheta\,\sigma_3+\sin\vartheta\,\sigma_1)$, with
$\cos\vartheta=m/\omega_k$ and $\sin\vartheta=\kappa_k/\omega_k$.
The real half-angle eigenvectors have off-diagonal connection
$\dot\vartheta/2=m\dot\kappa_k/(2\omega_k^2)$ and zero diagonal
connection.  At zero momentum the homogeneous Dirac components evolve
diagonally; they are not obtained by counting an identically zero sine
mode as a second independent spatial mode.
\end{proof}
\src{R20}{thm:v20v37-twolevel}

The coefficient \eqref{eq:expanding-branch-coupling} is an exact
positive/negative-frequency mixing diagnostic.  Identifying it with a
particle-production probability requires a specified quantum vacuum and
in/out particle definition.  Those additional objects are not supplied
by the classical two-level reduction.

\begin{proposition}[Scalar adiabatic residual and detuning drift]
\label{prop:expanding-scalar-detuning}
Put $\Theta=2a+c$ and $\Omega_q^2=q^2e^{-2c}+m_h^2>0$.
The quadratic Higgs mode amplitude obeys
\begin{equation}
 \ddot h_q+\dot\Theta\dot h_q+\Omega_q^2h_q=0.
 \label{eq:expanding-higgs-mode}
\end{equation}
The trial amplitude
$h_q^{\rm ad}=e^{-\Theta/2}\Omega_q^{-1/2}
 e^{-\ii\int^t\Omega_q(s)\dd s}$ has exact residual
\begin{equation}
 (\partial_t^2+\dot\Theta\partial_t+\Omega_q^2)h_q^{\rm ad}
 =\left[-\frac{\ddot\Theta}{2}-\frac{\dot\Theta^2}{4}
 +\frac{3\dot\Omega_q^2}{4\Omega_q^2}
 -\frac{\ddot\Omega_q}{2\Omega_q}\right]h_q^{\rm ad}.
 \label{eq:expanding-higgs-residual}
\end{equation}
For the Yukawa kinematic mismatch
$\delta_Y(c)=\Omega_q(c)-\omega_{k_0}(c)-\omega_{k_-}(c)$,
\begin{equation}
 \frac{\partial\delta_Y}{\partial c}
 =e^{-2c}\left(
 \frac{k_0^2}{\omega_{k_0}}+\frac{k_-^2}{\omega_{k_-}}
                         -\frac{q^2}{\Omega_q}\right).
 \label{eq:expanding-yukawa-detuning}
\end{equation}
\end{proposition}
\begin{proof}
Removing the dilution factor $e^{-\Theta/2}$ from
\eqref{eq:expanding-higgs-mode} leaves an oscillator with the additional
potential $-\ddot\Theta/2-\dot\Theta^2/4$.  Differentiating the remaining
adiabatic amplitude gives the two frequency-derivative terms in
\eqref{eq:expanding-higgs-residual}.  Finally,
$\partial_c\sqrt{k^2e^{-2c}+m^2}=-k^2e^{-2c}/\omega_k$ proves
\eqref{eq:expanding-yukawa-detuning}.
\end{proof}
\src{R20}{thm:v20v37-wkb,thm:v20v37-drift}

The detuning derivative is the longitudinal pressure-weight difference
between a boson and a positive-energy fermion pair, with the signs fixed
by $\delta_Y=\Omega_q-\omega_{k_0}-\omega_{k_-}$.  It is not the same
occupation derivative as \eqref{eq:kinetic-pressure-drift} without a
change of channel orientation and particle interpretation.  At the
kinematic level, pressure matching makes the resonance stationary under
longitudinal expansion; otherwise $\dot\delta_Y=\dot c\,\partial_c\delta_Y$.

\scope{The metric in \eqref{eq:expanding-dirac-two-level} is prescribed,
and \eqref{eq:expanding-higgs-residual} is an adiabatic residual, not a
uniform error estimate for a coupled solution.  Nonlinear Higgs terms,
backreaction of the evolving occupations, and a kinetic equation with
redshift must still be assembled in the common stress convention.
An externally prescribed resonance sweep is therefore a reduced
conversion model, not a construction of a quantum kinetic--Einstein
continuum.}

\section{Role of the explicit system in the continuum programme}
The assembled equations exhibit a common mass mechanism, exact gauge and
Yukawa source transfer, on-shell total-stress conservation, and
constraint-compatible Einstein equations.  The selected matter sectors
make this content quantitative: both participating level shifts enter
the detuning, spatial vertices identify the resonances that actually
couple, and the canonical triad distinguishes dynamical pump evolution
from an external drive.  Its oscillator--doublet quantization provides
finite excitation blocks, not an inferred irreversible decay width.
The independent CAR normalization in Section~\ref{sec:canonical-box}
and the fixed-parent density theorem supply the next quantum interface;
the one-parent ultraviolet construction in
Chapter~\ref{ch:renormalized-conversion} then establishes a survival
law for its specified continuum Hamiltonian.

The extensions preserve this same order of construction.  Shared modes
are assembled through a common finite Hamiltonian rather than independent
copies of each triad.  Phase averaging yields an explicit second-order
response and a kinetic closure with an exact entropy identity, while
finite-circle normalization identifies the candidate density factor.
The canonical occupation equations retain Pauli blocking and an explicit
correlation remainder rather than reusing the classical collision
polynomial.  The expanding-metric calculation locates the redshift and branch-mixing
terms.  The two-polarization metric supplies the nonlinear geometric
response to a retained transverse stress.  Each calculation specifies a
part of the dynamics that a continuum construction would have to retain.

Their remaining requirements are different.  A finite polynomial rank
test is not a general non-integrability theorem; a second-order rate
coefficient is not a controlled kinetic limit; a signed classical Dirac
level is not already a positive-energy antiparticle occupation; and the
exact transverse Einstein identities do not supply the uncomputed
non-diagonal fermion system.  The fixed-background pressure identity and
the prescribed expanding-mode equations likewise do not yet form a
self-consistent kinetic--Einstein solution.  The role of these results is
to test and resolve the physical source mechanism, not to replace the
uniform analytic hypotheses of Theorem~\ref{thm:main}.

% ==================== 17_outlook.tex ====================
\chapter{Renormalized quantum conversion and spectral decay}
\label{ch:renormalized-conversion}
The canonical normalization of Section~\ref{sec:canonical-box} determines
a continuum coupling for an explicitly projected quantum system.  This
chapter constructs its ultraviolet limit before discussing irreversible
behavior.  The projection retains one rest parent and its opposite-momentum
fermion--antifermion pairs.  It is the one-excitation block of the
pair-creation interaction, not an invariant sector of the full
normal-ordered Yukawa action.  Within that declared model, self-adjointness,
spectral classification, and a weak-coupling decay limit are established
without replacing the exact dynamics by an assumed rate equation.

\section{The one-parent Hamiltonian and its exact memory equation}
\label{sec:one-parent-memory}
Fix $m,M>0$, $y\in\R\setminus\{0\}$, and write
\begin{equation}
 \omega(p)=2\sqrt{m^2+p^2},\qquad
 v(p)=\frac{yp}{\sqrt{4\pi M}\sqrt{m^2+p^2}},\qquad
 v_\Lambda=\chi_\Lambda v.
 \label{eq:one-parent-profile}
\end{equation}
Here $\chi_\Lambda$ are real even compactly supported cutoffs with
$0\le\chi_\Lambda\le1$ and $\chi_\Lambda\to1$ pointwise.  On
$\mathcal H=\C\oplus L^2(\R,\dd p)$ the cutoff Hamiltonian is
\begin{equation}
 H_\Lambda\binom A b=
 \binom{\mu_\Lambda A+\langle v_\Lambda,b\rangle}
       {\omega b+v_\Lambda A},
 \qquad D(H_\Lambda)=\C\oplus D(\omega).
 \label{eq:one-parent-cutoff}
\end{equation}
The scalar product is conjugate-linear in its first argument.  Since
$v_\Lambda\in L^2$, this is a bounded finite-rank perturbation of the
free self-adjoint operator for every real $\mu_\Lambda$.

Starting from the bare parent $(1,0)$, elimination of the pair amplitude
is exact:
\begin{equation}
 \begin{split}
 b_\Lambda(t,p)&=-\ii v_\Lambda(p)\int_0^t
              e^{-\ii\omega(p)(t-s)}A_\Lambda(s)\,\dd s,\\
 \dot A_\Lambda(t)&=-\ii\mu_\Lambda A_\Lambda(t)
       -\int_0^t K_\Lambda(t-s)A_\Lambda(s)\,\dd s,\\
 K_\Lambda(u)&=\int_{\R}|v_\Lambda(p)|^2e^{-\ii\omega(p)u}\,\dd p.
 \end{split}
 \label{eq:one-parent-volterra}
\end{equation}
The two signed roots of $w=\omega(p)$ give the spectral density
\begin{equation}
 J_\Lambda(w)=c\sqrt{1-\frac{4m^2}{w^2}}\,
       \chi_\Lambda\!\left(\sqrt{w^2/4-m^2}\right)^2,
 \quad w>2m,\qquad c=\frac{y^2}{4\pi M}.
 \label{eq:one-parent-spectral-density}
\end{equation}
Thus $K_\Lambda(u)=\int J_\Lambda(w)e^{-\ii wu}\dd w$.
For $M>2m$, $2\pi J(M)=y^2k_*/M^2$, precisely the rest-parent
benchmark of \eqref{eq:canonical-yukawa-width}.  At a fixed cutoff,
however, the parent survival probability is
$1-\norm{v_\Lambda}^2t^2+o(t^2)$ near zero, not an exact exponential.
\src{R20}{thm:v20v46-memory,thm:v20v46-density}

\section{A subtraction-defined self-adjoint ultraviolet limit}
\label{sec:one-parent-renormalization}
The uncut density $J(w)$ tends to $c$ as $w\to\infty$, so $v\notin L^2$.
The formal self-energy
$\Sigma_\Lambda(z)=\int J_\Lambda(w)/(z-w)\dd w$ has a logarithmic
negative divergence at every fixed real $z_0<2m$.  The convergent object is
\begin{equation}
 S(z)=\int_{2m}^\infty J(w)
       \left(\frac1{z-w}-\frac1{z_0-w}\right)\dd w,
 \qquad \mu_\Lambda=\mu_R-\Sigma_\Lambda(z_0).
 \label{eq:one-parent-subtraction}
\end{equation}
The difference of the two denominators is $O(w^{-2})$.  The counterterm
therefore removes the divergence while retaining a specified real
parameter $\mu_R$.

\begin{theorem}[Norm-resolvent construction and dressed domain]
\label{thm:one-parent-renormalized}
For the cutoffs above, $H_\Lambda$ converges in norm resolvent sense to
a self-adjoint $H_R$, independent of the cutoff profile at fixed
$(z_0,\mu_R)$.  For $z\notin\R$, set
\begin{equation}
 r(z)=(z-\omega)^{-1},\quad h_z=r(z)v\in L^2,\quad
 \ell_z(f)=\langle h_{\bar z},f\rangle,\quad
 D(z)=z-\mu_R-S(z).
 \label{eq:renormalized-resolvent-data}
\end{equation}
Then
\begin{equation}
 (z-H_R)^{-1}\binom a f=
 \binom{D(z)^{-1}(a+\ell_z(f))}
       {r(z)f+h_zD(z)^{-1}(a+\ell_z(f))}.
 \label{eq:renormalized-block-resolvent}
\end{equation}
Its domain is
\begin{equation}
 D(H_R)=\{(A,b):b-h_zA\in D(\omega)\},
 \label{eq:renormalized-dressed-domain}
\end{equation}
independently of the nonreal $z$.  The bare parent $(1,0)$ is a valid
Hilbert-space state but does not belong to $D(H_R)$ when $y\ne0$.
\end{theorem}
\begin{proof}
Schur elimination gives \eqref{eq:renormalized-block-resolvent} at
finite cutoff, with $v,h,S,D$ replaced by their cutoff versions.
Dominated convergence gives $h_{\Lambda,z}\to h_z$ in $L^2$ and
$S_\Lambda(z)\to S(z)$.  The imaginary part obeys
\[
 \operatorname{Im}D(z)=\operatorname{Im}z
 \left(1+\int_{2m}^\infty\frac{J(w)}{|z-w|^2}\,\dd w\right),
 \qquad |D(z)|\ge|\operatorname{Im}z|.
\]
All four resolvent blocks consequently converge in norm.  The limit
satisfies the resolvent and adjoint identities.  It is injective:
a zero output has $A=0$, then $r(z)f=0$, then $a=0$.  Its adjoint is
also injective, so its range is dense.  These identities define the
self-adjoint limit.  The range formula gives
\eqref{eq:renormalized-dressed-domain}; conversely every vector in
that set is obtained by choosing $f=(z-\omega)(b-h_zA)$ and
$a=D(z)A-\ell_z(f)$.  Finally
$\omega h_z=-v+zh_z\notin L^2$, which excludes the bare parent from
the operator domain.
\end{proof}
\src{R20}{thm:v20v47-limit,thm:v20v47-operator,prop:v20v47-domain}

Changing the reference point to another $z_1<2m$ changes the parameter
to $\mu_R^{(1)}=\mu_R+S(z_1)$ and the subtracted function to
$S^{(1)}(z)=S(z)-S(z_1)$.  The denominator, operator, and physical
unitary evolution are unchanged.  The absence of the bare parent from
the operator domain does not prevent its unitary evolution; it means
that the cutoff short-time second moment cannot be passed to the
limit as a finite number.\src{R20}{cor:v20v47-reference}

\section{Bound states, continuum density, and threshold tails}
\label{sec:parent-spectrum}
Let $a=2m$, and define
\begin{equation}
 \mu_{\rm crit}=a-S(a),\qquad
 s(x)=\operatorname{PV}\int_a^\infty J(w)
       \left(\frac1{x-w}-\frac1{z_0-w}\right)\dd w
       \quad(x>a).
 \label{eq:parent-critical-parameter}
\end{equation}
The finite boundary value $S(a)$ follows from the square-root vanishing
of $J$ at $a$.

\begin{theorem}[Complete parent spectral measure]
\label{thm:parent-spectral-measure}
There is exactly one subthreshold eigenvalue $x_b<a$ when
$\mu_R<\mu_{\rm crit}$ and none when $\mu_R\ge\mu_{\rm crit}$.
The parent weight of that eigenvalue is
\begin{equation}
 Z_b=\left(1+\int_a^\infty\frac{J(w)}{(x_b-w)^2}\,\dd w\right)^{-1}.
 \label{eq:parent-bound-weight}
\end{equation}
The rest of the parent spectral measure is absolutely continuous,
with density
\begin{equation}
 \rho(x)=\frac{J(x)}{(x-\mu_R-s(x))^2+\pi^2J(x)^2},\qquad x>a.
 \label{eq:parent-ac-density}
\end{equation}
There is no threshold atom, including at $\mu_R=\mu_{\rm crit}$,
and no singular continuous component.  Thus the parent survival
amplitude is exactly
\begin{equation}
 A_R(t)=\mathbf1_{\{\mu_R<\mu_{\rm crit}\}}Z_be^{-\ii tx_b}
                   +\int_a^\infty e^{-\ii tx}\rho(x)\,\dd x.
 \label{eq:parent-spectral-survival}
\end{equation}
Its limiting survival probability is $Z_b^2$ on the bound-state branch
and zero otherwise.  At subtraction point $z_0=0$,
$\mu_{\rm crit}=2m+c(1+\pi/2)$.
\end{theorem}
\begin{proof}
Below threshold,
$D'(x)=1+\int_a^\infty J(w)/(x-w)^2\dd w>1$,
$D(x)\to-\infty$ as $x\to-\infty$, and
$D(a-)=\mu_{\rm crit}-\mu_R$.  This proves existence and uniqueness
of the root and gives its residue.  The boundary identity
$S(x+\ii0)=s(x)-\ii\pi J(x)$ and Stieltjes inversion give
\eqref{eq:parent-ac-density}.  On every compact interval inside
$(a,\infty)$ its denominator has strictly positive imaginary part,
which excludes an atom or singular component there.  At the critical
endpoint, $\operatorname{Im}D(a+\ii\eta)$ is bounded below by a
positive multiple of $\sqrt\eta$; hence the possible atomic mass,
which is the limit of $-\eta\operatorname{Im}D(a+\ii\eta)^{-1}$,
vanishes.  There can be no remaining singular continuous measure
supported on the single endpoint.  The continuum density is integrable
because it is part of a probability measure.  The Riemann--Lebesgue
lemma applied to \eqref{eq:parent-spectral-survival} gives the limiting
probability.  Substitution $w=a\cosh u$ in $S(a)$ at $z_0=0$ gives
$S(a)=-c(1+\pi/2)$.
\end{proof}
\src{R20}{thm:v20v48-bound,thm:v20v48-ac,thm:v20v48-noatom,cor:v20v48-survival,prop:v20v48-explicit}

The threshold also determines the nonexponential late-time tail.  Set
$j=c\sqrt{2/a}$ and $\delta=\mu_{\rm crit}-\mu_R$.  Expansion of the
explicit density gives
\begin{equation}
 \rho(a+u)=
 \begin{cases}
 (j/\delta^2)u^{1/2}+O(u^{3/2}),&\delta\ne0,\\
 (\pi^2j)^{-1}u^{-1/2}+O(u^{1/2}),&\delta=0.
 \end{cases}
 \label{eq:parent-threshold-density}
\end{equation}
The density away from the threshold has two integrable derivatives.
A smooth partition, two integrations by parts on that remote part,
and the endpoint Fourier integral therefore yield, as $t\to+\infty$,
\begin{equation}
 \begin{split}
 A_R(t)&=\mathbf1_{\{\delta>0\}}Z_be^{-\ii tx_b}
 +e^{-\ii at}\frac{j\sqrt\pi}{2\delta^2}
          e^{-3\pi\ii/4}t^{-3/2}+O(t^{-2}),\quad\delta\ne0,\\
 A_R(t)&=e^{-\ii at}\frac{e^{-\pi\ii/4}}{\pi^{3/2}j}
                       t^{-1/2}+O(t^{-1}),\quad\delta=0.
 \end{split}
 \label{eq:parent-threshold-tails}
\end{equation}
The atom and continuum terms are additive.  For $\delta<0$ the leading
probability is $\pi j^2t^{-3}/(4\delta^4)$; at criticality it is
$t^{-1}/(\pi^3j^2)$.  These fixed-coupling laws concern this spectral
model and do not assert a universal decay exponent for the full theory.
\src{R20}{thm:v20v49-density,thm:v20v49-tails}

\section{A spectral weak-coupling limit}
\label{sec:parent-weak-coupling}
A genuine exponential law appears in a limit different from
$t\to\infty$ at fixed coupling.  Fix $E_0>a$, choose $y_0\ne0$, and
use $y_\lambda=\lambda y_0$, $\mu_R=E_0$.  The ultraviolet construction
is performed first.  Let $J_0,S_0,s_0$ denote the corresponding
reference-coupling functions, and put $J_*=J_0(E_0)>0$ and
$b=s_0(E_0)$.

\begin{theorem}[Uniform absolute-error survival limit]
\label{thm:parent-weak-coupling}
For sufficiently small $\lambda>0$, the parent measure has no bound
atom.  Its rescaled density
$r_\lambda(u)=\lambda^2\rho_\lambda(E_0+\lambda^2u)$, extended by zero
below threshold, satisfies
\begin{equation}
 r_\lambda\longrightarrow
 r_*(u)=\frac{J_*}{(u-b)^2+\pi^2J_*^2}
 \quad\hbox{in }L^1(\R).
 \label{eq:parent-cauchy-limit}
\end{equation}
Consequently
\begin{equation}
 \sup_{\tau\in\R}\left|
 e^{\ii E_0\tau/\lambda^2}A_\lambda(\tau/\lambda^2)
       -e^{-\ii b\tau-\pi J_*|\tau|}\right|\longrightarrow0.
 \label{eq:parent-uniform-decay}
\end{equation}
For forward rescaled time the survival probabilities converge, in
uniform absolute error, to $e^{-2\pi J_*\tau}$.
\end{theorem}
\begin{proof}
Since $E_0>a$ is fixed, the threshold criterion excludes a subthreshold
atom for small $\lambda$.  Substitution into
\eqref{eq:parent-ac-density} gives pointwise convergence to the
Cauchy density in \eqref{eq:parent-cauchy-limit}.  Both densities are
nonnegative with integral one.  Dominated convergence for
$\min(r_\lambda,r_*)\le r_*$ therefore gives
$\norm{r_\lambda-r_*}_1=2-2\int\min(r_\lambda,r_*)\to0$.
The difference of their Fourier transforms is bounded, at every
$\tau$, by this same $L^1$ difference.  The Fourier transform of the
Cauchy density is the displayed exponential.  The probability
conclusion follows since both amplitudes have modulus at most one.
\end{proof}
\src{R20}{thm:v20v51-L1,thm:v20v51-exponential}

The alternative choice $\mu_R=E_0-\lambda^2b$ removes the leading
frequency shift.  At $E_0=M$, the decay rate is
$2\pi J_*=y_0^2k_*/M^2$.  Uniform absolute error does not give a
uniform relative error when the amplitude is very small, so the
fixed-coupling algebraic tails in \eqref{eq:parent-threshold-tails}
are consistent with \eqref{eq:parent-uniform-decay}.  Nor does a
compressed survival limit identify the full unitary group with an
irreversible semigroup or establish a many-body kinetic equation.
\src{R20}{cor:v20v51-shift}

\section{Common-profile parents, coherent decay, and dark subspaces}
\label{sec:matrix-parent-decay}
Several parent states can couple to the same bath profile.  Let
$C:\C^n\to\C^q$ and $K=C^*C$.  Take $q$ identical continuum channels,
common profile $v_0$, and degenerate parent energy $E_0>a$.  The coupling
is $\lambda v_0C$ and the parent counterterm is the full matrix
$-\lambda^2\Sigma_{0,\Lambda}(z_0)K$.  Discarding its off-diagonal
entries would define a different model.

\begin{theorem}[Matrix decay and the invariant dark sector]
\label{thm:matrix-parent-limit}
The common-profile model has a self-adjoint norm-resolvent limit with
compressed parent resolvent
\begin{equation}
 \bigl[(z-E_0)I-\lambda^2S_0(z)K\bigr]^{-1}.
 \label{eq:matrix-parent-resolvent}
\end{equation}
Its rotating, rescaled parent propagator converges in operator norm,
uniformly for $\tau\in\R$, to
\begin{equation}
 \exp\{-\ii bK\tau-\pi J_*K|\tau|\}.
 \label{eq:matrix-common-profile-limit}
\end{equation}
If a bounded Hermitian splitting $\lambda^2B$ is added to the parent
block, the forward limit on every compact interval $0\le\tau\le T$ is
\begin{equation}
 e^{\tau\mathcal A},\qquad
 \mathcal A=-\ii(B+bK)-\gamma K,\quad\gamma=\pi J_*>0.
 \label{eq:split-parent-generator}
\end{equation}
No commutation of $B$ and $K$ is required.  The nondecaying subspace is
\begin{equation}
 \mathcal W=\bigcap_{j=0}^{n-1}\ker(KB^j),
 \label{eq:invariant-dark-subspace}
\end{equation}
the largest $B$-invariant subspace contained in $\ker K$.  It is an
exact reducing subspace of the full model with evolution generated
by $E_0I+\lambda^2B$.  The limiting semigroup is unitary on $\mathcal W$
and decays exponentially on $\mathcal W^\perp$.
\end{theorem}
\begin{proof}
A singular-value decomposition of $C$ reduces the unsplit model to
finitely many scalar constructions from
Theorem~\ref{thm:one-parent-renormalized}; zero singular values give
uncoupled parent states.  The same decomposition and
Theorem~\ref{thm:parent-weak-coupling} prove
\eqref{eq:matrix-common-profile-limit}.  For the splitting, the exact
compressed Duhamel equation in rescaled time is
\[
 T_{\lambda,B}(\tau)=T_{\lambda,0}(\tau)
 -\ii\int_0^\tau T_{\lambda,0}(\tau-s)B T_{\lambda,B}(s)\,\dd s.
\]
If the unsplit error on $[0,T]$ is $\varepsilon_\lambda$, subtraction
of the limiting equation and Gronwall's inequality bound the split
error by $\varepsilon_\lambda(1+\norm B T)e^{\norm B T}$.
Cayley--Hamilton proves that \eqref{eq:invariant-dark-subspace} is
$B$-invariant and maximal.  Hermiticity of $B,K$ makes it reducing.
For any eigenvector $u$ of $\mathcal A$,
$\operatorname{Re}\langle u,\mathcal Au\rangle=-\gamma\langle u,Ku\rangle$.
An eigenvalue on the imaginary axis therefore has $Ku=0$ and $u$ a
$B$-eigenvector, hence $u\in\mathcal W$.  All eigenvalues on the finite
complement have strictly negative real part, which gives exponential
decay there.  Since $Cu=0$ on $\mathcal W$, its exact decoupling from
the bath follows before taking the limit.
\end{proof}
\src{R20}{thm:v20v52-construction,thm:v20v52-limit,thm:v20v53-limit,thm:v20v53-dark}

For example, $C=(1,1)$ has one bright and one dark parent combination.
From initial parent $e_1$, the limiting retained parent probability
approaches $1/2$, whereas the probability to return to that same initial
vector approaches $1/4$.  Retention and return are different observables.
The matrix theorem concerns a common profile, identical thresholds,
finite parent multiplicity, and the specified splitting scale.  It is
not a classification of all multichannel Yukawa baths.


\chapter{Exclusion-preserving collective limits and renormalized energy}
\label{ch:collective-limits}
The one-parent construction is a controlled quantum continuum model,
but it does not supply a finite-density fermion kinetics.  Two further
questions must be separated: whether many pair modes can be replaced
by bosons, and whether a collective evolution that retains exclusion
has an ultraviolet limit.  The first replacement has a quantitative
obstruction; the second can be constructed for a specified
opposite-momentum collective system.  Neither assertion removes the
local stress and common-regulator requirements of the gravitational
closure theorem.

\section{Small channel entries do not imply small coherent feedback}
\label{sec:weighted-coherence}
In the finite one-parent sector write
$|\psi\rangle=A|H\rangle+\sum_i b_i|i\rangle$,
$P=|A|^2$, $p_i=|b_i|^2$, and $P+\sum_i p_i=1$.  Each $|i\rangle$
contains one disjoint fermion pair and no parent.  The exact diagonal
blocking error and off-diagonal pair covariance are
\begin{equation}
 E_i=(2P-1)p_i+p_i^2,\qquad Y_{ij}=\bar b_i b_j.
 \label{eq:one-parent-correlation-entries}
\end{equation}
For couplings $g_i$ define
\begin{equation}
 R_i=|g_i|^2E_i-\sum_{j\ne i}g_i\bar g_jY_{ij},\quad
 d_g=\sum_i|g_i|^2p_i,\quad s_g=\sum_i\bar g_i b_i.
 \label{eq:weighted-residual-definition}
\end{equation}

\begin{proposition}[The aggregate correlation remainder]
\label{prop:weighted-coherence}
One has the exact identity and bounds
\begin{equation}
 \begin{split}
 \sum_iR_i&=2P d_g+\sum_i|g_i|^2p_i^2-|s_g|^2,\\
 \sum_i|g_i|^2|E_i|&\le g_{\max}^2(1-P),\\
 0\le\sum_iR_i+|s_g|^2&\le2g_{\max}^2(1-P),
 \qquad g_{\max}=\max_i|g_i|.
 \end{split}
 \label{eq:weighted-residual-identity}
\end{equation}
For $r$ equal resonant channels $g_i=G/\sqrt r$ and an initially bare
parent, in the rotating frame
\begin{equation}
 A(t)=\cos(Gt),\qquad b_i(t)=-\frac{\ii}{\sqrt r}\sin(Gt).
 \label{eq:collective-one-parent-example}
\end{equation}
Thus every $p_i$ and $Y_{ij}$ is $O(r^{-1})$, while
$\sum_iR_i\to-G^2(1-P)$ at every fixed time.
\end{proposition}
\begin{proof}
The one-excitation occupation identities give
\eqref{eq:one-parent-correlation-entries}.  Summing the second term
of $R_i$ yields $-|s_g|^2+\sum_i|g_i|^2p_i$.  This proves the first
line.  Since $0\le p_i\le1-P$,
$|2P-1+p_i|\le1$, and
$2P\sum_i p_i+\sum_i p_i^2\le2(1-P)$, proving the bounds.
Only the normalized collective pair state couples to the parent in
the equal-channel example, giving the two-dimensional rotation
\eqref{eq:collective-one-parent-example}.  Here
$|s_g|^2=G^2(1-P)$ and the two remaining terms tend to zero.
\end{proof}
\src{R20}{thm:v20v56-residual,prop:v20v56-collective}

In the canonical box $g_{\max}=O(L^{-1/2})$ at a fixed ultraviolet
cutoff.  This suppresses the displayed number-correlation error, but
not the coherently weighted off-diagonal sum.  Indeed elimination of
the pair amplitudes gives
$s_g(t)=-\ii\int_0^t K(t-u)A(u)\dd u$, exactly the memory term of
Section~\ref{sec:one-parent-memory}.  Entrywise decorrelation is
therefore not the norm required for a kinetic closure.

\section{Exact leakage from the exclusion-preserving sector}
\label{sec:hard-pair-leakage}
Consider $r$ disjoint two-state pair modes.  In total excitation $q$,
a basis is indexed by an occupied subset $S$ and parent occupation
$n=q-|S|$.  Embed this space by $J_q$ into a model of $r$ pair bosons,
retaining only occupations zero and one in each pair mode, and set
$\Pi=J_qJ_q^*$.  Let $H_{\rm bos}$ be the bosonic conversion Hamiltonian,
$H_{\rm pair}=J_q^*H_{\rm bos}J_q$, and
$T_q=(I-\Pi)H_{\rm bos}J_q$.

\begin{theorem}[Leakage norm, dilute comparison, and collision obstruction]
\label{thm:hard-pair-leakage}
The leakage operator is diagonal after squaring:
\begin{equation}
 T_q^*T_q|S\rangle
 =2(q-|S|)\sum_{i\in S}|g_i|^2|S\rangle.
 \label{eq:hard-pair-leakage-diagonal}
\end{equation}
Writing $|g|_{(i)}^2$ for the decreasingly ordered weights,
\begin{equation}
 \norm{T_q}^2=2\max_{0\le k\le\min(q,r)}
       (q-k)\sum_{i=1}^k|g|_{(i)}^2
 \le2\lfloor q^2/4\rfloor g_{\max}^2.
 \label{eq:hard-pair-leakage-norm}
\end{equation}
For all real $t$,
\begin{equation}
 \norm{e^{-\ii tH_{\rm bos}}J_q-J_qe^{-\ii tH_{\rm pair}}}
 \le |t|\norm{T_q}.
 \label{eq:hard-pair-duhamel}
\end{equation}
If $g_{\max}\le\lambda C/\sqrt L$, then on
$|t|\le T/\lambda^2$ the right side is bounded by
$CT\sqrt{2\lfloor q^2/4\rfloor}/(\lambda\sqrt L)$.
Thus $q/(\lambda\sqrt L)\to0$ is a sufficient regime for this
comparison, not a fixed-density conclusion.

For equal resonant $g_i=G/\sqrt r$, start with $q$ parents and no pairs.
At $t_* =\pi/(2G)$ the bosonic probability of remaining in the embedded
exclusion-preserving sector is
\begin{equation}
 \frac{(r)_q}{r^q}\quad(q\le r),\qquad 0\quad(q>r),
 \label{eq:pair-collision-probability}
\end{equation}
where $(r)_q=r(r-1)\cdots(r-q+1)$.  The distance from that bosonic state
to any embedded pair-model state is at least the square root of one
minus this probability.  In particular the lower bound tends to one
when $q(q-1)/r\to\infty$.
\end{theorem}
\begin{proof}
Leakage occurs only when conversion creates a second boson in an
already occupied pair mode.  Its matrix element is
$g_i\sqrt{q-|S|}\sqrt2$.  Different choices of $(S,i)$ have distinct
final repeated-occupation states, proving
\eqref{eq:hard-pair-leakage-diagonal}.  Maximizing over the subset
proves \eqref{eq:hard-pair-leakage-norm}; differentiating the
intertwined propagators and integrating gives
\eqref{eq:hard-pair-duhamel}.  In the equal-channel bosonic model the
parent swaps with the normalized collective bath oscillator at $t_*$.
Its $q$ quanta are distributed among $r$ modes with equal weights.
The probability that all labels are distinct is precisely
\eqref{eq:pair-collision-probability}.  Orthogonal projection gives
the distance bound, and
$(r)_q/r^q\le\exp[-q(q-1)/(2r)]$ proves the asymptotic assertion.
\end{proof}
\src{R20}{thm:v20v58-leakage,cor:v20v58-scaling,thm:v20v58-collision}

For $q/r\to\rho>0$, the expected fraction of pair modes containing at
least two bosons at $t_*$ tends to $1-(1+\rho)e^{-\rho}>0$.
If $G=\lambda G_0$, this violation occurs at physical time
$O(\lambda^{-1})$, or rescaled time $\lambda^2t_*=O(\lambda)$.
It rules out a uniform comparison on a kinetic window containing zero
for this equal-resonance family.  It does not prove failure at every
fixed positive rescaled time, nor identify the dispersive canonical
spectrum with an equal-frequency bath.

\section{A compact spectrum with an exclusion-preserving spin at each energy}
\label{sec:compact-spin-spectrum}
An alternative to pair bosonization retains the two-state constraint.
On a compact parameter set $K\subset\R^2$, with
$\xi=(\omega,g)$, a probability measure $\nu$, and real $g$, consider
\begin{equation}
 \begin{split}
 \ii\dot\alpha&=\Omega\alpha+\int_K g s(\xi)\,\dd\nu(\xi),\\
 \ii\partial_t s(\xi)&=\omega s(\xi)-2g\alpha z(\xi),\\
 \partial_t z(\xi)&=2g\operatorname{Im}(\alpha\bar s(\xi)).
 \end{split}
 \label{eq:compact-spin-equations}
\end{equation}
The continuous initial spin fields satisfy $|s|^2+z^2\le1/4$.
Writing $x=z+1/2$ makes this exactly $|s|^2\le x(1-x)$.

\begin{theorem}[Compact-spectrum evolution and quadrature stability]
\label{thm:compact-spin-quadrature}
Equation~\eqref{eq:compact-spin-equations} has a unique global
$C^1$ solution in $\C\times C(K;\R^3)$, preserving each spin length
and the excitation and energy
\begin{equation}
 \begin{split}
 \rho_\nu&=|\alpha|^2+\int_K(z+\tfrac12)\,\dd\nu,\\
 h_\nu&=\Omega|\alpha|^2+
 \int_K\{\omega(z+\tfrac12)+2g\operatorname{Re}(\alpha\bar s)\}\,\dd\nu.
 \end{split}
 \label{eq:compact-spin-invariants}
\end{equation}
Let $g_* =\sup_K|g|$, $A_* =\sqrt{|\alpha_0|^2+1}$, and suppose
$w_0=(\operatorname{Re}s_0,\operatorname{Im}s_0,z_0)$ is Lipschitz
with constant $L_0$.  Two measures $\nu,\widetilde\nu$ with the same
initial fields satisfy
\begin{equation}
 \sup_{|t|\le T}\left\{|\alpha_\nu-\alpha_{\widetilde\nu}|+
                \norm{w_\nu-w_{\widetilde\nu}}_\infty\right\}
 \le B_TTe^{g_*T}W_1(\nu,\widetilde\nu),
 \label{eq:compact-spin-wasserstein}
\end{equation}
where $B_T=1/2+g_*[L_0+(1/2+A_*)T]$ and $W_1$ is the
one-Wasserstein distance on $K$.
\end{theorem}
\begin{proof}
The vector field is locally Lipschitz on the stated Banach space.
Direct differentiation gives \eqref{eq:compact-spin-invariants} and
the pointwise spin-length identity.  They imply
$|w|\le1/2$ and $|\alpha|\le A_*$, hence global existence in both
time directions.  For a prescribed parent trajectory the spin equation
is a real skew-symmetric rotation.  Comparison at two parameters with
its orthogonal propagator gives
$\operatorname{Lip}w(t)\le L_0+(1/2+A_*)|t|$, and therefore
$\operatorname{Lip}(gs)\le B_T$.  Let $a(t)$ be the parent difference
and $b(t)$ the supremum spin difference.  Variation of constants and
the Lipschitz test-function bound for $W_1$ give
$a(t)\le g_*\int_0^t b+B_TtW_1$ and
$b(t)\le g_*\int_0^t a$.  Adding and applying Gronwall proves
\eqref{eq:compact-spin-wasserstein}; time reversal gives negative times.
\end{proof}
\src{R20}{thm:v20v63-global,lem:v20v63-lipschitz,thm:v20v63-stability}

For atomic measures these are finite coupled-spin equations.  Their
quantum collective limit is taken with a fixed number of energy
groups, each population tending to infinity, before refining the
parameter measure.  Symmetrically ordered polynomial moments then
transport along the classical flow on fixed physical time intervals.
A limiting phase-space measure is required for general initial data;
its mean is not silently replaced by one deterministic trajectory.
The fixed-mode many-copy theorem in
Section~\ref{sec:many-copy-transport} gives a momentum-resolved version
of this principle.  Compactness of $K$ and boundedness of $g_*$ are
essential in \eqref{eq:compact-spin-wasserstein}; this estimate by itself
has no ultraviolet-uniform content.
\src{R20}{thm:v20v61-transport,thm:v20v62-transport,cor:v20v63-sequential}

\section{Physical-volume normalization and the ultraviolet source}
\label{sec:collective-volume-normalization}
Counting fractions of pair modes is different from counting states per
spatial volume.  For the canonical rest-parent profile set
\begin{equation}
 \dd\pi(p)=\frac{\dd p}{2\pi},\quad
 \omega(p)=2\sqrt{m^2+p^2},\quad
 h(p)=\frac{yp}{\sqrt{2M}\sqrt{m^2+p^2}}.
 \label{eq:collective-physical-profile}
\end{equation}
There are $r_{L,\Lambda}=2\lfloor\Lambda L/(2\pi)\rfloor+1$ signed
box modes, so $r_{L,\Lambda}/L\to d_\Lambda=\Lambda/\pi$.
At fixed $\Lambda$, their empirical probability measure tends to
$\dd p/(2\Lambda)$.  Rewriting $h(p)/\sqrt L$ as a coupling divided by
$\sqrt r$ therefore gives $G_\Lambda(p)=\sqrt{d_\Lambda}h(p)$.
With $\beta=\sqrt{d_\Lambda}\alpha$, the equations per unit volume are
\begin{equation}
 \begin{split}
 \ii\dot\beta&=\Omega_\Lambda\beta+
                     \int_{-\Lambda}^{\Lambda}hs\,\dd\pi,\\
 \ii\partial_t s&=\omega s+h\beta(1-2x),\qquad
 \partial_t x=2h\operatorname{Im}(\beta\bar s).
 \end{split}
 \label{eq:physical-collective-cutoff}
\end{equation}
The ultraviolet weights are explicitly
\begin{equation}
 \begin{split}
 B_\Lambda=\int_{-\Lambda}^{\Lambda}h^2\,\dd\pi
 &=\frac{y^2}{2\pi M}[\Lambda-m\arctan(\Lambda/m)],\\
 I_\Lambda=\int_{-\Lambda}^{\Lambda}\frac{h^2}{\omega}\,\dd\pi
 &=\frac{y^2}{4\pi M}\left[
 \operatorname{arsinh}(\Lambda/m)-\frac\Lambda{\sqrt{\Lambda^2+m^2}}\right].
 \end{split}
 \label{eq:collective-ultraviolet-weights}
\end{equation}
Excitation bounds control the instantaneous source only by
$\sqrt{B_\Lambda\mathcal N}$; a positive free pair-energy bound
$\mathcal E_{\rm pair}=\int\omega x\dd\pi$ improves this to
$\sqrt{I_\Lambda\mathcal E_{\rm pair}}$.  The latter still diverges
as $\sqrt{\log\Lambda}$.  The source construction includes physical
spin data with bounded pair energy and diverging instantaneous source,
so this is an actual ultraviolet obstruction rather than merely an
inconvenient estimate.  The following construction incorporates the
singular linear part into a self-adjoint propagator instead of assuming
that this instantaneous integral defines a bounded vector field.
\src{R20}{prop:v20v64-volume,thm:v20v64-source}

\section{Global ultraviolet mild dynamics with exclusion}
\label{sec:global-collective-mild}
Let $\mathcal H=\C\oplus L^2(\dd\pi)$, choose a real $\Omega_R$, and
use $h_\Lambda=\chi_\Lambda h$ and
$\Omega_\Lambda=\Omega_R+\int h_\Lambda^2/\omega\dd\pi$.
The linear block $H_\Lambda$ acting on $(\beta,s)$ is
\eqref{eq:one-parent-cutoff} in the measure $\dd\pi$, with subtraction
point zero.  The change $v=h/\sqrt{2\pi}$ identifies its limit with
Theorem~\ref{thm:one-parent-renormalized}.  Denote this self-adjoint
limit by $H_R$.
On the real Hilbert space
$\mathcal X=\mathcal H\oplus L^2(\dd\pi;\R)$ set
\begin{equation}
 \begin{split}
 U_R(t)(\beta,s,x)&=(e^{-\ii tH_R}(\beta,s),x),\\
 F(\beta,s,x)&=(0,\,2\ii h\beta x,\,
                         2h\operatorname{Im}(\beta\bar s)).
 \end{split}
 \label{eq:collective-mild-data}
\end{equation}
The nonlinear term is locally Lipschitz because $h\in L^\infty$;
it does not require $h\in L^2$.

\begin{theorem}[Global physical mild flow and cutoff convergence]
\label{thm:collective-global-mild}
Suppose $\beta_0\in\C$, $x_0\in L^1(\dd\pi)$, and almost everywhere
\begin{equation}
 0\le x_0\le1,\qquad |s_0|^2\le x_0(1-x_0),\qquad
 \mathcal N_0=|\beta_0|^2+\int x_0\,\dd\pi<\infty.
 \label{eq:collective-physical-data}
\end{equation}
There is a unique global mild solution
\begin{equation}
 Y(t)=U_R(t)Y_0+\int_0^t U_R(t-u)F(Y(u))\,\dd u.
 \label{eq:collective-mild-solution}
\end{equation}
It preserves the Pauli region, the pointwise deficit
\begin{equation}
 d(t,p)=x(t,p)-x(t,p)^2-|s(t,p)|^2=d_0(p)\ge0,
 \label{eq:collective-deficit}
\end{equation}
and the excitation $|\beta(t)|^2+\int x(t)\dd\pi=\mathcal N_0$.
In particular $\norm{Y(t)}_{\mathcal X}^2\le2\mathcal N_0$.
For the corresponding cutoff flows with the same physical initial data,
\begin{equation}
 \sup_{|t|\le T}\left(
 \norm{Y_\Lambda(t)-Y(t)}_{\mathcal X}
       +\norm{x_\Lambda(t)-x(t)}_{L^1(\dd\pi)}\right)\longrightarrow0
 \label{eq:collective-mild-cutoff-limit}
\end{equation}
for every finite $T$.  The limiting flow is independent of the cutoff
profile at the fixed subtraction convention.
\end{theorem}
\begin{proof}
The nonlinear maps with $h_\Lambda$ in place of $h$ have common
Lipschitz bounds on bounded subsets of $\mathcal X$ and converge
strongly on each fixed input.  The norm-resolvent linear convergence
gives strong convergence of their unitary groups uniformly on compact
time intervals.  Picard iteration and Gronwall therefore give a unique
local mild limit and local cutoff convergence.  At cutoff, the spin
rotation preserves $x-x^2-|s|^2$, and direct cancellation between
parent and pair equations preserves excitation.  These yield the
uniform bound $\norm{Y_\Lambda(t)}_{\mathcal X}^2\le2\mathcal N_0$.
Almost-everywhere subsequences of the $L^2$ convergence preserve the
Pauli region and deficit.  The identity
$x=x^2+|s|^2+d_0$ improves convergence to $L^1$ through
\[
 \norm{x_\Lambda-x}_1\le
 (\norm{x_\Lambda}_2+\norm x_2)\norm{x_\Lambda-x}_2
 +(\norm{s_\Lambda}_2+\norm s_2)\norm{s_\Lambda-s}_2.
\]
Thus excitation is not lost to an ultraviolet tail.  The common
bounded ball gives continuation and convergence on every compact time
interval, in both time directions, proving the global assertions.
\end{proof}
\src{R20}{lem:v20v66-local,lem:v20v66-cutoff,thm:v20v66-global}

\scope{This is a nonlinear classical collective ultraviolet evolution
for one parent amplitude and disjoint opposite-momentum pairs.  It is
not a simultaneous microscopic quantum, spatial-volume, and ultraviolet
limit.  Its mild equation is meaningful even when $\int hs\dd\pi$ is
not an instantaneous classical source or the state is not in $D(H_R)$.}

\section{Dressed energy and conservation beyond the bare domain}
\label{sec:collective-dressed-energy}
Define $k=h/\omega\in L^2(\dd\pi)$ and $\eta=s+k\beta$.
Although $k$ is square-integrable, $\sqrt\omega k$ is not.  The correct
finite-energy condition consequently involves $\eta$, not the bare
pair amplitude alone.

\begin{theorem}[Closed renormalized form and conserved nonlinear energy]
\label{thm:collective-energy}
The form
\begin{equation}
 q_R(\beta,s)=\Omega_R|\beta|^2+\norm{s+k\beta}_{L^2(\omega\dd\pi)}^2,
 \quad D(q_R)=\{(\beta,s):s+k\beta\in L^2(\omega\dd\pi)\}
 \label{eq:collective-renormalized-form}
\end{equation}
is densely defined, closed, and semibounded, and represents the
self-adjoint $H_R$ above.  Its operator domain has
$\eta\in D(\omega)$ and
\begin{equation}
 H_R(\beta,s)=(\Omega_R\beta+\langle k,\omega\eta\rangle,
                                            \omega\eta).
 \label{eq:collective-dressed-operator}
\end{equation}
For physical data of Theorem~\ref{thm:collective-global-mild} with
$\eta_0,x_0\in L^2(\omega\dd\pi)$ and $d_0\in L^1(\omega\dd\pi)$,
the global flow conserves
\begin{equation}
 \mathcal E_R=
 \Omega_R|\beta|^2+\norm\eta_{L^2(\omega\dd\pi)}^2
 +\norm x_{L^2(\omega\dd\pi)}^2+\int\omega d_0\,\dd\pi.
 \label{eq:collective-renormalized-energy}
\end{equation}
It is bounded below by $-\max(0,-\Omega_R)\mathcal N_0$.
The dressed pair field and occupation are continuous in their weighted
$L^2$ norms.  Sharp-cutoff recovery data give weighted strong
convergence at each fixed time.
\end{theorem}
\begin{proof}
The dressing $(\beta,s)\mapsto(\beta,s+k\beta)$ is bounded and
invertible on $\mathcal H$.  Pulling back the positive weighted form,
and adding the bounded term $\Omega_R|\beta|^2$, proves closedness,
density, and semiboundedness.  Testing the form against arbitrary
variations gives \eqref{eq:collective-dressed-operator}; its resolvent
is \eqref{eq:renormalized-block-resolvent} at $z_0=0$.
At cutoff the conserved energy is
\[
 (\Omega_R+I_\Lambda)|\beta|^2+
 \int\omega x\,\dd\pi+2\operatorname{Re}\int h_\Lambda\beta\bar s\,\dd\pi.
\]
Using $x=|s|^2+x^2+d_0$ completes its square into the cutoff version
of \eqref{eq:collective-renormalized-energy}.  For a sharp cutoff
$\chi_\Lambda=\mathbf1_{\{|p|\le\Lambda\}}$, choose recovery initial
data $(\beta_0,\chi_\Lambda s_0,\chi_\Lambda x_0)$.  Their deficit is
$\chi_\Lambda d_0$ and their dressed field is
$\chi_\Lambda\eta_0$, so their energies converge to $\mathcal E_R(0)$.
The strong unweighted convergence from the mild theorem and lower
semicontinuity of weighted norms give
$\mathcal E_R(t)\le\mathcal E_R(0)$.  In particular the time-$t$ data
remain in the energy class.  Applying the same argument to those data
for time $-t$, and using uniqueness of the global autonomous flow,
gives the reverse inequality.  Equality, weighted weak continuity,
and convergence of the sum of the squared weighted norms imply
weighted strong continuity and the stated fixed-time recovery
convergence.  The lower bound follows directly from the positive terms
and $|\beta|^2\le\mathcal N_0$.
\end{proof}
\src{R20}{thm:v20v67-form,lem:v20v67-recovery,thm:v20v67-conservation}

The energy class is nonempty with $\beta\ne0$: choose a tail
$s=-k\beta$, modify it on a compact interval so that $|s|\le1/2$, and
put $x=(1-\sqrt{1-4|s|^2})/2$, $d_0=0$.  The dressed field then has
compact support and $x=O(|p|^{-2})$.  In contrast, the bare parent
$s=x=0$, $\beta\ne0$, is admissible for the global mild flow but has
infinite renormalized energy.  Mild existence and finite energy are
therefore genuinely different conclusions.

\section{Why global energy is not a local stress construction}
\label{sec:collective-locality-obstruction}
The collective sector above retains a zero-momentum parent and
opposite-momentum pairs.  Let $P_{\rm coll}$ project onto it.  If
$A_k$ transfers nonzero total momentum, covariance under translations
gives
\begin{equation}
 P_{\rm coll}A_kP_{\rm coll}=0\quad(k\ne0).
 \label{eq:collective-momentum-projection}
\end{equation}
Nevertheless
\begin{equation}
 PABP=PAP\,PBP+PA(I-P)BP
 \label{eq:projected-two-source-defect}
\end{equation}
shows why two source insertions cannot be recovered from their
individual projections.  For a single free boson in the zero-momentum
state, the density mode $n_k=\sum_p a_{p+k}^\dagger a_p$ obeys
$Pn_kP=0$ but $Pn_{-k}n_kP=P$.  Intermediate nonzero momentum states
are essential to the response.

There is a second limitation.  For momentum-complete fermion pairs
$P_{pu}^\dagger=c_p^\dagger d_u^\dagger$, the CAR give
\begin{equation}
 [P_{pu},P_{vw}^\dagger]
 =\delta_{pv}\delta_{uw}-\delta_{pv}d_w^\dagger d_u
                              -\delta_{uw}c_v^\dagger c_p.
 \label{eq:overlapping-pair-algebra}
\end{equation}
Pairs sharing a constituent do not form independent two-state spins;
for example $P_{pu}^\dagger P_{pw}^\dagger=0$.  A complete vertex
$\sum_{q,p}g_{q,p}a_qc_p^\dagger d_{q-p}^\dagger+\mathrm{h.c.}$ can
also convert constituents of two opposite-momentum pairs into a moving
parent and an unmatched pair.  Thus the restricted collective sector
need not be invariant even inside total momentum zero.
\src{R20}{thm:v20v68-transfer,thm:v20v68-overlap}

These exact algebraic facts determine the next interface.  Global
conservation of \eqref{eq:collective-renormalized-energy} does not
construct a local stress tensor, its metric derivatives, or the
Einstein equation.  The next chapter restores all retained momentum
pairings through a single fermionic covariance matrix rather than
assigning an independent spin to each overlapping pair.


\chapter{Momentum-resolved covariance and many-copy dynamics}
\label{ch:momentum-covariance}
The overlap identity \eqref{eq:overlapping-pair-algebra} requires a
common fermionic algebra when several retained parent momenta are
present.  A finite momentum-complete system supplies an exact hierarchy,
a Pauli-preserving semiclassical covariance flow, and a controlled
many-copy limit on fixed time intervals.  These results address the
correlation closure of that system.  The spatial ultraviolet limit,
finite-species identification, local stress renormalization, and
Einstein coupling remain separate interfaces.

\section{A finite momentum-complete Hamiltonian and its hierarchy}
\label{sec:momentum-quantum-hierarchy}
Let $a_a$ be finitely many parent oscillators, $c_i$ fermion modes, and
$d_r$ antifermion modes.  The labels have momenta $q_a,p_i,u_r$ and
energies $\Omega_a,E_i,\widetilde E_r$.  Impose
$g_{a;ir}=0$ unless $q_a=p_i+u_r$ and define
\begin{equation}
 \begin{split}
 H={}&\sum_a\Omega_a a_a^\dagger a_a
 +\sum_iE_i c_i^\dagger c_i+\sum_r\widetilde E_r d_r^\dagger d_r\\
 &+\sum_{a,i,r}\left(g_{a;ir}a_ac_i^\dagger d_r^\dagger
                   +\bar g_{a;ir}a_a^\dagger d_rc_i\right).
 \end{split}
 \label{eq:momentum-complete-hamiltonian}
\end{equation}
The excitation
$\mathcal Q=\sum_a a_a^\dagger a_a+\tfrac12\sum_i c_i^\dagger c_i
+\tfrac12\sum_r d_r^\dagger d_r$ is conserved; every fixed-excitation
sector is finite-dimensional.  Define
\begin{equation}
 B_{ab}=\langle a_b^\dagger a_a\rangle,\quad
 C_{ij}=\langle c_j^\dagger c_i\rangle,\quad
 D_{rs}=\langle d_s^\dagger d_r\rangle,\quad
 T_{a;ir}=\langle a_ac_i^\dagger d_r^\dagger\rangle.
 \label{eq:momentum-hierarchy-variables}
\end{equation}
Here $D$ is the antifermion density, not a Dirac operator.

\begin{proposition}[Exact density and conversion hierarchy]
\label{prop:momentum-exact-hierarchy}
The quantum evolution satisfies
\begin{equation}
 \begin{split}
 \ii\dot C_{ij}&=(E_i-E_j)C_{ij}
 +\sum_{a,r}(g_{a;ir}T_{a;jr}-\bar g_{a;jr}\bar T_{a;ir}),\\
 \ii\dot D_{rs}&=(\widetilde E_r-\widetilde E_s)D_{rs}
 +\sum_{a,i}(g_{a;ir}T_{a;is}-\bar g_{a;is}\bar T_{a;ir}),\\
 \ii\dot B_{ab}&=(\Omega_a-\Omega_b)B_{ab}
 -\sum_{i,r}g_{b;ir}T_{a;ir}
 +\sum_{i,r}\bar g_{a;ir}\bar T_{b;ir}.
 \end{split}
 \label{eq:momentum-exact-densities}
\end{equation}
With $P_{ir}=d_rc_i$ and
\[
 Q_{ir;js}=\langle P_{ir}^\dagger P_{js}\rangle,\quad
 M^c_{ab;ji}=\langle a_b^\dagger a_a c_i^\dagger c_j\rangle,\quad
 M^d_{ab;sr}=\langle a_b^\dagger a_a d_r^\dagger d_s\rangle,
\]
the conversion equation is
\begin{equation}
 \begin{split}
 \ii\dot T_{a;ir}={}&-(E_i+\widetilde E_r-\Omega_a)T_{a;ir}
 +\sum_{j,s}\bar g_{a;js}Q_{ir;js}-\sum_b\bar g_{b;ir}B_{ab}\\
 &+\sum_{b,s}\bar g_{b;is}M^d_{ab;sr}
 +\sum_{b,j}\bar g_{b;jr}M^c_{ab;ji}.
 \end{split}
 \label{eq:momentum-exact-conversion}
\end{equation}
The matrices $C,D$ satisfy $0\preceq C,D\preceq I$ at all times.
Excitation, fermion charge, and total momentum are conserved.
\end{proposition}
\begin{proof}
Evaluate the Heisenberg commutators on the finite-excitation sector.
The bilinear CAR and \eqref{eq:overlapping-pair-algebra} give
\eqref{eq:momentum-exact-densities} and
\eqref{eq:momentum-exact-conversion}; the $Q$ term comes from the
parent commutator and the $M^c,M^d$ terms from the shared-constituent
pair commutator.  For any vector $v$, the CAR imply
$0\le\langle c(v)^\dagger c(v)\rangle\le\norm v^2$, and similarly
for $d$, proving the matrix Pauli bounds.  Each vertex removes one
parent and creates one particle and one antiparticle with the same
total momentum.  This proves the three conservation statements.
\end{proof}
\src{R20}{thm:v20v69-densities,thm:v20v69-conversion,prop:v20v69-conservation}

Replacing $Q_{ir;js}$ by $C_{ji}D_{sr}$,
$M^c_{ab;ji}$ by $B_{ab}C_{ji}$, and
$M^d_{ab;sr}$ by $B_{ab}D_{sr}$ leaves explicit fourth-order
correlation defects in \eqref{eq:momentum-exact-conversion}.
Those replacements are not a general Wick rule when anomalous pair
covariances are present.  They also do not eliminate the overlap
terms.  A more economical closure retains the full anomalous
covariance in one particle--hole matrix.

\section{A common Pauli-preserving covariance flow}
\label{sec:pauli-covariance-flow}
Let $n_c,n_d$ be the numbers of fermion and antifermion modes and
introduce the nonredundant CAR vector
\begin{equation}
 f=(c_1,\ldots,c_{n_c},d_1^\dagger,\ldots,d_{n_d}^\dagger)^T.
 \label{eq:particle-hole-car-vector}
\end{equation}
Its covariance is
\begin{equation}
 \Gamma_{ij}=\langle f_j^\dagger f_i\rangle,\qquad
 \Gamma=\begin{pmatrix}C&\kappa\\\kappa^*&I-D^T\end{pmatrix},
 \quad \kappa_{ir}=\langle d_rc_i\rangle.
 \label{eq:full-particle-hole-covariance}
\end{equation}
There is no redundant doubling in this definition.  Write
$E=\diag(E_i)$, $\widetilde E=\diag(\widetilde E_r)$ and, for
classical parent coordinates $\alpha_a$,
\begin{equation}
 \Delta_{ir}(\alpha)=\sum_a g_{a;ir}\alpha_a,\qquad
 \mathsf h(\alpha)=
 \begin{pmatrix}E&\Delta(\alpha)\\\Delta(\alpha)^*&-\widetilde E^T\end{pmatrix}.
 \label{eq:covariance-one-body-hamiltonian}
\end{equation}
The coupled finite-mode equations are
\begin{equation}
 \ii\dot\alpha_a=\Omega_a\alpha_a+
                    \sum_{i,r}\bar g_{a;ir}\kappa_{ir},\qquad
 \ii\dot\Gamma=[\mathsf h(\alpha),\Gamma].
 \label{eq:pauli-covariance-flow}
\end{equation}

\begin{theorem}[Global covariance evolution and conserved quantities]
\label{thm:pauli-covariance-flow}
For every $\alpha_0$ and Hermitian $0\preceq\Gamma_0\preceq I$,
\eqref{eq:pauli-covariance-flow} has a unique global smooth solution.
It preserves the entire spectrum of $\Gamma$ and hence the matrix
Pauli bounds, including exclusion between overlapping pairs.  It
conserves
\begin{equation}
 \begin{split}
 \mathcal N_{\rm sc}&=\sum_a|\alpha_a|^2+\tfrac12(\Tr C+\Tr D),\\
 \mathcal Q_{\rm f}&=\Tr C-\Tr D,\\
 \mathcal E_{\rm sc}&=\sum_a\Omega_a|\alpha_a|^2+
 \Tr(EC)+\Tr(\widetilde E D)+2\operatorname{Re}\Tr(\Delta^*\kappa).
 \end{split}
 \label{eq:covariance-invariants}
\end{equation}
With the declared momentum support of the vertex, it also conserves
$\sum_aq_a|\alpha_a|^2+\sum_i p_iC_{ii}+\sum_r u_rD_{rr}$.
The anomalous block equation is
\begin{equation}
 \ii\dot\kappa=E\kappa+\kappa\widetilde E^T
                  +\Delta(I-D^T)-C\Delta.
 \label{eq:matrix-pauli-blocking}
\end{equation}
\end{theorem}
\begin{proof}
The finite polynomial vector field gives local existence and
uniqueness.  Solving $\ii\dot U=\mathsf h(\alpha)U$ along this
solution gives a unitary propagator and
$\Gamma(t)=U(t)\Gamma_0U(t)^*$.  This proves isospectrality and the
Pauli bounds.  Block multiplication yields
\eqref{eq:matrix-pauli-blocking}; tracing the diagonal blocks and
using the parent equation proves conservation of excitation and
charge.  Direct differentiation of the energy cancels the parent
variation against the off-diagonal covariance variation; equivalently
the finite flow is Hamiltonian for the displayed energy.  The same
calculation with momentum weights cancels vertex by vertex because
$q_a=p_i+u_r$.  Excitation bounds every parent amplitude, while
$0\preceq\Gamma\preceq I$ bounds the covariance.  The solution remains
in a compact set, giving global smooth continuation.
\end{proof}
\src{R20}{thm:v20v71-pauli,thm:v20v71-conservation}

For a prescribed classical $\alpha(t)$, covariance evolution under the
quadratic fermion Hamiltonian is exact for arbitrary states, not only
quasifree states.  What requires justification is replacing quantum
parents by those classical coordinates.  Let
$\alpha_a=\langle a_a\rangle$ in the exact quantum model, write
$\widehat{\mathsf h}=\mathsf h(\alpha)+\delta\mathsf h$, and let the
off-diagonal block of $\delta\mathsf h$ be
$\sum_a g_a(a_a-\alpha_a)$.  Then
\begin{equation}
 \begin{split}
 \ii\dot\Gamma&=[\mathsf h(\alpha),\Gamma]+\mathcal R,\\
 \mathcal R_{ij}&=\sum_k\left[
 \langle\delta\mathsf h_{ik}f_j^\dagger f_k\rangle
 -\langle\delta\mathsf h_{kj}f_k^\dagger f_i\rangle\right].
 \end{split}
 \label{eq:covariance-mixed-residual}
\end{equation}
Thus \eqref{eq:pauli-covariance-flow} suppresses a specific
parent--fermion correlation remainder.  Number-state parents with
empty fermions have $\alpha=\kappa=0$ initially, for which a
single-trajectory closure would incorrectly predict no conversion.
A phase-space distribution, or a limit with concentrated initial
empirical data, is needed.  The following theorem supplies both
possibilities for a declared many-copy scaling.

\section{Fixed-time transport in a many-copy quantum limit}
\label{sec:many-copy-transport}
Keep the spatial modes and vertex coefficients fixed.  Introduce $N$
fermion copies, indexed by $\ell$, with CAR across all copies; even
bilinears in distinct copies commute.  The parents are shared.  Define
\begin{equation}
 \begin{split}
 H_N={}&\sum_a\Omega_a a_a^\dagger a_a+
 \sum_{\ell=1}^N\left(\sum_i E_i c_{\ell i}^\dagger c_{\ell i}
                  +\sum_r\widetilde E_r d_{\ell r}^\dagger d_{\ell r}\right)\\
 &+\frac1{\sqrt N}\sum_{\ell,a,i,r}
 \left(g_{a;ir}a_ac_{\ell i}^\dagger d_{\ell r}^\dagger
              +\bar g_{a;ir}a_a^\dagger d_{\ell r}c_{\ell i}\right).
 \end{split}
 \label{eq:many-copy-hamiltonian}
\end{equation}
The excitation
$\mathcal Q_N=\sum_a a_a^\dagger a_a+\tfrac12\sum_{\ell,i}
 c_{\ell i}^\dagger c_{\ell i}+\tfrac12\sum_{\ell,r}
 d_{\ell r}^\dagger d_{\ell r}$ is conserved.  Set
\begin{equation}
 \widehat\alpha_a=\frac{a_a}{\sqrt N},\qquad
 \widehat\Gamma_{ij}=\frac1N\sum_{\ell=1}^N
                         f_{\ell j}^\dagger f_{\ell i}.
 \label{eq:many-copy-empirical-coordinates}
\end{equation}
The upper-right block is $\widehat\kappa$.  These empirical operators
satisfy
\begin{equation}
 \begin{split}
 [\widehat\alpha_a,\widehat\alpha_b^\dagger]&=\delta_{ab}/N,
 \qquad [\widehat\alpha_a,\widehat\Gamma_{ij}]=0,\\
 [\widehat\Gamma_{ij},\widehat\Gamma_{kl}]&=
 \frac1N(\delta_{il}\widehat\Gamma_{kj}
                   -\delta_{kj}\widehat\Gamma_{il}),\\
 \ii\dot{\widehat\alpha}_a&=\Omega_a\widehat\alpha_a+
                            \sum_{i,r}\bar g_{a;ir}\widehat\kappa_{ir},\\
 \ii\dot{\widehat\Gamma}&=
               [\mathsf h(\widehat\alpha),\widehat\Gamma].
 \end{split}
 \label{eq:many-copy-exact-algebra}
\end{equation}
The last equation uses operator products; it is not yet a closed
expectation equation.  For every fixed complex vector $v$ the CAR give
$0\le v^*\widehat\Gamma v\le\norm v^2I$.  These directional operator
inequalities are sufficient below.  No stronger positivity claim about
an operator-entry matrix on a tensor product is needed.
\src{R20}{thm:v20v72-algebra}

Let $x$ denote the independent real coordinates of $(\alpha,\Gamma)$,
and let $\operatorname{Sym}_N(p)$ be the completely symmetrically
ordered operator associated with a real-coordinate polynomial $p$.
For $R>0$ put
\begin{equation}
 K_R=\left\{(\alpha,\Gamma):\Gamma=\Gamma^*,\ 0\preceq\Gamma\preceq I,
 \quad\sum_a|\alpha_a|^2+\tfrac12(\Tr C+\Tr D)\le R\right\}.
 \label{eq:many-copy-compact-state-space}
\end{equation}
This compact set is invariant under the global flow $\Phi_t$ of
Theorem~\ref{thm:pauli-covariance-flow}.

\begin{theorem}[Transport of empirical polynomial moments]
\label{thm:many-copy-transport}
Let $\rho_N$ be density matrices supported on
$\mathcal Q_N\le RN$.  Suppose a probability measure $\mu_0$ on $K_R$
satisfies, for every polynomial $p$,
\begin{equation}
 \Tr(\rho_N\operatorname{Sym}_N(p))\longrightarrow\int_{K_R}p\,\dd\mu_0.
 \label{eq:many-copy-initial-moments}
\end{equation}
Then for every such $p$ and every finite $T$,
\begin{equation}
 \sup_{|t|\le T}\left|
 \Tr(e^{-\ii tH_N}\rho_Ne^{\ii tH_N}\operatorname{Sym}_N(p))
       -\int_{K_R}p(\Phi_t(x))\,\dd\mu_0(x)\right|\longrightarrow0.
 \label{eq:many-copy-moment-transport}
\end{equation}
For $\mu_0=\delta_{x_0}$, the limiting moments concentrate on
$\Phi_t(x_0)$, and the empirical version of
\eqref{eq:covariance-mixed-residual} tends to zero on compact time
intervals.  For a nonconcentrated $\mu_0$, the limit is the transported
measure $(\Phi_t)_*\mu_0$, not in general the trajectory starting from
its mean.
\end{theorem}
\begin{proof}
A word of fixed degree $k$ in the real empirical coordinates changes
total excitation by at most $k$.  On the resulting sectors,
$\norm{a_a/\sqrt N}\le\sqrt{R+k/N}$ and
$\norm{a_a^\dagger/\sqrt N}\le\sqrt{R+(k+1)/N}$ as maps into the full
Hilbert space.  Fermionic empirical coordinates are bounded uniformly
in $N$.  Successive application bounds each fixed word uniformly;
excitation conservation makes the same bounds available for all $t$.
No compressed oscillator is inserted into the Heisenberg equations.

The commutators in \eqref{eq:many-copy-exact-algebra} show that exchanging
adjacent factors in any fixed word has restricted norm $O(N^{-1})$.
Differentiating a word by the product rule therefore gives
\begin{equation}
 \frac{\dd}{\dd t}\Tr(\rho_N(t)\operatorname{Sym}_N(p))
 =\Tr(\rho_N(t)\operatorname{Sym}_N(F\cdot\nabla p))+O(N^{-1}),
 \label{eq:many-copy-generator-bound}
\end{equation}
where $F$ is the real polynomial vector field of
\eqref{eq:pauli-covariance-flow}.  For fixed $p,R$ and the fixed
Hamiltonian coefficients, the error is uniform on compact times.
This proves the generator approximation without assuming moment
factorization.

The word bounds and \eqref{eq:many-copy-generator-bound} give boundedness
and equicontinuity of all polynomial moments.  A diagonal
Arzel\`a--Ascoli argument gives subsequential limits $L_t(p)$.
Reordering products identifies
$L_t(\bar p p)$ with a limit of nonnegative operator expectations, so
$L_t$ is normalized and positive on squares.  The coordinate bounds
also give $L_t(x_j^2|p|^2)\le M^2L_t(|p|^2)$ for one fixed sufficiently
large $M$.  In the polynomial Hilbert-space representation, multiplication
by each coordinate is therefore bounded self-adjoint; these multiplication
operators commute.  Their joint spectral measure represents $L_t$ as
a probability measure on a compact box.  The tested directional CAR
inequalities force $0\preceq\Gamma\preceq I$ on its support.  Testing
the excitation inequality, whose symmetric-ordering correction is
$O(N^{-1})$, forces the support into $K_R$.

Pass to the limit in the integrated generator identity to obtain
\[
 \int p\,\dd\mu_t-\int p\,\dd\mu_0
   =\int_0^t\int F\cdot\nabla p\,\dd\mu_s\,\dd s.
\]
Polynomial approximation of smooth tests and their first derivatives
on the compact set gives the weak transport equation.  The globally
smooth finite-dimensional flow is unique on $K_R$, so this equation
has the unique solution $\mu_t=(\Phi_t)_*\mu_0$.  Every subsequence has
the same limit, proving convergence of the whole sequence, uniformly
on compact times.  For a point measure the first and second moments
factorize in the limit; the finite sums of mixed empirical covariances
in \eqref{eq:covariance-mixed-residual} therefore vanish.
\end{proof}
\src{R20}{thm:v20v72-transport}

\subsection{Nonempty initial-state class and limit order}
The hypotheses admit concrete states with nonzero pair covariance.
Diagonalize any $0\preceq\Gamma_0\preceq I$ in the particle--hole CAR
space, take independent occupations with those eigenvalue probabilities,
and rotate back.  This gives an even finite-mode state with covariance
$\Gamma_0$.  The product of $N$ identical copies has empirical polynomial
moments converging to the point $\Gamma_0$: terms with a repeated copy
index occupy an $O(N^{-1})$ fraction at each fixed degree.  Shared
parents in coherent states of amplitudes $\sqrt N\alpha_0$ have the
corresponding point moments.  To enforce the strict excitation support,
choose $R>\sum_a|\alpha_{0,a}|^2+(n_c+n_d)/2$ and project parent number
onto
$\sum_a a_a^\dagger a_a\le\lfloor N(R-(n_c+n_d)/2)\rfloor$.
The removed Poisson tail is exponentially small, so fixed polynomial
moments are unchanged in the limit.  Mixtures produce nonpoint initial
measures as well.

\scope{The parameter $N$ is an auxiliary number of identical fermion
copies, not a count of physical Standard Model generations or species.
The theorem fixes the spatial momentum set, the couplings, and the
physical time interval before $N\to\infty$.  It proves neither a
finite-species closure nor a result on kinetic times growing with $N$.
The ultraviolet limit of Chapter~\ref{ch:collective-limits} belongs to
a different, restricted collective model; the two limits have not been
shown to commute.  A common local stress packet and the ten hypotheses
of Theorem~\ref{thm:main} are still needed for the
Standard Model--Einstein endpoint.}


\chapter{Synthesis: what the programme now reduces to}
\label{ch:outlook}
\section{One construction, several genuinely different completions}
The finite renewal construction is not merely an analogy with spacetime and internal symmetry. Under its explicit operational hypotheses it produces a persistent carrier, positive all-order word data, a calibrated local geometric branch, the faithful Standard Model representation, and a common action with exact source identities. These are strong finite statements. They make it possible to ask a definite continuum question rather than postulate the answer as the meaning of a tensor.

The analytic programme then divides into several interacting tasks. The chiral sector requires a rough and globally glued physical line with exact zero-mode handling. The probabilistic sector requires a normalized positive law, joint compactness, and locally normal physical limits. The renormalization sector requires a bounded common trajectory and source-preserving summable comparison. The geometric sector requires nondegenerate ordered contractions and vanishing Bianchi/intertwining defects. The reconstruction sector requires a physical reflection, common domains, locality, and an honest clock. The dynamical sector requires the actual Einstein residual to vanish. None can be inferred solely from success in another.

The fixed-spacing construction in
Theorem~\ref{thm:fixed-spacing-dlr} closes a concrete intermediate existence
problem: the positive conformal and Higgs variables, including the restored
common conformal mode, coexist with compact gauge links in a homogeneous
Gibbs state.  Its proof supplies local moment bounds rather than merely a
formal stable exponent.  The dependence of those bounds on $\kappa_g$
explains why this thermodynamic result does not yet solve the canonical
scaling problem.  Equations~\eqref{eq:third-shell-bound} and
\eqref{eq:normalized-mass-derivative} similarly expose the quantities that
must become summable or uniform before a fixed-scale identity can be
transported cofinally.

The boundary construction is similarly resolved into concrete inputs.
The full-face identity removes a specific rotational loss before block
mixing.  A faithful old-clock lift and actual detail dynamics replace
one-step erasure, with an explicit angle and commutator cost when the
pieces do not commute.  Stable Higgs reference forms and elliptic
holonomy propagation give genuine finite heat mechanisms.  None of
these mechanisms permits a missing chiral, gravitational, or interface
factor to be discarded, and none makes a metric-dependent vacuum
normalization invisible to the source.

\section{Constructed model limits and their relation to the endpoint}
The new continuum constructions close specified model problems rather
than merely naming further estimates.  Theorem~\ref{thm:one-parent-renormalized}
constructs a self-adjoint ultraviolet Hamiltonian from the canonical
rest-parent vertex, including its dressed domain.  Its spectral measure
and threshold behavior give an exact survival representation.  Taking
weak coupling only after that construction gives the uniform
absolute-error exponential law of
Theorem~\ref{thm:parent-weak-coupling}.  The common-profile extension
retains matrix interference and distinguishes the invariant dark
subspace from the instantaneous null space of a decay matrix.

At finite density, exclusion and coherent feedback cannot be discarded
because their individual matrix entries become small.  The aggregate
identity in Proposition~\ref{prop:weighted-coherence} and the exact
leakage norm in Theorem~\ref{thm:hard-pair-leakage} quantify those two
obstructions.  The alternative construction retains the Pauli region
and yields the global ultraviolet mild flow of
Theorem~\ref{thm:collective-global-mild}.  Its dressed energy theorem
separates a nonempty finite-energy class from bare parent data that
remain meaningful only at the mild level.

Momentum completion changes the fermionic algebra: shared constituents
require a single covariance rather than independent pair spins.
Theorem~\ref{thm:pauli-covariance-flow} preserves that covariance's full
spectrum, and Theorem~\ref{thm:many-copy-transport} derives its
fixed-time moment transport under an explicit many-copy scaling.
A point initial phase-space measure gives a deterministic closure;
a nonpoint measure remains a transported distribution.  The spatial
cutoff and physical time are fixed in that limit, and the auxiliary
copy number does not identify physical Standard Model species.

The analytic chain is sharpened without conflating these models with it.
Trace-class proper-time convergence must be differentiated to supply
stress, while vanishing constrained and Ward components must be joined
by a vanishing reduced Einstein residual.  Theorems
\ref{thm:proper-time-c1} and \ref{thm:three-defect} express those
requirements quantitatively.  Global energy conservation in a projected
momentum sector is not a local source construction, as
\eqref{eq:projected-two-source-defect} demonstrates.

\section{The obstructions are constructive information}
The conformal calculation rules out the unrestricted real Euclidean density as a coercive starting point. The annihilating-detail calculation rules out a convenient but unphysical refinement clock. The heat-link example shows that exact projectivity can concentrate away from every fixed smooth chart. The cumulant theorem shows that strong mixing can yield a free limit unless a separately protected interaction survives. The compact constraint and quartic calculations show that low-mode positivity and a finite Fock seed do not automatically persist in an unrenormalized nonlinear extension.

These statements do not have the same logical force as a general impossibility theorem for Standard Model--Einstein emergence. They identify failure of specific constructions under specific scalings. Their constructive use is to constrain the next regulator family: solve or otherwise treat the conformal direction, retain physical detail dynamics, work in a topology supporting typical rough fields, preserve the full chiral source line, and measure a genuinely non-quasi-free response after the complete comparator and contact subtraction.

\section{The strongest available conditional endpoint}
Theorem~\ref{thm:main} is the central compression of the continuum programme. It assembles a sequence of explicit sufficient inputs rather than hiding the remaining work in a single phrase such as ``take the continuum limit.'' Its conclusion concerns a regular physical branch with a reconstructed positive theory, conserved compiled stress, a critical Einstein metric, and a nonzero separated interaction witness. A more ambitious theory with an operator-valued quantum metric, all topological sectors, or global Lorentzian evolution would require additional definitions and estimates not supplied by that theorem.

A useful next result should therefore close a \emph{named quantitative row} on a \emph{single actual family}. Examples include a uniform rough chiral-line/reflection construction, a no-escape estimate for the positive reduced metric law, a common renormalized source comparison, or a nonzero separated excess with a controlled limit. Repeating finite algebraic identities with an added cutoff index would not close any of these rows.

\section{Why the explicit dynamics belongs in the monograph}
The compact and plane-symmetric systems show that the intended endpoint has calculable content before the entire quantum continuum is available. They distinguish massless and massive spinors, phase-dependent colour exchange, Higgs-generated masses, Yukawa backreaction, local anisotropic stress, averaged radiation, compact constraint balance, and finite-time focusing of a normal congruence. The fully assembled toy system conserves its total stress on the matter equations and sources all five Einstein combinations in one convention. Within its declared matter reductions, the two conversion channels share a canonical three-wave structure, while their different spatial vertices determine distinct selection rules and coupling constants. Frequency matching, a nonzero matrix element, and the interpretation of a quantized reduced block remain separate statements. The finite-circle formulas make this particularly explicit: a marginal frequency match can have a vanishing vertex, while the local density factor does not by itself control an infinite-volume rate.

Beyond isolated conversion, shared-mode Hamiltonians identify the
additional dynamical problem, and the classical kinetic closure separates
its conservation and entropy identities from the unproved long-time
limit.  Positive-energy occupation stresses provide an anisotropic
cosmological comparison, while the exact expanding-mode reduction locates
resonance drift without solving the statistical backreaction.  The
hyperbolic-plane polarization system in
Theorem~\ref{thm:polarization-sigma} makes the transverse geometric
response nonlinear and explicit.  These results enlarge the calculable
physical content while retaining their distinct classical, adiabatic,
kinetic-model, and finite-quantum scopes.  The subsequent spectral and
collective chapters go beyond finite quantum blocks, but only through
the explicitly constructed ultraviolet and population limits.  Their
existence and convergence conclusions are not transferred to a larger
field content without the corresponding uniform and source estimates.

The corrected results matter precisely because the programme is quantitative. A momentum constraint changes the quartic coefficients. A stress-sign convention changes a claimed secular cosmological effect into a static spatial modulation. These are not cosmetic alterations. A continuum construction must preserve the same sensitivity to source identity, domains, projections, and normalization at every scale.

The resulting mathematical position is therefore specific: a finite structural and dynamical reconstruction, a broad collection of exact analytic modules and obstruction theorems, a sufficient regular-branch continuum chain, and a set of explicit coupled sectors that test the physical source mechanism. The remaining achievement is to realize the full chain cofinally on one physical family, without discarding the massless sectors, replacing the clock, or normalizing away the Einstein source.


% ==================== 18_appendices.tex ====================
\appendix
\chapter{Additional analytic interfaces}
\label{app:interfaces}
\section{Relative bulk--boundary Ward cohomology}
Boundary repair is not merely a choice of a convenient extension operator. Let $r^*:\mathcal B^\bullet\to\mathcal E^\bullet$ be the restriction of the bulk Ward complex to the boundary complex. The shifted mapping cone defines the relative complex $\mathcal C_{\rm rel}$. Its cohomology fits into
\begin{equation}
\begin{aligned}
 \cdots\longrightarrow H^{p-1}(\mathcal B)
 &\xrightarrow{r^*}H^{p-1}(\mathcal E)
 \longrightarrow H^p(\mathcal C_{\rm rel})\\
 &\longrightarrow H^p(\mathcal B)
 \xrightarrow{r^*}H^p(\mathcal E)\longrightarrow\cdots.
\end{aligned}
 \label{eq:relative-sequence}
\end{equation}
A bulk class admits a relative lift if and only if its restriction class vanishes. A relative class with zero bulk component is a boundary class modulo restrictions of bulk classes. The long exact sequence therefore identifies the actual obstruction before a norm estimate for a lift is attempted. \src{R2}{thm:v2v92-long-exact}

When the obstruction vanishes, a quantitative continuum theorem still needs a bounded trace lift on the physical domain. Common boundary-angle estimates can be protected by private lift directions. Arbitrary shared couplings need not preserve that protection. A common elliptic spectral gap also does not, without uniform coefficient and boundary-symbol assumptions, control chiral Poisson operators near glancing directions. These are the domain-level hypotheses used in Chapter~\ref{ch:closure}.

\section{An assembled Ward numerator can remove a pinch}
For Hilbert spaces $X,Y$, take a smooth compactly supported operator-valued numerator $A(k,\omega)$ in the source's local two-denominator chart and define
\begin{equation}
 \mathcal P_\epsilon[A](k)
 =\int\frac{A(k,\omega)}{(\omega-\ii\epsilon)
                             (\omega-k+\ii\epsilon)}\,\dd\omega.
 \label{eq:pinch}
\end{equation}

\begin{theorem}[Complete-numerator pinch criterion]
\label{thm:pinch}
On the stated local chart, there is a smooth operator-valued remainder $S_A$ such that
\begin{equation}
 \mathcal P_\epsilon[A]\longrightarrow
 -\frac{2\pi\ii A(0,0)}{k-\ii0}+S_A(k)
 \label{eq:pinch-residue}
\end{equation}
in distributions. In particular, $A(0,0)=0$ removes the singular part.
\end{theorem}
\begin{proof}[Proof mechanism]
Localize near the two crossing poles and subtract the value of the complete numerator at their common intersection. The constant part has the explicit two-pole integral in~\eqref{eq:pinch-residue}. The subtracted numerator cancels a transverse factor; terms away from the crossing have uniformly differentiable integrals. Their sum is the smooth remainder.
\end{proof}
\src{R2}{thm:v2v85-residue}
The evaluated quantity is the \emph{assembled} numerator. Bounding or cancelling individual diagrams before forming the common Ward sum can give the wrong residue. A BRST-exact numerator may be smooth after pushforward when the full identity forces its value at the pinch to vanish; this does not license multiplying undefined boundary distributions without the stated nonpinching/domain hypotheses.

\section{Spectral sections and pseudoinverse stability}
Let $A_n$ be the self-adjoint family used to define a retained spectral section and a physical complement. A uniform separating gap and convergence in the required resolvent topology control the Riesz projection and the pseudoinverse on the complement. They also control the homotopy operator used in the Ward or Hodge complex. If the gap collapses, strong-resolvent convergence alone does not give a uniform pseudoinverse bound. The low spectrum must instead be retained as a soft packet and transported explicitly. This is the same principle used for physical Dirac crossings in Chapter~\ref{ch:chiral}. \cite{R2,R3,R12}

\section{Subtraction order and derivative order}
The heat/proper-time analysis separates bulk and boundary ultraviolet orders. On the source's smooth Gaussian branch, subtraction through local weight four leaves a boundary remainder of order $\Lambda^{-1}$ and a closed four-dimensional bulk remainder of order $\Lambda^{-2}$ under its stated symbol bounds. Leaving the marginal weight-four term unsubtracted retains a logarithm. A first metric variation of an interacting expectation is not just the expectation of the varied insertion: it also has the covariance with the varied Gaussian action from~\eqref{eq:normalized-response}. These facts explain why convergence of one unvaried determinant or energy does not automatically prove convergence of its stress response. \cite{R1,R2}

\section{Topology-bearing references and local rough charts}
In a bundle atlas with partition of unity $\rho_j$ and transition functions $g_{ij}$, the source's topology-bearing reference connection is locally
\begin{equation}
 A_i^\sharp=-\sum_j\rho_j\,\dd g_{ij}\,g_{ij}^{-1},
 \qquad
 A_i-A_i^\sharp=\sum_j\rho_j g_{ij}A_jg_{ij}^{-1}.
 \label{eq:topology-reference}
\end{equation}
This separates a fixed transition-function carrier from the variable local connection. Uhlenbeck-type small-energy estimates can then be used on the declared chart without pretending that a nontrivial bundle is globally represented by a single small potential. The reference retains its transition and metric derivatives in the determinant/source comparison. \cite{R4}

\chapter{Result map and remaining hypotheses}
\label{app:map}
The table distinguishes the mathematical object obtained from the additional input required to use it in the full continuum chain. ``Finite'' and ``conditional'' refer to the hypotheses in the text, not to an independent verification of external executable certificates.

\footnotesize
\begin{longtable}{@{}>{\raggedright\arraybackslash}p{0.19\textwidth}>{\raggedright\arraybackslash}p{0.34\textwidth}>{\raggedright\arraybackslash}p{0.41\textwidth}@{}}
\toprule
Module & Strong conclusion retained & Additional continuum requirement\\
\midrule
\endfirsthead
\toprule
Module & Strong conclusion retained & Additional continuum requirement\\
\midrule
\endhead
\bottomrule
\endfoot
Predictive carrier & Complete history--future quotient; all-order positive word construction (Ch.~\ref{ch:finite}) & Cofinal word, support, and physical source transport; not just a finite covariance shadow.\\[1pt]
Local dimension & Conditional $K_4/A_3$ and $3+1$ kinematics (Thm.~\ref{thm:k4}) & Persistent rank margins, vanishing source losses, independently calibrated clock.\\[1pt]
Internal Standard Model & Faithful $\GSM$, charges, three structural generations, common finite Dirac writer (Thm.~\ref{thm:finite-sm}) & Numerical parameters and physical continuum matching remain separate.\\[1pt]
Common action & Exact finite gauge/BRST/source identities and Higgs vacuum (Thm.~\ref{thm:action}) & A common positive quantum regulator and complete differentiated limit.\\[1pt]
Chiral topology & Smooth faithful phase transport; curved index (Thms.~\ref{thm:smooth-phase}, \ref{thm:curved-index}) & Rough support, zero-mode saturation, sector and scale gluing, physical reflection.\\[1pt]
Proper-time source transfer & $C^1$ trace-class transport of subtracted heat integrals (Thm.~\ref{thm:proper-time-c1}) & Realized rough-field subtraction, short-time derivative majorants, retained local and boundary variations.\\[1pt]
Hard/soft elimination & Exact determinant and Berezin Schur factorization (Prop.~\ref{prop:hard-soft}) & Uniform hard inverse, soft packet transport, renormalized source bounds.\\[1pt]
Wilsonian trajectory & Exact shell identities and stable irrelevant graph (Thm.~\ref{thm:stable-graph}) & A bounded physical relevant/marginal base and an invariant local projector.\\[1pt]
Localized shell control & Stable one-sided polynomial tilts and differentiated third-order remainder (Thm.~\ref{thm:localized-shell}) & Summable localized constants, matched lower-order terms, and large-field tails.\\[1pt]
Massless decay & Exact multipole threshold for Coulomb row summability (Thm.~\ref{thm:coulomb}) & Actual endpoint cancellations and infrared control; no replacement by a massive norm.\\[1pt]
Gravity positivity & Fixed-volume conformal obstruction; conditional reduced positive sectors (Ch.~\ref{ch:positive-gravity}) & One physical positive reduction or contour with uniform source and reflection control.\\[1pt]
Constraint coarea & Positive regular-locus law with exact Jacobian (Thm.~\ref{thm:coarea}) & Singular strata, orbit stabilizers, exceptional-set tails, refinement compatibility.\\[1pt]
Ordered geometry & Exact Cartan identities and local torsion selector (Thms.~\ref{thm:cartan}, \ref{thm:torsion-selector}) & Uniform inverse charts and ordered contraction convergence.\\[1pt]
Einstein divergence & Exact Bianchi/intertwining defect identity (Thm.~\ref{thm:einstein-contraction}) & Both realized defects must vanish; uncontracted Bianchi alone is insufficient.\\[1pt]
Constraint-angle control & Exact compression error and three-defect residual bound (Thm.~\ref{thm:three-defect}) & Uniform pseudoinverse floors, angle margin, and common distributional source landing.\\[1pt]
Rigid coarse modes & Exact full-face moment and periodic translation-only kernel (Prop.~\ref{prop:full-face-rigid}) & Fine/coarse Schur margin and cutoff-scaled coercivity on the physical family.\\[1pt]
Block mixing & $M$-matrix covariance and physical-distance estimates (Thm.~\ref{thm:covariance}) & Uniform physical rate and amplitude for the intended massive family.\\[1pt]
Higgs--conformal sign & Incidence-sharp constant, fixed-graph massless reduction (Prop.~\ref{prop:massless-sign}) & Uniform massless sign and normalized mass derivative on a cofinal family.\\[1pt]
Supported conformal response & Dirichlet inverse and normalized source estimate (Thm.~\ref{thm:supported-response}) & Exterior form control for an unrestricted inverse; uniform massless sign and derivative bounds.\\[1pt]
Trees and cycles & Depth-independent radial certificate and conditional cycle margin (Thms.~\ref{thm:tree-margin}, \ref{thm:cycle-restoration}) & Physical coefficient inequalities and complete coupled response estimates.\\[1pt]
Joint probability & Normalized coercivity, nuclear/Sobolev tightness, local normality (Ch.~\ref{ch:states}) & No escape in metric, energy, and topological sectors; moment landing.\\[1pt]
Fixed-spacing existence & Uniform local moments and a homogeneous joint Gibbs law (Thms.~\ref{thm:fixed-spacing-moments}, \ref{thm:fixed-spacing-dlr}) & Fixed spacing only; ultraviolet constants, chiral extension, reflection, and Einstein identification remain separate.\\[1pt]
BRST limit & Zero-mode-safe compactness channels and projective physical states & Common graph domains, anomaly-class control, vanishing channel defects.\\[1pt]
Cofinal comparison & Exact lifted references and summable local expectations (Thm.~\ref{thm:summable}) & Complete source ledger and independent direct-route consistency.\\[1pt]
Reflection & Closed Gram positivity and conditional-half square (Thms.~\ref{thm:half-square}, \ref{thm:rp-limit}) & Actual coupled finite positivity and full reconstruction hypotheses.\\[1pt]
Physical clock & Faithful old/detail product and noncommuting split bound (Thm.~\ref{thm:faithful-clock}, Prop.~\ref{prop:split-clock}) & Local reflected cofinal fibre transport, actual detail generator, and clock-scaled commutator control.\\[1pt]
Stable Higgs reference & Exact Wick separation and finite reference heat bound (Prop.~\ref{prop:stable-wick-reference}) & Physical running coefficients and uniform lower-degree and derivative interaction bounds.\\[1pt]
Holonomy propagation & Regular quotient heat trace; obstruction to bare continuous multiplicity (Prop.~\ref{prop:collective-heat}) & Cutoff-derived interface conditions, stratified control, and reflection sewing.\\[1pt]
Boundary nuclearity & Conditional trace-norm interacting heat limit (Thm.~\ref{thm:nuclearity}) & Every physical sector, collective coordinates, uniform relative form margin.\\[1pt]
Interaction & Elementary Gaussianization and exact composite comparator (Ch.~\ref{ch:stress}) & Nonzero separated complete excess with moment and follower convergence.\\[1pt]
Stress & Exact partition compiler and distributional conservation (Thms.~\ref{thm:response-compiler}, \ref{thm:ward-descent}) & Complete contacts, phase/collar control, enhanced-field compactness.\\[1pt]
Einstein equation & Conditional variational landing (Thm.~\ref{thm:einstein-landing}) & Separate matter subgradient and vanishing transverse critical residual.\\[1pt]
Causal continuation & Retarded inverse and subsidiary constraint propagation & Common physical boundary domain and local nonlinear hyperbolic estimates.\\[1pt]
Gauge charts & Local nonlinear reconstruction and exact heat-link projectivity & Typical support compatibility; the explicit heat law excludes the fixed smooth chart.\\[1pt]
Compact TT gravity & Positive finite seed, exact free spectrum, constrained branches & Nonlinear and renormalized cofinal dynamics; no pathwise substitution for state constraints.\\[1pt]
Quartic gravity & York-corrected finite coefficients and bare self-energy diagnostic & Uniform remainder and interacting operator construction before spectral interpretation.\\[1pt]
Classical waves & Balanced coefficients, $H^2$ continuity, congruence focusing & Higher-order/global spacetime conclusions require further estimates.\\[1pt]
Coupled toy sector & Exact assembled matter equations, stress, Einstein system (Ch.~\ref{ch:coupled}) & Full Standard Model content and quantum cutoff removal are not included.\\[1pt]
Resonant conversion & Explicit vertex selection, population quadrature, and finite quantum blocks (Ch.~\ref{ch:coupled}) & Declared flat or rotating-wave projections; the canonical CAR benchmark uses an independent normalization.\\[1pt]
Two polarizations & Exact weighted sigma model and transverse stress map (Thm.~\ref{thm:polarization-sigma}, Prop.~\ref{prop:polarization-sources}) & Remaining non-diagonal Einstein and fermion equations, full stress, and coupled control.\\[1pt]
Shared-mode dynamics & Finite Hamiltonian, three charges, and a scoped polynomial test (Sec.~\ref{sec:shared-triads}) & Further integrals or dynamical analysis; no general non-integrability conclusion.\\[1pt]
Statistical conversion & Second-order phase response and entropy identity of the closed model (Sec.~\ref{sec:kinetic-conversion}) & Uniform kinetic-time remainder and phase closure; canonical CAR dynamics use a distinct collision factor.\\[1pt]
Finite-volume triads & Coupling scaling and simple-root counting Jacobian (Secs.~\ref{sec:finite-volume-triads}--\ref{sec:triad-density}) & Threshold and joint-limit control; quantum decay uses the separately normalized fixed-parent channel.\\[1pt]
Canonical quantum vertex & CAR box normalization, two parity channels, and fixed-parent phase space (Secs.~\ref{sec:canonical-box}, \ref{sec:fixed-parent-density}) & Declared $1+1$ benchmark; no direct four-dimensional width or full-action invariant truncation.\\[1pt]
Pauli conversion & Exact blocking and correlation defect; entropy of a declared rate model (Sec.~\ref{sec:canonical-pauli}) & Weighted cross-coherence control and a justified quantum kinetic limit.\\[1pt]
One-parent continuum & Self-adjoint norm-resolvent limit, spectral measure, threshold tails (Thms.~\ref{thm:one-parent-renormalized}, \ref{thm:parent-spectral-measure}) & One-excitation projected interaction; a larger coupled field algebra requires a new construction.\\[1pt]
Spectral decay limit & Uniform absolute-error weak-coupling survival and matrix dark-sector reduction (Thms.~\ref{thm:parent-weak-coupling}, \ref{thm:matrix-parent-limit}) & Ultraviolet limit precedes weak coupling; no many-body kinetic or full-group semigroup assertion.\\[1pt]
Collective comparison & Weighted coherence and exact exclusion leakage (Prop.~\ref{prop:weighted-coherence}, Thm.~\ref{thm:hard-pair-leakage}) & Entrywise smallness and fixed-volume coupling scaling do not establish finite-density bosonization.\\[1pt]
Compact pair spectrum & Global Pauli spin flow and Wasserstein quadrature bound (Thm.~\ref{thm:compact-spin-quadrature}) & Fixed compact parameter set; physical-volume ultraviolet constants are not uniform.\\[1pt]
Collective ultraviolet flow & Global physical mild solution and occupation $L^1$ convergence (Thm.~\ref{thm:collective-global-mild}) & Restricted opposite-pair classical model, not microscopic quantum ultraviolet removal.\\[1pt]
Dressed collective energy & Closed renormalized form and reversible energy conservation (Thm.~\ref{thm:collective-energy}) & Finite dressed-energy data; global energy does not determine a local stress tensor.\\[1pt]
Momentum covariance & Global isospectral Pauli flow and explicit mixed residual (Thm.~\ref{thm:pauli-covariance-flow}) & A specified parent--fermion correlation limit; shared pairs cannot be treated as independent spins.\\[1pt]
Many-copy transport & Fixed-time empirical polynomial-moment convergence (Thm.~\ref{thm:many-copy-transport}) & Fixed spatial modes and auxiliary copy scaling; finite species, growing times, and ultraviolet interchange remain separate.\\[1pt]
Expanding matter & Exact Dirac mode, scalar residual, and detuning drift (Sec.~\ref{sec:expanding-resonances}) & Self-consistent redshift, kinetic occupations, and Einstein backreaction.\\[1pt]
\end{longtable}
\normalsize


\chapter{Mathematical scope, corrections, and normalization}
\label{app:sources}
\section{Scope of the cited constructions}
The reference GT denotes the foundational Grand Tensor construction, and
R1--R20 denote successive expositions of the continuum programme.  The
analytic results cited to R17 include the differentiated shell theorem,
fixed-spacing Gibbs construction, supported response estimate, and
massless sign reduction.  The explicit dynamics cited to R20 include its
general-wavenumber and finite-volume results, shared-mode dynamics,
classical kinetic response, expanding modes, two-polarization geometry,
and the canonical, spectral, collective, and momentum-covariance
constructions through V72.  In particular, the retained material
contains one-parent norm-resolvent construction and weak-coupling
spectral decay, exclusion-preserving global mild evolution and dressed
energy, and fixed-time many-copy moment transport.  The analytic details
from R14--R16 include quantitative constraint-angle reduction,
differentiated proper-time trace-class transport, faithful clock
completion, stable boundary reference forms, regular holonomy
propagation, and full-face rigid control.  The foundational references are
retained with the hypotheses of the corresponding reconstruction branch.

The proofs in this monograph distinguish finite algebra, fixed-spacing
probability, conditional uniform transport, and projected dynamics.
Computer-algebra or numerical certificates cited by an underlying
manuscript are not treated as independently verified objects when their
executable files are not available.  No external certificate is required
for the elementary matrix, trigonometric, and population identities
derived explicitly here; those derivations do not establish the omitted
full-field or cutoff-uniform hypotheses.

\section{Corrections that determine the retained mathematics}
The bounded conformal Ward estimate is used as the supported Dirichlet result of Theorem~\ref{thm:supported-response}, not as a global inverse of the unrestricted Witten operator. The fixed-spacing Gibbs existence theorem retains the full common conformal mode and does not assert uniqueness or uniform ultraviolet control. Stable one-sided polynomial tilts replace absolute exponential-moment assumptions that fail for Gaussian polynomials of degree greater than two. The elementary strong-mixing branch is allowed to Gaussianize; a degree-incomplete Gaussian stress comparator is not used as evidence of interaction. Fixed smooth gauge charts are separated from the support of the explicit heat-link projective law. Global-gap statements are restricted to the intended massive sector and are not imposed on a photon-containing endpoint.

For the compact-gravity quartic sector, the York-completed momentum constraint replaces the earlier diagonal ansatz. Equation~\eqref{eq:quartic-york}, the mode-pair table, and the corrected $K_*(3/4)=4$ self-energy diagnostic use the later coefficients. For the gauge--fermion sector, Eqs.~\eqref{eq:corrected-pi}--\eqref{eq:corrected-b} replace the withdrawn uniform polarization growth. The condensate and its reported colour self-energy remain unchanged.

Two further distinctions determine the scope of the explicit dynamics. First, the displayed gauged Dirac equation depends on $b$ through $e^{-a-b}f$, despite an adjacent source sentence saying that it contains no $b$; the equation and the qualified paired-sector conclusion are retained. Second, the homogeneous Higgs oscillator in the two-mode exchange model is a projected, linearized Higgs subsystem, whereas the full assembled Higgs equation contains nonlinear self-interaction and spatial modes. Exact conservation of the reduced system is not promoted to exact closure of those omitted full-field modes.

Focusing refers to the geodesic normal congruence and synchronous development, not automatically to curvature blow-up or maximal-spacetime inextendibility. A perturbative occupation coefficient is not automatically the spectrum of a fully constructed interacting Hamiltonian. Finite sampled or symbolic checks are not used as universal all-parameter results beyond the displayed proved scope. The general-wavenumber statements classify the two retained vertices, not all cubic interactions of the full action. The oscillator--doublet spectrum belongs to an isolated projected quantum model and is not an irreversible continuum decay law.  The continuum decay law in Chapter~\ref{ch:renormalized-conversion} instead follows from a separately defined one-parent Hamiltonian, its ultraviolet limit, and its spectral weak-coupling scaling.

The shared-mode polynomial exclusion is restricted to the displayed
couplings $G_a=1$, $G_b=2$ and degree at most four.  The finite coefficient
rank does not prove non-integrability or a universal coupling statement.
The finite-circle detuning measure is discrete; the inverse-relative-
velocity formula is a simple-root continuum counting Jacobian, not an
already constructed finite-volume density or physical decay width.
The kinetic closure is distinguished from its second-order derivation,
and its entropy identity is not a proof of convergence to equilibrium.

The energy convention also limits the gravitational reading of the
kinetics.  The Yukawa negative-frequency amplitude carries a signed
classical level energy; replacing it by a positive-energy antifermion
occupation requires a particle--hole and quantum normalization not
provided by the amplitude system.  Section~\ref{sec:canonical-box}
supplies that normalization for an independent canonical $1+1$ model,
without identifying the two phase spaces.  Accordingly, the positive-energy
mixture and the Yukawa detuning drift retain their different channel
orientations explicitly.  Similarly, the colour contribution $2Q\Pi_f$
is a signed multiplicative source, not an unconditional cross-polarization
growth law.  The scalar's vanishing direct transverse traceless source
does not remove its indirect gravitational or gauge--Higgs effects.

The canonical rest-parent vertex has two orthogonal standing parity
channels.  The earlier one-channel projection is retained as a classical
reduction, not interpreted as the complete canonical quantum vertex.
Likewise, the fixed-daughter derivative in
\eqref{eq:relative-velocity-jacobian} and the fixed-parent derivative in
\eqref{eq:fixed-parent-jacobian} answer different counting questions;
they are not interchanged.  The canonical occupation polynomial uses
Pauli blocking, whereas the classical triad polynomial and its
Rayleigh--Jeans entropy retain their classical meaning.

The threshold-tail formula is written with the spectral atom and
continuum Fourier term added, as required by
\eqref{eq:parent-spectral-survival}.  This restores the missing additive
separator in the earlier displayed threshold expansion; no multiplicative
bound-state factor is attached to the continuum tail.  The weak-coupling
error is uniform in absolute amplitude, not in relative late-time error.
The bounded noncommuting parent splitting has a compact-rescaled-time
limit, not the all-time estimate proved for the unsplit common-profile
case.

The collective ultraviolet energy theorem uses sharp-cutoff recovery
data in its dressed energy class; bare parent data belong to the mild
state space but not that energy class.  The momentum-covariance theorem
uses the full particle--hole matrix for shared fermions.  Its derivation
from many-copy quantum moments fixes the spatial modes and time interval
before the copy limit, and does not supply the local stress-response
products excluded by \eqref{eq:projected-two-source-defect}.

\section{Notation and normalization}
\begin{center}
\begin{tabularx}{\textwidth}{@{}lX@{}}
\toprule
Symbol & Meaning\\
\midrule
$a_n,h_n$ & Microscopic mesh in continuum estimates; $\ell_n$ is physical block size.\\
$\omega_n,\mu_n$ & Physical functional and, where defined, its positive bosonic law.\\
$S_n,W_n$ & Complete regulated action and $-\log Z_n$; normalization/source derivatives are retained.\\
$Q,\Delta_Q$ & BRST differential and its Hodge operator; distinct from a quartic balance form also denoted $Q$ in Ch.~\ref{ch:compact}.\\
$E,\theta,G$ & Ordered coefficient coframe, coframe one-form, and metric writer in Ch.~\ref{ch:cartan}.\\
$T,\mathsf T$ & Compiled mean stress and enhanced random/source stress packet, respectively.\\
$\Pi,P$ & Declared orthogonal projections; their subspaces are specified locally.\\
$\kappa$ & Gravitational coupling in the general continuum chain. It is absorbed into $T$ in Ch.~\ref{ch:coupled}.\\
$a,b,c$ & Plane-metric functions in the explicit coupled chapters, not ultraviolet mesh parameters.\\
$P,Q,\mathsf M$ & Local two-polarization coordinates and transverse unit-determinant metric matrix in Sec.~\ref{sec:two-polarizations}.\\
$G_\tau,\varrho$ & Canonical triad coupling and a declared detuning density; distinct from the Einstein tensor and spatial energy density.\\
$\varepsilon_N$ & Amplitude of a unit-normalized spatial spinor; not an independently fixed quantum one-particle normalization.\\
$M,m;\ \Omega_q,E_p$ & Scalar and fermion masses, and their positive frequencies, in the canonical $1+1$ benchmark.\\
$\mu_R,S(z)$ & Renormalized parent parameter and subtracted self-energy in Ch.~\ref{ch:renormalized-conversion}.\\
$\beta,s,x;\ \eta$ & Physical-volume collective parent, pair coherence, occupation, and dressed pair field $\eta=s+(h/\omega)\beta$.\\
$\Gamma,\kappa$ & Full particle--hole covariance and its anomalous block in Ch.~\ref{ch:momentum-covariance}; this $\kappa$ is not a gravitational coupling.\\
$N,\mathcal Q_N$ & Auxiliary fermion copy number and total excitation in the many-copy model; distinct from physical species.\\
$W_\pm$ & Metric factors $e^{-2a\pm2b}$; $W=\omega-m$ is a separate spinor amplitude ratio.\\
$\pi$ & $\tfrac12(W_+T_{yy}-W_-T_{xx})$ in the final plane system. Its sign is fixed.\\
\bottomrule
\end{tabularx}
\end{center}

The finite Standard Model potential in Chapter~\ref{ch:action} is $\lambda_H(\norm H^2-v_H^2)^2$. The unitary toy potential in Chapter~\ref{ch:coupled} is $\lambda((v+h)^2-v^2)^2/4$. Their displayed Hessians use their own radial-coordinate normalization. The mode amplitudes and slow times in the resonant systems also follow their declared local conventions; no numerical coefficient is transferred between these conventions without translation.

\backmatter
\begingroup\small
\begin{thebibliography}{R20}
\setlength{\itemsep}{0pt}
\setlength{\parsep}{0pt}
\setlength{\parskip}{0pt}
\addcontentsline{toc}{chapter}{Research sources}
\bibitem[GT]{GT} A.\ P\'elissier. \emph{Renewal Grand Tensor Monograph}, V075. Research manuscript.
\bibitem[R1]{R1} \emph{Renewal Standard Model--Einstein Continuum Research}, extended first cycle, V157. Research manuscript.
\bibitem[R2]{R2} \emph{Renewal Standard Model--Einstein Continuum Research}, second cycle, V100. Research manuscript.
\bibitem[R3]{R3} \emph{Renewal Standard Model--Einstein Continuum Research}, third cycle, V100. Research manuscript.
\bibitem[R4]{R4} \emph{Renewal Standard Model--Einstein Continuum Research}, fourth cycle, V100. Research manuscript.
\bibitem[R5]{R5} \emph{Renewal Standard Model--Einstein Continuum Research}, fifth cycle, V100. Research manuscript.
\bibitem[R6]{R6} \emph{Renewal Standard Model--Einstein Continuum Research}, sixth cycle, V100. Research manuscript.
\bibitem[R7]{R7} \emph{Renewal Standard Model--Einstein Continuum Research}, seventh cycle, V100. Research manuscript.
\bibitem[R8]{R8} \emph{Renewal Standard Model--Einstein Continuum Research}, eighth cycle, V100. Research manuscript.
\bibitem[R9]{R9} \emph{Renewal Standard Model--Einstein Continuum Research}, ninth cycle, V100. Research manuscript.
\bibitem[R10]{R10} \emph{Renewal Standard Model--Einstein Continuum Research}, tenth cycle, V100. Research manuscript.
\bibitem[R11]{R11} \emph{Renewal Standard Model--Einstein Continuum Research}, eleventh cycle, V100. Research manuscript.
\bibitem[R12]{R12} \emph{Renewal Standard Model--Einstein Continuum Research}, twelfth cycle, V100. Research manuscript.
\bibitem[R13]{R13} \emph{Renewal Standard Model--Einstein Continuum Research}, thirteenth cycle, V100. Research manuscript.
\bibitem[R14]{R14} \emph{Renewal Standard Model--Einstein Continuum Research}, fourteenth cycle, V100. Research manuscript.
\bibitem[R15]{R15} \emph{Renewal Standard Model--Einstein Continuum Research}, fifteenth cycle, V100. Research manuscript.
\bibitem[R16]{R16} \emph{Renewal Standard Model--Einstein Continuum Research}, sixteenth cycle, V100. Research manuscript.
\bibitem[R17]{R17} \emph{Renewal Standard Model--Einstein Continuum Research}, seventeenth cycle, V100. Research manuscript.
\bibitem[R18]{R18} \emph{Renewal Standard Model--Einstein Continuum Research}, eighteenth cycle, V100. Research manuscript.
\bibitem[R19]{R19} \emph{Renewal Standard Model--Einstein Continuum Research}, nineteenth cycle, V100. Research manuscript.
\bibitem[R20]{R20} \emph{Renewal Standard Model--Einstein Continuum Research}, twentieth cycle, V72. Research manuscript.
\end{thebibliography}
\endgroup
\end{document}
